\documentclass[11pt,a4paper]{report}
\usepackage[T1]{fontenc}\usepackage[utf8]{inputenc}
\usepackage[margin=2.3cm]{geometry}
\usepackage{amsmath,amssymb,bm,booktabs,longtable,array,siunitx,xcolor,listings,enumitem}
\usepackage[hidelinks]{hyperref}
% Version, date and every enumerated option list come from tools/buildDocs.py, which reads
% toolboxVersion.m and Options.validateValue; the complete API comes from the same generator,
% reading doc/functionlist.csv and the help blocks. auditPackage.py refuses either file when it is
% stale, BY CONTENT, so this document cannot state a version, an option or a signature that the
% toolbox would not honour.
% GENERATED by tools/buildDocs.py from the authorities named beside each macro.
% DO NOT EDIT: auditPackage.py compares this file by content and refuses a stale one.
% Regenerate with: python tools/buildDocs.py
\newcommand{\ngfversion}{0.11.0}                 % +NonGravForces/toolboxVersion.m
\newcommand{\ngfdate}{16 September 2026}          % toolboxVersion.m .date
\newcommand{\ngfmodel}{Claude Opus 5 (Anthropic)}      % toolboxVersion.m .model
\newcommand{\ngfunitcount}{56}      % rows of doc/functionlist.csv
%% The admissible option values, from Options.validateValue's `sets` struct:
\newcommand{\ngfoptdragdensity}{\texttt{exponential}, \texttt{msis00}, \texttt{msis21}}
\newcommand{\ngfoptdragwind}{\texttt{none}, \texttt{corotation}, \texttt{hwm14}}
\newcommand{\ngfoptdraggsi}{\texttt{constantCd}, \texttt{sentman}, \texttt{maxwell}}
\newcommand{\ngfoptdragaccommodation}{\texttt{fixed}, \texttt{pilinski2010}, \texttt{pilinski2011}, \texttt{walker2014}, \texttt{sesam}}
\newcommand{\ngfoptsrpeclipse}{\texttt{cylindrical}, \texttt{conical}, \texttt{solaarsCF}}
\newcommand{\ngfopterpbackend}{\texttt{none}, \texttt{pointSource}, \texttt{knocke}, \texttt{ceres}}
\newcommand{\ngfoptthermalmodel}{\texttt{none}, \texttt{steadyState}, \texttt{transient}}
\newcommand{\ngfoptgeomselfShadow}{\texttt{none}, \texttt{lut}, \texttt{raycast}}
\newcommand{\ngfoptsunsource}{\texttt{rs.analytic}, \texttt{rs.de440}}
\newcommand{\ngfoptpartialsmethod}{\texttt{analytic}, \texttt{finiteDiff}}

\sisetup{per-mode=symbol,detect-all}
\lstset{basicstyle=\ttfamily\small,columns=fullflexible,frame=single,breaklines=true,
        keepspaces=true,backgroundcolor=\color{gray!8},showstringspaces=false}
\newcommand{\vv}[1]{\bm{#1}}
\newcommand{\ngf}{\texttt{NonGravForces}}

\title{\ngf: User Guide\\[6pt]
\large Non-gravitational forces and torques on LEO satellites\\[4pt]
\normalsize Theory, installation, processing, the integrator interface, and the complete API\\[2pt]
\normalsize Version \ngfversion\ (concept~3)}
\author{Prepared by Claude Opus 5 (Anthropic) for Matthias Weigelt}
\date{\ngfdate \quad\textperiodcentered\quad MATLAB R2026a Update 4}

\begin{document}
\maketitle

\begin{abstract}
\ngf\ computes the non-gravitational acceleration, the surface and body torques, and their partial
derivatives for a spacecraft described by a panel model: atmospheric drag and lift, solar radiation
pressure, Earth albedo and infrared, and thermal re-radiation. Its reference case is GRACE-FO
accelerometer processing at the \SI{e-10}{\metre\per\second\squared} level.

This guide is written for someone who has not used the toolbox before. It installs it from nothing,
explains the physics it implements and the approximations it makes, walks through a GRACE-type
processing chain step by step, shows how to drive an orbit integrator through the compiled
interface, and gives worked examples with the output they actually produce on the reference day.
It ends with what the toolbox cannot do --- stated plainly, because a limitation discovered by a
user is a defect and a limitation written down is a specification --- and with the complete API,
which is generated from the code so that it cannot drift from it.
\end{abstract}

\tableofcontents
\clearpage

%=======================================================================================
\chapter{Introduction}

\section{What this toolbox is}
\ngf\ computes the non-gravitational acceleration $\vv a_{ng}$, the torques $\vv\tau$, and their
partials, given a spacecraft's state, attitude and environment. One rule governs every other:

\begin{quote}
\textbf{\ngf\ computes forces and torques, and applies them; it implements nothing it can call.}
\end{quote}

There is no gravity field in it, no equations of motion, no integrator and no least-squares
machinery: since 0.13.0 the calibration application composes \texttt{timeSeriesPreparation.run}
and \texttt{lsqvce.fit} instead of solving its own problem. Everything else is delegated to a
named dependency, by qualified call, never copied or wrapped:

\begin{table}[h]\centering\footnotesize
\begin{tabular}{@{}lp{9.6cm}@{}}\toprule
Toolbox & What it owns\\\midrule
\texttt{referenceSystem} & Time scales, Earth orientation, ITRS$\leftrightarrow$GCRS with
 derivatives, geodesy, quaternions, attitude interpolation, Sun and Moon states\\
\texttt{gracefoReading} & Level-1B ASCII reading, quality flags, GROOPS instrument export\\
\texttt{gjToolbox} (odeSuite) & Gauss--Jackson integration. Optional for the toolbox, required by
 the gate: \texttt{app.DynamicsSource} implements that toolbox's frozen dynamics interface\\
\texttt{timeSeriesPreparation} & The \texttt{Arc}, the L1B CRN kernel, mask, align, QC, noise,
 observability, bootstrap and \texttt{run}. Optional for the toolbox, required by the gate: the
 default accelerometer filter and the calibration's instrument stages\\
\texttt{lsqvce} & Least squares, variance components, whitening, robust reweighting. Optional for
 the toolbox, required by the gate: every number of \texttt{app.fitCalibration}\\
\bottomrule
\end{tabular}
\end{table}

\noindent This partition is the single most important design decision in the project. It is what
makes the toolbox smaller than either of its two rejected predecessors while covering more physics.

\section{The layer map}
Nine layers; a layer calls only lower layers plus the dependencies.

\begin{lstlisting}
L8 app      calibrate, fitCalibration, Calibration, syntheticPair, forceBudget, densityInversion,
            DynamicsSource (integrator entry, native C)
L7 qc       residuals, hampel, detectImpulses, detectJumps, eventTable, screen
L6 sensor   Sensor (abstract), Electrostatic, AtomInterferometer, DragFree
L5 Model    accel | torque | series | partials | budget      <- single point of access
L4 kernel   gsiCoefficients, accommodation, dragPanelForce, radiationPanelForce,
            thermalRecoil, assemble, bodyTorques, partials, singleEpoch (C/MEX)
L3 geom     MacroModel, graceFO, grace, boxWing, selfShadow, flowShadow, thermalState,
            AttitudeProvider, checkAttitudeSense
L2 env      density, wind, eclipseFactor, earthRadiance, magneticField, radiationSources
L1 io       readACC1B/SCA1B/GNV1B/THR1B/MAS1B/MAG1B, massModel, spaceWeather,
            importCERES, importGroopsGrid, igrfCoefficients
L0 core     const, toolboxVersion, dependencyReport, fail, provenance, Options, buildMex
\end{lstlisting}

\noindent \ngfunitcount{} public units are catalogued in \texttt{doc/functionlist.csv}, which is the
authority: continuous integration checks it against the files on disk in both directions.

\section{Design principles a user will notice}
Most of the architecture is invisible from outside. Four decisions are not.

\begin{description}[leftmargin=1.4em,style=nextline]
\item[Configure once, evaluate many.] \texttt{Options} and \texttt{Model} are built once and then
 evaluated over many epochs. Everything is vectorised over epochs; a $3\times K$ array is the
 normal shape of a result.
\item[Sources, not effects.] \texttt{env} produces a list of radiation sources (Sun, Earth tiles)
 and a gas state; the kernel does not know which effect it is computing. Albedo and sunlight go
 through the same panel code with different fluxes, which is why they cannot disagree.
\item[Never degrade silently.] There is no fallback between fidelity tiers. An unknown option field
 is an error; a missing optional dependency is an error at the point of use, naming its setup
 function; a model evaluated outside the range it was fitted in warns and says by how much. A
 \texttt{NaN} where an error belonged is the worst outcome available and is treated as a defect.
\item[One raise point.] Every error is built by \texttt{NonGravForces.fail}, so every identifier has
 the form \texttt{NonGravForces:<unit>:<mnemonic>} and every message ends with what the caller
 should change, after ``\texttt{ -> }''.
\end{description}

%=======================================================================================
\chapter{Installation}
\label{ch:install}

\section{What you need first}
\begin{itemize}[leftmargin=1.4em]
\item \textbf{MATLAB R2026a} or newer. The toolbox uses \emph{base MATLAB and the Parallel Computing
 Toolbox only}. It deliberately calls nothing from Signal Processing, Statistics, Optimization,
 Curve Fitting, Symbolic, Aerospace or Mapping --- not in code and not in an example, because an
 example that needs a product the toolbox refuses to need has misled the reader who copied it.
 Continuous integration scans for this.
\item \textbf{A C compiler}, for the compiled density and kernel paths. On Windows this is the
 MinGW-w64 support package. MATLAB routinely forgets where it is; \texttt{NonGravForces.buildMex}
 locates it itself and runs \texttt{mex -setup C} for the session, so you should not have to.
\item \textbf{Git}, and roughly \SI{1}{\giga\byte} of disk if you fetch the optional data.
\end{itemize}

\noindent You do \emph{not} need \LaTeX, GROOPS, or any of the third-party atmosphere models: the
last are fetched from their maintainers at pinned versions when you ask for them.

\section{Getting the code}
Clone the toolbox and its two mandatory dependencies \emph{as siblings}, in one directory:

\begin{lstlisting}[language=bash]
cd ~/MATLAB
git clone https://github.com/astrobayer/NonGravForces.git
git clone https://github.com/astrobayer/referenceSystem.git
git clone https://github.com/astrobayer/gracefoReading.git
# optional, and required only to drive an integrator or to run the full gate:
git clone <your odeSuite remote> odeSuite
\end{lstlisting}

\noindent The sibling layout matters: the gate and the CI job both look for the dependencies next to
the repository.

\section{Setting up, in order}
Run each dependency's setup first, then the toolbox's. Order matters, because the toolbox's setup
probes what is already on the path.

\begin{lstlisting}
run('~/MATLAB/referenceSystem/setupReferenceSystem.m')
run('~/MATLAB/gracefoReading/setupGraceFOReading.m')
run('~/MATLAB/odeSuite/gjToolboxSetup.m')        % only if you cloned it
cd('~/MATLAB/NonGravForces')
setupNonGravForces
\end{lstlisting}

\noindent \texttt{setupNonGravForces} adds the toolbox root, \texttt{tests/} and \texttt{tools/} to
the path and registers the documentation. It prints a line for everything it does.

\section{What setup fetches, and what it refuses to}
\textbf{Setup is the only place in the toolbox that touches the network}, it does so only when
something is missing, and it says so when it does. \texttt{NonGravForces.buildMex} never fetches:
it refuses and names setup. An installer may reach the network; a build step may not.

\begin{table}[h]\centering\footnotesize
\begin{tabular}{@{}llp{7.2cm}@{}}\toprule
Option & Default & What it controls\\\midrule
\texttt{'FetchMsis'} & \texttt{"ifMissing"} & NRLMSISE-00 and NRLMSIS 2.1 source, cloned at a
 pinned commit from its maintainer. Not redistributed here: it carries no redistribution grant.\\
\texttt{'FetchData'} & \texttt{"ifMissing"} & IERS Earth-orientation and CelesTrak space-weather
 files, and GRACE-FO Level-1B from the GFZ ISDC archive.\\
\texttt{'FetchKernels'} & \texttt{"ifMissing"} & The DE440 planetary ephemeris, pinned by SHA-256.\\
\texttt{'FetchGroops'} & \texttt{"never"} & The GROOPS CERES grids. Default \emph{never}, because
 two of the files are \SI{1.25}{\giga\byte} \emph{each}.\\
\texttt{'DataPool'} & --- & Where the large external data live, if you keep them outside the
 repository.\\
\bottomrule
\end{tabular}
\end{table}

\noindent Every fetch is verified by content, not by HTTP status: a status code is not a file, and
one of these URLs answers \texttt{200} with an HTML error page. Every downloaded file is hashed on
\emph{every} call, not only when it is first downloaded, because a check that runs only on download
protects only the first run.

To install without any network access at all:

\begin{lstlisting}
setupNonGravForces('FetchMsis', "never", 'FetchData', "never", 'FetchKernels', "never")
\end{lstlisting}

\noindent Everything except the \texttt{msis00} and \texttt{msis21} density backends works without
any of it.

\section{Building the compiled parts}
\begin{lstlisting}
NonGravForces.buildMex                 % all targets
NonGravForces.buildMex('Check', true)  % report only: what is built, and from which sources
\end{lstlisting}

\noindent There are five targets: \texttt{msis}, \texttt{msis21}, \texttt{hwm14}, \texttt{kernel}
(the single-epoch force kernel) and \texttt{rhs} (the integrator's right-hand side, which needs
\texttt{gjToolbox}). \texttt{buildMex} reports an unbuildable target rather than raising, so a user
without \texttt{gjToolbox} still gets the other four.

\section{Verifying the installation}
\begin{lstlisting}
runNgfTests                      % the suite; 'ExcludeSlow', false adds the speed tag
NonGravForces.dependencyReport   % what is present, which version, and what is missing
doc NonGravForces                % the browsable documentation
\end{lstlisting}

\noindent If you want exactly what continuous integration runs --- dependencies installed from the
sibling folders, the suite, and a single verdict --- run \texttt{matlab -batch ciGate}. It raises
\texttt{NonGravForces:ciGate:red} on any failure, so a batch call exits non-zero. It restores your
saved MATLAB path afterwards: dependency setups call \texttt{savepath}, and a gate has no business
rewriting your environment.

\section{When something goes wrong}
\begin{description}[leftmargin=1.4em,style=nextline]
\item[\texttt{buildMex} says it cannot find a compiler.] Install the MinGW-w64 support package from
 the Add-On Explorer. The toolbox does not attempt to install it for you.
\item[A MEX file cannot be overwritten on Windows.] A mapped DLL is locked, and so is every ancestor
 folder. Restart MATLAB and build before calling anything else; \texttt{clear mex} was measured not
 to release it once an OpenMP thread pool is live.
\item[A test is \emph{incomplete} rather than failing.] It skipped because a fixture is missing.
 Fixtures are committed, never fetched, precisely so that this cannot happen --- if it does, the
 checkout is damaged.
\item[\texttt{density} refuses with \texttt{inputDrivers}.] The \texttt{msis} backends need space
 weather. Pass a CelesTrak \texttt{spaceWeather} struct, or explicit
 \texttt{'Drivers', struct('f107',\ldots,'f107a',\ldots,'ap',\ldots)}.
\end{description}


%=======================================================================================
\chapter{Theory and implementation}
\label{ch:theory}

\noindent The sections of this chapter are shared, verbatim, with the standalone reference
document \texttt{NonGravForces\_Theory\_and\_Implementation.pdf}: both documents input one
source, so the theory exists once and cannot differ between them. Equations are given in the
form the code implements, and every approximation is named with its magnitude.

% SHARED SECTION, \input by both NonGravForces_Theory_and_Implementation.tex and
% NonGravForces_UserGuide.tex. It carries no preamble and opens no document, so that the
% theory exists ONCE and both documents render the same words (R1).
\section{Scope, notation and the force hierarchy}
\subsection{Scope}
\texttt{NonGravForces} computes the non-gravitational acceleration $\vv a_{ng}$, the surface and body torques $\vv\tau$, and their partial derivatives, for a spacecraft described by a panel (macro) model, given its state, attitude and environment. It computes nothing else: time scales, Earth orientation, frame transformations, geodetic coordinates, ephemerides and attitude kinematics are the \texttt{referenceSystem} toolbox; Level-1B reading is \texttt{gracefoReading}; orbit integration is \texttt{gjToolbox}, which \texttt{app.DynamicsSource} serves through that toolbox's frozen dynamics interface. Parameter estimation is NOT delegated: the calibration problem is linear and separates by axis, so \texttt{app.calibrate} carries its own equilibrated QR (owner ruling of 2026-09-09; \texttt{lsqvce} was a dependency until then). This partition (HANDOVER Part~A.2) is the single most important design decision and it is what makes the toolbox smaller than either predecessor while covering more physics.

\subsection{Notation}
$\vv r,\vv v$ position and velocity in GCRS [\si{\metre}, \si{\metre\per\second}]; $\vv q$ attitude quaternion BODY$\to$GCRS (scalar-first Hamilton, passive convention of \texttt{referenceSystem}); $R_{b}$ the rotation matrix; $m$ spacecraft mass; panel $k$ with area $A_k$, outward unit normal $\hat{\vv n}_k$, centroid $\vv c_k$ (body frame), optical triplets $(\alpha,\rho_s,\rho_d)$ in the visible band and $(\varepsilon,\rho_s,\rho_d)$ in the infrared, energy accommodation $\alpha_E$. $\hat{\vv s}$ unit vector to the Sun, $\hat{\vv u}=\vv v_r/|\vv v_r|$ the relative-flow direction, $\theta_k$ the incidence angle with $\gamma_k=\cos\theta_k=\hat{\vv n}_k\!\cdot\!(-\hat{\vv u})$ for drag and $\hat{\vv n}_k\!\cdot\!\hat{\vv s}$ for radiation. SI throughout; angles in radians; $K$ epochs are stored as $3\times K$ arrays.

\subsection{Force hierarchy at the GRACE-FO altitude ($\sim$\SI{490}{\kilo\metre})}
\begin{table}[h]\centering\footnotesize
\begin{tabular}{@{}lllp{4.4cm}@{}}\toprule
Effect & along-track & cross-track & source / note\\\midrule
atmospheric drag & \num{5e-8} (solar min) \ldots\ \num{5e-7} (max) & $\sim\num{e-9}$ & Vielberg \& Kusche 2020; Wöske et al.\ 2019\\
solar radiation pressure & \SI{13}{nm/s^2} RMS & \SI{21}{nm/s^2} & Siemes et al.\ 2023, GRACE-A 2008\\
thermal re-radiation (TRP) & \SI{0.8}{nm/s^2} & \SI{6.4}{nm/s^2} & second-largest radiation term\\
Earth albedo & \SI{0.2}{nm/s^2} & \SI{0.2}{nm/s^2} & $<\SI{1}{nm/s^2}$ combined with IR\\
Earth infrared & \SI{0.2}{nm/s^2} & \SI{0.1}{nm/s^2} & \\
thruster impulses & up to \num{e-6} & & THR1B; Bandikova et al.\ 2019\\
CoM-offset coupling & $\dot{\vv\omega}\times\Delta\vv r+\vv\omega\times(\vv\omega\times\Delta\vv r)$ & & $|\Delta\vv r|\le\SI{100}{\micro\metre}$ (JPL D-56935)\\
antenna thrust, Lorentz & $\lesssim\num{e-11}$ & & hooks only\\
accelerometer noise floor & \SI{1e-10}{m\,s^{-2}\,Hz^{-1/2}} & $\sim10\times$ worse & Kornfeld et al.\ 2019\\\bottomrule
\end{tabular}
\caption{Magnitudes in \si{\metre\per\second\squared} unless stated. Every unit of the toolbox is validated against these orders of magnitude before any finer oracle.}
\end{table}

\section{Atmospheric drag and lift}
\subsection{Free-molecular flow and the per-panel force}
At \SI{400}{}--\SI{500}{\kilo\metre} the Knudsen number is $\gg1$; molecules do not interact with each other near the surface and the force is the integral over the incident Maxwellian of the momentum exchanged at each panel. With mass density $\rho$, relative velocity $\vv v_r$ and per-panel pressure and shear coefficients $C_{p,k}, C_{\tau,k}$,
\begin{equation}
\vv F_{drag}=-\tfrac12\rho\,|\vv v_r|^2\sum_k A_k\left[C_{p,k}\,\hat{\vv n}_k+C_{\tau,k}\,\hat{\vv t}_k\right],\qquad
\hat{\vv t}_k=\frac{\hat{\vv u}-(\hat{\vv u}\!\cdot\!\hat{\vv n}_k)\hat{\vv n}_k}{|\hat{\vv u}-(\hat{\vv u}\!\cdot\!\hat{\vv n}_k)\hat{\vv n}_k|},
\label{eq:drag}
\end{equation}
where the sign convention makes $C_p>0$ push the panel inward along $-\hat{\vv n}_k$ for a windward panel. The projection onto $\hat{\vv u}$ is drag, the remainder is lift plus side force. The kernel returns the per-panel $\vv F_k$ (principle P2); $\vv F$ and $\vv\tau=\sum_k(\vv c_k-\vv r_{cm})\times\vv F_k$ are reductions.

\subsection{Relative velocity}
\begin{equation}
\vv v_r=\vv v-\vv\omega_\oplus\times\vv r-\vv v_{wind},
\end{equation}
with the co-rotation term formed in \texttt{referenceSystem.applyRotation} (it carries $\vv\omega\times\vv r$ exactly, never by hand) and $\vv v_{wind}$ from HWM14 (Drob et al.\ 2015; quiet-time RMS \SI{40}{}--\SI{80}{\metre\per\second}) as an opt-in. At \SI{7.6}{\kilo\metre\per\second} the co-rotation term is \SI{\sim470}{\metre\per\second} at the equator and enters cross-track drag at the \SI{6}{\percent} level; winds of \SI{100}{\metre\per\second} enter at \SI{1.3}{\percent} of the along-track force but dominate the cross-track signal (Doornbos 2011).

\subsection{Sentman's model (exact form implemented)}
Let $s=|\vv v_r|/\sqrt{2\mathcal RT_\infty/M}$ be the speed ratio with $M$ the mean molar mass of the local gas (\emph{not} air: at \SI{500}{\kilo\metre} atomic oxygen with $M\approx\SI{16}{g/mol}$ dominates, so $s\approx8$; using $R_{air}=\SI{287}{J\,kg^{-1}\,K^{-1}}$ as v1 did biases $s$ by $\sqrt{29/16}=1.35$), $\gamma=\cos\theta$, $\ell=\sin\theta$, and
\begin{equation}
P=\frac{e^{-\gamma^2s^2}}{s},\quad Q=1+\frac1{2s^2},\quad G=\frac1{2s^2},\quad Z=1+\operatorname{erf}(\gamma s),\quad
\frac{v_{re}}{v_{in}}=\sqrt{\tfrac12\Big[1+\alpha_E\Big(\frac{4\mathcal RT_w}{M|\vv v_r|^2}-1\Big)\Big]}.
\end{equation}
\textbf{Which pair is which.} Sentman's published coefficients, and Doornbos (2011, eqs 3.55--3.56), are \emph{drag} and \emph{lift} --- along and across $\vv v_r$:
\begin{align}
C_D&=\frac{P}{\sqrt\pi}+\gamma QZ+\frac{\gamma}{2}\frac{v_{re}}{v_{in}}\big(\gamma\sqrt\pi Z+P\big),\label{eq:sentmanCD}\\
C_L&=\ell GZ+\frac{\ell}{2}\frac{v_{re}}{v_{in}}\big(\gamma\sqrt\pi Z+P\big).\label{eq:sentmanCL}
\end{align}
The toolbox does not use that pair. Its kernel is written in the plate's own frame, \emph{normal} and \emph{tangential} (eq.~\eqref{eq:drag}), which is what a per-panel force and a torque both need, and it obtains those from the same momentum integrals:
\begin{align}
C_p&=\underbrace{\frac{\gamma P}{\sqrt\pi}+\Big(\gamma^2+\frac1{2s^2}\Big)Z}_{\text{incident}}+\underbrace{\frac12\frac{v_{re}}{v_{in}}\big(\gamma\sqrt\pi Z+P\big)}_{\text{re-emitted}},\label{eq:sentmanCp}\\
C_\tau&=\ell\Big(\frac{P}{\sqrt\pi}+\gamma Z\Big).\label{eq:sentmanCtau}
\end{align}
The two pairs are one rotation apart, $C_D=C_p\gamma+C_\tau\ell$ and $C_L=C_p\ell-C_\tau\gamma$, and the toolbox asserts both projections against eqs.~\eqref{eq:sentmanCD}--\eqref{eq:sentmanCL}. \textbf{This document labelled eqs.~\eqref{eq:sentmanCD}--\eqref{eq:sentmanCL} $C_p$ and $C_\tau$ until 0.10.0}, which is the same confusion that produced defect D1 below; nothing in the code ever read it, and it is corrected here. Note in particular that the \emph{tangential} coefficient carries no $1/s^2$ term: the lone $Z/(2s^2)$ belongs to the lift. The limit that settles it is edge-on ($\gamma\to0$): $C_p\to1/(2s^2)$ and $C_\tau\to\ell/(s\sqrt\pi)$, and the hyperthermal ram limit is $C_p\to2\gamma^2$, $C_\tau\to2\gamma\ell$, $C_D\to2\gamma$.

\paragraph{How good is the closed form?} Incomplete accommodation is treated in three ways in the
literature, and this toolbox implements the first: (a) the energy-flux form above, where
$v_{re}/v_{in}$ carries $(1-\alpha_E)/2$ of the incident energy; (b) the kinetic-temperature form
$T_{k,r}=T_{k,i}(1-\alpha_E)+\alpha_E T_w$ with $T_{k,i}=mU^2/3k$, which carries $(1-\alpha_E)\,2/3$
(Walker et al.\ 2014, eqs 10--12); and (c) sampling the Cercignani--Lampis--Lord kernels, which is the
physically correct reflected velocity distribution and requires a DSMC rather than a closed form.
At $\alpha_E=1$ (a) and (b) are the SAME NUMBER, and eq.~\eqref{eq:sentmanCp} reproduces Walker's
published flat-plate coefficient to \num{1.6e-15} there --- the only check of this chain against
arithmetic written outside this project. Below full accommodation they separate: measured on the
plate, (a) reads \SI{1.7}{\percent} below (b) at $\alpha_E=0.95$, \SI{2.4}{\percent} at $0.90$ and
\SI{5.6}{\percent} at $\alpha_E=0$. Walker et al.\ measure (b) against (c) for the GRACE geometry at
\SI{\sim1.4}{\percent} for $\alpha_E=0.9$ and \SI{\sim1.7}{\percent} at $\alpha_E=0$. Both are of order
\SIrange{1}{2}{\percent} of $C_D$ at the $\alpha_E\approx0.85$--$1$ of a real LEO surface, inside the
\SIrange{5}{10}{\percent} that $\alpha_E$ itself is uncertain by --- so the closed form is kept, and
the deviation is stated instead of assumed to be zero.

Diffuse re-emission is isotropic about the normal, so it contributes to $C_p$ and \emph{not} to $C_\tau$ --- which is why $C_\tau$ carries no $v_{re}/v_{in}$ while $C_L$ does. The re-emission ratio carries the wall temperature $T_w$ and the accommodation coefficient $\alpha_E$; complete accommodation ($\alpha_E=1$) gives diffuse re-emission at $T_w$. \textbf{The gas force is therefore not antiparallel to $\vv v_r$}: $C_L\ne0$ off the ram direction, \SI{0.34}{\degree} on GRACE-FO (measured). Only \texttt{constantCd} is antiparallel, by construction. \textbf{Defect D1 of v1}: the exponential group and the whole shear term were divided by $s$ instead of $s^2$, an error of $s\approx8$ in the shear on inclined panels; the toolbox pins this as a regression test with the correct value from an independent kinetic-theory integration (\texttt{tools/pyxcheck/sentman\_check.py}).

\subsection{Other gas--surface interaction models}
\emph{Maxwell}: a fraction $\sigma$ diffuse, $1-\sigma$ specular. \emph{Cercignani--Lampis--Lord (CLL)}: separate normal and tangential accommodation $(\alpha_n,\sigma_t)$ with quasi-specular lobes; used by Walker et al.\ (2014) and Mehta et al.\ (2014). \emph{Schamberg}: quasi-specular lobe of adjustable width. All are \texttt{drag.gsi} options; none is in a preset except Sentman, because the density error (\S\ref{sec:density}) exceeds their differences except below \SI{300}{\kilo\metre}.

\subsection{Energy accommodation}
Published values: March et al.\ (2021) optimise a \emph{constant} $\alpha_E$ per mission against attitude manoeuvres and wind analyses and obtain $0.85$ (CHAMP) and $0.82$ (GOCE), recommending $\alpha_E=0.85$ and superseding the $0.93$ of Doornbos (2011); $0.89$ from CHAMP yaw (Pilinski \& Argrow 2013). They fit rather than compute \emph{because} the in-situ composition and temperature a physical model needs are not measured on any of these satellites. The SESAM model (Pilinski et al.\ 2010, 2013) makes $\alpha_E$ a function of the atomic-oxygen partial pressure $p_O$ through a Langmuir isotherm $\Theta=Kp_O/(1+Kp_O)$, $\alpha_E=\Theta\,\alpha_{ads}+(1-\Theta)\,\alpha_{clean}$, reproducing measurements to \SI{\sim3}{\percent}. The toolbox implements \ngfoptdragaccommodation. \texttt{pilinski2010} is the Langmuir isotherm in $n_O T$ (Pilinski 2011, \S7.3); \texttt{walker2014} carries the adsorbate coverage with a clean-substrate $\alpha_s=3\mu/(1+\mu)^2$ (Paul et al.\ 2022, eqs 10--13). \texttt{sesam} is the model of Pilinski et al.\ (2013) itself: the same mixture, but with $K_L$ a function of the incident kinetic energy through an initial sticking coefficient (their eq.~26b, the fraction of the incident Maxwellian falling below a binding energy of \SI{5.7}{\electronvolt}), and with $P_O$ a \emph{ram} pressure $\tfrac12\rho_O V_r^2 C_{D,\mathrm{in},O}$ rather than $n_O k T$ --- at \SI{520}{\kilo\metre} a factor of 140, so the two parameterisations' Langmuir constants are not interchangeable either. Its substrate coefficient uses Goodman's $2.4\mu/(1+\mu)^2$ where \texttt{walker2014} ships $3.0$; the sources disagree and each option follows its own. \textbf{Not implemented from that paper}: the DSMC lookup tables and the tumbling-sphere shape factor, which map $\alpha$ to a drag coefficient --- this toolbox has its own panel geometry and Sentman kernel and takes only $\alpha$ from SESAM (HANDOVER O-16). A document that offers an option the code refuses is worse than one that offers none, which is why this list is generated from \texttt{Options.validateValue}.

\textbf{All three presets take the fitted $\alpha_E=0.85$, and the adsorption models are opt-ins.} That is a ruling (HANDOVER O-42, 2026-09-17), not an approximation, and the reason is measured. Evaluated at CHAMP's altitude and CHAMP's solar activity (\SI{300}{}--\SI{350}{\kilo\metre}, $F_{10.7}=70$), \texttt{pilinski2010} returns $0.86$--$0.95$, straddling the $0.85$--$0.93$ March et al.\ (2021) report for SESAM there: the adsorption family reproduces the literature where the literature was fitted. The same code returns $0.11$--$0.17$ at GRACE-FO's \SI{520}{\kilo\metre} in deep solar minimum, and $0.63$--$0.83$ at that altitude at $F_{10.7}=180$, because the atomic-oxygen partial pressure driving $\Theta$ falls by two orders of magnitude between them ($\Theta=0.83$ at \SI{300}{\kilo\metre}, $0.05$ at \SI{520}{\kilo\metre}). So these models predict a factor-of-five solar-cycle swing in GRACE-FO's accommodation, and \SI{2020}{}-01-01---the day this toolbox is validated on---sits at its extreme low end. On that day the choice is worth \SI{17.5}{\percent} of the drag (\texttt{walker2014} \SI{8.09e-9}{\metre\per\second\squared} against \SI{6.88e-9}{} for $\alpha_E=0.85$; $0.93$ gives \SI{6.62e-9}{}), which for accelerometer calibration is absorbed by the estimated scale factor and propagates into any inverted density. Both positions are defensible; the presets take the one the missions are actually processed with, and the other is one option away.

\subsection{Density and composition}\label{sec:density}
Backends, as \texttt{drag.density} admits them: \ngfoptdragdensity\ --- exponential (sizing only), NRLMSISE-00 (Picone et al.\ 2002; the Brodowski C port, whose 17-case DOCUMENTATION table is the bit-exact oracle) and NRLMSIS~2.1 (Emmert et al.\ 2021, NRL's unmodified Fortran, checked against its own 200 reference cases). DTM2020 (Bruinsma \& Boniface 2021) and JB2008 (Bowman et al.\ 2008) are \emph{not} implemented and were removed from the option set in 0.10.0, because an option that cannot be evaluated is a promise the code does not keep; their space-weather drivers stay readable in \texttt{io.spaceWeather}. Inputs: $F_{10.7}$ (\emph{observed}, previous day, and 81-day mean; v1 fed adjusted same-day values, $\pm\SI{3.4}{\percent}$), $F_{30}$, $S_{10.7}, M_{10.7}, Y_{10.7}$ (JB2008, from SET), $A_p$ 3-hourly history (the \texttt{ap7} vector), $K_p$/Hp, $D_{st}$ for storms. Uncertainty: \SI{10}{}--\SI{30}{\percent} between models and data sets; NRLMSIS~2.1 is \SI{\sim10}{\percent} below NRLMSISE-00 at GRACE altitudes (Guo \& Zhang 2024). Composition is required by the GSI ($M$, and $p_O$ for SESAM). The models are evaluated at the \emph{geodetic} height on WGS-84: a spherical Earth gives \SI{-21.4}{\kilo\metre} at the poles and a density factor of 1.40 on a polar mission (v1 defect D2); the conversion is \texttt{referenceSystem.cartesianToGeodetic} with the ellipsoid passed explicitly (its default is GRS80).

\subsection{Macro model and self-shadowing}
GRACE-FO (JPL D-56935 Table~5): eight plates, side panels at \SI{50}{\degree} with normals printed as $\sin50^\circ=0.766044$, $\cos50^\circ=0.642787$ (their squares sum to $1-\num{5e-7}$; the toolbox accepts $10^{-5}$ and documents it), total area \SI{16.9163284}{\metre\squared}, mass \SI{601.214}{kg} (C) / \SI{601.209}{kg} (D), inertia tensor Table~4. The body is not convex: the flanks expose their undersides, so $\sum_kA_k\hat{\vv n}_k\ne0$ and self-shadowing is physical. Wöske et al.\ (2019) quote $\sim\SI{10}{\percent}$ force error of an 8-plate model against finite-element ray tracing. Self-shadowing options (\texttt{geom.selfShadow}): \ngfoptgeomselfShadow\ --- none, an equiangular CUBE-MAP look-up over incidence directions with $O(1)$ bilinear retrieval (\texttt{geom.lutResolution} cells per cube-face edge, convergence-tested), or an exact ray cast. The table records OCCLUSION ONLY: whether a plate faces the source is the kernels' business, and folding the two together is what deleted the edge-on shear before 0.10.0 (\S\ref{sec:flow}). The S-band boom is a cylinder (projection area \SI{0.0461901}{\metre\squared}, \SI{0.27}{\percent}); it is a separate element kind, not a plate.

\subsection{Occlusion is not facing, and a gas is not a beam}\label{sec:flow}
A look-up table that returns $f=0$ for a plate the source cannot see is a rule about a PARALLEL BEAM.
Applied to the ambient gas it is wrong twice, and both errors shipped until 0.10.0 (review findings
F1 and F2). First, a thermal gas arrives from every direction of the front hemisphere with the
drifting-Maxwellian weight, so an edge-on plate ($\gamma=0$) still receives flux and still carries
Sentman's shear $C_\tau\to\ell/(s\sqrt\pi)$ --- setting it to zero deleted that shear on
$\sim\SI{15}{\metre\squared}$ of GRACE-FO plate. Second, the same column was being used for the
Earth's radiation, which zeroed the heating of every panel facing away from the \emph{Sun}, and the
nadir panel's steady-state temperature read \SI{0}{\kelvin} at 11 of 120 real epochs.

The separation is now explicit: \texttt{geom.selfShadow} answers only ``is anything in the way'', and
whether a plate faces a source is the kernels' business. For the gas, \texttt{geom.flowShadow}
computes per plate, epoch and species three ratios of occlusion-weighted to unweighted directional
integrals of the incident flux over the front hemisphere --- number flux (for the re-emitted
momentum), normal momentum and tangential momentum --- with closed-form speed moments and a
flow-aligned product quadrature over directions. The limits that pin it: a plate with nothing in
front of it returns exactly 1, bitwise; and $s\to\infty$ reduces to the parallel-beam occlusion along
$\vv v_r$, which is what \texttt{constantCd} uses. Not modelled, and bounded by the oracle:
re-emitted molecules striking a second plate. Measured on 120 real epochs, the correction moves the
drag by \SI{1.68e-9}{\metre\per\second\squared} RMS, almost all along track, and the thermal recoil
by \SI{1.77e-9}{\metre\per\second\squared} RMS, radial.

\subsection{Implementation notes}
Vectorise over $K$ epochs: $\gamma$ is a $N\times K$ array, $\operatorname{erf}$ and $\exp$ are evaluated once per array; the density backend accepts $K$ points and returns $\rho$, species number densities and $T_\infty$; the MSIS inner loop is the one place where hand-written C (MEX) is sanctioned for post-processing. NaN density is an error (\texttt{numericNaN}), never a zero force.

\section{Solar radiation pressure}
\subsection{Irradiance and pressure}
Total solar irradiance $S_0=\SI{1361}{W/m^2}$ at \SI{1}{au} (IAU 2015 Resolution B3; SORCE/TIM; TSIS-1: $1361.5\pm5.9$), $P_0=S_0/c=\SI{4.5398e-6}{Pa}$. The older \SI{1367}{W/m^2} (Behzadpour et al.\ 2021, Table~1) is available as \texttt{Options.srp.tsi} for legacy comparisons and is never a constant (v1 used both values in two routines, \SI{0.44}{\percent} inconsistency; v0.1.1 repeated the class of error between its oracle and its constant). Scale with $(\mathrm{au}/|\vv r_s-\vv r|)^2$; the Sun position comes from \texttt{referenceSystem.geocentricBodyState} (analytic SOFA tier, RMS \SI{8.3e3}{\metre}; DE440 tier through MICE), never from a second implementation.

\subsection{Per-panel optical model}
For a panel with $\gamma=\hat{\vv n}_k\!\cdot\!\hat{\vv s}>0$,
\begin{equation}
\vv F_k=-\nu\,P\,A_k\,\gamma\Big[(\alpha+\rho_d)\,\hat{\vv s}+\Big(2\rho_s\gamma+\tfrac23\rho_d\Big)\hat{\vv n}_k\Big],\qquad \alpha+\rho_s+\rho_d=1,
\label{eq:srp}
\end{equation}
with $\nu\in[0,1]$ the shadow function and $2/3$ the Lambertian recoil factor. The same kernel serves Earth albedo (visible triplet) and Earth infrared (infrared triplet: $\varepsilon,\rho_s,\rho_d$) through the \emph{source list} $\{\hat{\vv d}_j,P_j,\text{band}_j\}$ (principle P3): SRP is one source, point-source Earth radiation two, a tiled Earth $2M$. Thermal re-emission of the absorbed part is \emph{not} folded into eq.~\eqref{eq:srp} by default (\S\ref{sec:trp}).

\subsection{Shadow function}
Apparent radii $a=\arcsin(R_s/|\vv r_s-\vv r|)$, $b=\arcsin(R_\oplus/|\vv r|)$, angular separation $c$ between $-\hat{\vv r}$ and $\widehat{\vv r_s-\vv r}$. Cylindrical model: $\nu=0$ iff $\vv r\!\cdot\!\hat{\vv s}<0$ and $|\vv r-(\vv r\!\cdot\!\hat{\vv s})\hat{\vv s}|<R_\oplus$. Conical model (Montenbruck \& Gill 2000, \S3.4.2): full light for $c\ge a+b$, umbra for $c\le b-a$, annular ($\nu=1-b^2/a^2$) for $c\le a-b$, else
\begin{equation}
x=\frac{c^2+a^2-b^2}{2c},\quad y=\sqrt{a^2-x^2},\quad
A=a^2\arccos\frac xa+b^2\arccos\frac{c-x}{b}-cy,\quad \nu=1-\frac{A}{\pi a^2}.
\end{equation}
The conical function is continuous but only $C^0$ at the umbra edge; for orbit integration a $C^1$ transition is required or the integrator collapses its step. SOLAARS-CF (Robertson 2015; Zhang et al.\ 2019) adds Earth oblateness and atmospheric extinction/refraction as a tabulated attenuation curve; Siemes et al.\ (2023) quantify the model-version effect at \SI{6}{}--\SI{9}{\percent} of the radiation-pressure acceleration. Lunar eclipses use the same geometry with $R_{moon}$ and are an opt-in. Tests: $\nu\in[0,1]$, monotone across the penumbra, symmetric in the geometry, and the worked example of Montenbruck \& Gill; a numerical disc integration in Python is the independent oracle.

\section{Earth radiation pressure}
\subsection{Knocke--Ries--Tapley (1988)}
Albedo and emissivity as zonal series with an annual term, $t_0$ = 22 December:
\begin{equation}
a(\phi,t)=a_0+a_1\cos\omega(t-t_0)\,P_1(\sin\phi)+a_2P_2(\sin\phi),\quad
e(\phi,t)=e_0+e_1\cos\omega(t-t_0)\,P_1(\sin\phi)+e_2P_2(\sin\phi),
\end{equation}
$a_0=0.34,a_1=0.10,a_2=0.29$; $e_0=0.68,e_1=-0.07,e_2=-0.18$; $\omega=2\pi/\SI{365.25}{d}$. Each visible surface element $dA$ at distance $\varrho$ with outgoing angle $\vartheta$ is a Lambertian source: albedo flux $dP=\dfrac{a\,S_0\cos\theta_s\,\cos\vartheta\,dA}{\pi c\,\varrho^2}$ (only for $\cos\theta_s>0$, gated by the shadow function of the element), infrared flux $dP=\dfrac{e\,S_0\cos\vartheta\,dA}{4\pi c\,\varrho^2}$. The visible cap is tiled ($12\times24$ in \texttt{standard}, $24\times48$ in \texttt{precise}); a few hundred tiles suffice at LEO (convergence test).
\subsection{CERES}
Observed top-of-atmosphere fluxes (EBAF climatology, SYN1deg hourly) replace the zonal model at the nm/s$^2$ level (Vielberg \& Kusche 2020; Wang et al.\ 2023); anisotropic reflection through the CERES angular distribution models (Bhattarai et al.\ 2019) is \emph{not} modelled: Lambertian reflection is the model, stated in the unit's ACCURACY block. An \texttt{erp.anisotropy} option was accepted and silently ignored until 0.10.0 and has been removed (review finding F4). The backends are \ngfopterpbackend. The netCDF variable names must be read from a real file (open item O-3).

\section{Thermal re-radiation}\label{sec:trp}
A panel at temperature $T_k$ with emissivity $\varepsilon_k$ emits $\varepsilon_k\sigma T_k^4A_k$ and recoils with
$\vv F_k=-\tfrac23\,\varepsilon_k\sigma T_k^4A_k\,\hat{\vv n}_k/c$. In steady state $T_k$ follows from the balance $\alpha_kQ_{sun,k}+\alpha_kQ_{alb,k}+\varepsilon_kQ_{ir,k}+q_{int,k}=\varepsilon_k\sigma T_k^4$; the transient model adds $C_k\dot T_k$ with a per-panel heat capacity (Wöske 2019 thesis, Ch.~4; Rievers \& Lämmerzahl 2011). For GRACE, TRP is the second-largest radiation term after SRP (Wöske et al.\ 2019; Vielberg \& Kusche 2020; Siemes et al.\ 2023: \SI{6.4}{nm/s^2} cross-track RMS), so ``folding it into SRP'' (the $2/3$ factor on absorbed flux) is a documented, deliberately crude option, not the default. Metamorphic test: in steady state the total absorbed power equals the total emitted power.

\section{Other effects and artefacts}
\textbf{Thruster forces}: cold-gas attitude thrusters (THR1B) produce linear accelerations up to \SI{e-6}{\metre\per\second\squared} that survive the \SI{35}{mHz} CRN filter (Meyer et al.\ 2012; Bandikova et al.\ 2019 model them as impulse responses per thruster pair). The toolbox models them as impulses with the CRN response and uses THR1B epochs as \emph{priors} for the impulse detector (\S\ref{sec:qc}). \textbf{Antenna thrust} (K-band, GPS): $\lesssim\SI{e-11}{\metre\per\second\squared}$, hook only. \textbf{Lorentz force}: negligible at GRACE-FO, hook only. \textbf{Artefacts the model does not explain}: twangs (up to \SI{2e-5}{\metre\per\second\squared} radial; Flury et al.\ 2008; Peterseim 2014), heater switching, magnetorquer spikes (Peterseim et al.\ 2012). These are removed by quality control, never modelled.

\section{Torques and the centre-of-mass coupling}
Surface torque $\vv\tau=\sum_k(\vv c_k-\vv r_{cm})\times\vv F_k$ from the same per-panel array (v1's separate torque implementation lost the vis/IR split: \SI{8}{\percent}). Gravity gradient $\vv\tau_{gg}=\frac{3\mu}{r^3}\hat{\vv r}\times(\mathbf I\hat{\vv r})$ with $\hat{\vv r}$ in the body frame. Magnetic $\vv\tau_{m}=\vv m\times\vv B$ with residual and commanded dipoles (MAG1B) and $\vv B$ from IGRF-14 (IAGA 2024; degree 13 main field, degree 8 secular variation; oracle: the igrf14 test cases) evaluated in ITRS and rotated to the body frame.

The accelerometer proof mass sits at $\Delta\vv r$ from the centre of mass ($\le\SI{100}{\micro\metre}$ per axis by design, \SI{\sim50}{\micro\metre} achieved; GRACE-C CoM calibration agrees with JPL to \SI{0.5}{\micro\metre}, Hu et al.\ 2022) and senses
\begin{equation}
\vv a_{sensed}=\vv a_{ng}+\dot{\vv\omega}\times\Delta\vv r+\vv\omega\times(\vv\omega\times\Delta\vv r),
\end{equation}
with $\vv\omega,\dot{\vv\omega}$ from SCA1B through \texttt{referenceSystem.angularVelocityFromRotationRate} after smoothing (Klinger \& Mayer-Gürr 2016; Wöske 2019). This is a property of the \emph{instrument}, not of the physics, and lives in the sensor layer (\S\ref{sec:sensors}).

\section{GRACE-FO data and the measurement model}\label{sec:sensors}
\subsection{Products and frames}
ACC1A/ACT1A (\SI{10}{Hz}) $\to$ ACC1B/ACT1B (\SI{1}{Hz}) via time-tag correction and the CRN low-pass filter (\SI{\sim35}{mHz}, Wu et al.\ 2006); ACT (McCullough et al.\ 2019) substitutes thruster spikes and, for GRACE-D from 2018-06-21, transplants GRACE-C data (Bandikova et al.\ 2019); ACH (2022) hybridises. SCA1B quaternions are ICRF$\to$SRF; the toolbox stores BODY$\to$GCRS and conjugates exactly once in the attitude provider, verified by physics (nadir and line-of-sight), never by a flag. GNV1B is Earth-fixed (ITRS) and is rotated with \texttt{rotationItrsToGcrs} (v1 had no such path and treated GPS tags as UT1: \SI{18}{\second} of Earth rotation). MAS1B holds cold-gas tank masses (\SI{\sim30}{kg}), not spacecraft mass; the handbook constant is the default.

\subsection{Calibration model}
$\vv a_{cal}=\vv b(t)+\mathbf S\,\vv a_{meas}$ with a bias vector (cubic B-spline, 6-hour knots, Klinger \& Mayer-Gürr 2016) and a full $3\times3$ scale/coupling matrix; typical scales $0.83$--$1.05$ (GRACE-C along-track $1.05\pm\num{2e-4}$, Siemes et al.\ 2023), biases $\sim\SI{e-6}{\metre\per\second\squared}$ with thermal excursions of hundreds of nm/s$^2$. Observability: along-track bias/scale well determined (\SI{0.5}{nm/s^2}, 1--3\,\%), cross-track poorly at solar minimum (\SI{22}{nm/s^2}, \SI{\sim15}{\percent}; van Helleputte 2009). The estimation is delegated to \texttt{lsqvce} (ARMA whitening from the sensor noise model, iterative snooping, variance components); \texttt{NonGravForces.app.calibrate} builds the observation vector and the design blocks, including Heaviside columns for detected jumps.

\subsection{The sensor abstraction}
The physics core returns the true dynamics record $\{\vv a_{ng},\vv\tau,\vv\omega,\dot{\vv\omega},t\}$; a \texttt{Sensor} maps it to an observable and provides the partials with respect to its own parameters and its one-sided noise ASD. \textbf{Electrostatic}: bias, scale matrix, CoM coupling, CRN filter; ASD $10^{-10}\sqrt{1+(f/0.5)^4+0.005/f}\ \si{m\,s^{-2}\,Hz^{-1/2}}$ (Kornfeld et al.\ 2019). \textbf{Atom interferometer} (CARIOQA pathfinder, MAGIC/NGGM hybrids): phase $\Delta\phi=k_{eff}\,a\,T^2$ with $k_{eff}=4\pi/\lambda$, interrogation time $T$ (\SI{1}{\second}), cycle time (\SI{4}{\second}, dead time), phase ambiguity modulo $2\pi$ requiring hybridisation with an electrostatic channel (Zahzam et al.\ 2022; Trimeche et al.\ 2019), and a rotation input because nadir pointing implies mrad/s rates whose Coriolis phase must be compensated (Lévèque et al.\ 2022); projection-noise floor \SI{e-10}{m\,s^{-2}\,Hz^{-1/2}}, target \num{e-12}. \textbf{Drag-free}: residual acceleration after control (LISA-Pathfinder heritage; NGGM \SI{e-11}{m\,s^{-2}\,Hz^{-1/2}} class). Optomechanical or MEMS instruments are further subclasses; calibration and quality control are written against the abstract interface, so no core code changes.

\section{Partial derivatives}
For orbit determination and calibration the toolbox returns $\partial\vv a/\partial\vv r$ (through $\partial\rho/\partial h$ and the shadow function), $\partial\vv a/\partial\vv v$ (drag, exact from eq.~\eqref{eq:drag} since $C_p,C_\tau$ depend on $\vv v_r$ through $s$ and $\gamma$), and parameter partials $\partial\vv a/\partial C_D=\vv a_{drag}/C_D$, $\partial\vv a/\partial C_R$, $\partial\vv a/\partial\alpha_E$. Analytic where clean; \texttt{partials.method} admits \ngfoptpartialsmethod\ and nothing else. \textbf{Complex-step differentiation is impossible in this chain and was removed}: MATLAB's \texttt{erf} refuses a complex argument and \texttt{abs}/\texttt{vecnorm} discard the imaginary part, so no complex step survives Sentman's coefficients or any norm on the way. The verifier is Richardson-extrapolated central differences instead, and the accuracy claim is stated against that. One evaluation serves both consumers (\texttt{Model.partials}).

\section{Quality control: outliers, impulses, jumps}\label{sec:qc}
Three event classes with different signatures need three detectors, all operating on whitened residuals $r(t)=y_{meas}-(b+\mathbf S\,y_{model})$ from a preliminary calibration:
\begin{itemize}[nosep]
\item \textbf{Outliers} (isolated samples): Hampel filter, moving median and $1.4826\cdot$MAD over a \SI{61}{\second} window, threshold $4\sigma$; robust to 50\,\% contamination, so thruster spikes and twangs cannot inflate the scale (a moving mean/std would). Implemented and verified against an independent numpy oracle (\texttt{qc.hampel}).
\item \textbf{Impulses} (CRN-shaped pulses, \SI{\sim30}{\second} wide, simultaneous in three axes): matched filter with the CRN impulse response; THR1B epochs are assigned first, unassigned pulses are twang/impact candidates; an impact (not a twang) also changes $\vv\omega$ in SCA1B.
\item \textbf{Jumps} (steps in level or slope): PELT change-point detection (Killick et al.\ 2012) with BIC penalty on piecewise-constant and piecewise-linear models; eclipse entries/exits (from the shadow function) and mass-trim events (MAS1B) are priors. Jumps are \emph{estimated} as Heaviside columns in the calibration, never cut.
\end{itemize}
Output: a canonical event table (\texttt{qc.eventTable}: $t_{start},t_{end}$, class, axis, amplitude, $\sigma$, source, confidence) and a logical mask. Limits: impulses shorter than the CRN width are unresolvable in L1B (use ACC1A); jumps below \SI{\sim1}{nm/s^2} are indistinguishable from $1/f$ bias drift without temperature data; along-track thresholds must adapt to geomagnetic activity because the residual carries the 10--30\,\% density error. Metamorphic tests: time shift, affine scaling, time reversal, planted events recovered with amplitude to $3\sigma$.

\section{Implementation}
\subsection{Architecture}
Nine layers (L0 core, L1 io, L2 env, L3 geom, L4 kernel, L5 \texttt{Model}, L6 sensor, L7 qc, L8 app), \ngfunitcount{} units, one catalogue (\texttt{doc/functionlist.csv}) that generates the function reference and drives the stub and audit tests. A layer calls only lower layers plus the four dependencies. \texttt{Model} is the single point of access (\texttt{accel | torque | series | partials | budget}; \texttt{rhs} was removed in LP-C---an integrator binds \texttt{app.DynamicsSource} instead); \texttt{Options} is a value class with presets \texttt{fast/standard/precise} and validated dotted overrides; extended models are opt-ins that no preset enables, so preset results never move when a model lands.

\subsection{Workflow}
\begin{lstlisting}
run('.../referenceSystem/setupReferenceSystem.m'); run('.../gracefoReading/setupGraceFOReading.m')
setupNonGravForces; runNgfTests
[ep,r,v]  = NonGravForces.io.readGNV1B(gnv);      % ITRS -> GCRS via referenceSystem
[~,q]     = NonGravForces.io.readSCA1B(sca);      % ICRF->SRF conjugated once
[~,y]     = NonGravForces.io.readACC1B(act);
macro = NonGravForces.geom.graceFO("C");
opt   = NonGravForces.Options("standard","drag.wind","hwm14");
M     = NonGravForces.Model(macro, opt, 'Attitude', att, 'Mass', mass);
T     = M.series(ep, r, v, q);                    % a_ng, tau, provenance
sens  = NonGravForces.sensor.Electrostatic(...);
[mask, events] = NonGravForces.qc.screen(y, sens.observe(T), thr, eclipse, mass);
R = NonGravForces.app.calibrate(y, T, sens, 'Events', events, 'Mask', mask);
\end{lstlisting}

\subsection{Error handling}
Every error passes through \texttt{NonGravForces.fail(unit, mnemonic, message)}; the identifier is \texttt{NonGravForces:<unit>:<mnemonic>}, the mnemonic's first token is one of \texttt{input, dep, data, numeric, unsupported, notImplemented}, and the message ends with a correction hint after ``\texttt{ -> }''; \texttt{fail} refuses anything else. Contracts are \texttt{arguments} blocks (shape, class, range) plus semantic checks at the top of the body; nothing is validated twice (the macro model validates once, kernels trust its packed arrays). No fidelity fallback exists; unknown option fields error; absent optional dependencies error at the point of use naming their setup function; dependency errors propagate unwrapped. Every result carries \texttt{info.provenance} (toolbox and dependency versions, option hash, files, warnings, UTC time). NaN policy is stated per unit: a NaN density is an error because it would silently zero the drag; a NaN sample in quality control is never flagged and never poisons a neighbour. Numerical guards report their diagnostic (condition number, residual, epoch index). Tests assert identifiers, not messages, and pair every refusal with the case that must succeed.

\subsection{Testing}
Seven categories per public unit, counted mechanically: contract, reference (an authority outside the project: published tables, independent Python, real data), precision (a bound with the reason for its size in a comment), vectorization ($K$ at once equals one at a time, and chunked/parallel series are bit-identical), robustness (NaN, gaps, empty, flags), metamorphic (rotation covariance, affine/shift/reversal invariance, mass and density scaling, $\nu\in[0,1]$ monotone), speed (ratios against a built-in, absolute times recorded). A passing bound leaves its number in \texttt{tests/last\_run.txt}; the gate is failures + incomplete + audit shortfall + written log, not the pass count. This session: 74 tests, GREEN in R2026a Update 4; two failures were authoring errors found by decomposition (a typed oracle value; a 0/1-based index), one was a guessed bound corrected by measurement (interpreted loop 382$\times$ slower than \texttt{movmedian}; algorithmically equivalent).

\subsection{Performance}
$3\times K$ column layout; environment evaluated in bulk; \texttt{pagemtimes} for attitude; the panel loop over $N\approx8$ with $K$-vectorised arithmetic; self-shadow LUT instead of per-epoch ray casting; \texttt{series} chunks $K$ and optionally runs the chunks under \texttt{parfor} (Parallel Computing Toolbox, feature-detected) with bit-identical results either way; the single-epoch C/MEX kernel exists only for the integrator right-hand side, where MATLAB call overhead (5--15\,\si{\micro\second}) dominates, and links \texttt{referenceSystem}'s \texttt{rsApiRotation}; no GPU (eight panels give a GPU nothing to do).

\section{Audit of the predecessors (summary; full table in HANDOVER Part C)}
\textbf{v1} (\texttt{SatelliteNonGravToolbox}): Sentman shear $\times8$; air gas constant; spherical Earth (\SI{-21.4}{\kilo\metre}, density $\times1.40$); two solar constants; GPS as UT1, no ITRS$\to$GCRS; adjusted same-day $F_{10.7}$; silent density fallback; drag routine computes its own Sun; five optical implementations, torque without vis/IR split; single-epoch interfaces; hand-rolled tests; duplicated frame helpers. \textbf{v0.1.1} (\texttt{NonGravToolbox}): the physics specification is correct; the flaws are architectural: gravity and equations of motion inside a non-gravitational toolbox, a compatibility layer to a defective v1, a Sun ``selector'' that is a wrapper, a 4171-line handover mixing specification and diaries, a setup that does not add the tests, four instead of seven test categories, a 1367 oracle against a 1361 constant, no sensor layer, no quality control, error handling by convention only. Concept 3 keeps v0.1.1's dependency discipline, principles, help template and $\circ$/$\boxtimes$ rule, removes the scope creep, and adds the sensor and QC layers.

\section{Recommended defaults and test oracles}
\begin{table}[h]\centering\footnotesize
\begin{tabular}{@{}lp{6.6cm}p{5cm}@{}}\toprule
Quantity & Default & Authority\\\midrule
density & NRLMSISE-00 (\texttt{standard}), NRLMSIS 2.1 (\texttt{precise}) & Picone 2002; Emmert 2021\\
wind & co-rotation; HWM14 opt-in & Drob 2015\\
GSI & Sentman, fitted $\alpha_E=0.85$ in every preset & Sentman 1961; March 2021\\
$S_0$ & \SI{1361}{W/m^2}; 1367 override & IAU 2015 B3; SORCE/TIM\\
Lambert factor & $2/3$ & standard\\
shadow & conical; SOLAARS-CF in \texttt{precise} & Montenbruck \& Gill 2000; Robertson 2015\\
Earth radiation & Knocke $12\times24$; CERES $24\times48$ & Knocke 1988; Wielicki 1996\\
TRP & steady-state per panel; transient opt-in & Wöske 2019\\
magnetic field & IGRF-14 & IAGA 2024\\
ephemeris & SOFA analytic; DE440 in \texttt{precise} & \texttt{referenceSystem}\\
ellipsoid & WGS-84, passed explicitly & \\
GRACE-FO mass, geometry & \SI{601.214}{kg}; JPL D-56935 Tables 4--5 & Wen 2019\\\bottomrule
\end{tabular}
\end{table}
\begin{table}[h]\centering\footnotesize
\begin{tabular}{@{}p{9.6cm}p{6cm}@{}}\toprule
Oracle & Used for\\\midrule
CODATA 2018, IAU 2015 B3 & \texttt{const} (verified)\\
JPL D-56935 Tables 4--5 & \texttt{geom.graceFO} (verified: areas, optics, inertia, normals)\\
numpy Hampel (\texttt{tools/pyxcheck}) & \texttt{qc.hampel} (verified, bit-for-bit flags)\\
Brodowski NRLMSISE-00 DOCUMENTATION table (17 cases) & density backend\\
NRL/CCMC NRLMSIS 2.x and HWM14 test outputs; CNES DTM2020; SET JB2008 & extended density and wind\\
igrf14 test cases; Zenodo 10.5281/zenodo.14012303 & magnetic field\\
Montenbruck \& Gill \S3.4.2 example; Python disc integration & shadow function\\
Knocke 1988 Table 1; Orekit \texttt{KnockeRediffusedForceModel} & Earth radiation\\
Python kinetic-theory integration; v1 $\times8$ regression & Sentman coefficients\\
Pilinski 2013 Fig.~3 & SESAM\\
Siemes 2023 Table 3; Vielberg \& Kusche 2020 monthly means & radiation magnitudes per axis\\
Hu 2022 CoM calibration; Kornfeld 2019 ASD & electrostatic sensor\\
Lévèque 2022 CARIOQA parameters & atom-interferometer sensor\\
real ACT1B/GNV1B/SCA1B/THR1B/MAS1B 2020-01-01 GRACE-C & readers, end-to-end calibration (scale $\approx1.05$)\\\bottomrule
\end{tabular}
\end{table}

%=======================================================================================
\chapter{Practical implementation}
\label{ch:practical}

\section{\texttt{Options}: the fidelity ladder}
An \texttt{Options} object is a validated, immutable option set. Three presets fix the default tier
of every model family:

\begin{table}[h]\centering\footnotesize
\begin{tabular}{@{}llll@{}}\toprule
 & \texttt{fast} & \texttt{standard} & \texttt{precise}\\\midrule
\texttt{drag.density} & exponential & NRLMSISE-00 & NRLMSIS 2.1\\
\texttt{drag.wind} & co-rotation & co-rotation & co-rotation (HWM14 opt-in)\\
\texttt{drag.gsi} & \texttt{constantCd} & Sentman & Sentman\\
\texttt{drag.accommodation} & fixed 0.85 & fixed 0.85 & fixed 0.85\\
\texttt{srp.eclipse} & conical & conical & SOLAARS-CF\\
\texttt{erp.backend} & none & Knocke & CERES\\
\texttt{erp.grid} & --- & $12\times24$ & $24\times48$\\
\texttt{thermal.model} & none & steady state & steady state (transient opt-in)\\
\texttt{geom.selfShadow} & none & LUT 16 & LUT 32\\
\texttt{sun.source} & analytic & analytic & DE440\\
\bottomrule
\end{tabular}
\end{table}

\begin{lstlisting}
opt = NonGravForces.Options("standard");
opt = NonGravForces.Options("precise", "drag.wind", "hwm14");   % an opt-in extension
opt = opt.set("srp.tsi", 1367);        % a COPY: Options is a value class
\end{lstlisting}

\noindent Two rules make this safe to rely on. \textbf{Every admissible value evaluates}: a value
the toolbox cannot compute is not offered, so an option that validates will not fail three layers
down. And \textbf{extended models are opt-ins that no preset enables} --- HWM14 winds, the transient
thermal model, body torques --- so preset results never move when a new model lands.

An option a model does not read is not silently ignored when it contradicts the choice:
\texttt{Model} refuses a non-empty \texttt{drag.alpha} beside a computed accommodation, a
\texttt{maxwell} gas without a specular fraction, and \texttt{walker2014} without a surface mass.

\section{\texttt{Model}: the single point of access}
\begin{lstlisting}
M = NonGravForces.Model(macro, opt, 'Attitude', att, 'Eop', eop, 'SpaceWeather', sw);

[a, info]      = M.accel(epoch, r, v, q);      % 3xK m/s^2, body frame by default
[tau, info]    = M.torque(epoch, r, v, q);     % 3xK N m
[a, tau, info] = M.evaluate(epoch, r, v, q);   % both from ONE environment pass
T              = M.series(epoch, r, v, q);     % timetable, chunked, optionally parfor
[J, names]     = M.partials(epoch, r, v, q, 'WithRespectTo', ["r","v","cdScale"]);
B              = M.budget(epoch, r, v, q);     % per-effect table
\end{lstlisting}

\noindent \texttt{accel} and \texttt{torque} are two reductions of \emph{the same} per-panel force
array: $\vv F=\sum_k\vv F_k$ and $\vv\tau=\sum_k(\vv c_k-\vv r_{cm})\times\vv F_k$. They cannot
disagree about the physics because they are not two implementations of it. If you want both, call
\texttt{evaluate} and pay for the environment once.

\texttt{info} carries the environment the evaluation actually saw --- density, species, gas
temperature, relative velocity, the shadow fractions, the panel temperatures --- and
\texttt{info.provenance}, which records the toolbox and dependency versions, the option hash, the
files read and the warnings raised. A result that cannot say where it came from is not a result.

\section{Three evaluation paths, all first class}
This is a design ruling, not an optimisation, and it is worth understanding before you choose one.

\begin{description}[leftmargin=1.4em,style=nextline]
\item[The vectorised MATLAB path.] \texttt{Model.accel} and \texttt{Model.series} over an arc. This
 is normative: everything else is verified against it, never the reverse.
\item[The compiled single-epoch kernel.] \texttt{kernel.singleEpoch}, a hand-written C port of the
 force kernel, for a pointwise right-hand side.
\item[The native chain.] \texttt{app.DynamicsSource} evaluates the \emph{whole} chain in C over
 tables built once per run, for an integrator's inner loop (Chapter~\ref{ch:integrator}).
\end{description}

\noindent Neither replaces the other and none is justified by speed. \texttt{Model} gains no
\texttt{'Backend'} option, because that would make one a mode of another. The price of this choice
is that they must agree, and that agreement is asserted: the native evaluator matches
\texttt{Model.accel} to $2.5\times10^{-14}$ relative at a table node, in seven configurations chosen
to reach every branch.

\section{Errors, warnings and what they mean}
Every error identifier is \texttt{NonGravForces:<unit>:<mnemonic>}, where the mnemonic's first token
is one of \texttt{input}, \texttt{dep}, \texttt{data}, \texttt{numeric}, \texttt{unsupported} or
\texttt{notImplemented}, and every message ends with the correction after ``\texttt{ -> }''. Errors
from a dependency propagate unwrapped: a \texttt{referenceSystem:} identifier means that toolbox
said it, and re-badging it here would hide where the fault is.

Only two units may \emph{warn}. \texttt{dependencyReport} reports what is installed.
\texttt{kernel.accommodation} warns when a computed accommodation model is evaluated outside the
range it was fitted in:

\begin{lstlisting}
Warning: accommodation walker2014: 600 of 600 epochs lie outside the atomic-oxygen partial
pressure of 1e-5 Pa and above, which is below about 300 km. The returned coefficient is an
extrapolation of a fitted model -> read info.insideFit, choose a fixed alpha (March et al. 2021
recommend 0.85), or silence this with warning('off', 'NonGravForces:accommodation:outsideFit').
\end{lstlisting}

\noindent \textbf{This warning will fire on GRACE-FO}, and that is its purpose --- see
\S\ref{sec:lim-accommodation}. Silence it with \texttt{warning('off', \ldots)} once you have decided
the extrapolation is what you want.

%=======================================================================================
\chapter{Driving an orbit integrator}
\label{ch:integrator}

\section{What this interface is for}
An integrator calls its right-hand side many times per step, at times it chooses, non-monotonically,
and often repeatedly at the same $t$. That is a different problem from evaluating an arc: the
per-call cost dominates, and anything that can be computed once per run should be.

\texttt{app.DynamicsSource} subclasses \texttt{gjToolbox.DynamicsSource}, that toolbox's frozen
dynamics interface. It builds time tables once at \texttt{prepare} --- Earth orientation, attitude
nodes, space-weather bins, the Sun and Moon, the CERES month --- and then evaluates the whole force
chain in C against them.

\section{Using it}
\begin{lstlisting}
M   = NonGravForces.Model(macro, NonGravForces.Options("standard"), ...
          'Attitude', att, 'Eop', eop, 'SpaceWeather', sw);

src = NonGravForces.app.DynamicsSource(M, 'Epoch', t0_tai, ...
          'Margin_s', 30, 'TableStep_s', 60);

src = src.prepare(struct('Span', [0 5400], 'StencilSpan', [0 5400], 'MaxColumns', 1));

a   = src.accel(t, r, v);        % t seconds from 'Epoch'; r, v in GCRS
\end{lstlisting}

\noindent \texttt{'Epoch'} is the TAI datetime that $t=0$ means. \texttt{'Margin\_s'} widens the
tables beyond the span you declare, because an integrator's stencil reaches outside the interval it
is integrating. \texttt{'TableStep\_s'} is the table spacing: \SI{60}{\second} costs
$8\times10^{-11}$ relative between nodes, and the error falls as the fourth power of the step.

\section{Native or vectorised}
\begin{lstlisting}
src = NonGravForces.app.DynamicsSource(M, 'Epoch', t0, 'Native', false);   % MATLAB binding
\end{lstlisting}

\noindent With \texttt{'Native'} true (the default when the \texttt{rhs} MEX is built) the chain runs
in C. With false it runs the MATLAB path over the same tables. The two agree to rounding, and the
MATLAB binding exists so that a user without a compiler can still integrate, and so that a parity
difference can be localised.

\texttt{src.diagnostics(t, r, v)} returns what the native evaluator computed on the way to one
acceleration --- the three effects separately, the density, the species, the gas temperature, the
relative velocity, the sub-satellite point, the eclipse factor, the mass and the attitude matrix. It
is how you find \emph{where} two implementations differ rather than only that they do.

\section{What it refuses, and why that is a feature}
\begin{itemize}[leftmargin=1.4em]
\item \texttt{thermal.model = "transient"} --- it carries temperatures from epoch to epoch, and an
 integrator calls the right-hand side non-monotonically and repeatedly at one $t$. There is no
 correct answer, so there is no answer.
\item \texttt{geom.selfShadow = "raycast"} --- a ray cast per direction is not an inner loop.
\item \texttt{erp.backend = "ceresClimatology"} --- not for want of memory: the native path already
 carries three EBAF grids of that size for \texttt{erp.backend = "ceres"}. It is refused because it
 would buy no physics. EBAF is also $180\times360$ at one degree and also a measured exitance, so
 over an integration span the two differ in source and reference level rather than in what they can
 express; what the climatology uniquely offers is an epoch range of 1900--2101 and provenance shared
 with GROOPS, and both are post-processing concerns. Use \texttt{ceres} for an integrated arc, or
 \texttt{knocke} if you do not hold the months.
\item An attitude provider built with \texttt{"squad"} --- the native evaluator interpolates by
 slerp.
\end{itemize}

\noindent Each of these \emph{refuses at construction}, immediately, rather than degrading. The
alternative was worse than it looks: the option name is mapped to a C enum by position, and a name
the C does not know would map to zero --- which is \texttt{"none"}. A model that computes Earth
radiation would have quietly become one that does not, on the native path only.
%=======================================================================================
\chapter{A GRACE-type processing chain, step by step}
\label{ch:walkthrough}

This chapter builds a processing chain from nothing, using only the fixtures committed to the
repository, so you can run it as you read. \textbf{Every number printed below was produced by
running this code}, not composed for the document.

The excerpt is 600 seconds of GRACE-C on 2020-01-01. That is deliberately short: it is enough to
show the chain and, as \S\ref{sec:cal-short} explains, too short to determine a scale factor ---
which is itself worth seeing.

\section{Step 1: the auxiliary data}
\begin{lstlisting}
root = 'C:\path\to\NonGravForces';
f = @(varargin) fullfile(root, 'data', varargin{:});

eop = referenceSystem.readEop(f('eop', 'finals2000A_2019-09_2020-03.all'));
sw  = NonGravForces.io.spaceWeather(f('spaceweather', 'SW-All_2019-09_2020-03.csv'));
\end{lstlisting}

\noindent Earth orientation comes from \texttt{referenceSystem}, not from here. Space weather is the
CelesTrak file; \texttt{io.spaceWeather} returns a struct carrying a \texttt{query} function, and
the \texttt{msis} density backends take it directly.

\section{Step 2: orbit and attitude}
\begin{lstlisting}
[ep, r, v] = NonGravForces.io.readGNV1B(f('samples','GNV1B_2020-01-01_C_04.txt'), 'Eop', eop);
[eq, q]    = NonGravForces.io.readSCA1B(f('samples','SCA1B_2020-01-01_C_04.txt'));
att        = NonGravForces.geom.AttitudeProvider(eq, q);
\end{lstlisting}

\noindent \textbf{Two traps live in this step, and a reader should know they are already handled.}
GNV1B is Earth-fixed (ITRS), not inertial --- \texttt{readGNV1B} rotates it through
\texttt{referenceSystem.rotationItrsToGcrs}, which also produces the velocity correctly, and a GPS
time tag is not UT1. And SCA1B is ICRF$\to$SRF, so it must be conjugated exactly \emph{once} to give
BODY$\to$GCRS. Doing it twice, or not at all, produces an attitude that looks plausible and is
wrong; \texttt{geom.checkAttitudeSense} exists to catch precisely that.

\section{Step 3: the spacecraft}
\begin{lstlisting}
macro = NonGravForces.geom.graceFO("C");
\end{lstlisting}

\noindent \texttt{geom.graceFO} builds the handbook panel model. For the classic GRACE-A/B satellites,
\texttt{geom.grace("A")} builds the model of the GRACE Product Specification (Bettadpur 2012,
Sec.~I.5.3), whose surface table is the same as GRACE-FO's but whose dimensions, origin and
\SI{487.2}{\kilo\gram} launch mass are its own. Use \texttt{geom.boxWing} for a
generic satellite, or \texttt{geom.MacroModel} directly if you have your own panels.

\textbf{Do not take the mass from MAS1B.} That product is cold-gas \emph{tank} mass, about
\SI{30}{\kilo\gram} in total, not spacecraft mass; using it gives accelerations 20 to 40 times too
large. \texttt{io.massModel} defaults to the handbook constant and asserts that the result exceeds
\SI{500}{\kilo\gram}. \emph{Classic GRACE is the opposite case:} its MAS1B carries the spacecraft
mass (\SI{476.85}{\kilo\gram} for GRACE-A on 2008-01-01), which \texttt{io.massModel} does not serve;
pass it to \texttt{Model} as \texttt{'Mass'} (HANDOVER O-46).

\section{Step 4: the options and the model}
\begin{lstlisting}
opt = NonGravForces.Options("standard");
M   = NonGravForces.Model(macro, opt, 'Attitude', att, 'Eop', eop, 'SpaceWeather', sw);
\end{lstlisting}

\section{Step 5: evaluate, and look at the environment it saw}
\begin{lstlisting}
idx = 1:200;
e   = referenceSystem.daySecondToEpoch([ep.day(idx) ep.sec(idx)], ep.scale);
[a, info] = M.accel(e, r(:, idx), v(:, idx));
\end{lstlisting}

\noindent \texttt{info} is not a diagnostic afterthought: it carries the density, the species, the
relative velocity, the shadow fractions, the panel temperatures and the accommodation the
evaluation actually used, plus the provenance of all of it.

\section{Step 6: which effect carries the arc}
\begin{lstlisting}
B = M.budget(e, r(:, idx), v(:, idx));
\end{lstlisting}

\section{Step 7: compare with the measured accelerations}
\begin{lstlisting}
[epa, aMeas] = NonGravForces.io.readACC1B(f('samples','ACT1B_2020-01-01_C_04.txt'));
[tf, loc]    = ismember(ep.sec(idx), epa.sec);
am           = aMeas(:, loc(tf));
\end{lstlisting}

\noindent Compare \emph{shape}, not level. The measured along-track is dominated by instrument bias
drift over an excerpt this short, so its correlation means nothing; cross-track and radial are the
informative axes.

\section{Step 8: calibrate}
\label{sec:cal-short}
\begin{lstlisting}
acc = timeSeriesPreparation.Arc(referenceSystem.epochTo(epa, "gpsNanoseconds"), aMeas.', auxACT1B);
res = NonGravForces.app.calibrate(acc, M, 'Orbit', struct('epoch', ep, 'r', r, 'v', v), ...
                                  'Attitude', att, 'Fit', {'FixCrossTrack', false});
res.calibration.summaryTable()
\end{lstlisting}

\noindent Since 0.13.0 (N-02) \texttt{app.calibrate} takes the observation as a
\texttt{timeSeriesPreparation.Arc} (here built from the reader's triple with the ACT1B column names,
\texttt{auxACT1B}) and the reference as a \texttt{Model}. The Model is evaluated along the orbit per
effect and interpolated onto the 10\,Hz grid of the CRN filter; the instrument stages (mask, CRN,
align, QC) run through \texttt{timeSeriesPreparation.run}; the estimate is \texttt{lsqvce.fit}
through \texttt{app.fitCalibration}, every deterministic column filtered by the same CRN kernel as
the observation. The result is the run's struct plus \texttt{calibration} (an
\texttt{app.Calibration}) and \texttt{attribution}. The matrix interface of 0.12.0
(\texttt{calibrate(yMeas, aModel, epoch)}) is gone.

\section{What this chain actually prints}
\begin{lstlisting}
--- Step 1: the auxiliary data -------------------------------------------
EOP rows 213, space weather family celestrak

--- Step 2: orbit and attitude -------------------------------------------
GNV1B epochs 600 (GPS), |r| = 6885.4 km, |v| = 7.5970 km/s
SCA1B quaternions 600, provider method slerp

--- Step 3: the spacecraft ----------------------------------------------
panels 8, area sum 16.9163 m^2, mass 601.214 kg
silhouette defect |sum A_k n_k| = 1.406e-01 m^2
   (zero only for a CLOSED body; the 8-panel GRACE-FO model is an open box-wing,
    so a non-zero value here is expected and is not a fault)

--- Step 4: the options and the model ------------------------------------
preset standard: density msis00, gsi sentman, accommodation fixed (alpha 0.85), erp knocke, thermal steadyState

--- Step 5: evaluate over an arc -----------------------------------------
epochs 200, |a| RMS 2.4555e-08 m/s^2
   drag      7.6358e-09   (31.1 % of the total RMS)
   radiation 2.3439e-08
   thermal   5.5046e-09
info fields: frame, K, rho, nu, mass_kg, vrel_body, sources, temperature_K, flowShadow, accommodation, lla, backends, drag, radiation, thermal, provenance
density 4.5107e-14 kg/m^3, |v_rel| 7594.2 m/s

--- Step 6: the per-effect budget ----------------------------------------
      effect        rms_norm
    ___________    __________

    "drag"         7.6358e-09
    "radiation"    2.3439e-08
    "thermal"      5.5046e-09
    "total"        2.4555e-08


--- Step 7: compare with the measured accelerations ----------------------
cross-track correlation with ACT1B  0.9722
radial      correlation with ACT1B  0.9954

--- Step 8: calibrate ----------------------------------------------------
(printed by the 0.12.0 estimator; not re-run for 0.13.0 -- see below)
\end{lstlisting}

\noindent \textbf{The lesson of this step survives the rewrite, the numbers did not.} With the
0.12.0 estimator, 200 seconds of this excerpt gave scale factors of 0.16, 0.70 and 0.85 at
condition numbers 57, 46 and 7: a scale factor is determined by the \emph{spread} of the modelled
signal across the arc, and minutes of a 90-minute orbit carry almost none. The 0.13.0 estimator
reports the same fact in its gates (\texttt{res.calibration.gates.cond}, \texttt{.vif}) and refuses
a full scale matrix the arc cannot separate; its numbers for this excerpt are written to
\texttt{tests/last\_run.txt} by \texttt{tNgfCalibrate/calibrateRealDay} on every gate run rather
than copied here, where nothing would compare them. \textbf{Use hours rather than minutes, and read
the gates before believing \texttt{scale}.}

Note also the silhouette defect of \SI{0.14}{\metre\squared} in step~3. That quantity is zero for a
\emph{closed} body; the 8-panel GRACE-FO model is an open box-wing, so a non-zero value there is
expected and is not a fault.
%=======================================================================================
\chapter{Worked examples}
\label{ch:examples}

Four runnable demonstrations ship in \texttt{demos/}. The test suite executes every one of them on
every run --- not for their numbers, which the unit tests already bound, but because \textbf{a demo
nobody executes rots}, exactly as one of this project's own documentation pages once did. What
follows is what they print.

\section{Which effect carries the arc, and what a modelling choice costs}
\begin{lstlisting}
demoForceBudget
\end{lstlisting}
\begin{lstlisting}
Per-effect budget over 10 epochs, body frame:

      model         effect        rms_norm      share
    __________    ___________    __________    _______

    "fast"        "drag"         1.0287e-07     1.1502
    "fast"        "radiation"    2.5663e-08    0.28695
    "fast"        "thermal"               0          0
    "fast"        "total"        8.9433e-08          1
    "standard"    "drag"         8.4956e-09    0.32603
    "standard"    "radiation"    2.5001e-08    0.95945
    "standard"    "thermal"      5.9707e-09    0.22914
    "standard"    "total"        2.6057e-08          1

The share column is formed from rms_norm, so it can sum to more than one when two
effects partly cancel in the total. That is a property of the arc, not an error,
which is why the total row is reported beside them.
\end{lstlisting}

\noindent The \texttt{fast} preset overstates the drag by an order of magnitude here, because its
exponential atmosphere is a stand-in rather than a model, and it carries no thermal term at all. The
point is not that \texttt{fast} is wrong: a fidelity tier is a statement about which physics you are
prepared to omit, and this is what omitting it costs on this arc.

Note the \texttt{share} column summing above one. That is a property of the arc rather than an
error --- the shares are formed from each effect's RMS norm, and two effects that partly cancel in
the total will exceed it. The total row is printed beside them so that the reader can see it.

\section{Recovering a planted density scale}
\begin{lstlisting}
demoDensityInversion
\end{lstlisting}
\begin{lstlisting}
planted scale        1.7000
recovered, worst     1.700000  (relative error 3.92e-16)
model density        8.8736e-13 kg/m^3
inverted density     1.5085e-12 kg/m^3
residual across drag 5.305e-23 m/s^2  (zero: the scaling was along the drag direction)

Against real data the perpendicular residual is the honest quality flag: it is what
the density cannot explain, and it grows when the attitude, the radiation model or
the calibration is wrong.
\end{lstlisting}

\noindent The inversion is a \emph{division}. The gas force is exactly linear in the species
densities, so a single model evaluation gives the whole sensitivity, and the planted scale comes
back to four parts in $10^{16}$ --- which is rounding, not skill. The interesting output is the
residual across the drag direction: it is what the density cannot explain, and it grows when the
attitude, the radiation model or the calibration is wrong. \texttt{app.densityInversion} does not
calibrate and says so.

\section{Recovering a planted bias and scale}
\label{sec:ex-cal}
\begin{lstlisting}
demoGraceFOCalibration
\end{lstlisting}
\noindent The demo prints, per axis, the planted bias and scale beside the recovered ones, the
scale's bootstrap and formal sigma, the EIV attenuation bound and the stages \texttt{run} executed.
Its truth is \texttt{app.syntheticPair}: a planted scale, a drifting bias, coloured noise, spikes
and a gap, passed through the real Level-1B CRN chain. This is the contrast promised in
\S\ref{sec:cal-short}: on an arc built to determine the parameters, the same estimator recovers
them, and the recovery is asserted to within five bootstrap sigma by
\texttt{tNgfCalibrate/calibratePrecision}.

The two sigmas are reported side by side on purpose. The CRN-filtered residual is correlated, so
the formal white-noise sigma is optimistic; the Calibration carries the moving-block bootstrap
covariance, and \texttt{cal.uncertainty.inflation} is their ratio.

\section{A cold-atom accelerometer}
\begin{lstlisting}
demoCAI
\end{lstlisting}
\begin{lstlisting}
readings per cycle   29 valid of 600 samples (dead time 1.0 s)
drag phase           1.2232 rad  = 3.0389e-07 m/s^2 equivalent
wrapped phase        1.2575 rad  (what the instrument returns)
fringe integer       0  (what the readout cannot tell you)

Coriolis phase       8.8550e+00 rad = 1.409e+00 fringes
as an acceleration   2.20e-06 m/s^2, against a drag of 3.0e-07

One millimetre per second of atom velocity outweighs the signal several times over. The
rotation input is not a refinement to this instrument; it has to be measured or compensated.
\end{lstlisting}

\noindent The headline is not the drag. \textbf{The Coriolis phase is seven times the signal}: one
millimetre per second of atom velocity relative to the instrument, at GRACE-FO's
\SI{1.1}{\milli\radian\per\second} nadir-pointing rate, produces an apparent
\SI{2.2e-6}{\metre\per\second\squared} against a drag of \SI{3e-7}{}. For this instrument the
rotation input is not a refinement --- it has to be measured or compensated.

Two further details are visible in the output. Only 29 of 600 samples carry a reading, because the
interrogation time is dead time and \emph{nothing is interpolated across it}; and the readout
returns a wrapped phase, so the fringe integer is reported separately as the thing the instrument
cannot tell you.
%=======================================================================================
\chapter{Limitations}
\label{ch:limits}

Everything in this chapter is a property of the toolbox as shipped, measured where a number is
given. A limitation a user discovers is a defect; a limitation written down is a specification.

\section{Energy accommodation: the presets use a fitted constant}
\label{sec:lim-accommodation}
All three presets set \texttt{drag.accommodation = "fixed"} with $\alpha_E=0.85$, the value March
et al.\ (2021) recommend after optimising a \emph{constant} per mission against attitude manoeuvres
and wind analyses. The adsorption models --- \texttt{pilinski2010}, \texttt{pilinski2011},
\texttt{walker2014}, \texttt{sesam} --- are available by name and are not defaults.

This is a ruling, not an approximation, and the reason is measured. At CHAMP's altitude and solar
activity (\SIrange{300}{350}{\kilo\metre}, $F_{10.7}=70$) \texttt{pilinski2010} returns
\numrange{0.86}{0.95}, straddling the \numrange{0.85}{0.93} that March et al.\ report for SESAM
there: \textbf{the adsorption family reproduces the published range where the published range was
measured}. The same code returns \numrange{0.11}{0.17} at GRACE-FO's \SI{520}{\kilo\metre} in deep
solar minimum, and \numrange{0.63}{0.83} at that altitude at $F_{10.7}=180$, because the
atomic-oxygen partial pressure that drives the Langmuir coverage falls by two orders of magnitude
between them.

So these models predict a factor-of-five swing in GRACE-FO's accommodation over the solar cycle, and
2020-01-01 --- the day this toolbox is validated on --- sits at its extreme low end. On that day the
choice is worth \SI{17.5}{\percent} of the drag:

\begin{table}[h]\centering\footnotesize
\begin{tabular}{@{}lll@{}}\toprule
\texttt{drag.accommodation} & $\alpha_E$ & drag RMS [\si{\metre\per\second\squared}]\\\midrule
\texttt{fixed} 0.85 (shipped) & 0.850 & \num{6.8826e-9}\\
\texttt{fixed} 0.93 (Doornbos 2011, superseded) & 0.930 & \num{6.6158e-9}\\
\texttt{walker2014} & \numrange{0.116}{0.164} & \num{8.0886e-9}\\
\texttt{pilinski2010} & \numrange{0.051}{0.098} & \num{8.1882e-9}\\
\texttt{pilinski2011} & \numrange{0.034}{0.067} & \num{8.2164e-9}\\
\bottomrule
\end{tabular}
\end{table}

\noindent For accelerometer calibration that is not a refinement: it is absorbed by the estimated
scale factor and propagates into any density inverted from it. If you select a computed model, the
toolbox warns once per call that it is extrapolating, and \texttt{info.insideFit} says which epochs.

\section{Gas--surface interaction: CLL is refused}
\texttt{drag.gsi} offers \texttt{sentman}, \texttt{maxwell} and \texttt{constantCd}. The
Cercignani--Lampis--Lord model is \emph{not} offered and is refused by name. Walker et al.\ (2014)
publish closed forms for \emph{whole bodies} --- sphere, cylinder, and a flat plate at normal
incidence only --- and defer arbitrary orientation to a source this project does not hold. This
kernel needs per-panel $C_p$ and $C_\tau$ at arbitrary incidence, which no equation in that paper
provides. Offering a CLL that silently used the normal-incidence form would be worse than refusing.

\section{Earth radiation: no diurnal cycle}
\texttt{erp.backend = "ceres"} uses monthly CERES EBAF fluxes; \texttt{"ceresClimatology"} uses the
1\textdegree\ GROOPS monthly climatology. \textbf{Neither carries a diurnal cycle}, whatever the
file is named --- the GROOPS file is called SYN1deg but is a monthly mean. The albedo is applied to
the \emph{instantaneous} insolation, which restores the day side; the longwave is a monthly mean and
is not corrected for local time. The hourly SYN1deg product would fix this and is not implemented.

\texttt{"ceresClimatology"} additionally has \textbf{no native evaluator}: it reads a 64\,800-cell
grid, which is not table material for an integrator, and \texttt{app.DynamicsSource} refuses it at
construction.

\section{Thermal re-radiation: the literature disagrees with itself}
The presets model thermal re-radiation instantaneously (\texttt{thermal.model = "steadyState"}).
Two peer-reviewed results disagree about whether that helps:

\begin{itemize}[leftmargin=1.4em]
\item Hładczuk et al.\ (2024), GRACE-FO: adding instantaneous re-radiation to a panel model takes
 the cross-track residual from \SI{4.22}{} to \SI{3.04}{\nano\metre\per\second\squared}, and thermal
 inertia takes it to \SI{1.70}{}. It helps.
\item Vielberg \& Kusche (2020), GRACE-A: instantaneous re-radiation fits ``more than \SI{39}{\percent}
 worse'' to calibrated cross-track accelerations, and ``did not lead to improvements''. It hurts.
\end{itemize}

\noindent The presets follow the first: it is the right satellite, four years later, and its
Section~3 repairs the cross-track calibration that the second names as a possible cause of its own
result --- it also names neglected horizontal winds, which this toolbox models and that study did
not. \textbf{It is an open question in the field.} Set \texttt{thermal.model = "none"} to take the
other position.

There is also no exact per-axis TRP oracle for GRACE-FO in the literature this project holds. What
is verified is the magnitude: \SI{5.263}{\nano\metre\per\second\squared} cross-track here, between
the \SI{4.22}{} that Hładczuk et al.\ leave when they omit the term and the \SI{6.4}{} Siemes et al.\
(2023) publish for GRACE-A.

\section{Cost}
\textbf{The cost per epoch depends on how many epochs you ask for in one call, far more than on which
preset you choose}, and that is the most useful thing to know about this toolbox's performance.
Measured by \texttt{tools/benchmark.m} on 2026-09-18 (v0.12.0, R2026a, median of five after five
warm-up calls); the table it writes is committed at \texttt{tools/benchmark\_reference.csv}, and the
audit checks every figure below against it.

\begin{table}[h]\centering\footnotesize
\begin{tabular}{@{}lrrr@{}}\toprule
 & \multicolumn{3}{c}{ms per epoch, for a call of $K$ epochs}\\\cmidrule(l){2-4}
path & $K=1$ & $K=10$ & $K=100$\\\midrule
\texttt{Model.accel}, \texttt{standard} & 15.9 & 2.99 & 0.798\\
\texttt{Model.accel}, \texttt{fast} & 4.00 & 0.695 & 0.076\\
\texttt{Model.budget}, \texttt{fast} & 4.30 & 0.762 & 0.084\\
\texttt{app.forceBudget}, \texttt{fast} & 5.50 & 0.826 & 0.094\\\midrule
\multicolumn{4}{@{}l}{\emph{pointwise by construction --- these take one state and do not batch:}}\\
\texttt{kernel.singleEpoch} (compiled) & \multicolumn{3}{r}{0.456}\\
\texttt{app.DynamicsSource.accel} (native) & \multicolumn{3}{r}{0.544}\\
\bottomrule
\end{tabular}
\end{table}

\noindent \texttt{Model.series} over the 600-epoch fixture measures \SI{0.886}{\milli\second} per
epoch, so \textbf{a day at \SI{1}{\hertz} costs about 1 minute 16 seconds}.

Every path in that table gets dramatically cheaper per epoch between $K=1$ and $K=100$ ---
$20\times$ for \texttt{standard}, and $50$ to $60\times$ for the \texttt{fast} rows, whose
per-epoch work is smaller and so amortises further --- and the reason is a large fixed cost per
\emph{call}. Fitting $\text{cost}(K)=a+bK$ to the
\texttt{standard} rows gives $a\approx\SI{15.2}{\milli\second}$ of setup and
$b\approx\SI{0.65}{\milli\second}$ of actual per-epoch work; for \texttt{fast},
$a\approx\SI{4.0}{\milli\second}$ and $b\approx\SI{0.037}{\milli\second}$. \textbf{Asking for one
epoch buys you almost entirely setup.} The presets separate by about $4\times$ at $K=1$ and
$10.5\times$ at $K=100$, and on the marginal cost --- the number that describes the physics rather
than the call --- by $17.6\times$, which is where self-shadowing, Earth radiation and the thermal
model actually show up.

\textbf{So: pass many epochs at once, or use \texttt{Model.series}.} For the one caller that cannot
batch by nature --- an integrator's right-hand side, which asks for one state at a time --- the
native evaluator of Chapter~\ref{ch:integrator} answers in \SI{0.544}{\milli\second}, which is
\num{29} times cheaper than \texttt{Model.accel} at $K=1$ and is the whole reason that path exists.

\begin{quote}\footnotesize
This section read \SI{88}{\milli\second} per epoch and \SI{2.1}{\hour} for a day until 2026-09-18. It
was written in good faith from an obligation record that had been superseded two days earlier and
never closed. It is left mentioned rather than silently replaced, because a reader who saw the old
figure elsewhere deserves to know which number this project stands behind.
\end{quote}

\section{The accelerometer filter}
Level-1B accelerations are the \SI{10}{Hz} ACT1A stream through the CRN low-pass of the L1B ATBD
(JPL D-27672, App.~B). Since N-04 (2026-10-14) \texttt{sensor.Electrostatic} applies that filter by
default, with the kernel taken from \texttt{timeSeriesPreparation.crnKernel("ACC")} and never built
here: 1407 taps on the \SI{10}{Hz} grid, unit gain at $f_0 = \SI{0.37}{mHz}$, DC gain
$1 + 1.1\times10^{-7}$. The default therefore needs \SI{10}{Hz} dynamics with \texttt{dyn.epoch}, which
it checks; \texttt{'OutputRate', 1} samples the result at \SI{1}{Hz} as ACT1B is formed. Pass
\texttt{[]} as the kernel for an unfiltered sensor at any rate, or any odd-length FIR of your own.
The first and last 703 samples (\SI{70.3}{s}) of an arc are marked invalid in \texttt{y.mask}.

\section{What is not modelled at all}
\begin{itemize}[leftmargin=1.4em]
\item \textbf{Gravity}, in any form. $T=\partial\vv g/\partial\vv r$ enters
 \texttt{sensor.Electrostatic}'s CoM coupling and must be supplied by the caller in
 \texttt{dyn.gravityGradient}; without it that term is omitted and \texttt{info.gradient} says
 \texttt{"none"}. For GRACE-FO's few-hundred-micrometre offset it is of order
 \SI{e-12}{\metre\per\second\squared}.
\item \textbf{Antenna thrust and Lorentz forces} ($\lesssim\SI{e-11}{\metre\per\second\squared}$):
 hooks only.
\item \textbf{A GROOPS cross-validation.} The reader and writer of GROOPS instrument files have
 never been checked against a real GROOPS installation.
\end{itemize}

\section{Data the toolbox does not ship}
Level-1B products, CERES grids, the DE440 ephemeris and the atmosphere model sources are fetched or
supplied by you, never redistributed. One consequence is worth stating rather than discovering:
\textbf{the demos do not run from an installed copy without the data}, because they need the
2020-01-01 GRACE-C arc.

%=======================================================================================
\appendix
\chapter{Complete API}
\label{ch:api}
% GENERATED by tools/buildDocs.py from doc/functionlist.csv (which units exist, and
% each one's H1) and the help block of each unit (everything else).
% DO NOT EDIT: auditPackage.py compares this file by content and refuses a stale one.
% Regenerate with: python tools/buildDocs.py
%
% \input-able from a report-class document; it opens no document and no chapter.

\noindent This reference is generated from the catalogue \texttt{doc/functionlist.csv} and from the help block of each unit, by \texttt{tools/buildDocs.py}. It is therefore what \texttt{help NonGravForces.\textless{}unit\textgreater{}} prints, and cannot drift from it: \texttt{auditPackage.py} re-renders this file and refuses a stale one. 64 units, toolbox version 0.18.0.

\section{L0 \texttt{core} --- constants, options, provenance, uniform errors and the MEX build}

\subsection{\texttt{NonGravForces.Options}}
\label{api:NonGravForces.Options}
\noindent Validated option set for NonGravForces.Model, with fidelity presets.

\begin{lstlisting}
opt = NonGravForces.Options                       % preset "standard"
opt = NonGravForces.Options("precise")
opt = NonGravForces.Options("fast", "drag.density", "msis00")
opt = opt.set("srp.tsi", 1367)                    % legacy TSI, explicit
s   = opt.toStruct()                              % plain struct, for hashing/display
\end{lstlisting}

\subsubsection*{METHODS}
\begin{lstlisting}
opt = set(obj, name, value)        A COPY with ONE dotted field replaced (chain calls for
                                   several); this is a value class, so the original is
                                   unchanged. Unknown names error; text becomes a string.
s   = toStruct(obj)                The same option set as a plain struct, with one field per
                                   family (.preset .drag .srp .erp .thermal .torques .geom
                                   .sun .series .partials). It exists so that provenance can
                                   HASH an option set and so that display is readable; it is
                                   not a second way to configure anything, and nothing reads
                                   it back.
\end{lstlisting}
Configure once, evaluate many (principle P5). Every field is validated at construction; an unknown field name errors, and EVERY ADMISSIBLE VALUE EVALUATES: a value this toolbox cannot compute is not offered (rule R4, HANDOVER LP-C finding F5). Presets fix the DEFAULT tier of every model family; extended models (HWM14 winds, body torques) are opt-ins no preset enables, so preset results never move when one lands. The deliberate moves: 0.17.0's standard preset takes the combined observed density scale (owner ruling 2026-10-29, HANDOVER O-19), and since 2026-10-31 the transient thermal model (owner ruling 2026-10-30): the steady state's recoil led ACT1B's by 8-18 s on y and z over 2020-01, and the transient walls take that to 2-4 s (tools/alignmentMonth.m, HANDOVER "The checkAlignment lags of the 0.17.0 baseline"). A transient Model integrates the walls over the epochs of ONE call, from the steady state at its first epoch (Model.accel 'T0\_K' carries them on; Model.series does so between its chunks). app.DynamicsSource serves it only on its native path, from a temperature table built with 'ThermalOrbit'. The precise preset follows since 2026-10-31 (owner ruling): its steady-state walls led ACT1B by 12-20 s on y and z over five days of 2020-01, the transient walls take that to 2-7 s (release-checks run 41213060789).

\begin{lstlisting}
FIDELITY LADDER
  field                   fast          standard        precise
  drag.density            exponential   msis00          msis21  (+dtm2020 opt-in)
  drag.densityScale       none          combined        none    (+swarm, driver opt-in)
  drag.wind               corotation    corotation      corotation (+hwm14 opt-in)
  drag.gsi                constantCd    sentman         sentman
  drag.accommodation      fixed 0.85    fixed 0.85      fixed 0.85 (March et al. 2021)
  srp.eclipse             conical       conical         solaarsCF
  erp.backend             none          knocke          ceres
  erp.grid                -             12x24           24x48
  thermal.model           none          transient       transient   (+steadyState opt-in)
  geom.selfShadow         none          lut 16          lut 32
  geom.flowShadow         exact         table           exact
  sun.source              rs.analytic   rs.analytic     rs.de440
Common to all three: drag.wind corotation, drag.Tw_K 300, drag.cdScale 1, drag.cd 2.2,
drag.specularFraction [], drag.surfaceMass_amu [], srp.tsi const.S0 (1361),
srp.lunarEclipse false, srp.crScale 1, thermal.internalPower_W 0, torques false,
series.useParallel false, series.chunk 8640, partials.method "analytic".
\end{lstlisting}

\subsubsection*{INPUT ARGUMENTS}
\begin{longtable}{@{}>{\footnotesize\raggedright\arraybackslash\ttfamily}p{3.00cm}>{\footnotesize\raggedright\arraybackslash}p{1.50cm}>{\footnotesize\raggedright\arraybackslash}p{2.20cm}>{\footnotesize\raggedright\arraybackslash}p{1.50cm}>{\footnotesize\raggedright\arraybackslash}p{6.40cm}@{}}
\toprule
{\footnotesize\textbf{Name}} & {\footnotesize\textbf{Size}} & {\footnotesize\textbf{Type}} & {\footnotesize\textbf{Units}} & {\footnotesize\textbf{Description}} \\
\midrule
\endfirsthead
\toprule
{\footnotesize\textbf{Name}} & {\footnotesize\textbf{Size}} & {\footnotesize\textbf{Type}} & {\footnotesize\textbf{Units}} & {\footnotesize\textbf{Description}} \\
\midrule
\endhead
\midrule
\endfoot
\bottomrule
\endlastfoot
preset & 1x1 & string & - & "fast", "standard" (default) or "precise" \\
overrides & 1x2N & cell & - & optional "group.field", value pairs, applied by set() in order \\
\end{longtable}

\subsubsection*{NAME-VALUE ARGUMENTS}
The overrides above: any dotted field of PROPERTIES, e.g. "drag.density", "msis00". Numeric ranges: drag.alpha and drag.specularFraction [0, 1] or [], drag.surfaceMass\_amu (0, 1e6] or [], drag.Tw\_K [1, 2000], drag.cdScale and srp.crScale [0.1, 10], drag.cd [0.5, 10], srp.tsi [1300, 1400], thermal.internalPower\_W [0, 1e4], series.chunk [1, 1e7], geom.lutResolution integer [1, 256]; erp.grid [nRings nSectors] non-negative integers; srp.lunarEclipse, torques.*, series.useParallel logical scalars.


\subsubsection*{OUTPUT ARGUMENTS}
\begin{longtable}{@{}>{\footnotesize\raggedright\arraybackslash\ttfamily}p{3.00cm}>{\footnotesize\raggedright\arraybackslash}p{1.50cm}>{\footnotesize\raggedright\arraybackslash}p{2.20cm}>{\footnotesize\raggedright\arraybackslash}p{7.90cm}@{}}
\toprule
{\footnotesize\textbf{Name}} & {\footnotesize\textbf{Size}} & {\footnotesize\textbf{Type}} & {\footnotesize\textbf{Description}} \\
\midrule
\endfirsthead
\toprule
{\footnotesize\textbf{Name}} & {\footnotesize\textbf{Size}} & {\footnotesize\textbf{Type}} & {\footnotesize\textbf{Description}} \\
\midrule
\endhead
\midrule
\endfoot
\bottomrule
\endlastfoot
opt & 1x1 & NonGravForces.\allowbreak{}Options & the validated option set (value class) \\
\end{longtable}

\subsubsection*{PROPERTIES (read-only after construction; change through set())}
\begin{lstlisting}
preset      string    "fast" | "standard" | "precise"
drag        struct    .density  "exponential" | "msis00" | "msis21" | "dtm2020" (CNES licence:
                                academic, non-commercial use only; THIRD-PARTY-NOTICES.md)
                      .wind     "none" | "corotation" | "hwm14"
                      .gsi      "constantCd" | "sentman" | "maxwell"
                      .accommodation "fixed" | "pilinski2010" | "pilinski2011" | "walker2014" | "sesam"
                      .alpha    accommodation coefficient, read by "fixed" only ([] otherwise)
                      .surfaceMass_amu  substrate particle mass, read by "walker2014" only
                                        (sesam fixes its own at 65 amu)
                      .specularFraction specular fraction, read by "maxwell" only
                      .Tw_K     wall temperature [K]
                      .cdScale  estimated scale on the whole drag force
                      .cd       the coefficient constantCd uses (2.2 for a convex LEO body;
                                Montenbruck & Gill 2005)
                      .densityScale "none" | "swarm" | "driver" | "combined": the model density times an
                                OBSERVED daily factor from Swarm's POD-only densities ("swarm",
                                env.densityScale) or a factor PREDICTED from F30 and Ap for any
                                date ("driver", env.densityDriverScale), or Swarm's where it has a
                                day and the predicted one elsewhere ("combined", env.densityScale
                                with 'Fallback'), passed to Model as 'DensityScale' (HANDOVER O-19).
                                "combined" in the standard preset since 2026-10-29 (owner): a
                                standard Model evaluates only with 'DensityScale', never
                                silently unscaled; set "none" for the unscaled model
srp         struct    .tsi [W/m^2] .eclipse "cylindrical"|"conical"|"solaarsCF"
                      .lunarEclipse (logical) .crScale (estimated scale on the whole
                      radiation force; the optics are per panel, there is no lumped Cr)
erp         struct    .backend "none"|"pointSource"|"knocke"|"ceres"|"ceresClimatology"
                      .grid [nRings nSectors].
                      Earth radiance is Lambertian in every backend (no angular
                      distribution models: the field that pretended otherwise is gone)
thermal     struct    .model "none"|"steadyState"|"transient" .internalPower_W
torques     struct    .gravityGradient .magnetic (logical)
geom        struct    .selfShadow "none"|"lut"|"raycast" .lutResolution (cube-map cells per face edge)
                      .flowShadow "exact"|"table": the gas stream's shadow ratios computed per epoch,
                      or read from geom.flowShadow's table (Method "table"; drag changes by <= 7e-13
                      m/s^2 on GRACE-FO, measured 2026-10-17). "table" needs selfShadow "lut"; with
                      selfShadow "none" or drag.gsi constantCd there is no stream shadow to tabulate
sun         struct    .source "rs.analytic"|"rs.de440"
series      struct    .useParallel (logical) .chunk (double, epochs)
partials    struct    .method "analytic" | "finiteDiff"
\end{lstlisting}
A field a model does not read is NOT silently ignored when it contradicts the choice: Model refuses a non-empty drag.alpha beside a non-fixed accommodation (and fixed without one), a maxwell gas without drag.specularFraction, walker2014 without drag.surfaceMass\_amu (and a surface mass beside any other accommodation), geom.flowShadow "table" beside selfShadow "raycast", and internal power on a macro model with no shares.


\subsubsection*{ERRORS}
\begin{description}[leftmargin=1.2em,style=nextline,itemsep=0pt]
\item[\texttt{NonGravForces:\allowbreak{}Options:\allowbreak{}inputPreset}] unknown preset
\item[\texttt{NonGravForces:\allowbreak{}Options:\allowbreak{}inputField}] unknown dotted field name (or "preset"), or overrides not in pairs
\item[\texttt{NonGravForces:\allowbreak{}Options:\allowbreak{}inputValue}] value outside its admissible set
\end{description}

\subsection{\texttt{NonGravForces.const}}
\label{api:NonGravForces.const}
\noindent Constants owned by NonGravForces (everything geodetic lives in referenceSystem).

\begin{lstlisting}
c = NonGravForces.const
\end{lstlisting}

\subsubsection*{OUTPUT ARGUMENTS}
\begin{longtable}{@{}>{\footnotesize\raggedright\arraybackslash\ttfamily}p{3.00cm}>{\footnotesize\raggedright\arraybackslash}p{1.50cm}>{\footnotesize\raggedright\arraybackslash}p{2.20cm}>{\footnotesize\raggedright\arraybackslash}p{1.50cm}>{\footnotesize\raggedright\arraybackslash}p{6.40cm}@{}}
\toprule
{\footnotesize\textbf{Name}} & {\footnotesize\textbf{Size}} & {\footnotesize\textbf{Type}} & {\footnotesize\textbf{Units}} & {\footnotesize\textbf{Description}} \\
\midrule
\endfirsthead
\toprule
{\footnotesize\textbf{Name}} & {\footnotesize\textbf{Size}} & {\footnotesize\textbf{Type}} & {\footnotesize\textbf{Units}} & {\footnotesize\textbf{Description}} \\
\midrule
\endhead
\midrule
\endfoot
\bottomrule
\endlastfoot
c.\allowbreak{}S0 & 1x1 & double & W/\allowbreak{}m\textasciicircum{}2 & Total solar irradiance at 1 au, 1361 (IAU 2015 B3; SORCE/TIM, TSIS-1). Legacy 1367 is an Options field, not a constant. \\
c.\allowbreak{}au & 1x1 & double & m & Astronomical unit, IAU 2012 (149597870700). referenceSystem publishes none (verified 2026-08). \\
c.\allowbreak{}Rsun & 1x1 & double & m & Nominal solar radius, IAU 2015 B3. \\
c.\allowbreak{}Rmoon & 1x1 & double & m & Mean lunar radius, IAU WGCCRE 2015 (1737.4 km); lunar-eclipse geometry. \\
c.\allowbreak{}sigmaSB & 1x1 & double & W/\allowbreak{}m\textasciicircum{}2/\allowbreak{}K\textasciicircum{}4 & Stefan-Boltzmann, CODATA 2018 exact. \\
c.\allowbreak{}kB & 1x1 & double & J/\allowbreak{}K & Boltzmann, exact. \\
c.\allowbreak{}NA & 1x1 & double & 1/\allowbreak{}mol & Avogadro, exact. \\
c.\allowbreak{}Rgas & 1x1 & double & J/\allowbreak{}mol/\allowbreak{}K & Molar gas constant, exact. \\
c.\allowbreak{}amu & 1x1 & double & kg & Atomic mass unit, CODATA 2018. \\
c.\allowbreak{}eV & 1x1 & double & J & Electron volt, exact (SI 2019). \\
c.\allowbreak{}torr & 1x1 & double & Pa & Torr, exact (101325/760); SESAM's Langmuir parameters are published in torr\textasciicircum{}-1. \\
c.\allowbreak{}speciesName & 1x8 & string & - & He O N2 O2 Ar H N O\_anomalous \\
c.\allowbreak{}speciesMass & 1x8 & double & kg/\allowbreak{}mol & Molar masses of the MSIS species \\
c.\allowbreak{}lambertDiffuse & 1x1 & double & - & 2/3, Lambertian recoil factor \\
c.\allowbreak{}graceFO & 1x1 & struct & - & Handbook (JPL D-56935 Tab.4): .massC .massD [kg] \\
c.\allowbreak{}grace & 1x1 & struct & - & Launch mass (Bettadpur 2012, GRACE 327-720 Sec. I.5.3.3): .massA .massB [kg] \\
\end{longtable}

\subsubsection*{ERRORS}
(none) The function takes no input and cannot fail.


\subsection{\texttt{NonGravForces.toolboxVersion}}
\label{api:NonGravForces.toolboxVersion}
\noindent Semantic version of NonGravForces - the single authority.

\begin{lstlisting}
v = NonGravForces.toolboxVersion            % struct
s = NonGravForces.toolboxVersion('string')  % '0.18.0'
\end{lstlisting}

\subsubsection*{INPUT ARGUMENTS}
\begin{longtable}{@{}>{\footnotesize\raggedright\arraybackslash\ttfamily}p{3.00cm}>{\footnotesize\raggedright\arraybackslash}p{1.50cm}>{\footnotesize\raggedright\arraybackslash}p{2.20cm}>{\footnotesize\raggedright\arraybackslash}p{1.50cm}>{\footnotesize\raggedright\arraybackslash}p{6.40cm}@{}}
\toprule
{\footnotesize\textbf{Name}} & {\footnotesize\textbf{Size}} & {\footnotesize\textbf{Type}} & {\footnotesize\textbf{Units}} & {\footnotesize\textbf{Description}} \\
\midrule
\endfirsthead
\toprule
{\footnotesize\textbf{Name}} & {\footnotesize\textbf{Size}} & {\footnotesize\textbf{Type}} & {\footnotesize\textbf{Units}} & {\footnotesize\textbf{Description}} \\
\midrule
\endhead
\midrule
\endfoot
\bottomrule
\endlastfoot
format & 1x1 & string & - & "struct" (default) or "string". \\
\end{longtable}

\subsubsection*{OUTPUT ARGUMENTS}
\begin{longtable}{@{}>{\footnotesize\raggedright\arraybackslash\ttfamily}p{3.00cm}>{\footnotesize\raggedright\arraybackslash}p{1.50cm}>{\footnotesize\raggedright\arraybackslash}p{2.20cm}>{\footnotesize\raggedright\arraybackslash}p{1.50cm}>{\footnotesize\raggedright\arraybackslash}p{6.40cm}@{}}
\toprule
{\footnotesize\textbf{Name}} & {\footnotesize\textbf{Size}} & {\footnotesize\textbf{Type}} & {\footnotesize\textbf{Units}} & {\footnotesize\textbf{Description}} \\
\midrule
\endfirsthead
\toprule
{\footnotesize\textbf{Name}} & {\footnotesize\textbf{Size}} & {\footnotesize\textbf{Type}} & {\footnotesize\textbf{Units}} & {\footnotesize\textbf{Description}} \\
\midrule
\endhead
\midrule
\endfoot
\bottomrule
\endlastfoot
v & 1x1 & struct & - & .major .minor .patch (uint16), .semantic (string, composed), .date (string, ISO), .concept (uint8, 3 = this redesign), .model (string: authoring model of this revision). With "string": char '0.18.0'. \\
\end{longtable}

\subsubsection*{ERRORS}
\begin{description}[leftmargin=1.2em,style=nextline,itemsep=0pt]
\item[\texttt{MATLAB:\allowbreak{}validators:\allowbreak{}mustBeMember}] Unknown format keyword.
\end{description}

\subsection{\texttt{NonGravForces.dependencyReport}}
\label{api:NonGravForces.dependencyReport}
\noindent Presence and version of the five ecosystem toolboxes and MICE.

\begin{lstlisting}
deps = NonGravForces.dependencyReport
deps = NonGravForces.dependencyReport('Verbose', false)
\end{lstlisting}

\subsubsection*{NAME-VALUE ARGUMENTS}
\begin{longtable}{@{}>{\footnotesize\raggedright\arraybackslash\ttfamily}p{3.00cm}>{\footnotesize\raggedright\arraybackslash}p{1.50cm}>{\footnotesize\raggedright\arraybackslash}p{2.20cm}>{\footnotesize\raggedright\arraybackslash}p{1.50cm}>{\footnotesize\raggedright\arraybackslash}p{6.40cm}@{}}
\toprule
{\footnotesize\textbf{Name}} & {\footnotesize\textbf{Size}} & {\footnotesize\textbf{Type}} & {\footnotesize\textbf{Units}} & {\footnotesize\textbf{Description}} \\
\midrule
\endfirsthead
\toprule
{\footnotesize\textbf{Name}} & {\footnotesize\textbf{Size}} & {\footnotesize\textbf{Type}} & {\footnotesize\textbf{Units}} & {\footnotesize\textbf{Description}} \\
\midrule
\endhead
\midrule
\endfoot
\bottomrule
\endlastfoot
Verbose & 1x1 & logical & - & Print a table (default true). \\
\end{longtable}

\subsubsection*{OUTPUT ARGUMENTS}
\begin{longtable}{@{}>{\footnotesize\raggedright\arraybackslash\ttfamily}p{3.00cm}>{\footnotesize\raggedright\arraybackslash}p{1.50cm}>{\footnotesize\raggedright\arraybackslash}p{2.20cm}>{\footnotesize\raggedright\arraybackslash}p{1.50cm}>{\footnotesize\raggedright\arraybackslash}p{6.40cm}@{}}
\toprule
{\footnotesize\textbf{Name}} & {\footnotesize\textbf{Size}} & {\footnotesize\textbf{Type}} & {\footnotesize\textbf{Units}} & {\footnotesize\textbf{Description}} \\
\midrule
\endfirsthead
\toprule
{\footnotesize\textbf{Name}} & {\footnotesize\textbf{Size}} & {\footnotesize\textbf{Type}} & {\footnotesize\textbf{Units}} & {\footnotesize\textbf{Description}} \\
\midrule
\endhead
\midrule
\endfoot
\bottomrule
\endlastfoot
deps & 1x6 & struct & - & One element each for referenceSystem, gracefoReading (required), gjToolbox, timeSeriesPreparation, lsqvce and mice (optional), in that order: .name .required (logical) .present (logical) .version (string, "" if absent) .minimum (string) .comparable (logical: the two are dotted numbers and can be ordered) .belowMinimum (logical: present, comparable and older than .minimum) .setup (string: setup function) .role (string) \\
\end{longtable}

\subsubsection*{ERRORS}
None. Warning NonGravForces:dependencyReport:depMissing per absent

REQUIRED dependency (an absent optional one is only reported).


\subsection{\texttt{NonGravForces.fail}}
\label{api:NonGravForces.fail}
\noindent Build (and by default throw) a NonGravForces error with a uniform identifier.

\begin{lstlisting}
NonGravForces.fail(unit, mnemonic, message, args...)          throws
ME = NonGravForces.fail(unit, mnemonic, message, args...)     returns
\end{lstlisting}

\subsubsection*{INPUT ARGUMENTS}
\begin{longtable}{@{}>{\footnotesize\raggedright\arraybackslash\ttfamily}p{3.00cm}>{\footnotesize\raggedright\arraybackslash}p{1.50cm}>{\footnotesize\raggedright\arraybackslash}p{2.20cm}>{\footnotesize\raggedright\arraybackslash}p{1.50cm}>{\footnotesize\raggedright\arraybackslash}p{6.40cm}@{}}
\toprule
{\footnotesize\textbf{Name}} & {\footnotesize\textbf{Size}} & {\footnotesize\textbf{Type}} & {\footnotesize\textbf{Units}} & {\footnotesize\textbf{Description}} \\
\midrule
\endfirsthead
\toprule
{\footnotesize\textbf{Name}} & {\footnotesize\textbf{Size}} & {\footnotesize\textbf{Type}} & {\footnotesize\textbf{Units}} & {\footnotesize\textbf{Description}} \\
\midrule
\endhead
\midrule
\endfoot
\bottomrule
\endlastfoot
unit & 1x1 & string & - & Function or class name, e.g. "eclipseFactor" \\
mnemonic & 1x1 & string & - & camelCase reason, first token from the taxonomy above, e.g. "inputShape" \\
message & 1x1 & string & - & sprintf-style message. Must contain the correction hint after the marker ' -\textgreater{} '. \\
args & any & - & - & sprintf arguments for message \\
\end{longtable}

\subsubsection*{OUTPUT ARGUMENTS}
\begin{longtable}{@{}>{\footnotesize\raggedright\arraybackslash\ttfamily}p{3.00cm}>{\footnotesize\raggedright\arraybackslash}p{1.50cm}>{\footnotesize\raggedright\arraybackslash}p{2.20cm}>{\footnotesize\raggedright\arraybackslash}p{1.50cm}>{\footnotesize\raggedright\arraybackslash}p{6.40cm}@{}}
\toprule
{\footnotesize\textbf{Name}} & {\footnotesize\textbf{Size}} & {\footnotesize\textbf{Type}} & {\footnotesize\textbf{Units}} & {\footnotesize\textbf{Description}} \\
\midrule
\endfirsthead
\toprule
{\footnotesize\textbf{Name}} & {\footnotesize\textbf{Size}} & {\footnotesize\textbf{Type}} & {\footnotesize\textbf{Units}} & {\footnotesize\textbf{Description}} \\
\midrule
\endhead
\midrule
\endfoot
\bottomrule
\endlastfoot
ME & 1x1 & MException & - & Only when an output is requested; then nothing is thrown. \\
\end{longtable}

\subsubsection*{ERRORS}
\begin{description}[leftmargin=1.2em,style=nextline,itemsep=0pt]
\item[\texttt{NonGravForces:\allowbreak{}fail:\allowbreak{}inputMnemonic}] mnemonic does not start with a taxonomy token
\item[\texttt{NonGravForces:\allowbreak{}fail:\allowbreak{}inputHint}] message lacks the ' -\textgreater{} ' hint marker
\item[\texttt{MATLAB:\allowbreak{}validators:\allowbreak{}mustBeNonzeroLengthText}] unit, mnemonic or message empty
\end{description}
These two identifiers are raised with error(), not through fail itself.


\subsection{\texttt{NonGravForces.provenance}}
\label{api:NonGravForces.provenance}
\noindent Metadata record that travels with every NonGravForces result (rule R2).

\begin{lstlisting}
p = NonGravForces.provenance("Model.accel")
p = NonGravForces.provenance("io.readACC1B", 'Files', files, 'Options', opts)
\end{lstlisting}

\subsubsection*{INPUT ARGUMENTS}
\begin{longtable}{@{}>{\footnotesize\raggedright\arraybackslash\ttfamily}p{3.00cm}>{\footnotesize\raggedright\arraybackslash}p{1.50cm}>{\footnotesize\raggedright\arraybackslash}p{2.20cm}>{\footnotesize\raggedright\arraybackslash}p{1.50cm}>{\footnotesize\raggedright\arraybackslash}p{6.40cm}@{}}
\toprule
{\footnotesize\textbf{Name}} & {\footnotesize\textbf{Size}} & {\footnotesize\textbf{Type}} & {\footnotesize\textbf{Units}} & {\footnotesize\textbf{Description}} \\
\midrule
\endfirsthead
\toprule
{\footnotesize\textbf{Name}} & {\footnotesize\textbf{Size}} & {\footnotesize\textbf{Type}} & {\footnotesize\textbf{Units}} & {\footnotesize\textbf{Description}} \\
\midrule
\endhead
\midrule
\endfoot
\bottomrule
\endlastfoot
unit & 1x1 & string & - & Producing unit, qualified inside the package \\
\end{longtable}

\subsubsection*{NAME-VALUE ARGUMENTS}
\begin{longtable}{@{}>{\footnotesize\raggedright\arraybackslash\ttfamily}p{3.00cm}>{\footnotesize\raggedright\arraybackslash}p{1.50cm}>{\footnotesize\raggedright\arraybackslash}p{2.20cm}>{\footnotesize\raggedright\arraybackslash}p{1.50cm}>{\footnotesize\raggedright\arraybackslash}p{6.40cm}@{}}
\toprule
{\footnotesize\textbf{Name}} & {\footnotesize\textbf{Size}} & {\footnotesize\textbf{Type}} & {\footnotesize\textbf{Units}} & {\footnotesize\textbf{Description}} \\
\midrule
\endfirsthead
\toprule
{\footnotesize\textbf{Name}} & {\footnotesize\textbf{Size}} & {\footnotesize\textbf{Type}} & {\footnotesize\textbf{Units}} & {\footnotesize\textbf{Description}} \\
\midrule
\endhead
\midrule
\endfoot
\bottomrule
\endlastfoot
Files & 1xN & string & - & Data files consumed (default empty) \\
Options & 1x1 & struct & - & Resolved options; hashed (default empty) \\
Warnings & 1xN & string & - & Warning identifiers issued (default empty) \\
Note & 1x1 & string & - & Free text (default "") \\
\end{longtable}

\subsubsection*{OUTPUT ARGUMENTS}
\begin{longtable}{@{}>{\footnotesize\raggedright\arraybackslash\ttfamily}p{3.00cm}>{\footnotesize\raggedright\arraybackslash}p{1.50cm}>{\footnotesize\raggedright\arraybackslash}p{2.20cm}>{\footnotesize\raggedright\arraybackslash}p{1.50cm}>{\footnotesize\raggedright\arraybackslash}p{6.40cm}@{}}
\toprule
{\footnotesize\textbf{Name}} & {\footnotesize\textbf{Size}} & {\footnotesize\textbf{Type}} & {\footnotesize\textbf{Units}} & {\footnotesize\textbf{Description}} \\
\midrule
\endfirsthead
\toprule
{\footnotesize\textbf{Name}} & {\footnotesize\textbf{Size}} & {\footnotesize\textbf{Type}} & {\footnotesize\textbf{Units}} & {\footnotesize\textbf{Description}} \\
\midrule
\endhead
\midrule
\endfoot
\bottomrule
\endlastfoot
p & 1x1 & struct & - & .unit .toolboxVersion .dependencies (name, version; cached per session) .optionHash (string, 8 hex) .files .warnings .note .created (datetime UTC) \\
\end{longtable}

\subsubsection*{ERRORS}
None of its own; MATLAB:validators:* for arguments of the wrong class or size.


\subsection{\texttt{NonGravForces.buildMex}}
\label{api:NonGravForces.buildMex}
\noindent Compile the MEX gateways in mex/ into their package private folders, or report their state.

\begin{lstlisting}
report = NonGravForces.buildMex                          % build every target
report = NonGravForces.buildMex('Check', true)            % report only, nothing compiled
report = NonGravForces.buildMex('Target', "kernel")       % one target, scalar report
\end{lstlisting}

\subsubsection*{NAME-VALUE ARGUMENTS}
\begin{longtable}{@{}>{\footnotesize\raggedright\arraybackslash\ttfamily}p{3.00cm}>{\footnotesize\raggedright\arraybackslash}p{1.50cm}>{\footnotesize\raggedright\arraybackslash}p{2.20cm}>{\footnotesize\raggedright\arraybackslash}p{1.50cm}>{\footnotesize\raggedright\arraybackslash}p{6.40cm}@{}}
\toprule
{\footnotesize\textbf{Name}} & {\footnotesize\textbf{Size}} & {\footnotesize\textbf{Type}} & {\footnotesize\textbf{Units}} & {\footnotesize\textbf{Description}} \\
\midrule
\endfirsthead
\toprule
{\footnotesize\textbf{Name}} & {\footnotesize\textbf{Size}} & {\footnotesize\textbf{Type}} & {\footnotesize\textbf{Units}} & {\footnotesize\textbf{Description}} \\
\midrule
\endhead
\midrule
\endfoot
\bottomrule
\endlastfoot
Check & 1x1 & logical & - & true: report only (default false). \\
Target & 1x1 & string & - & "all" (default), "msis", "msis21", "hwm14", "dtm2020", "rhs" or "kernel". "all" returns one report per target, in that order. \\
MsisSource & 1x1 & string & - & Directory holding the NRLMSISE-00 C source. Default: the NGF\_NRLMSISE00 environment variable if set, else mex/nrlmsise00. \\
Msis21Source & 1x1 & string & - & Directory holding the NRLMSIS 2.1 package. Default: the NGF\_NRLMSIS21 environment variable if set, else mex/nrlmsis21. \\
Hwm14Source & 1x1 & string & - & Directory holding the HWM14 files. Default: the NGF\_HWM14 environment variable if set, else mex/hwm14. \\
Dtm2020Source & 1x1 & string & - & Directory holding the DTM2020 files. Default: the NGF\_DTM2020 environment variable if set, else mex/dtm2020. \\
GjToolboxSource & 1x1 & string & - & Directory holding gj\_native.h (odeSuite's csrc). Default: the NGF\_GJTOOLBOX environment variable, else the csrc folder of the gjToolbox on the path. \\
\end{longtable}

\subsubsection*{OUTPUT ARGUMENTS}
\begin{longtable}{@{}>{\footnotesize\raggedright\arraybackslash\ttfamily}p{3.00cm}>{\footnotesize\raggedright\arraybackslash}p{1.50cm}>{\footnotesize\raggedright\arraybackslash}p{2.20cm}>{\footnotesize\raggedright\arraybackslash}p{1.50cm}>{\footnotesize\raggedright\arraybackslash}p{6.40cm}@{}}
\toprule
{\footnotesize\textbf{Name}} & {\footnotesize\textbf{Size}} & {\footnotesize\textbf{Type}} & {\footnotesize\textbf{Units}} & {\footnotesize\textbf{Description}} \\
\midrule
\endfirsthead
\toprule
{\footnotesize\textbf{Name}} & {\footnotesize\textbf{Size}} & {\footnotesize\textbf{Type}} & {\footnotesize\textbf{Units}} & {\footnotesize\textbf{Description}} \\
\midrule
\endhead
\midrule
\endfoot
\bottomrule
\endlastfoot
report & 1xT & struct & - & One element per target (T = 1 for a named Target, 6 for "all"), each with .target (string), .file (string), .present (logical), .upToDate (logical: present and digest matches the sources), .sourceDigest (string, SHA-256 of the sources), .binaryDigest (string recorded at build time, "" if none), .compiler (string, "" when not built now), .built (logical: compiled in this call), .sources (1xN string), .sourcesPresent (logical: every source file exists; false is a REPORT under 'Check' and an error otherwise), .sourceDir (string, where the third-party source was looked for; gjToolbox's csrc for rhs; "" for the kernel, which has none, and for an unbuildable rhs). The digest covers whatever source was actually used, so a different NRLMSISE-00 release marks the binary stale and forces a rebuild -- which is the mechanism that makes tracking upstream safe rather than merely convenient. \\
\end{longtable}

\subsubsection*{ERRORS}
\begin{description}[leftmargin=1.2em,style=nextline,itemsep=0pt]
\item[\texttt{NonGravForces:\allowbreak{}buildMex:\allowbreak{}depCompiler}] no C compiler selected and no MinGW-w64 support package found; or no gfortran beside the selected compiler (msis21, hwm14, dtm2020), or it does not run, or its static runtime is missing, or mex -setup C did not select the MinGW-w64 found
\item[\texttt{NonGravForces:\allowbreak{}buildMex:\allowbreak{}depMissing}] the NRLMSISE-00, NRLMSIS 2.1, HWM14 or DTM2020 source is absent, or Target "rhs" without gjToolbox or its gj\_native.h
\item[\texttt{NonGravForces:\allowbreak{}buildMex:\allowbreak{}dataMissing}] a source file of mex/ is missing
\item[\texttt{NonGravForces:\allowbreak{}buildMex:\allowbreak{}numericBuild}] gfortran refused a file (its output is in the message)
\item[\texttt{MATLAB:\allowbreak{}mex:\allowbreak{}*}] compiler errors propagate unwrapped
\end{description}

\section{L1 \texttt{io} --- Level-1B products, the mass model, space weather and external grids}

\subsection{\texttt{NonGravForces.io.readACC1B}}
\label{api:io.readACC1B}
\noindent ACC1B/ACT1B/ACH1B reader over gracefoReading: 3xK SRF accelerations, GPS epochs, aux.

\begin{lstlisting}
[epoch, a_srf, aux] = NonGravForces.io.readACC1B(file)
[epoch, a_srf, aux] = NonGravForces.io.readACC1B(files, 'TimeRange', tr, 'QualityBitMask', bits)
\end{lstlisting}

\subsubsection*{INPUT ARGUMENTS}
\begin{longtable}{@{}>{\footnotesize\raggedright\arraybackslash\ttfamily}p{3.00cm}>{\footnotesize\raggedright\arraybackslash}p{1.50cm}>{\footnotesize\raggedright\arraybackslash}p{2.20cm}>{\footnotesize\raggedright\arraybackslash}p{1.50cm}>{\footnotesize\raggedright\arraybackslash}p{6.40cm}@{}}
\toprule
{\footnotesize\textbf{Name}} & {\footnotesize\textbf{Size}} & {\footnotesize\textbf{Type}} & {\footnotesize\textbf{Units}} & {\footnotesize\textbf{Description}} \\
\midrule
\endfirsthead
\toprule
{\footnotesize\textbf{Name}} & {\footnotesize\textbf{Size}} & {\footnotesize\textbf{Type}} & {\footnotesize\textbf{Units}} & {\footnotesize\textbf{Description}} \\
\midrule
\endhead
\midrule
\endfoot
\bottomrule
\endlastfoot
files & 1xF & string & - & One or more ACC1B/ACT1B/ACH1B files of ONE satellite, chronological; .gz accepted. \\
\end{longtable}

\subsubsection*{NAME-VALUE ARGUMENTS}
\begin{longtable}{@{}>{\footnotesize\raggedright\arraybackslash\ttfamily}p{3.00cm}>{\footnotesize\raggedright\arraybackslash}p{1.50cm}>{\footnotesize\raggedright\arraybackslash}p{2.20cm}>{\footnotesize\raggedright\arraybackslash}p{1.50cm}>{\footnotesize\raggedright\arraybackslash}p{6.40cm}@{}}
\toprule
{\footnotesize\textbf{Name}} & {\footnotesize\textbf{Size}} & {\footnotesize\textbf{Type}} & {\footnotesize\textbf{Units}} & {\footnotesize\textbf{Description}} \\
\midrule
\endfirsthead
\toprule
{\footnotesize\textbf{Name}} & {\footnotesize\textbf{Size}} & {\footnotesize\textbf{Type}} & {\footnotesize\textbf{Units}} & {\footnotesize\textbf{Description}} \\
\midrule
\endhead
\midrule
\endfoot
\bottomrule
\endlastfoot
TimeRange & 1x2 & double\textbar{}\allowbreak{}datetime & - & Inclusive window forwarded to gracefoReading.readL1b: decimal MJD, or datetime on the GPS scale. Default [] (whole files). \\
QualityBitMask & 1xQ & double & - & qualflg bits (0-based, bit 0 rightmost) whose records are removed; forwarded. Default [] (nothing removed). \\
\end{longtable}

\subsubsection*{OUTPUT ARGUMENTS}
\begin{longtable}{@{}>{\footnotesize\raggedright\arraybackslash\ttfamily}p{3.00cm}>{\footnotesize\raggedright\arraybackslash}p{1.50cm}>{\footnotesize\raggedright\arraybackslash}p{2.20cm}>{\footnotesize\raggedright\arraybackslash}p{1.50cm}>{\footnotesize\raggedright\arraybackslash}p{6.40cm}@{}}
\toprule
{\footnotesize\textbf{Name}} & {\footnotesize\textbf{Size}} & {\footnotesize\textbf{Type}} & {\footnotesize\textbf{Units}} & {\footnotesize\textbf{Description}} \\
\midrule
\endfirsthead
\toprule
{\footnotesize\textbf{Name}} & {\footnotesize\textbf{Size}} & {\footnotesize\textbf{Type}} & {\footnotesize\textbf{Units}} & {\footnotesize\textbf{Description}} \\
\midrule
\endhead
\midrule
\endfoot
\bottomrule
\endlastfoot
epoch & 1x1 & struct & - & referenceSystem epoch struct, K entries, scale "GPS": .day (MJD), .sec (s of day). \\
a\_\allowbreak{}srf & 3xK & double & m/\allowbreak{}s\textasciicircum{}2 & lin\_accl\_x, \_y, \_z in the SRF. \\
aux & 1x1 & struct & - & .angular\_rad\_s2 (3xK, ang\_accl\_*), .residual\_m\_s2 (3xK acl\_*\_res, or [] if the product has none), .qualflgBits (Kx8 logical), .product, .satellite ("C"\textbar{}"D"), .files, .numRecords, .sampling\_s (median spacing), .obsNames, .units, .provenance. \\
\end{longtable}

\subsubsection*{ERRORS}
\begin{description}[leftmargin=1.2em,style=nextline,itemsep=0pt]
\item[\texttt{NonGravForces:\allowbreak{}readACC1B:\allowbreak{}dataProduct}] a file is not an ACC1B/ACT1B/ACH1B product
\item[\texttt{NonGravForces:\allowbreak{}readACC1B:\allowbreak{}dataSatellite}] the files mix GRACE-C and GRACE-D
\item[\texttt{NonGravForces:\allowbreak{}readACC1B:\allowbreak{}dataEmpty}] no record survives TimeRange/QualityBitMask
\item[\texttt{NonGravForces:\allowbreak{}readACC1B:\allowbreak{}dataColumn}] lin\_accl\_* or ang\_accl\_* missing from the header
\item[\texttt{gracefoReading:\allowbreak{}*}] propagated unwrapped (missing file, malformed record, column mismatch)
\end{description}

\subsection{\texttt{NonGravForces.io.readSCA1B}}
\label{api:io.readSCA1B}
\noindent SCA1B reader: ICRF-\textgreater{}SRF quaternions conjugated ONCE to BODY-\textgreater{}GCRS, 4xK scalar first.

\begin{lstlisting}
[epoch, q, aux] = NonGravForces.io.readSCA1B(file)
[epoch, q, aux] = NonGravForces.io.readSCA1B(files, 'TimeRange', tr, 'QualityBitMask', bits)
\end{lstlisting}

\subsubsection*{INPUT ARGUMENTS}
\begin{longtable}{@{}>{\footnotesize\raggedright\arraybackslash\ttfamily}p{3.00cm}>{\footnotesize\raggedright\arraybackslash}p{1.50cm}>{\footnotesize\raggedright\arraybackslash}p{2.20cm}>{\footnotesize\raggedright\arraybackslash}p{1.50cm}>{\footnotesize\raggedright\arraybackslash}p{6.40cm}@{}}
\toprule
{\footnotesize\textbf{Name}} & {\footnotesize\textbf{Size}} & {\footnotesize\textbf{Type}} & {\footnotesize\textbf{Units}} & {\footnotesize\textbf{Description}} \\
\midrule
\endfirsthead
\toprule
{\footnotesize\textbf{Name}} & {\footnotesize\textbf{Size}} & {\footnotesize\textbf{Type}} & {\footnotesize\textbf{Units}} & {\footnotesize\textbf{Description}} \\
\midrule
\endhead
\midrule
\endfoot
\bottomrule
\endlastfoot
files & 1xF & string & - & One or more SCA1B files of ONE satellite, chronological; .gz accepted. \\
\end{longtable}

\subsubsection*{NAME-VALUE ARGUMENTS}
\begin{longtable}{@{}>{\footnotesize\raggedright\arraybackslash\ttfamily}p{3.00cm}>{\footnotesize\raggedright\arraybackslash}p{1.50cm}>{\footnotesize\raggedright\arraybackslash}p{2.20cm}>{\footnotesize\raggedright\arraybackslash}p{1.50cm}>{\footnotesize\raggedright\arraybackslash}p{6.40cm}@{}}
\toprule
{\footnotesize\textbf{Name}} & {\footnotesize\textbf{Size}} & {\footnotesize\textbf{Type}} & {\footnotesize\textbf{Units}} & {\footnotesize\textbf{Description}} \\
\midrule
\endfirsthead
\toprule
{\footnotesize\textbf{Name}} & {\footnotesize\textbf{Size}} & {\footnotesize\textbf{Type}} & {\footnotesize\textbf{Units}} & {\footnotesize\textbf{Description}} \\
\midrule
\endhead
\midrule
\endfoot
\bottomrule
\endlastfoot
TimeRange & 1x2 & double\textbar{}\allowbreak{}datetime & - & Forwarded to readL1b: decimal MJD or GPS datetime. Default []. \\
QualityBitMask & 1xQ & double & - & qualflg bits removing a record; forwarded. Default []. \\
\end{longtable}

\subsubsection*{OUTPUT ARGUMENTS}
\begin{longtable}{@{}>{\footnotesize\raggedright\arraybackslash\ttfamily}p{3.00cm}>{\footnotesize\raggedright\arraybackslash}p{1.50cm}>{\footnotesize\raggedright\arraybackslash}p{2.20cm}>{\footnotesize\raggedright\arraybackslash}p{1.50cm}>{\footnotesize\raggedright\arraybackslash}p{6.40cm}@{}}
\toprule
{\footnotesize\textbf{Name}} & {\footnotesize\textbf{Size}} & {\footnotesize\textbf{Type}} & {\footnotesize\textbf{Units}} & {\footnotesize\textbf{Description}} \\
\midrule
\endfirsthead
\toprule
{\footnotesize\textbf{Name}} & {\footnotesize\textbf{Size}} & {\footnotesize\textbf{Type}} & {\footnotesize\textbf{Units}} & {\footnotesize\textbf{Description}} \\
\midrule
\endhead
\midrule
\endfoot
\bottomrule
\endlastfoot
epoch & 1x1 & struct & - & referenceSystem epoch struct, K entries, "GPS". \\
q\_\allowbreak{}body2gcrs & 4xK & double & - & Unit quaternions BODY-\textgreater{}GCRS, [q0;q1;q2;q3], canonical sign q0 \textgreater{}= 0 (as referenceSystem.invertQuaternion returns). \\
aux & 1x1 & struct & - & .frame "BODY-\textgreater{}GCRS", .convention (referenceSystem.attitudeConvention()), .qualRss (1xK, qual\_rss column), .scaId (1xK), .normDeviationMax (max \textbar{}norm-1\textbar{} of the file quaternions), .qualflgBits, .product, .satellite, .files, .numRecords, .sampling\_s, .obsNames, .units, .provenance. \\
\end{longtable}

\subsubsection*{ERRORS}
\begin{description}[leftmargin=1.2em,style=nextline,itemsep=0pt]
\item[\texttt{NonGravForces:\allowbreak{}readSCA1B:\allowbreak{}dataProduct}] a file is not an SCA1B product
\item[\texttt{NonGravForces:\allowbreak{}readSCA1B:\allowbreak{}dataSatellite}] the files mix GRACE-C and GRACE-D
\item[\texttt{NonGravForces:\allowbreak{}readSCA1B:\allowbreak{}dataEmpty}] no record survives the filters
\item[\texttt{NonGravForces:\allowbreak{}readSCA1B:\allowbreak{}dataColumn}] a quaternion column is missing
\item[\texttt{NonGravForces:\allowbreak{}readSCA1B:\allowbreak{}dataQuaternion}] a file quaternion is not of unit norm
\item[\texttt{gracefoReading:\allowbreak{}*, referenceSystem:\allowbreak{}*}] propagated unwrapped
\end{description}

\subsection{\texttt{NonGravForces.io.readGNV1B}}
\label{api:io.readGNV1B}
\noindent GNV1B reader: Earth-fixed state rotated to the GCRS through referenceSystem, 3xK.

\begin{lstlisting}
[epoch, r, v, aux] = NonGravForces.io.readGNV1B(file, 'Eop', eop)
[epoch, r, v, aux] = NonGravForces.io.readGNV1B(files, 'Eop', eopFile, 'TimeRange', tr)
[epoch, r, v, aux] = NonGravForces.io.readGNV1B(file, 'Frame', "ITRS")
\end{lstlisting}

\subsubsection*{INPUT ARGUMENTS}
\begin{longtable}{@{}>{\footnotesize\raggedright\arraybackslash\ttfamily}p{3.00cm}>{\footnotesize\raggedright\arraybackslash}p{1.50cm}>{\footnotesize\raggedright\arraybackslash}p{2.20cm}>{\footnotesize\raggedright\arraybackslash}p{1.50cm}>{\footnotesize\raggedright\arraybackslash}p{6.40cm}@{}}
\toprule
{\footnotesize\textbf{Name}} & {\footnotesize\textbf{Size}} & {\footnotesize\textbf{Type}} & {\footnotesize\textbf{Units}} & {\footnotesize\textbf{Description}} \\
\midrule
\endfirsthead
\toprule
{\footnotesize\textbf{Name}} & {\footnotesize\textbf{Size}} & {\footnotesize\textbf{Type}} & {\footnotesize\textbf{Units}} & {\footnotesize\textbf{Description}} \\
\midrule
\endhead
\midrule
\endfoot
\bottomrule
\endlastfoot
files & 1xF & string & - & One or more GNV1B files of ONE satellite, chronological; .gz accepted. \\
\end{longtable}

\subsubsection*{NAME-VALUE ARGUMENTS}
\begin{longtable}{@{}>{\footnotesize\raggedright\arraybackslash\ttfamily}p{3.00cm}>{\footnotesize\raggedright\arraybackslash}p{1.50cm}>{\footnotesize\raggedright\arraybackslash}p{2.20cm}>{\footnotesize\raggedright\arraybackslash}p{1.50cm}>{\footnotesize\raggedright\arraybackslash}p{6.40cm}@{}}
\toprule
{\footnotesize\textbf{Name}} & {\footnotesize\textbf{Size}} & {\footnotesize\textbf{Type}} & {\footnotesize\textbf{Units}} & {\footnotesize\textbf{Description}} \\
\midrule
\endfirsthead
\toprule
{\footnotesize\textbf{Name}} & {\footnotesize\textbf{Size}} & {\footnotesize\textbf{Type}} & {\footnotesize\textbf{Units}} & {\footnotesize\textbf{Description}} \\
\midrule
\endhead
\midrule
\endfoot
\bottomrule
\endlastfoot
Eop & 1x1 & struct\textbar{}\allowbreak{}string & - & EOP series from referenceSystem.readEop, or the path of a series file (forwarded as 'EopFile'). REQUIRED for Frame "GCRS"; ignored for "ITRS". Default [] (none). \\
Frame & 1x1 & string & - & "GCRS" (default) or "ITRS". \\
TimeRange & 1x2 & double\textbar{}\allowbreak{}datetime & - & Forwarded to readL1b. Default []. \\
QualityBitMask & 1xQ & double & - & Forwarded to readL1b. Default []. \\
\end{longtable}

\subsubsection*{OUTPUT ARGUMENTS}
\begin{longtable}{@{}>{\footnotesize\raggedright\arraybackslash\ttfamily}p{3.00cm}>{\footnotesize\raggedright\arraybackslash}p{1.50cm}>{\footnotesize\raggedright\arraybackslash}p{2.20cm}>{\footnotesize\raggedright\arraybackslash}p{1.50cm}>{\footnotesize\raggedright\arraybackslash}p{6.40cm}@{}}
\toprule
{\footnotesize\textbf{Name}} & {\footnotesize\textbf{Size}} & {\footnotesize\textbf{Type}} & {\footnotesize\textbf{Units}} & {\footnotesize\textbf{Description}} \\
\midrule
\endfirsthead
\toprule
{\footnotesize\textbf{Name}} & {\footnotesize\textbf{Size}} & {\footnotesize\textbf{Type}} & {\footnotesize\textbf{Units}} & {\footnotesize\textbf{Description}} \\
\midrule
\endhead
\midrule
\endfoot
\bottomrule
\endlastfoot
epoch & 1x1 & struct & - & referenceSystem epoch struct, K entries, "GPS". \\
r & 3xK & double & m & Position in aux.frame. \\
v & 3xK & double & m/\allowbreak{}s & Velocity in aux.frame (GCRS: includes Rdot r). \\
aux & 1x1 & struct & - & .frame ("GCRS"\textbar{}"ITRS"), .sigma\_r\_m (3xK, xpos\_err.., ITRS as printed), .sigma\_v\_m\_s (3xK), .rotation (info struct of rotationItrsToGcrs; [] for ITRS), .eopSeries (eop.meta.series or the file path; "" for ITRS), .qualflgBits, .product, .satellite, .files, .numRecords, .sampling\_s, .obsNames, .units, .provenance. \\
\end{longtable}

\subsubsection*{ERRORS}
\begin{description}[leftmargin=1.2em,style=nextline,itemsep=0pt]
\item[\texttt{NonGravForces:\allowbreak{}readGNV1B:\allowbreak{}inputEop}] Frame "GCRS" without a usable Eop
\item[\texttt{NonGravForces:\allowbreak{}readGNV1B:\allowbreak{}dataProduct}] a file is not a GNV1B product
\item[\texttt{NonGravForces:\allowbreak{}readGNV1B:\allowbreak{}dataSatellite}] the files mix satellites
\item[\texttt{NonGravForces:\allowbreak{}readGNV1B:\allowbreak{}dataEmpty}] no record survives the filters
\item[\texttt{NonGravForces:\allowbreak{}readGNV1B:\allowbreak{}dataColumn}] a position/velocity column is missing
\item[\texttt{NonGravForces:\allowbreak{}readGNV1B:\allowbreak{}dataFrame}] coord\_ref is not 'E' for every record
\item[\texttt{referenceSystem:\allowbreak{}*, gracefoReading:\allowbreak{}*}] propagated unwrapped (a pre-1972 EOP series is trimmed with a warning, not refused; see the note above)
\end{description}

\subsection{\texttt{NonGravForces.io.readTHR1B}}
\label{api:io.readTHR1B}
\noindent THR1B thruster-activation reader: timeSeriesPreparation.qc.eventTable impulse priors plus per-thruster on-times.

\begin{lstlisting}
[events, aux] = NonGravForces.io.readTHR1B(file)
[events, aux] = NonGravForces.io.readTHR1B(files, 'TimeRange', tr)
\end{lstlisting}

\subsubsection*{INPUT ARGUMENTS}
\begin{longtable}{@{}>{\footnotesize\raggedright\arraybackslash\ttfamily}p{3.00cm}>{\footnotesize\raggedright\arraybackslash}p{1.50cm}>{\footnotesize\raggedright\arraybackslash}p{2.20cm}>{\footnotesize\raggedright\arraybackslash}p{1.50cm}>{\footnotesize\raggedright\arraybackslash}p{6.40cm}@{}}
\toprule
{\footnotesize\textbf{Name}} & {\footnotesize\textbf{Size}} & {\footnotesize\textbf{Type}} & {\footnotesize\textbf{Units}} & {\footnotesize\textbf{Description}} \\
\midrule
\endfirsthead
\toprule
{\footnotesize\textbf{Name}} & {\footnotesize\textbf{Size}} & {\footnotesize\textbf{Type}} & {\footnotesize\textbf{Units}} & {\footnotesize\textbf{Description}} \\
\midrule
\endhead
\midrule
\endfoot
\bottomrule
\endlastfoot
files & 1xF & string & - & One or more THR1B files of ONE satellite; .gz accepted. \\
\end{longtable}

\subsubsection*{NAME-VALUE ARGUMENTS}
\begin{longtable}{@{}>{\footnotesize\raggedright\arraybackslash\ttfamily}p{3.00cm}>{\footnotesize\raggedright\arraybackslash}p{1.50cm}>{\footnotesize\raggedright\arraybackslash}p{2.20cm}>{\footnotesize\raggedright\arraybackslash}p{1.50cm}>{\footnotesize\raggedright\arraybackslash}p{6.40cm}@{}}
\toprule
{\footnotesize\textbf{Name}} & {\footnotesize\textbf{Size}} & {\footnotesize\textbf{Type}} & {\footnotesize\textbf{Units}} & {\footnotesize\textbf{Description}} \\
\midrule
\endfirsthead
\toprule
{\footnotesize\textbf{Name}} & {\footnotesize\textbf{Size}} & {\footnotesize\textbf{Type}} & {\footnotesize\textbf{Units}} & {\footnotesize\textbf{Description}} \\
\midrule
\endhead
\midrule
\endfoot
\bottomrule
\endlastfoot
TimeRange & 1x2 & double\textbar{}\allowbreak{}datetime & - & Forwarded to readL1b. Default []. \\
\end{longtable}

\subsubsection*{OUTPUT ARGUMENTS}
\begin{longtable}{@{}>{\footnotesize\raggedright\arraybackslash\ttfamily}p{3.00cm}>{\footnotesize\raggedright\arraybackslash}p{1.50cm}>{\footnotesize\raggedright\arraybackslash}p{2.20cm}>{\footnotesize\raggedright\arraybackslash}p{1.50cm}>{\footnotesize\raggedright\arraybackslash}p{6.40cm}@{}}
\toprule
{\footnotesize\textbf{Name}} & {\footnotesize\textbf{Size}} & {\footnotesize\textbf{Type}} & {\footnotesize\textbf{Units}} & {\footnotesize\textbf{Description}} \\
\midrule
\endfirsthead
\toprule
{\footnotesize\textbf{Name}} & {\footnotesize\textbf{Size}} & {\footnotesize\textbf{Type}} & {\footnotesize\textbf{Units}} & {\footnotesize\textbf{Description}} \\
\midrule
\endhead
\midrule
\endfoot
\bottomrule
\endlastfoot
events & Nx8 & table & - & timeSeriesPreparation.qc.eventTable schema; N records with a firing. \\
aux & 1x1 & struct & - & .epoch (GPS epoch struct of ALL K records), .thruster (1xT string, header names without "on\_time\_"), .onTime\_s (TxK), .fired (TxK logical), .count (TxK cumulative thrust counts, thrust\_count\_* in header order), .accumDur\_s (struct .attCtrl 1xK .orbCtrl 1xK, cumulative), .numRecordsNoFiring, .qualflgBits, .product, .satellite, .files, .numRecords, .sampling\_s, .obsNames, .units, .provenance. \\
\end{longtable}

\subsubsection*{ERRORS}
\begin{description}[leftmargin=1.2em,style=nextline,itemsep=0pt]
\item[\texttt{NonGravForces:\allowbreak{}readTHR1B:\allowbreak{}dataProduct}] a file is not a THR1B product
\item[\texttt{NonGravForces:\allowbreak{}readTHR1B:\allowbreak{}dataSatellite}] the files mix satellites
\item[\texttt{NonGravForces:\allowbreak{}readTHR1B:\allowbreak{}dataEmpty}] no record survives TimeRange
\item[\texttt{NonGravForces:\allowbreak{}readTHR1B:\allowbreak{}dataColumn}] no on\_time\_* variable in the header
\item[\texttt{NonGravForces:\allowbreak{}readTHR1B:\allowbreak{}dataTimeRef}] time\_ref is not 'G' (GPS) for every record
\item[\texttt{NonGravForces:\allowbreak{}readTHR1B:\allowbreak{}depMissing}] timeSeriesPreparation (the event table's owner) not on the path
\item[\texttt{gracefoReading:\allowbreak{}*, referenceSystem:\allowbreak{}*}] propagated unwrapped
\end{description}

\subsection{\texttt{NonGravForces.io.readMAS1B}}
\label{api:io.readMAS1B}
\noindent MAS1B reader: cold-gas TANK masses (kg) with GPS epochs. NOT the spacecraft mass.

\begin{lstlisting}
[epoch, tank, aux] = NonGravForces.io.readMAS1B(file)
[epoch, tank, aux] = NonGravForces.io.readMAS1B(files, 'TimeRange', tr)
\end{lstlisting}

\subsubsection*{INPUT ARGUMENTS}
\begin{longtable}{@{}>{\footnotesize\raggedright\arraybackslash\ttfamily}p{3.00cm}>{\footnotesize\raggedright\arraybackslash}p{1.50cm}>{\footnotesize\raggedright\arraybackslash}p{2.20cm}>{\footnotesize\raggedright\arraybackslash}p{1.50cm}>{\footnotesize\raggedright\arraybackslash}p{6.40cm}@{}}
\toprule
{\footnotesize\textbf{Name}} & {\footnotesize\textbf{Size}} & {\footnotesize\textbf{Type}} & {\footnotesize\textbf{Units}} & {\footnotesize\textbf{Description}} \\
\midrule
\endfirsthead
\toprule
{\footnotesize\textbf{Name}} & {\footnotesize\textbf{Size}} & {\footnotesize\textbf{Type}} & {\footnotesize\textbf{Units}} & {\footnotesize\textbf{Description}} \\
\midrule
\endhead
\midrule
\endfoot
\bottomrule
\endlastfoot
files & 1xF & string & - & One or more MAS1B files of ONE satellite; .gz accepted. \\
\end{longtable}

\subsubsection*{NAME-VALUE ARGUMENTS}
\begin{longtable}{@{}>{\footnotesize\raggedright\arraybackslash\ttfamily}p{3.00cm}>{\footnotesize\raggedright\arraybackslash}p{1.50cm}>{\footnotesize\raggedright\arraybackslash}p{2.20cm}>{\footnotesize\raggedright\arraybackslash}p{1.50cm}>{\footnotesize\raggedright\arraybackslash}p{6.40cm}@{}}
\toprule
{\footnotesize\textbf{Name}} & {\footnotesize\textbf{Size}} & {\footnotesize\textbf{Type}} & {\footnotesize\textbf{Units}} & {\footnotesize\textbf{Description}} \\
\midrule
\endfirsthead
\toprule
{\footnotesize\textbf{Name}} & {\footnotesize\textbf{Size}} & {\footnotesize\textbf{Type}} & {\footnotesize\textbf{Units}} & {\footnotesize\textbf{Description}} \\
\midrule
\endhead
\midrule
\endfoot
\bottomrule
\endlastfoot
TimeRange & 1x2 & double\textbar{}\allowbreak{}datetime & - & Forwarded to readL1b. Default []. \\
\end{longtable}

\subsubsection*{OUTPUT ARGUMENTS}
\begin{longtable}{@{}>{\footnotesize\raggedright\arraybackslash\ttfamily}p{3.00cm}>{\footnotesize\raggedright\arraybackslash}p{1.50cm}>{\footnotesize\raggedright\arraybackslash}p{2.20cm}>{\footnotesize\raggedright\arraybackslash}p{1.50cm}>{\footnotesize\raggedright\arraybackslash}p{6.40cm}@{}}
\toprule
{\footnotesize\textbf{Name}} & {\footnotesize\textbf{Size}} & {\footnotesize\textbf{Type}} & {\footnotesize\textbf{Units}} & {\footnotesize\textbf{Description}} \\
\midrule
\endfirsthead
\toprule
{\footnotesize\textbf{Name}} & {\footnotesize\textbf{Size}} & {\footnotesize\textbf{Type}} & {\footnotesize\textbf{Units}} & {\footnotesize\textbf{Description}} \\
\midrule
\endhead
\midrule
\endfoot
\bottomrule
\endlastfoot
epoch & 1x1 & struct & - & referenceSystem epoch struct, K entries, "GPS". \\
tank & 1x1 & struct & kg & .total\_kg (1xK: tank1 + tank2 where both are present, else mass\_tnk), .tank\_kg (2xK: gas\_mass\_tnk1, gas\_mass\_tnk2), .thruster\_kg (2xK: gas\_mass\_thr1, \_thr2), .tankTotal\_kg (1xK mass\_tnk), .thrusterTotal\_kg (1xK mass\_thr), .sigma\_kg (1xK mass\_tnk\_err), .present (struct of logicals per column: any record has it). NaN where the flag says absent. \\
aux & 1x1 & struct & - & .prodFlagBits, .qualflgBits, .product, .satellite, .files, .numRecords, .sampling\_s, .obsNames, .units, .provenance. \\
\end{longtable}

\subsubsection*{ERRORS}
\begin{description}[leftmargin=1.2em,style=nextline,itemsep=0pt]
\item[\texttt{NonGravForces:\allowbreak{}readMAS1B:\allowbreak{}dataProduct}] a file is not a MAS1B product
\item[\texttt{NonGravForces:\allowbreak{}readMAS1B:\allowbreak{}dataSatellite}] the files mix satellites
\item[\texttt{NonGravForces:\allowbreak{}readMAS1B:\allowbreak{}dataEmpty}] no record survives TimeRange
\item[\texttt{NonGravForces:\allowbreak{}readMAS1B:\allowbreak{}dataColumn}] gas\_mass\_tnk1 or gas\_mass\_tnk2 missing from the header
\item[\texttt{NonGravForces:\allowbreak{}readMAS1B:\allowbreak{}dataMissing}] no record has a usable total (both tank columns and mass\_tnk absent)
\item[\texttt{NonGravForces:\allowbreak{}readMAS1B:\allowbreak{}dataRange}] a total is \textgreater{}= 100 kg or \textless{}= 0 kg
\item[\texttt{NonGravForces:\allowbreak{}readMAS1B:\allowbreak{}dataTimeRef}] time\_ref is not 'G' for every record
\item[\texttt{gracefoReading:\allowbreak{}*}] propagated unwrapped
\end{description}

\subsection{\texttt{NonGravForces.io.readMAG1B}}
\label{api:io.readMAG1B}
\noindent MAG1B reader: measured magnetic field and magnetorquer currents in the SRF.

\begin{lstlisting}
[epoch, field, aux] = NonGravForces.io.readMAG1B(file)
[epoch, field, aux] = NonGravForces.io.readMAG1B(files, 'TimeRange', tr)
\end{lstlisting}

\subsubsection*{INPUT ARGUMENTS}
\begin{longtable}{@{}>{\footnotesize\raggedright\arraybackslash\ttfamily}p{3.00cm}>{\footnotesize\raggedright\arraybackslash}p{1.50cm}>{\footnotesize\raggedright\arraybackslash}p{2.20cm}>{\footnotesize\raggedright\arraybackslash}p{1.50cm}>{\footnotesize\raggedright\arraybackslash}p{6.40cm}@{}}
\toprule
{\footnotesize\textbf{Name}} & {\footnotesize\textbf{Size}} & {\footnotesize\textbf{Type}} & {\footnotesize\textbf{Units}} & {\footnotesize\textbf{Description}} \\
\midrule
\endfirsthead
\toprule
{\footnotesize\textbf{Name}} & {\footnotesize\textbf{Size}} & {\footnotesize\textbf{Type}} & {\footnotesize\textbf{Units}} & {\footnotesize\textbf{Description}} \\
\midrule
\endhead
\midrule
\endfoot
\bottomrule
\endlastfoot
files & 1xF & string & - & One or more MAG1B files of ONE satellite; .gz accepted. \\
\end{longtable}

\subsubsection*{NAME-VALUE ARGUMENTS}
\begin{longtable}{@{}>{\footnotesize\raggedright\arraybackslash\ttfamily}p{3.00cm}>{\footnotesize\raggedright\arraybackslash}p{1.50cm}>{\footnotesize\raggedright\arraybackslash}p{2.20cm}>{\footnotesize\raggedright\arraybackslash}p{1.50cm}>{\footnotesize\raggedright\arraybackslash}p{6.40cm}@{}}
\toprule
{\footnotesize\textbf{Name}} & {\footnotesize\textbf{Size}} & {\footnotesize\textbf{Type}} & {\footnotesize\textbf{Units}} & {\footnotesize\textbf{Description}} \\
\midrule
\endfirsthead
\toprule
{\footnotesize\textbf{Name}} & {\footnotesize\textbf{Size}} & {\footnotesize\textbf{Type}} & {\footnotesize\textbf{Units}} & {\footnotesize\textbf{Description}} \\
\midrule
\endhead
\midrule
\endfoot
\bottomrule
\endlastfoot
TimeRange & 1x2 & double\textbar{}\allowbreak{}datetime & - & Forwarded to readL1b. Default []. \\
\end{longtable}

\subsubsection*{OUTPUT ARGUMENTS}
\begin{longtable}{@{}>{\footnotesize\raggedright\arraybackslash\ttfamily}p{3.00cm}>{\footnotesize\raggedright\arraybackslash}p{1.50cm}>{\footnotesize\raggedright\arraybackslash}p{2.20cm}>{\footnotesize\raggedright\arraybackslash}p{1.50cm}>{\footnotesize\raggedright\arraybackslash}p{6.40cm}@{}}
\toprule
{\footnotesize\textbf{Name}} & {\footnotesize\textbf{Size}} & {\footnotesize\textbf{Type}} & {\footnotesize\textbf{Units}} & {\footnotesize\textbf{Description}} \\
\midrule
\endfirsthead
\toprule
{\footnotesize\textbf{Name}} & {\footnotesize\textbf{Size}} & {\footnotesize\textbf{Type}} & {\footnotesize\textbf{Units}} & {\footnotesize\textbf{Description}} \\
\midrule
\endhead
\midrule
\endfoot
\bottomrule
\endlastfoot
epoch & 1x1 & struct & - & referenceSystem epoch struct, K entries, "GPS". \\
field & 1x1 & struct & - & .B\_srf\_T (3xK, TESLA -- the file is microtesla and this converts, because the toolbox is SI throughout) .torquerNet\_A (3xK, AMPERE, the A + B sum per axis) .torquerBranch\_A (6xK, ampere, rows 1A 2A 3A 1B 2B 3B, as read) .magnitude\_T (1xK, \textbar{}B\textbar{}) \\
aux & 1x1 & struct & - & .epoch (the epoch struct again), .calibration (.magField 3xK, .torquer 1xK, unapplied; NaN when the columns are absent), .qualityBits (8xK logical, bit 0 first), .qualflgBits (Kx8, as the other readers), .torquerBranchNames (1x6 string), .product, .satellite, .files, .numRecords, .sampling\_s, .obsNames, .units, .provenance. \\
\end{longtable}

\subsubsection*{ERRORS}
\begin{description}[leftmargin=1.2em,style=nextline,itemsep=0pt]
\item[\texttt{NonGravForces:\allowbreak{}readMAG1B:\allowbreak{}dataProduct}] a file is not a MAG1B product
\item[\texttt{NonGravForces:\allowbreak{}readMAG1B:\allowbreak{}dataSatellite}] the files mix satellites
\item[\texttt{NonGravForces:\allowbreak{}readMAG1B:\allowbreak{}dataEmpty}] no record survives TimeRange
\item[\texttt{NonGravForces:\allowbreak{}readMAG1B:\allowbreak{}dataColumn}] a MfvX/Y/Z\_RAW or torque column is missing
\item[\texttt{NonGravForces:\allowbreak{}readMAG1B:\allowbreak{}dataTimeRef}] time\_ref is not 'G' for every record
\item[\texttt{NonGravForces:\allowbreak{}readMAG1B:\allowbreak{}dataRange}] \textbar{}B\textbar{} above 100 microtesla or a branch current above 1 A -- an Earth field at 490 km reaches about 50 microtesla and the measured currents reach 55 mA, so either bound means the columns were misread
\item[\texttt{gracefoReading:\allowbreak{}*}] propagated unwrapped
\end{description}

\subsection{\texttt{NonGravForces.io.massModel}}
\label{api:io.massModel}
\noindent Spacecraft mass model: handbook constant, or MAS1B-driven variation about it.

\begin{lstlisting}
mass = NonGravForces.io.massModel("C")
mass = NonGravForces.io.massModel("C", 'Mode', "mas1b", 'Epoch', epMas, 'GasMass_kg', tank.total_kg)
m_kg = mass.at(epoch)                        % 1xK
\end{lstlisting}

\subsubsection*{INPUT ARGUMENTS}
\begin{longtable}{@{}>{\footnotesize\raggedright\arraybackslash\ttfamily}p{3.00cm}>{\footnotesize\raggedright\arraybackslash}p{1.50cm}>{\footnotesize\raggedright\arraybackslash}p{2.20cm}>{\footnotesize\raggedright\arraybackslash}p{1.50cm}>{\footnotesize\raggedright\arraybackslash}p{6.40cm}@{}}
\toprule
{\footnotesize\textbf{Name}} & {\footnotesize\textbf{Size}} & {\footnotesize\textbf{Type}} & {\footnotesize\textbf{Units}} & {\footnotesize\textbf{Description}} \\
\midrule
\endfirsthead
\toprule
{\footnotesize\textbf{Name}} & {\footnotesize\textbf{Size}} & {\footnotesize\textbf{Type}} & {\footnotesize\textbf{Units}} & {\footnotesize\textbf{Description}} \\
\midrule
\endhead
\midrule
\endfoot
\bottomrule
\endlastfoot
satellite & 1x1 & string & - & "C" or "D". \\
\end{longtable}

\subsubsection*{NAME-VALUE ARGUMENTS}
\begin{longtable}{@{}>{\footnotesize\raggedright\arraybackslash\ttfamily}p{3.00cm}>{\footnotesize\raggedright\arraybackslash}p{1.50cm}>{\footnotesize\raggedright\arraybackslash}p{2.20cm}>{\footnotesize\raggedright\arraybackslash}p{1.50cm}>{\footnotesize\raggedright\arraybackslash}p{6.40cm}@{}}
\toprule
{\footnotesize\textbf{Name}} & {\footnotesize\textbf{Size}} & {\footnotesize\textbf{Type}} & {\footnotesize\textbf{Units}} & {\footnotesize\textbf{Description}} \\
\midrule
\endfirsthead
\toprule
{\footnotesize\textbf{Name}} & {\footnotesize\textbf{Size}} & {\footnotesize\textbf{Type}} & {\footnotesize\textbf{Units}} & {\footnotesize\textbf{Description}} \\
\midrule
\endhead
\midrule
\endfoot
\bottomrule
\endlastfoot
Mode & 1x1 & string & - & "constant" (default) or "mas1b". \\
Epoch & 1x1 & struct & - & referenceSystem epoch struct of the MAS1B samples (mas1b only; strictly increasing). \\
GasMass\_\allowbreak{}kg & 1xN & double & kg & Tank total at those epochs (mas1b only; finite; readMAS1B's tank.total\_kg). \\
Anchor\_\allowbreak{}kg & 1x1 & double & kg & Mass at the anchor (constant: the value; mas1b: the mass at the first sample). Default: the handbook constant. \\
\end{longtable}

\subsubsection*{OUTPUT ARGUMENTS}
\begin{longtable}{@{}>{\footnotesize\raggedright\arraybackslash\ttfamily}p{3.00cm}>{\footnotesize\raggedright\arraybackslash}p{1.50cm}>{\footnotesize\raggedright\arraybackslash}p{2.20cm}>{\footnotesize\raggedright\arraybackslash}p{1.50cm}>{\footnotesize\raggedright\arraybackslash}p{6.40cm}@{}}
\toprule
{\footnotesize\textbf{Name}} & {\footnotesize\textbf{Size}} & {\footnotesize\textbf{Type}} & {\footnotesize\textbf{Units}} & {\footnotesize\textbf{Description}} \\
\midrule
\endfirsthead
\toprule
{\footnotesize\textbf{Name}} & {\footnotesize\textbf{Size}} & {\footnotesize\textbf{Type}} & {\footnotesize\textbf{Units}} & {\footnotesize\textbf{Description}} \\
\midrule
\endhead
\midrule
\endfoot
\bottomrule
\endlastfoot
mass & 1x1 & struct & - & .kind ("constant"\textbar{}"mas1b"), .satellite, .value\_kg (the anchor), .source ("JPL D-56935 handbook" \textbar{} "caller" \textbar{} "MAS1B about \textless{}source\textgreater{}"), .span (1x2 decimal MJD in .scale, [-Inf Inf] for constant), .scale (epoch scale of the samples, "" for constant), .at (function handle: epoch struct -\textgreater{} 1xK kg), .provenance. \\
\end{longtable}

\subsubsection*{ERRORS}
\begin{description}[leftmargin=1.2em,style=nextline,itemsep=0pt]
\item[\texttt{NonGravForces:\allowbreak{}massModel:\allowbreak{}inputMas1b}] Mode "mas1b" without Epoch/GasMass\_kg, mismatched lengths, non-finite gas mass, or non-increasing epochs
\item[\texttt{NonGravForces:\allowbreak{}massModel:\allowbreak{}inputAnchor}] Anchor\_kg not finite or \textless{}= 500 kg
\item[\texttt{NonGravForces:\allowbreak{}massModel:\allowbreak{}inputEpoch}] (from .at) query epoch further than two sample intervals beyond the series, or not an epoch struct
\item[\texttt{NonGravForces:\allowbreak{}massModel:\allowbreak{}numericRange}] (from .at) a resulting mass \textless{}= 500 kg
\item[\texttt{referenceSystem:\allowbreak{}*}] propagated unwrapped (scale conversion)
\end{description}

\subsection{\texttt{NonGravForces.io.spaceWeather}}
\label{api:io.spaceWeather}
\noindent Space-weather provider: CelesTrak SW-All, SET SOLFSMY+DTCFILE, Kyoto Dst or CLS F30, with query(epoch).

\begin{lstlisting}
sw = NonGravForces.io.spaceWeather(swAllCsv)                       % family celestrak
sw = NonGravForces.io.spaceWeather([solfsmyTxt dtcfileTxt])        % family jb2008
sw = NonGravForces.io.spaceWeather([dst0712 dst0801 ...])          % family dst (Kyoto WDC months)
sw = NonGravForces.io.spaceWeather(clsRadioFluxTxt)                % family f30 (CLS)
d  = sw.query(epoch)                                               % drivers at K epochs
\end{lstlisting}

\subsubsection*{INPUT ARGUMENTS}
\begin{longtable}{@{}>{\footnotesize\raggedright\arraybackslash\ttfamily}p{3.00cm}>{\footnotesize\raggedright\arraybackslash}p{1.50cm}>{\footnotesize\raggedright\arraybackslash}p{2.20cm}>{\footnotesize\raggedright\arraybackslash}p{1.50cm}>{\footnotesize\raggedright\arraybackslash}p{6.40cm}@{}}
\toprule
{\footnotesize\textbf{Name}} & {\footnotesize\textbf{Size}} & {\footnotesize\textbf{Type}} & {\footnotesize\textbf{Units}} & {\footnotesize\textbf{Description}} \\
\midrule
\endfirsthead
\toprule
{\footnotesize\textbf{Name}} & {\footnotesize\textbf{Size}} & {\footnotesize\textbf{Type}} & {\footnotesize\textbf{Units}} & {\footnotesize\textbf{Description}} \\
\midrule
\endhead
\midrule
\endfoot
\bottomrule
\endlastfoot
files & 1xF & string & - & One CelesTrak SW-All CSV; or the SET SOLFSMY and DTCFILE text files (either order); or one or more Kyoto WDC Dst month files (dst\textless{}yymm\textgreater{}.for.request, any order, no day twice); or one CLS radio\_flux\_*\_observation.txt. \\
\end{longtable}

\subsubsection*{NAME-VALUE ARGUMENTS}
\begin{longtable}{@{}>{\footnotesize\raggedright\arraybackslash\ttfamily}p{3.00cm}>{\footnotesize\raggedright\arraybackslash}p{1.50cm}>{\footnotesize\raggedright\arraybackslash}p{2.20cm}>{\footnotesize\raggedright\arraybackslash}p{1.50cm}>{\footnotesize\raggedright\arraybackslash}p{6.40cm}@{}}
\toprule
{\footnotesize\textbf{Name}} & {\footnotesize\textbf{Size}} & {\footnotesize\textbf{Type}} & {\footnotesize\textbf{Units}} & {\footnotesize\textbf{Description}} \\
\midrule
\endfirsthead
\toprule
{\footnotesize\textbf{Name}} & {\footnotesize\textbf{Size}} & {\footnotesize\textbf{Type}} & {\footnotesize\textbf{Units}} & {\footnotesize\textbf{Description}} \\
\midrule
\endhead
\midrule
\endfoot
\bottomrule
\endlastfoot
Format & 1x1 & string & - & "auto" (default; identified from content), "celestrak", "jb2008", "dst" or "f30" (must then match). \\
\end{longtable}

\subsubsection*{OUTPUT ARGUMENTS}
\begin{longtable}{@{}>{\footnotesize\raggedright\arraybackslash\ttfamily}p{3.00cm}>{\footnotesize\raggedright\arraybackslash}p{1.50cm}>{\footnotesize\raggedright\arraybackslash}p{2.20cm}>{\footnotesize\raggedright\arraybackslash}p{1.50cm}>{\footnotesize\raggedright\arraybackslash}p{6.40cm}@{}}
\toprule
{\footnotesize\textbf{Name}} & {\footnotesize\textbf{Size}} & {\footnotesize\textbf{Type}} & {\footnotesize\textbf{Units}} & {\footnotesize\textbf{Description}} \\
\midrule
\endfirsthead
\toprule
{\footnotesize\textbf{Name}} & {\footnotesize\textbf{Size}} & {\footnotesize\textbf{Type}} & {\footnotesize\textbf{Units}} & {\footnotesize\textbf{Description}} \\
\midrule
\endhead
\midrule
\endfoot
\bottomrule
\endlastfoot
sw & 1x1 & struct & - & .format ("celestrak"\textbar{}"jb2008"\textbar{}"dst"\textbar{}"f30"), .mjd (Nx1, 0h UTC of each day), .span (1x2 MJD), .index (struct of Nx1 / Nx8 / Nx24 arrays, names below), .files, .query (handle), .provenance. celestrak .index: f107Observed f107Adjusted f107ObservedCentre81 f107ObservedLast81 apDaily ap3h (Nx8) kp3h (Nx8, /10) kpSum sunspotNumber dataType (string) jb2008 .index: f10 f10b s10 s10b m10 m10b y10 y10b ssrc (string) dtc (Nx24, NaN where DTCFILE has no day) dst .index: dst (Nx24 nT, base value applied, NaN for the WDC missing value 9999), dailyMean (Nx1 nT, as printed), version (Nx1 string) f30 .index: for each band b of f30 f15 f107 f8 f32 (CLS f30, f15, f10.7, f8, f3.2): b, bCorrected (CLS \_c, flare-corrected), bPrecision (\_p), bFlag (\_f); Nx1 sfu, NaN where CLS writes -1 f30 also sets .adjusted (true: "adjusted to 1 AU", false: "not adjusted"), read from the file header. \\
\end{longtable}

\subsubsection*{ERRORS}
\begin{description}[leftmargin=1.2em,style=nextline,itemsep=0pt]
\item[\texttt{NonGravForces:\allowbreak{}spaceWeather:\allowbreak{}inputFiles}] wrong number of files for the family
\item[\texttt{NonGravForces:\allowbreak{}spaceWeather:\allowbreak{}inputFormat}] a file cannot be opened, its content matches no known family, or not the requested one
\item[\texttt{NonGravForces:\allowbreak{}spaceWeather:\allowbreak{}dataColumn}] a required column is missing
\item[\texttt{NonGravForces:\allowbreak{}spaceWeather:\allowbreak{}dataMalformed}] a record does not parse (bad date, DTC row without 24 values, JD and YYYY/DDD disagree, a WDC Dst record not 120 columns, a CLS cont\_day that disagrees with its date, a day twice)
\item[\texttt{NonGravForces:\allowbreak{}spaceWeather:\allowbreak{}inputEpoch}] (query) not an epoch struct
\item[\texttt{NonGravForces:\allowbreak{}spaceWeather:\allowbreak{}dataRange}] (query) a needed day is not in the table
\item[\texttt{NonGravForces:\allowbreak{}spaceWeather:\allowbreak{}dataMissing}] (query) a needed value is empty
\item[\texttt{MATLAB:\allowbreak{}validators:\allowbreak{}mustBeFile}] a file does not exist
\item[\texttt{referenceSystem:\allowbreak{}*}] propagated unwrapped
\end{description}

\subsection{\texttt{NonGravForces.io.importCERES}}
\label{api:io.importCERES}
\noindent CERES EBAF TOA monthly fluxes (obs4MIPs netCDF) as the grids env.earthRadiance("ceres") evaluates.

\begin{lstlisting}
cer = NonGravForces.io.importCERES(files)
cer = NonGravForces.io.importCERES(files, 'Period', [datetime(2020,1,1) datetime(2020,1,31)])
[cer, aux] = NonGravForces.io.importCERES(___)
\end{lstlisting}

\subsubsection*{INPUT ARGUMENTS}
\begin{longtable}{@{}>{\footnotesize\raggedright\arraybackslash\ttfamily}p{3.00cm}>{\footnotesize\raggedright\arraybackslash}p{1.50cm}>{\footnotesize\raggedright\arraybackslash}p{2.20cm}>{\footnotesize\raggedright\arraybackslash}p{1.50cm}>{\footnotesize\raggedright\arraybackslash}p{6.40cm}@{}}
\toprule
{\footnotesize\textbf{Name}} & {\footnotesize\textbf{Size}} & {\footnotesize\textbf{Type}} & {\footnotesize\textbf{Units}} & {\footnotesize\textbf{Description}} \\
\midrule
\endfirsthead
\toprule
{\footnotesize\textbf{Name}} & {\footnotesize\textbf{Size}} & {\footnotesize\textbf{Type}} & {\footnotesize\textbf{Units}} & {\footnotesize\textbf{Description}} \\
\midrule
\endhead
\midrule
\endfoot
\bottomrule
\endlastfoot
files & 1x3 & string & - & The rsut, rlut and rsdt obs4MIPs netCDF files, any order \\
\end{longtable}

\subsubsection*{NAME-VALUE ARGUMENTS}
\begin{longtable}{@{}>{\footnotesize\raggedright\arraybackslash\ttfamily}p{3.00cm}>{\footnotesize\raggedright\arraybackslash}p{1.50cm}>{\footnotesize\raggedright\arraybackslash}p{2.20cm}>{\footnotesize\raggedright\arraybackslash}p{1.50cm}>{\footnotesize\raggedright\arraybackslash}p{6.40cm}@{}}
\toprule
{\footnotesize\textbf{Name}} & {\footnotesize\textbf{Size}} & {\footnotesize\textbf{Type}} & {\footnotesize\textbf{Units}} & {\footnotesize\textbf{Description}} \\
\midrule
\endfirsthead
\toprule
{\footnotesize\textbf{Name}} & {\footnotesize\textbf{Size}} & {\footnotesize\textbf{Type}} & {\footnotesize\textbf{Units}} & {\footnotesize\textbf{Description}} \\
\midrule
\endhead
\midrule
\endfoot
\bottomrule
\endlastfoot
Period & 1x2 & datetime & UTC & [start end], start \textless{}= end. Default: every month in the files. \\
\end{longtable}

\subsubsection*{OUTPUT ARGUMENTS}
\begin{longtable}{@{}>{\footnotesize\raggedright\arraybackslash\ttfamily}p{3.00cm}>{\footnotesize\raggedright\arraybackslash}p{1.50cm}>{\footnotesize\raggedright\arraybackslash}p{2.20cm}>{\footnotesize\raggedright\arraybackslash}p{1.50cm}>{\footnotesize\raggedright\arraybackslash}p{6.40cm}@{}}
\toprule
{\footnotesize\textbf{Name}} & {\footnotesize\textbf{Size}} & {\footnotesize\textbf{Type}} & {\footnotesize\textbf{Units}} & {\footnotesize\textbf{Description}} \\
\midrule
\endfirsthead
\toprule
{\footnotesize\textbf{Name}} & {\footnotesize\textbf{Size}} & {\footnotesize\textbf{Type}} & {\footnotesize\textbf{Units}} & {\footnotesize\textbf{Description}} \\
\midrule
\endhead
\midrule
\endfoot
\bottomrule
\endlastfoot
cer & 1x1 & struct & - & .product (source\_id), .version, .doi, .license, .time (1xT datetime, UTC, the files' mid-month nodes), .mjd (1xT, the same nodes as MJD UTC), .lat\_deg (1xNlat centres, ascending), .lon\_deg (1xNlon centres, ascending, east, a full circle), .latitude "geodetic", .referenceHeight\_m 20e3, .rsut, .rlut, .rsdt (NlatxNlonxT single, W m\textasciicircum{}-2) \\
aux & 1x1 & struct & - & .files (1x3, in the order rsut rlut rsdt), .trackingId (1x3), .timeIndex (1xT indices into the files' time axis), .monthsInFiles, .provenance \\
\end{longtable}

\subsubsection*{ERRORS}
\begin{description}[leftmargin=1.2em,style=nextline,itemsep=0pt]
\item[\texttt{NonGravForces:\allowbreak{}importCERES:\allowbreak{}inputFiles}] not three files, or not one each of rsut, rlut, rsdt
\item[\texttt{NonGravForces:\allowbreak{}importCERES:\allowbreak{}dataMissing}] a file does not exist, or a loaded value is a fill value, NaN or negative
\item[\texttt{NonGravForces:\allowbreak{}importCERES:\allowbreak{}dataFormat}] not an obs4MIPs CERES EBAF file (a global attribute missing), units not W m-2, dimensions not lon/lat/time, time units not "days since", non-uniform grid, or the files disagree
\item[\texttt{NonGravForces:\allowbreak{}importCERES:\allowbreak{}inputPeriod}] Period malformed or not bracketed by the files' months
\end{description}

\subsection{\texttt{NonGravForces.io.importGroopsGrid}}
\label{api:io.importGroopsGrid}
\noindent Read a GROOPS binary gridded-data file: global grid, ellipsoid, solid angles and values.

\begin{lstlisting}
grid = NonGravForces.io.importGroopsGrid(file)
grid = NonGravForces.io.importGroopsGrid(file, 'TimeRange', [mjd0 mjd1])
[grid, aux] = NonGravForces.io.importGroopsGrid(___)
\end{lstlisting}

\subsubsection*{INPUT ARGUMENTS}
\begin{longtable}{@{}>{\footnotesize\raggedright\arraybackslash\ttfamily}p{3.00cm}>{\footnotesize\raggedright\arraybackslash}p{1.50cm}>{\footnotesize\raggedright\arraybackslash}p{2.20cm}>{\footnotesize\raggedright\arraybackslash}p{1.50cm}>{\footnotesize\raggedright\arraybackslash}p{6.40cm}@{}}
\toprule
{\footnotesize\textbf{Name}} & {\footnotesize\textbf{Size}} & {\footnotesize\textbf{Type}} & {\footnotesize\textbf{Units}} & {\footnotesize\textbf{Description}} \\
\midrule
\endfirsthead
\toprule
{\footnotesize\textbf{Name}} & {\footnotesize\textbf{Size}} & {\footnotesize\textbf{Type}} & {\footnotesize\textbf{Units}} & {\footnotesize\textbf{Description}} \\
\midrule
\endhead
\midrule
\endfoot
\bottomrule
\endlastfoot
file & 1x1 & string & - & Path to the GROOPS binary file \\
\end{longtable}

\subsubsection*{NAME-VALUE ARGUMENTS}
\begin{longtable}{@{}>{\footnotesize\raggedright\arraybackslash\ttfamily}p{3.00cm}>{\footnotesize\raggedright\arraybackslash}p{1.50cm}>{\footnotesize\raggedright\arraybackslash}p{2.20cm}>{\footnotesize\raggedright\arraybackslash}p{1.50cm}>{\footnotesize\raggedright\arraybackslash}p{6.40cm}@{}}
\toprule
{\footnotesize\textbf{Name}} & {\footnotesize\textbf{Size}} & {\footnotesize\textbf{Type}} & {\footnotesize\textbf{Units}} & {\footnotesize\textbf{Description}} \\
\midrule
\endfirsthead
\toprule
{\footnotesize\textbf{Name}} & {\footnotesize\textbf{Size}} & {\footnotesize\textbf{Type}} & {\footnotesize\textbf{Units}} & {\footnotesize\textbf{Description}} \\
\midrule
\endhead
\midrule
\endfoot
\bottomrule
\endlastfoot
TimeRange & 1x2 & double & MJD & Read only the grids whose interval intersects [mjd0 mjd1]. Time series only; default empty, which reads every grid. \\
\end{longtable}

\subsubsection*{OUTPUT ARGUMENTS}
\begin{longtable}{@{}>{\footnotesize\raggedright\arraybackslash\ttfamily}p{3.00cm}>{\footnotesize\raggedright\arraybackslash}p{1.50cm}>{\footnotesize\raggedright\arraybackslash}p{2.20cm}>{\footnotesize\raggedright\arraybackslash}p{1.50cm}>{\footnotesize\raggedright\arraybackslash}p{6.40cm}@{}}
\toprule
{\footnotesize\textbf{Name}} & {\footnotesize\textbf{Size}} & {\footnotesize\textbf{Type}} & {\footnotesize\textbf{Units}} & {\footnotesize\textbf{Description}} \\
\midrule
\endfirsthead
\toprule
{\footnotesize\textbf{Name}} & {\footnotesize\textbf{Size}} & {\footnotesize\textbf{Type}} & {\footnotesize\textbf{Units}} & {\footnotesize\textbf{Description}} \\
\midrule
\endhead
\midrule
\endfoot
\bottomrule
\endlastfoot
grid & 1x1 & struct & - & .type, .version (string); .lon\_deg (1xNlon), .lat\_deg (1xNlat, GEOCENTRIC, decreasing -- see below); .radius\_m (1xNlat); .area\_sr (NlatxNlon); .values (NlatxNlonxT; for a griddedData T is its C value columns); .mjd (1x(T+1) node epochs BOUNDING THE GRIDS RETURNED, so after a 'TimeRange' seek it is the subrange and not the file's whole node list; [] for griddedData); .isTimeSeries (logical) \\
aux & 1x1 & struct & - & .file, .bytes, .ellipsoid (.a, .b), .pointCount, .gridCount, .gridsRead (1xT index of the grids actually read, into the FILE), .mjdFile (1x2 first and last node of the file, whatever was read), .areaSum\_sr (must be 4*pi), .provenance \\
\end{longtable}

\subsubsection*{ERRORS}
\begin{description}[leftmargin=1.2em,style=nextline,itemsep=0pt]
\item[\texttt{NonGravForces:\allowbreak{}importGroopsGrid:\allowbreak{}inputFile}] the file does not exist or cannot be opened
\item[\texttt{NonGravForces:\allowbreak{}importGroopsGrid:\allowbreak{}dataFormat}] not a GROOPS binary, the solid angles do not match the axes, the payload does not fill the file, or the node, column and grid counts disagree
\item[\texttt{NonGravForces:\allowbreak{}importGroopsGrid:\allowbreak{}unsupportedType}] a GROOPS type this reader does not serve
\item[\texttt{NonGravForces:\allowbreak{}importGroopsGrid:\allowbreak{}inputRange}] 'TimeRange' not two increasing finite MJD, given for a file that carries no time, or covered by no grid of the file
\end{description}

\subsection{\texttt{NonGravForces.io.igrfCoefficients}}
\label{api:io.igrfCoefficients}
\noindent IGRF-14 Schmidt semi-normalised spherical-harmonic coefficients, with the model's own time rule.

\begin{lstlisting}
coeff = NonGravForces.io.igrfCoefficients(file)
coeff = NonGravForces.io.igrfCoefficients(file, 'Epoch', ep)
coeff = NonGravForces.io.igrfCoefficients(file, 'Epoch', 2020.5)
\end{lstlisting}

\subsubsection*{INPUT ARGUMENTS}
\begin{longtable}{@{}>{\footnotesize\raggedright\arraybackslash\ttfamily}p{3.00cm}>{\footnotesize\raggedright\arraybackslash}p{1.50cm}>{\footnotesize\raggedright\arraybackslash}p{2.20cm}>{\footnotesize\raggedright\arraybackslash}p{1.50cm}>{\footnotesize\raggedright\arraybackslash}p{6.40cm}@{}}
\toprule
{\footnotesize\textbf{Name}} & {\footnotesize\textbf{Size}} & {\footnotesize\textbf{Type}} & {\footnotesize\textbf{Units}} & {\footnotesize\textbf{Description}} \\
\midrule
\endfirsthead
\toprule
{\footnotesize\textbf{Name}} & {\footnotesize\textbf{Size}} & {\footnotesize\textbf{Type}} & {\footnotesize\textbf{Units}} & {\footnotesize\textbf{Description}} \\
\midrule
\endhead
\midrule
\endfoot
\bottomrule
\endlastfoot
file & 1x1 & string & - & Path to igrf14coeffs.txt (IAGA V-MOD format). \\
\end{longtable}

\subsubsection*{NAME-VALUE ARGUMENTS}
\begin{longtable}{@{}>{\footnotesize\raggedright\arraybackslash\ttfamily}p{3.00cm}>{\footnotesize\raggedright\arraybackslash}p{1.50cm}>{\footnotesize\raggedright\arraybackslash}p{2.20cm}>{\footnotesize\raggedright\arraybackslash}p{1.50cm}>{\footnotesize\raggedright\arraybackslash}p{6.40cm}@{}}
\toprule
{\footnotesize\textbf{Name}} & {\footnotesize\textbf{Size}} & {\footnotesize\textbf{Type}} & {\footnotesize\textbf{Units}} & {\footnotesize\textbf{Description}} \\
\midrule
\endfirsthead
\toprule
{\footnotesize\textbf{Name}} & {\footnotesize\textbf{Size}} & {\footnotesize\textbf{Type}} & {\footnotesize\textbf{Units}} & {\footnotesize\textbf{Description}} \\
\midrule
\endhead
\midrule
\endfoot
\bottomrule
\endlastfoot
Epoch & 1x1 & struct\textbar{}\allowbreak{}double & - & A referenceSystem epoch struct or a decimal year. Given, coeff.g and coeff.h are the interpolated field at that time; absent, only the tables are returned and coeff.g is empty. \\
\end{longtable}

\subsubsection*{OUTPUT ARGUMENTS}
\begin{longtable}{@{}>{\footnotesize\raggedright\arraybackslash\ttfamily}p{3.00cm}>{\footnotesize\raggedright\arraybackslash}p{1.50cm}>{\footnotesize\raggedright\arraybackslash}p{2.20cm}>{\footnotesize\raggedright\arraybackslash}p{1.50cm}>{\footnotesize\raggedright\arraybackslash}p{6.40cm}@{}}
\toprule
{\footnotesize\textbf{Name}} & {\footnotesize\textbf{Size}} & {\footnotesize\textbf{Type}} & {\footnotesize\textbf{Units}} & {\footnotesize\textbf{Description}} \\
\midrule
\endfirsthead
\toprule
{\footnotesize\textbf{Name}} & {\footnotesize\textbf{Size}} & {\footnotesize\textbf{Type}} & {\footnotesize\textbf{Units}} & {\footnotesize\textbf{Description}} \\
\midrule
\endhead
\midrule
\endfoot
\bottomrule
\endlastfoot
coeff & 1x1 & struct & nT & .gTable, .hTable ((nMax+1)x(nMax+1)xE, one page per epoch) .epochs (1xE decimal years, 1900.0 ... 2025.0) .svG, .svH ((nMax+1)x(nMax+1), nT/year, degree \textless{}= svMax) .nMax (13), .svMax (8) .g, .h ((nMax+1)x(nMax+1) at 'Epoch', else []) .year (decimal year used, else NaN) \\
aux & 1x1 & struct & - & .file, .generation, .numCoefficients, .validFrom, .validTo, .header (the file's comment lines), .provenance \\
\end{longtable}

\subsubsection*{ERRORS}
\begin{description}[leftmargin=1.2em,style=nextline,itemsep=0pt]
\item[\texttt{NonGravForces:\allowbreak{}igrfCoefficients:\allowbreak{}dataMissing}] the file does not exist
\item[\texttt{NonGravForces:\allowbreak{}igrfCoefficients:\allowbreak{}dataFormat}] no "g/h n m ..." header line, its epoch columns do not parse as years, or a row whose field count disagrees with that header
\item[\texttt{NonGravForces:\allowbreak{}igrfCoefficients:\allowbreak{}dataColumn}] a row is neither g nor h, or n and m are not a valid Schmidt pair (m \textless{}= n, n \textgreater{}= 1)
\item[\texttt{NonGravForces:\allowbreak{}igrfCoefficients:\allowbreak{}inputRange}] 'Epoch' before the first epoch or more than five years past the last, or a numeric 'Epoch' that is not one finite decimal year, or an epoch struct holding more than one epoch
\item[\texttt{NonGravForces:\allowbreak{}igrfCoefficients:\allowbreak{}numericNaN}] a coefficient does not parse as a number
\end{description}

\subsection{\texttt{NonGravForces.io.readSwarmDensity}}
\label{api:io.readSwarmDensity}
\noindent Thermosphere density from VirES HAPI CSV: Swarm POD-only, CHAMP/GRACE/GRACE-FO accelerometer.

\begin{lstlisting}
[epoch, lla, rho, aux] = NonGravForces.io.readSwarmDensity(file)
[epoch, lla, rho, aux] = NonGravForces.io.readSwarmDensity(files, 'TimeRange', [mjd0 mjd1])
\end{lstlisting}

\subsubsection*{INPUT ARGUMENTS}
\begin{longtable}{@{}>{\footnotesize\raggedright\arraybackslash\ttfamily}p{3.00cm}>{\footnotesize\raggedright\arraybackslash}p{1.50cm}>{\footnotesize\raggedright\arraybackslash}p{2.20cm}>{\footnotesize\raggedright\arraybackslash}p{1.50cm}>{\footnotesize\raggedright\arraybackslash}p{6.40cm}@{}}
\toprule
{\footnotesize\textbf{Name}} & {\footnotesize\textbf{Size}} & {\footnotesize\textbf{Type}} & {\footnotesize\textbf{Units}} & {\footnotesize\textbf{Description}} \\
\midrule
\endfirsthead
\toprule
{\footnotesize\textbf{Name}} & {\footnotesize\textbf{Size}} & {\footnotesize\textbf{Type}} & {\footnotesize\textbf{Units}} & {\footnotesize\textbf{Description}} \\
\midrule
\endhead
\midrule
\endfoot
\bottomrule
\endlastfoot
files & 1xF & string & - & One or more HAPI CSV files of ONE dataset (SW\_OPER\_DNS\textless{}A\textbar{}B\textbar{}C\textgreater{}POD\_2\_, CH\_OPER\_DNS\_ACC\_2\_, GR\_OPER\_DNS\textless{}1\textbar{}2\textgreater{}ACC\_2\_, GF\_OPER\_DNS\textless{}1\textbar{}2\textgreater{}ACC\_2\_); concatenated and sorted by time, duplicates of an epoch kept once. \\
\end{longtable}

\subsubsection*{NAME-VALUE ARGUMENTS}
\begin{longtable}{@{}>{\footnotesize\raggedright\arraybackslash\ttfamily}p{3.00cm}>{\footnotesize\raggedright\arraybackslash}p{1.50cm}>{\footnotesize\raggedright\arraybackslash}p{2.20cm}>{\footnotesize\raggedright\arraybackslash}p{1.50cm}>{\footnotesize\raggedright\arraybackslash}p{6.40cm}@{}}
\toprule
{\footnotesize\textbf{Name}} & {\footnotesize\textbf{Size}} & {\footnotesize\textbf{Type}} & {\footnotesize\textbf{Units}} & {\footnotesize\textbf{Description}} \\
\midrule
\endfirsthead
\toprule
{\footnotesize\textbf{Name}} & {\footnotesize\textbf{Size}} & {\footnotesize\textbf{Type}} & {\footnotesize\textbf{Units}} & {\footnotesize\textbf{Description}} \\
\midrule
\endhead
\midrule
\endfoot
\bottomrule
\endlastfoot
TimeRange & 1x2 & double & MJD(UTC) & Keep records with mjd0 \textless{}= t \textless{}= mjd1. Default [] (all). \\
\end{longtable}

\subsubsection*{OUTPUT ARGUMENTS}
\begin{longtable}{@{}>{\footnotesize\raggedright\arraybackslash\ttfamily}p{3.00cm}>{\footnotesize\raggedright\arraybackslash}p{1.50cm}>{\footnotesize\raggedright\arraybackslash}p{2.20cm}>{\footnotesize\raggedright\arraybackslash}p{1.50cm}>{\footnotesize\raggedright\arraybackslash}p{6.40cm}@{}}
\toprule
{\footnotesize\textbf{Name}} & {\footnotesize\textbf{Size}} & {\footnotesize\textbf{Type}} & {\footnotesize\textbf{Units}} & {\footnotesize\textbf{Description}} \\
\midrule
\endfirsthead
\toprule
{\footnotesize\textbf{Name}} & {\footnotesize\textbf{Size}} & {\footnotesize\textbf{Type}} & {\footnotesize\textbf{Units}} & {\footnotesize\textbf{Description}} \\
\midrule
\endhead
\midrule
\endfoot
\bottomrule
\endlastfoot
epoch & 1x1 & struct & - & referenceSystem epoch struct, K entries, scale "UTC" \\
lla & 3xK & double & rad,rad,m & Geodetic latitude, longitude, height (GRS80, as the product states; see ACCURACY for WGS-84) \\
rho & 1xK & double & kg/\allowbreak{}m\textasciicircum{}3 & Thermosphere mass density (NaN where the product has its fill) \\
aux & 1x1 & struct & - & .rhoOrbitMean (1xK kg/m\textasciicircum{}3, the product's orbit average, NaN at the fill), .lst\_h (1xK local solar time, h), .valid (1xK logical, validity\_flag == 0 and the density not the fill), .satellite ("A"\textbar{}"B"\textbar{}"C" for Swarm, "CHAMP", "GRACE-A"\textbar{}"GRACE-B", "GRACE-C"\textbar{}"GRACE-D" for GRACE-FO 1\textbar{}2), .dataset, .header (decoded HAPI header of the first file), .files, .numRecords, .provenance \\
\end{longtable}

\subsubsection*{ERRORS}
\begin{description}[leftmargin=1.2em,style=nextline,itemsep=0pt]
\item[\texttt{NonGravForces:\allowbreak{}readSwarmDensity:\allowbreak{}inputRange}] TimeRange is not [mjd0 mjd1] with mjd0 \textless{}= mjd1
\item[\texttt{NonGravForces:\allowbreak{}readSwarmDensity:\allowbreak{}dataFormat}] no HAPI JSON header, an undecodable one, or a status other than 1200 (OK)
\item[\texttt{NonGravForces:\allowbreak{}readSwarmDensity:\allowbreak{}dataColumn}] the first parameter is not Timestamp, a required parameter is missing from the header, or a data line does not have the header's columns
\item[\texttt{NonGravForces:\allowbreak{}readSwarmDensity:\allowbreak{}dataUnits}] a parameter's units are not the ones converted here
\item[\texttt{NonGravForces:\allowbreak{}readSwarmDensity:\allowbreak{}dataSatellite}] the files are of different datasets, or of a dataset this reader does not know
\item[\texttt{NonGravForces:\allowbreak{}readSwarmDensity:\allowbreak{}dataEmpty}] no record survives TimeRange
\end{description}

\subsection{\texttt{NonGravForces.io.readIMU1B}}
\label{api:io.readIMU1B}
\noindent IMU1B reader: the filtered, accumulated angle of each GRACE-FO gyro, in radians, one record per sample.

\begin{lstlisting}
[epoch, angle_rad, aux] = NonGravForces.io.readIMU1B(file)
[epoch, angle_rad, aux] = NonGravForces.io.readIMU1B(files, 'TimeRange', tr, 'QualityBitMask', bits)
\end{lstlisting}

\subsubsection*{INPUT ARGUMENTS}
\begin{longtable}{@{}>{\footnotesize\raggedright\arraybackslash\ttfamily}p{3.00cm}>{\footnotesize\raggedright\arraybackslash}p{1.50cm}>{\footnotesize\raggedright\arraybackslash}p{2.20cm}>{\footnotesize\raggedright\arraybackslash}p{1.50cm}>{\footnotesize\raggedright\arraybackslash}p{6.40cm}@{}}
\toprule
{\footnotesize\textbf{Name}} & {\footnotesize\textbf{Size}} & {\footnotesize\textbf{Type}} & {\footnotesize\textbf{Units}} & {\footnotesize\textbf{Description}} \\
\midrule
\endfirsthead
\toprule
{\footnotesize\textbf{Name}} & {\footnotesize\textbf{Size}} & {\footnotesize\textbf{Type}} & {\footnotesize\textbf{Units}} & {\footnotesize\textbf{Description}} \\
\midrule
\endhead
\midrule
\endfoot
\bottomrule
\endlastfoot
files & 1xF & string & - & One or more IMU1B files of ONE satellite, chronological; .gz accepted. \\
\end{longtable}

\subsubsection*{NAME-VALUE ARGUMENTS}
\begin{longtable}{@{}>{\footnotesize\raggedright\arraybackslash\ttfamily}p{3.00cm}>{\footnotesize\raggedright\arraybackslash}p{1.50cm}>{\footnotesize\raggedright\arraybackslash}p{2.20cm}>{\footnotesize\raggedright\arraybackslash}p{1.50cm}>{\footnotesize\raggedright\arraybackslash}p{6.40cm}@{}}
\toprule
{\footnotesize\textbf{Name}} & {\footnotesize\textbf{Size}} & {\footnotesize\textbf{Type}} & {\footnotesize\textbf{Units}} & {\footnotesize\textbf{Description}} \\
\midrule
\endfirsthead
\toprule
{\footnotesize\textbf{Name}} & {\footnotesize\textbf{Size}} & {\footnotesize\textbf{Type}} & {\footnotesize\textbf{Units}} & {\footnotesize\textbf{Description}} \\
\midrule
\endhead
\midrule
\endfoot
\bottomrule
\endlastfoot
TimeRange & 1x2 & double\textbar{}\allowbreak{}datetime & - & Forwarded to readL1b: decimal MJD or GPS datetime. Default []. \\
QualityBitMask & 1xQ & double & - & qualflg bits removing a record; forwarded. Default []. \\
\end{longtable}

\subsubsection*{OUTPUT ARGUMENTS}
\begin{longtable}{@{}>{\footnotesize\raggedright\arraybackslash\ttfamily}p{3.00cm}>{\footnotesize\raggedright\arraybackslash}p{1.50cm}>{\footnotesize\raggedright\arraybackslash}p{2.20cm}>{\footnotesize\raggedright\arraybackslash}p{1.50cm}>{\footnotesize\raggedright\arraybackslash}p{6.40cm}@{}}
\toprule
{\footnotesize\textbf{Name}} & {\footnotesize\textbf{Size}} & {\footnotesize\textbf{Type}} & {\footnotesize\textbf{Units}} & {\footnotesize\textbf{Description}} \\
\midrule
\endfirsthead
\toprule
{\footnotesize\textbf{Name}} & {\footnotesize\textbf{Size}} & {\footnotesize\textbf{Type}} & {\footnotesize\textbf{Units}} & {\footnotesize\textbf{Description}} \\
\midrule
\endhead
\midrule
\endfoot
\bottomrule
\endlastfoot
epoch & 1x1 & struct & - & referenceSystem epoch struct, K entries, "GPS" (each record's own time) \\
angle\_\allowbreak{}rad & 1xK & double & rad & FiltAng of the record's gyro, deg2rad of the published value \\
aux & 1x1 & struct & - & .gyroId (1xK), .gyros (sorted ids present), .timeRef (1xK string, "G" = GPS; 1x0 when the file has no time\_ref), .angleRange\_rad ([-1 1]*deg2rad(5758), the header's valid range), .product, .satellite, .files, .numRecords, .sampling\_s (median spacing of ALL records, i.e. the gyro interleave, not one gyro's 0.125 s), .obsNames, .units, .qualflgBits, .provenance \\
\end{longtable}

\subsubsection*{ERRORS}
\begin{description}[leftmargin=1.2em,style=nextline,itemsep=0pt]
\item[\texttt{NonGravForces:\allowbreak{}readIMU1B:\allowbreak{}dataProduct}] a file is not an IMU1B product
\item[\texttt{NonGravForces:\allowbreak{}readIMU1B:\allowbreak{}dataSatellite}] the files mix GRACE-C and GRACE-D
\item[\texttt{NonGravForces:\allowbreak{}readIMU1B:\allowbreak{}dataEmpty}] no record survives the filters
\item[\texttt{NonGravForces:\allowbreak{}readIMU1B:\allowbreak{}dataColumn}] FiltAng or gyro\_id is missing
\item[\texttt{NonGravForces:\allowbreak{}readIMU1B:\allowbreak{}dataTimeRef}] a record is not tagged in GPS time (time\_ref \textasciitilde{}= "G")
\item[\texttt{gracefoReading:\allowbreak{}*, referenceSystem:\allowbreak{}*}] propagated unwrapped
\end{description}

\section{L2 \texttt{env} --- the environment: density, wind, shadow, Earth radiance, magnetic field}

\subsection{\texttt{NonGravForces.env.density}}
\label{api:env.density}
\noindent Thermosphere density and composition: exponential, NRLMSISE-00, NRLMSIS 2.1 or DTM2020 (compiled).

\begin{lstlisting}
[rho, species, T, info] = NonGravForces.env.density("msis21", epoch, lla, sw)
[rho, species, T, info] = NonGravForces.env.density("msis00", epoch, lla, sw)
[rho, species, T, info] = NonGravForces.env.density("dtm2020", epoch, lla, sw)
[rho, species, T, info] = NonGravForces.env.density("msis00", epoch, lla, [], 'Drivers', d)
rho = NonGravForces.env.density("exponential", epoch, lla, [], 'Exponential', p)
\end{lstlisting}

\subsubsection*{INPUT ARGUMENTS}
\begin{longtable}{@{}>{\footnotesize\raggedright\arraybackslash\ttfamily}p{3.00cm}>{\footnotesize\raggedright\arraybackslash}p{1.50cm}>{\footnotesize\raggedright\arraybackslash}p{2.20cm}>{\footnotesize\raggedright\arraybackslash}p{1.50cm}>{\footnotesize\raggedright\arraybackslash}p{6.40cm}@{}}
\toprule
{\footnotesize\textbf{Name}} & {\footnotesize\textbf{Size}} & {\footnotesize\textbf{Type}} & {\footnotesize\textbf{Units}} & {\footnotesize\textbf{Description}} \\
\midrule
\endfirsthead
\toprule
{\footnotesize\textbf{Name}} & {\footnotesize\textbf{Size}} & {\footnotesize\textbf{Type}} & {\footnotesize\textbf{Units}} & {\footnotesize\textbf{Description}} \\
\midrule
\endhead
\midrule
\endfoot
\bottomrule
\endlastfoot
model & 1x1 & string & - & "exponential" \textbar{} "msis00" \textbar{} "msis21" \textbar{} "dtm2020" \\
epoch & 1x1 & struct & - & referenceSystem epoch struct, K entries, any scale \\
lla & 3xK & double & rad,rad,m & Geodetic latitude, longitude, height (WGS-84) \\
sw & 1x1 & struct & - & NonGravForces.io.spaceWeather result of the celestrak family, or [] with 'Drivers' or for "exponential". \\
\end{longtable}

\subsubsection*{NAME-VALUE ARGUMENTS}
\begin{longtable}{@{}>{\footnotesize\raggedright\arraybackslash\ttfamily}p{3.00cm}>{\footnotesize\raggedright\arraybackslash}p{1.50cm}>{\footnotesize\raggedright\arraybackslash}p{2.20cm}>{\footnotesize\raggedright\arraybackslash}p{1.50cm}>{\footnotesize\raggedright\arraybackslash}p{6.40cm}@{}}
\toprule
{\footnotesize\textbf{Name}} & {\footnotesize\textbf{Size}} & {\footnotesize\textbf{Type}} & {\footnotesize\textbf{Units}} & {\footnotesize\textbf{Description}} \\
\midrule
\endfirsthead
\toprule
{\footnotesize\textbf{Name}} & {\footnotesize\textbf{Size}} & {\footnotesize\textbf{Type}} & {\footnotesize\textbf{Units}} & {\footnotesize\textbf{Description}} \\
\midrule
\endhead
\midrule
\endfoot
\bottomrule
\endlastfoot
Drivers & 1x1 & struct & - & .f107 (1xK, previous-day observed F10.7), .f107a (1xK, 81-day centred), .ap (1xK daily Ap, or 7xK NRLMSISE history). Overrides sw. dtm2020 instead: .f107 (previous day), .f107bar (81 days ending at the day), .kp3 (Kp 3 h back), .kp24 (24 h mean). \\
AnomalousOxygen & 1x1 & logical & - & true (default): the mass density includes anomalous oxygen (msis00: gtd7d; msis21: the model's default); false: it does not (gtd7). \\
Exponential & 1x1 & struct & - & .rho0\_kg\_m3, .h0\_m, .H\_m; required for "exponential", ignored otherwise. \\
LocalSolarTime\_\allowbreak{}h & 1xK & double & h & msis00 and dtm2020. Override of the local solar time handed to the model. Default [] = sec/3600 + lon/15, the consistent choice the reference recommends. The published test cases set an inconsistent 16 h and are reproduced only with this override. \\
\end{longtable}

\subsubsection*{OUTPUT ARGUMENTS}
\begin{longtable}{@{}>{\footnotesize\raggedright\arraybackslash\ttfamily}p{3.00cm}>{\footnotesize\raggedright\arraybackslash}p{1.50cm}>{\footnotesize\raggedright\arraybackslash}p{2.20cm}>{\footnotesize\raggedright\arraybackslash}p{1.50cm}>{\footnotesize\raggedright\arraybackslash}p{6.40cm}@{}}
\toprule
{\footnotesize\textbf{Name}} & {\footnotesize\textbf{Size}} & {\footnotesize\textbf{Type}} & {\footnotesize\textbf{Units}} & {\footnotesize\textbf{Description}} \\
\midrule
\endfirsthead
\toprule
{\footnotesize\textbf{Name}} & {\footnotesize\textbf{Size}} & {\footnotesize\textbf{Type}} & {\footnotesize\textbf{Units}} & {\footnotesize\textbf{Description}} \\
\midrule
\endhead
\midrule
\endfoot
\bottomrule
\endlastfoot
rho & 1xK & double & kg/\allowbreak{}m\textasciicircum{}3 & Total mass density. \\
species & 8xK & double & 1/\allowbreak{}m\textasciicircum{}3 & Number densities He O N2 O2 Ar H N O\_anomalous (const.speciesName order); NaN for "exponential". \\
T & 2xK & double & K & [exospheric temperature; temperature at h]; NaN for "exponential". \\
info & 1x1 & struct & - & .backend, .drivers (struct actually used: f107 f107a ap for the msis backends, f107 f107bar kp3 kp24 for dtm2020, and .source "spaceWeather" \textbar{} "caller" \textbar{} "none"), .anomalousOxygen (false for exponential and dtm2020), .msisInput (struct doy sec alt\_km lat\_deg lon\_deg, 1xK each, and lst\_h for msis00; both msis backends), .dtmInput (doy alt\_km lst\_h lat\_rad lon\_rad; dtm2020), .exponential (the parameters; exponential), .provenance. \\
\end{longtable}

\subsubsection*{ERRORS}
\begin{description}[leftmargin=1.2em,style=nextline,itemsep=0pt]
\item[\texttt{NonGravForces:\allowbreak{}density:\allowbreak{}inputShape}] lla not 3xK or K differs from epoch
\item[\texttt{NonGravForces:\allowbreak{}density:\allowbreak{}inputDrivers}] no sw and no Drivers; sw not of the celestrak family; Drivers or LocalSolarTime\_h malformed
\item[\texttt{NonGravForces:\allowbreak{}density:\allowbreak{}inputExponential}] "exponential" without a complete Exponential struct
\item[\texttt{NonGravForces:\allowbreak{}density:\allowbreak{}numericNaN}] NaN in lla or in the drivers
\item[\texttt{NonGravForces:\allowbreak{}density:\allowbreak{}depMissing}] the backend's MEX is not built
\item[\texttt{NonGravForces:\allowbreak{}density:\allowbreak{}inputOption}] LocalSolarTime\_h given to msis21, which forms local time itself
\item[\texttt{NonGravForces:\allowbreak{}density:\allowbreak{}inputRange}] dtm2020 below 150 km
\item[\texttt{NonGravForces:\allowbreak{}dtm2020Gateway:\allowbreak{}*}] path, coefficient file and screening errors of the gateway, unwrapped
\item[\texttt{NonGravForces:\allowbreak{}msis21Gateway:\allowbreak{}*}] path, parameter file and screening errors of the gateway, unwrapped
\item[\texttt{NonGravForces:\allowbreak{}spaceWeather:\allowbreak{}*}] from sw.query, unwrapped
\item[\texttt{referenceSystem:\allowbreak{}*}] propagated unwrapped
\end{description}

\subsection{\texttt{NonGravForces.env.densityScale}}
\label{api:env.densityScale}
\noindent Observed daily scale on a model density from Swarm's POD-only densities.

\begin{lstlisting}
S = NonGravForces.env.densityScale(files, sw)
S = NonGravForces.env.densityScale(files, sw, 'Model', "msis00", 'Stride', 10)
f = S.at(epoch, h_m)                        % 1xK factor on the model density
\end{lstlisting}

\subsubsection*{INPUT ARGUMENTS}
\begin{longtable}{@{}>{\footnotesize\raggedright\arraybackslash\ttfamily}p{3.00cm}>{\footnotesize\raggedright\arraybackslash}p{1.50cm}>{\footnotesize\raggedright\arraybackslash}p{2.20cm}>{\footnotesize\raggedright\arraybackslash}p{1.50cm}>{\footnotesize\raggedright\arraybackslash}p{6.40cm}@{}}
\toprule
{\footnotesize\textbf{Name}} & {\footnotesize\textbf{Size}} & {\footnotesize\textbf{Type}} & {\footnotesize\textbf{Units}} & {\footnotesize\textbf{Description}} \\
\midrule
\endfirsthead
\toprule
{\footnotesize\textbf{Name}} & {\footnotesize\textbf{Size}} & {\footnotesize\textbf{Type}} & {\footnotesize\textbf{Units}} & {\footnotesize\textbf{Description}} \\
\midrule
\endhead
\midrule
\endfoot
\bottomrule
\endlastfoot
files & 1xF & string & - & Swarm DNSxPOD HAPI CSV files (io.readSwarmDensity), any months, of B and of A and/or C. Files of one satellite are read together; an epoch in two files counts once. \\
sw & 1x1 & struct & - & io.spaceWeather result (celestrak family), the drivers of the model density. \\
\end{longtable}

\subsubsection*{NAME-VALUE ARGUMENTS}
\begin{longtable}{@{}>{\footnotesize\raggedright\arraybackslash\ttfamily}p{3.00cm}>{\footnotesize\raggedright\arraybackslash}p{1.50cm}>{\footnotesize\raggedright\arraybackslash}p{2.20cm}>{\footnotesize\raggedright\arraybackslash}p{1.50cm}>{\footnotesize\raggedright\arraybackslash}p{6.40cm}@{}}
\toprule
{\footnotesize\textbf{Name}} & {\footnotesize\textbf{Size}} & {\footnotesize\textbf{Type}} & {\footnotesize\textbf{Units}} & {\footnotesize\textbf{Description}} \\
\midrule
\endfirsthead
\toprule
{\footnotesize\textbf{Name}} & {\footnotesize\textbf{Size}} & {\footnotesize\textbf{Type}} & {\footnotesize\textbf{Units}} & {\footnotesize\textbf{Description}} \\
\midrule
\endhead
\midrule
\endfoot
\bottomrule
\endlastfoot
Model & 1x1 & string & - & env.density backend the factor is FOR: "msis00" (default), "msis21" or "dtm2020". A Model refuses a scale formed for another density model than its drag.density. \\
Stride & 1x1 & double & - & Every Stride-th valid 30 s record enters the sums (default 10, i.e. 5 min; the product's own resolution is \textasciitilde{}20 min). \\
MinRecords & 1x1 & double & - & Fewer used records in a day make that day's factor NaN, which S.at refuses (default 100, about a third of a day at Stride 10). \\
Fallback & 1x1 & struct & - & an env.densityDriverScale result for the same Model (default none). Given, S is of kind "combined" (Options drag.densityScale "combined"): Swarm's factor on every day that has one, the fallback's predicted factor on every other day -- never mixed within a day. S.source(epoch) says which per epoch. \\
\end{longtable}

\subsubsection*{OUTPUT ARGUMENTS}
\begin{longtable}{@{}>{\footnotesize\raggedright\arraybackslash\ttfamily}p{3.00cm}>{\footnotesize\raggedright\arraybackslash}p{1.50cm}>{\footnotesize\raggedright\arraybackslash}p{2.20cm}>{\footnotesize\raggedright\arraybackslash}p{1.50cm}>{\footnotesize\raggedright\arraybackslash}p{6.40cm}@{}}
\toprule
{\footnotesize\textbf{Name}} & {\footnotesize\textbf{Size}} & {\footnotesize\textbf{Type}} & {\footnotesize\textbf{Units}} & {\footnotesize\textbf{Description}} \\
\midrule
\endfirsthead
\toprule
{\footnotesize\textbf{Name}} & {\footnotesize\textbf{Size}} & {\footnotesize\textbf{Type}} & {\footnotesize\textbf{Units}} & {\footnotesize\textbf{Description}} \\
\midrule
\endhead
\midrule
\endfoot
\bottomrule
\endlastfoot
S & 1x1 & struct & - & .kind "swarm" ("combined" with 'Fallback'), .fallback (the 'Fallback' struct, or []), .day (1xD UTC MJD of each day), .satellite (1xP "A"\textbar{}"B"\textbar{}"C"), .factor (PxD), .height\_m (PxD mean geodetic height of the used records), .nUsed (PxD), .low (1xD) and .high (1xD) bracket factors, .hLow\_m .hHigh\_m (1xD), .model, .stride, .files, .provenance, and .at (function handle: epoch struct (any scale), 1xK heights [m] -\textgreater{} 1xK factors), .source (function handle: epoch -\textgreater{} 1xK "swarm" \textbar{} "driver" \textbar{} "none", the factor .at takes) \\
\end{longtable}

\subsubsection*{ERRORS}
\begin{description}[leftmargin=1.2em,style=nextline,itemsep=0pt]
\item[\texttt{NonGravForces:\allowbreak{}densityScale:\allowbreak{}inputSatellites}] no B, or neither A nor C, among the files; or a file that is not Swarm A/B/C (a TOLEOS density of CHAMP, GRACE or GRACE-FO)
\item[\texttt{NonGravForces:\allowbreak{}densityScale:\allowbreak{}inputModel}] Model is not a density backend with drivers
\item[\texttt{NonGravForces:\allowbreak{}densityScale:\allowbreak{}dataDay}] S.at at an epoch whose day has no factor (no file covers it, or fewer than MinRecords records) and no 'Fallback' was given
\item[\texttt{NonGravForces:\allowbreak{}densityScale:\allowbreak{}inputFallback}] 'Fallback' is not an env.densityDriverScale result, or was formed for another density model
\item[\texttt{NonGravForces:\allowbreak{}densityDriverScale:\allowbreak{}*}] propagated from the fallback on a day without Swarm
\item[\texttt{NonGravForces:\allowbreak{}densityScale:\allowbreak{}inputRange}] S.at at a height further outside the bracket than its width, or h not 1xK
\item[\texttt{NonGravForces:\allowbreak{}readSwarmDensity:\allowbreak{}*, NonGravForces:\allowbreak{}density:\allowbreak{}*}] propagated unwrapped
\end{description}

\subsection{\texttt{NonGravForces.env.densityDriverScale}}
\label{api:env.densityDriverScale}
\noindent Daily scale on NRLMSISE-00 from solar and geomagnetic drivers alone, for any date.

\begin{lstlisting}
S = NonGravForces.env.densityDriverScale(sw, cls)
S = NonGravForces.env.densityDriverScale(sw, cls, 'Offset', -0.041)   % GRACE-FO's own level
f = S.at(epoch, h_m)                        % 1xK factor on the model density
\end{lstlisting}

\subsubsection*{INPUT ARGUMENTS}
\begin{longtable}{@{}>{\footnotesize\raggedright\arraybackslash\ttfamily}p{3.00cm}>{\footnotesize\raggedright\arraybackslash}p{1.50cm}>{\footnotesize\raggedright\arraybackslash}p{2.20cm}>{\footnotesize\raggedright\arraybackslash}p{1.50cm}>{\footnotesize\raggedright\arraybackslash}p{6.40cm}@{}}
\toprule
{\footnotesize\textbf{Name}} & {\footnotesize\textbf{Size}} & {\footnotesize\textbf{Type}} & {\footnotesize\textbf{Units}} & {\footnotesize\textbf{Description}} \\
\midrule
\endfirsthead
\toprule
{\footnotesize\textbf{Name}} & {\footnotesize\textbf{Size}} & {\footnotesize\textbf{Type}} & {\footnotesize\textbf{Units}} & {\footnotesize\textbf{Description}} \\
\midrule
\endhead
\midrule
\endfoot
\bottomrule
\endlastfoot
sw & 1x1 & struct & - & io.spaceWeather result, celestrak family (SW-All): the daily Ap \\
cls & 1x1 & struct & - & io.spaceWeather result, f30 family, of CLS's ADJUSTED file (radio\_flux\_adjusted\_observation.txt): the daily F30 \\
\end{longtable}

\subsubsection*{NAME-VALUE ARGUMENTS}
\begin{longtable}{@{}>{\footnotesize\raggedright\arraybackslash\ttfamily}p{3.00cm}>{\footnotesize\raggedright\arraybackslash}p{1.50cm}>{\footnotesize\raggedright\arraybackslash}p{2.20cm}>{\footnotesize\raggedright\arraybackslash}p{1.50cm}>{\footnotesize\raggedright\arraybackslash}p{6.40cm}@{}}
\toprule
{\footnotesize\textbf{Name}} & {\footnotesize\textbf{Size}} & {\footnotesize\textbf{Type}} & {\footnotesize\textbf{Units}} & {\footnotesize\textbf{Description}} \\
\midrule
\endfirsthead
\toprule
{\footnotesize\textbf{Name}} & {\footnotesize\textbf{Size}} & {\footnotesize\textbf{Type}} & {\footnotesize\textbf{Units}} & {\footnotesize\textbf{Description}} \\
\midrule
\endhead
\midrule
\endfoot
\bottomrule
\endlastfoot
Model & 1x1 & string & - & density backend the factor is for; only "msis00" (default), the model the coefficients were fitted against \\
Offset & 1x1 & double & log f & added to log f: a satellite's own level (default 0, Swarm's; -0.041 for GRACE-FO 1, +0.036 CHAMP, +0.045 GRACE-A, run 40676249479) \\
\end{longtable}

\subsubsection*{OUTPUT ARGUMENTS}
\begin{longtable}{@{}>{\footnotesize\raggedright\arraybackslash\ttfamily}p{3.00cm}>{\footnotesize\raggedright\arraybackslash}p{1.50cm}>{\footnotesize\raggedright\arraybackslash}p{2.20cm}>{\footnotesize\raggedright\arraybackslash}p{1.50cm}>{\footnotesize\raggedright\arraybackslash}p{6.40cm}@{}}
\toprule
{\footnotesize\textbf{Name}} & {\footnotesize\textbf{Size}} & {\footnotesize\textbf{Type}} & {\footnotesize\textbf{Units}} & {\footnotesize\textbf{Description}} \\
\midrule
\endfirsthead
\toprule
{\footnotesize\textbf{Name}} & {\footnotesize\textbf{Size}} & {\footnotesize\textbf{Type}} & {\footnotesize\textbf{Units}} & {\footnotesize\textbf{Description}} \\
\midrule
\endhead
\midrule
\endfoot
\bottomrule
\endlastfoot
S & 1x1 & struct & - & .kind "driver", .model "msis00", .proxy "F30", .coef (struct of the transcribed coefficients, .a0 .Fbar .dF .ap .sinAnnual .cosAnnual .sinSemiannual .cosSemiannual .FbarSq .centre\_sfu .run), .offset, .drivers (function handle: epoch -\textgreater{} 8xK drivers in the order above), .at (function handle: epoch struct (any scale), 1xK heights [m] -\textgreater{} 1xK factors; the heights are checked for size and otherwise unused), .provenance \\
\end{longtable}

\subsubsection*{ERRORS}
\begin{description}[leftmargin=1.2em,style=nextline,itemsep=0pt]
\item[\texttt{NonGravForces:\allowbreak{}densityDriverScale:\allowbreak{}inputModel}] Model is not msis00
\item[\texttt{NonGravForces:\allowbreak{}densityDriverScale:\allowbreak{}inputFamily}] sw is not the celestrak family, or cls not f30, or cls is the flux not adjusted to 1 AU (.adjusted false)
\item[\texttt{NonGravForces:\allowbreak{}densityDriverScale:\allowbreak{}dataDay}] S.at at a UTC day whose 81 previous days of F30 or whose Ap are not all in the files
\item[\texttt{NonGravForces:\allowbreak{}densityDriverScale:\allowbreak{}inputRange}] h not 1xK, one height per epoch
\end{description}

\subsection{\texttt{NonGravForces.env.wind}}
\label{api:env.wind}
\noindent Velocity of the satellite relative to the air, GCRS: co-rotation, optional thermospheric wind.

\begin{lstlisting}
[vrel, info] = NonGravForces.env.wind("corotation", epoch, r, v, 'Eop', eop)
[vrel, info] = NonGravForces.env.wind("none", epoch, r, v)
[vrel, info] = NonGravForces.env.wind("corotation", epoch, r, v, 'Eop', eop, 'Wind_m_s', w)
[vrel, info] = NonGravForces.env.wind("hwm14", epoch, r, v, 'Eop', eop, 'SpaceWeather', sw)
[vrel, info] = NonGravForces.env.wind("hwm14", epoch, r, v, 'Eop', eop, 'Disturbed', false)
\end{lstlisting}

\subsubsection*{INPUT ARGUMENTS}
\begin{longtable}{@{}>{\footnotesize\raggedright\arraybackslash\ttfamily}p{3.00cm}>{\footnotesize\raggedright\arraybackslash}p{1.50cm}>{\footnotesize\raggedright\arraybackslash}p{2.20cm}>{\footnotesize\raggedright\arraybackslash}p{1.50cm}>{\footnotesize\raggedright\arraybackslash}p{6.40cm}@{}}
\toprule
{\footnotesize\textbf{Name}} & {\footnotesize\textbf{Size}} & {\footnotesize\textbf{Type}} & {\footnotesize\textbf{Units}} & {\footnotesize\textbf{Description}} \\
\midrule
\endfirsthead
\toprule
{\footnotesize\textbf{Name}} & {\footnotesize\textbf{Size}} & {\footnotesize\textbf{Type}} & {\footnotesize\textbf{Units}} & {\footnotesize\textbf{Description}} \\
\midrule
\endhead
\midrule
\endfoot
\bottomrule
\endlastfoot
model & 1x1 & string & - & "none" \textbar{} "corotation" \textbar{} "hwm14" \\
epoch & 1x1 & struct & - & referenceSystem epoch struct with K entries \\
r\_\allowbreak{}gcrs & 3xK & double & m & Satellite position, GCRS \\
v\_\allowbreak{}gcrs & 3xK & double & m/\allowbreak{}s & Satellite velocity, GCRS \\
\end{longtable}

\subsubsection*{NAME-VALUE ARGUMENTS}
\begin{longtable}{@{}>{\footnotesize\raggedright\arraybackslash\ttfamily}p{3.00cm}>{\footnotesize\raggedright\arraybackslash}p{1.50cm}>{\footnotesize\raggedright\arraybackslash}p{2.20cm}>{\footnotesize\raggedright\arraybackslash}p{1.50cm}>{\footnotesize\raggedright\arraybackslash}p{6.40cm}@{}}
\toprule
{\footnotesize\textbf{Name}} & {\footnotesize\textbf{Size}} & {\footnotesize\textbf{Type}} & {\footnotesize\textbf{Units}} & {\footnotesize\textbf{Description}} \\
\midrule
\endfirsthead
\toprule
{\footnotesize\textbf{Name}} & {\footnotesize\textbf{Size}} & {\footnotesize\textbf{Type}} & {\footnotesize\textbf{Units}} & {\footnotesize\textbf{Description}} \\
\midrule
\endhead
\midrule
\endfoot
\bottomrule
\endlastfoot
Eop & 1x1 & struct & - & referenceSystem.readEop result, or a file path; required by "corotation" and "hwm14" (no default: the rotation cannot be guessed, rule R5) \\
Wind\_\allowbreak{}m\_\allowbreak{}s & 3xK & double & m/\allowbreak{}s & Air velocity in GCRS to subtract in addition, default [] (none) \\
SpaceWeather & 1x1 & struct & - & io.spaceWeather (celestrak family): the current 3-hour ap at each epoch. "hwm14" only. \\
Ap & 1x1\textbar{}\allowbreak{}1xK & double & 2 nT & The current 3-hour ap, given instead of SpaceWeather. "hwm14" only. \\
Disturbed & 1x1 & logical & - & "hwm14": add DWM07 (default true, which needs an ap source); false: quiet-time winds only \\
\end{longtable}

\subsubsection*{OUTPUT ARGUMENTS}
\begin{longtable}{@{}>{\footnotesize\raggedright\arraybackslash\ttfamily}p{3.00cm}>{\footnotesize\raggedright\arraybackslash}p{1.50cm}>{\footnotesize\raggedright\arraybackslash}p{2.20cm}>{\footnotesize\raggedright\arraybackslash}p{1.50cm}>{\footnotesize\raggedright\arraybackslash}p{6.40cm}@{}}
\toprule
{\footnotesize\textbf{Name}} & {\footnotesize\textbf{Size}} & {\footnotesize\textbf{Type}} & {\footnotesize\textbf{Units}} & {\footnotesize\textbf{Description}} \\
\midrule
\endfirsthead
\toprule
{\footnotesize\textbf{Name}} & {\footnotesize\textbf{Size}} & {\footnotesize\textbf{Type}} & {\footnotesize\textbf{Units}} & {\footnotesize\textbf{Description}} \\
\midrule
\endhead
\midrule
\endfoot
\bottomrule
\endlastfoot
vrel & 3xK & double & m/\allowbreak{}s & Satellite velocity relative to the air, GCRS \\
info & 1x1 & struct & - & .model, .corotation\_m\_s (3xK, the term removed), .hwm14 (hwm14: .northEast\_m\_s 2xK, .gcrs\_m\_s 3xK, .ap 1xK or [], .disturbed, .geodetic 3xK [lat; lon; h] rad rad m), .speed\_m\_s (1xK, \textbar{}vrel\textbar{}), .windGiven (logical), .K, .provenance \\
\end{longtable}

\subsubsection*{ERRORS}
\begin{description}[leftmargin=1.2em,style=nextline,itemsep=0pt]
\item[\texttt{NonGravForces:\allowbreak{}wind:\allowbreak{}inputShape}] r, v or Wind\_m\_s not 3xK with K matching the epoch
\item[\texttt{NonGravForces:\allowbreak{}wind:\allowbreak{}inputEop}] "corotation" or "hwm14" without 'Eop'
\item[\texttt{NonGravForces:\allowbreak{}wind:\allowbreak{}inputAp}] "hwm14" disturbed without an ap source (or with a SpaceWeather not of the celestrak family), both sources, a negative or malformed Ap, an ap source with 'Disturbed' false, or an ap input to another model
\item[\texttt{NonGravForces:\allowbreak{}wind:\allowbreak{}depMissing}] "hwm14" MEX not built
\item[\texttt{NonGravForces:\allowbreak{}wind:\allowbreak{}inputRange}] "hwm14" at a position below the ellipsoid (the model does not return there)
\item[\texttt{NonGravForces:\allowbreak{}hwm14Gateway:\allowbreak{}*}] the gateway's refusals, unwrapped
\item[\texttt{NonGravForces:\allowbreak{}wind:\allowbreak{}numericNaN}] NaN/Inf in r, v or the supplied wind
\item[\texttt{referenceSystem:\allowbreak{}*}] propagated unwrapped
\end{description}

\subsection{\texttt{NonGravForces.env.eclipseFactor}}
\label{api:env.eclipseFactor}
\noindent Shadow function nu in [0,1]: cylindrical, conical or SOLAARS-CF Earth shadow, optional lunar shadow.

\begin{lstlisting}
nu = NonGravForces.env.eclipseFactor(r, r_sun)                       % conical
[nu, info] = NonGravForces.env.eclipseFactor(r, r_sun, 'Model', "cylindrical")
[nu, info] = NonGravForces.env.eclipseFactor(r, r_sun, 'Model', "solaarsCF", 'EclipticNormal', e)
[nu, info] = NonGravForces.env.eclipseFactor(r, r_sun, 'Moon', r_moon)  % Earth x Moon factor
\end{lstlisting}

\subsubsection*{INPUT ARGUMENTS}
\begin{longtable}{@{}>{\footnotesize\raggedright\arraybackslash\ttfamily}p{3.00cm}>{\footnotesize\raggedright\arraybackslash}p{1.50cm}>{\footnotesize\raggedright\arraybackslash}p{2.20cm}>{\footnotesize\raggedright\arraybackslash}p{1.50cm}>{\footnotesize\raggedright\arraybackslash}p{6.40cm}@{}}
\toprule
{\footnotesize\textbf{Name}} & {\footnotesize\textbf{Size}} & {\footnotesize\textbf{Type}} & {\footnotesize\textbf{Units}} & {\footnotesize\textbf{Description}} \\
\midrule
\endfirsthead
\toprule
{\footnotesize\textbf{Name}} & {\footnotesize\textbf{Size}} & {\footnotesize\textbf{Type}} & {\footnotesize\textbf{Units}} & {\footnotesize\textbf{Description}} \\
\midrule
\endhead
\midrule
\endfoot
\bottomrule
\endlastfoot
r & 3xK & double & m & Satellite position, Earth-centred (GCRS). \\
r\_\allowbreak{}sun & 3xK\textbar{}\allowbreak{}3x1 & double & m & Earth -\textgreater{} Sun vector (same for all K if 3x1). \\
\end{longtable}

\subsubsection*{NAME-VALUE ARGUMENTS}
\begin{longtable}{@{}>{\footnotesize\raggedright\arraybackslash\ttfamily}p{3.00cm}>{\footnotesize\raggedright\arraybackslash}p{1.50cm}>{\footnotesize\raggedright\arraybackslash}p{2.20cm}>{\footnotesize\raggedright\arraybackslash}p{1.50cm}>{\footnotesize\raggedright\arraybackslash}p{6.40cm}@{}}
\toprule
{\footnotesize\textbf{Name}} & {\footnotesize\textbf{Size}} & {\footnotesize\textbf{Type}} & {\footnotesize\textbf{Units}} & {\footnotesize\textbf{Description}} \\
\midrule
\endfirsthead
\toprule
{\footnotesize\textbf{Name}} & {\footnotesize\textbf{Size}} & {\footnotesize\textbf{Type}} & {\footnotesize\textbf{Units}} & {\footnotesize\textbf{Description}} \\
\midrule
\endhead
\midrule
\endfoot
\bottomrule
\endlastfoot
Model & 1x1 & string & - & "conical" (default) \textbar{} "cylindrical" \textbar{} "solaarsCF" \\
Moon & 3xK\textbar{}\allowbreak{}3x1 & double & m & Earth -\textgreater{} Moon vector; adds the lunar factor. Default [] (no lunar shadow). \\
EclipticNormal & 3xK\textbar{}\allowbreak{}3x1 & double & - & Unit normal of the ecliptic in the frame of r (e.g. r\_sun x v\_sun normalised, as radiationSources measures it). Required for solaarsCF, ignored otherwise. Default []. \\
\end{longtable}

\subsubsection*{OUTPUT ARGUMENTS}
\begin{longtable}{@{}>{\footnotesize\raggedright\arraybackslash\ttfamily}p{3.00cm}>{\footnotesize\raggedright\arraybackslash}p{1.50cm}>{\footnotesize\raggedright\arraybackslash}p{2.20cm}>{\footnotesize\raggedright\arraybackslash}p{1.50cm}>{\footnotesize\raggedright\arraybackslash}p{6.40cm}@{}}
\toprule
{\footnotesize\textbf{Name}} & {\footnotesize\textbf{Size}} & {\footnotesize\textbf{Type}} & {\footnotesize\textbf{Units}} & {\footnotesize\textbf{Description}} \\
\midrule
\endfirsthead
\toprule
{\footnotesize\textbf{Name}} & {\footnotesize\textbf{Size}} & {\footnotesize\textbf{Type}} & {\footnotesize\textbf{Units}} & {\footnotesize\textbf{Description}} \\
\midrule
\endhead
\midrule
\endfoot
\bottomrule
\endlastfoot
nu & 1xK & double & - & Fraction of the solar radiation, in [0, 1]. \\
info & 1x1 & struct & - & .model, .earth (1xK Earth factor), .moon (1xK or []), .radiusEarth\_m, .radiusMoon\_m, .radiusSun\_m, .solaars (struct rR\_Mm, rE\_Mm 1xK, oblatenessScale, clampedCount; solaarsCF only), .provenance. \\
\end{longtable}

\subsubsection*{ERRORS}
\begin{description}[leftmargin=1.2em,style=nextline,itemsep=0pt]
\item[\texttt{NonGravForces:\allowbreak{}eclipseFactor:\allowbreak{}inputShape}] r not 3xK; r\_sun, Moon or EclipticNormal not 3xK / 3x1
\item[\texttt{NonGravForces:\allowbreak{}eclipseFactor:\allowbreak{}numericNaN}] NaN/Inf in a position
\item[\texttt{NonGravForces:\allowbreak{}eclipseFactor:\allowbreak{}inputEclipticNormal}] solaarsCF without EclipticNormal, or not unit
\item[\texttt{referenceSystem:\allowbreak{}*}] propagated unwrapped
\end{description}

\subsection{\texttt{NonGravForces.env.earthRadiance}}
\label{api:env.earthRadiance}
\noindent Earth albedo and longwave emission on a cap of tiles around the sub-satellite point (Knocke or CERES), or as a point source.

\begin{lstlisting}
tiles = NonGravForces.env.earthRadiance("knocke", epoch, [12 24], r)
tiles = NonGravForces.env.earthRadiance("ceres", epoch, [24 48], r, 'Ceres', cer, 'RotationItrsToGcrs', R)
tiles = NonGravForces.env.earthRadiance("pointSource", epoch, [], r)
tiles = NonGravForces.env.earthRadiance("ceresClimatology", epoch, [24 48], r, 'Climatology', clm, 'RotationItrsToGcrs', R)
\end{lstlisting}

\subsubsection*{INPUT ARGUMENTS}
\begin{longtable}{@{}>{\footnotesize\raggedright\arraybackslash\ttfamily}p{3.00cm}>{\footnotesize\raggedright\arraybackslash}p{1.50cm}>{\footnotesize\raggedright\arraybackslash}p{2.20cm}>{\footnotesize\raggedright\arraybackslash}p{1.50cm}>{\footnotesize\raggedright\arraybackslash}p{6.40cm}@{}}
\toprule
{\footnotesize\textbf{Name}} & {\footnotesize\textbf{Size}} & {\footnotesize\textbf{Type}} & {\footnotesize\textbf{Units}} & {\footnotesize\textbf{Description}} \\
\midrule
\endfirsthead
\toprule
{\footnotesize\textbf{Name}} & {\footnotesize\textbf{Size}} & {\footnotesize\textbf{Type}} & {\footnotesize\textbf{Units}} & {\footnotesize\textbf{Description}} \\
\midrule
\endhead
\midrule
\endfoot
\bottomrule
\endlastfoot
backend & 1x1 & string & - & "knocke" \textbar{} "pointSource" \textbar{} "ceres" \textbar{} "ceresClimatology" \\
epoch & 1x1 & struct & - & referenceSystem epoch struct, K entries, any scale \\
grid & 1x2 & double & - & [nRings nSectors], positive integers (Options erp.grid; ignored for pointSource, pass []). \\
r & 3xK & double & m & Satellite position (Earth-centred, GCRS); the cap is centred on r/\textbar{}r\textbar{}. \\
\end{longtable}

\subsubsection*{NAME-VALUE ARGUMENTS}
\begin{longtable}{@{}>{\footnotesize\raggedright\arraybackslash\ttfamily}p{3.00cm}>{\footnotesize\raggedright\arraybackslash}p{1.50cm}>{\footnotesize\raggedright\arraybackslash}p{2.20cm}>{\footnotesize\raggedright\arraybackslash}p{1.50cm}>{\footnotesize\raggedright\arraybackslash}p{6.40cm}@{}}
\toprule
{\footnotesize\textbf{Name}} & {\footnotesize\textbf{Size}} & {\footnotesize\textbf{Type}} & {\footnotesize\textbf{Units}} & {\footnotesize\textbf{Description}} \\
\midrule
\endfirsthead
\toprule
{\footnotesize\textbf{Name}} & {\footnotesize\textbf{Size}} & {\footnotesize\textbf{Type}} & {\footnotesize\textbf{Units}} & {\footnotesize\textbf{Description}} \\
\midrule
\endhead
\midrule
\endfoot
\bottomrule
\endlastfoot
Ceres & 1x1 & struct & - & io.importCERES result; required by "ceres", refused by the others \\
Climatology & 1x1 & struct & - & .albedo and .longwave, each an io.importGroopsGrid result; required by "ceresClimatology", refused by the others \\
RotationItrsToGcrs & 3x3xK & double & - & ITRS-\textgreater{}GCRS rotation at the epochs; required by both gridded backends, refused by knocke and pointSource \\
\end{longtable}

\subsubsection*{OUTPUT ARGUMENTS}
\begin{longtable}{@{}>{\footnotesize\raggedright\arraybackslash\ttfamily}p{3.00cm}>{\footnotesize\raggedright\arraybackslash}p{1.50cm}>{\footnotesize\raggedright\arraybackslash}p{2.20cm}>{\footnotesize\raggedright\arraybackslash}p{1.50cm}>{\footnotesize\raggedright\arraybackslash}p{6.40cm}@{}}
\toprule
{\footnotesize\textbf{Name}} & {\footnotesize\textbf{Size}} & {\footnotesize\textbf{Type}} & {\footnotesize\textbf{Units}} & {\footnotesize\textbf{Description}} \\
\midrule
\endfirsthead
\toprule
{\footnotesize\textbf{Name}} & {\footnotesize\textbf{Size}} & {\footnotesize\textbf{Type}} & {\footnotesize\textbf{Units}} & {\footnotesize\textbf{Description}} \\
\midrule
\endhead
\midrule
\endfoot
\bottomrule
\endlastfoot
tiles & 1x1 & struct & - & .backend, .grid, .radius\_m, .M (tiles, 0 for pointSource), .n (3xMxK unit normals of the tile centres, GCRS), .area\_m2 (MxK), .lat\_rad (MxK geocentric latitude of the centres, GCRS; for ceres the ITRS geodetic latitude the grid was read at, for ceresClimatology the ITRS geocentric one), .capAngle\_rad (1xK), .albedo (MxK; 1xK for pointSource), .emissivity (same; [] for ceres and ceresClimatology), .exitanceLw\_W\_m2 (MxK, ceres and ceresClimatology; [] otherwise), .ceres (ceres: .product .version .months (the two bracketing nodes per epoch, 2xK datetime) .timeWeight (1xK)), .climatology (ceresClimatology: .referenceHeight\_m measured from the file, .gridIndex (1xK, into the grids PASSED), .month (1xK datetime of each grid's opening node), .nodes (1x2, first and last node of the grids passed), .gridCount), .daysSinceT0 (1xK), .coefficients, .provenance. \\
\end{longtable}

\subsubsection*{ERRORS}
\begin{description}[leftmargin=1.2em,style=nextline,itemsep=0pt]
\item[\texttt{NonGravForces:\allowbreak{}earthRadiance:\allowbreak{}inputClimatology}] the Climatology grids are missing, mismatched, not a time series, not uniform in longitude, not decreasing in latitude, not one constant height above the ellipsoid, given to a backend that does not read them, or ceresClimatology without a usable 'RotationItrsToGcrs'
\item[\texttt{NonGravForces:\allowbreak{}earthRadiance:\allowbreak{}inputGrid}] grid not two positive integers (every backend but pointSource)
\item[\texttt{NonGravForces:\allowbreak{}earthRadiance:\allowbreak{}inputShape}] r not 3xK or epoch not an epoch struct
\item[\texttt{NonGravForces:\allowbreak{}earthRadiance:\allowbreak{}numericNaN}] NaN/Inf in r
\item[\texttt{NonGravForces:\allowbreak{}earthRadiance:\allowbreak{}inputRange}] \textbar{}r\textbar{} \textless{}= R (inside the Earth, or below the reference level of ceres or ceresClimatology)
\item[\texttt{NonGravForces:\allowbreak{}earthRadiance:\allowbreak{}inputCeres}] ceres without a usable 'Ceres' or 'RotationItrsToGcrs', 'Ceres' given to another backend, or 'RotationItrsToGcrs' given to knocke or pointSource
\item[\texttt{NonGravForces:\allowbreak{}earthRadiance:\allowbreak{}dataRange}] an epoch outside the months the 'Ceres' struct holds, or outside the nodes of the 'Climatology' grids
\item[\texttt{referenceSystem:\allowbreak{}*}] propagated unwrapped
\end{description}

\subsection{\texttt{NonGravForces.env.magneticField}}
\label{api:env.magneticField}
\noindent IGRF main geomagnetic field at Earth-fixed positions, in tesla, Earth-fixed or rotated to any frame.

\begin{lstlisting}
B = NonGravForces.env.magneticField(coeff, rItrs)
B = NonGravForces.env.magneticField(coeff, rItrs, 'Rotation', M)          % e.g. ITRS -> body
[B, info] = NonGravForces.env.magneticField(coeff, rItrs, 'MaxDegree', 1)  % the dipole alone
\end{lstlisting}

\subsubsection*{INPUT ARGUMENTS}
\begin{longtable}{@{}>{\footnotesize\raggedright\arraybackslash\ttfamily}p{3.00cm}>{\footnotesize\raggedright\arraybackslash}p{1.50cm}>{\footnotesize\raggedright\arraybackslash}p{2.20cm}>{\footnotesize\raggedright\arraybackslash}p{1.50cm}>{\footnotesize\raggedright\arraybackslash}p{6.40cm}@{}}
\toprule
{\footnotesize\textbf{Name}} & {\footnotesize\textbf{Size}} & {\footnotesize\textbf{Type}} & {\footnotesize\textbf{Units}} & {\footnotesize\textbf{Description}} \\
\midrule
\endfirsthead
\toprule
{\footnotesize\textbf{Name}} & {\footnotesize\textbf{Size}} & {\footnotesize\textbf{Type}} & {\footnotesize\textbf{Units}} & {\footnotesize\textbf{Description}} \\
\midrule
\endhead
\midrule
\endfoot
\bottomrule
\endlastfoot
coeff & 1x1 & struct & nT & io.igrfCoefficients result evaluated at an epoch (.g, .h (nMax+1)x(nMax+1), indexed [n+1, m+1]; .nMax; .year) \\
rItrs & 3xK & double & m & Positions, Earth-fixed (ITRS), geocentric, \textbar{}r\textbar{} \textgreater{} 0 \\
\end{longtable}

\subsubsection*{NAME-VALUE ARGUMENTS}
\begin{longtable}{@{}>{\footnotesize\raggedright\arraybackslash\ttfamily}p{3.00cm}>{\footnotesize\raggedright\arraybackslash}p{1.50cm}>{\footnotesize\raggedright\arraybackslash}p{2.20cm}>{\footnotesize\raggedright\arraybackslash}p{1.50cm}>{\footnotesize\raggedright\arraybackslash}p{6.40cm}@{}}
\toprule
{\footnotesize\textbf{Name}} & {\footnotesize\textbf{Size}} & {\footnotesize\textbf{Type}} & {\footnotesize\textbf{Units}} & {\footnotesize\textbf{Description}} \\
\midrule
\endfirsthead
\toprule
{\footnotesize\textbf{Name}} & {\footnotesize\textbf{Size}} & {\footnotesize\textbf{Type}} & {\footnotesize\textbf{Units}} & {\footnotesize\textbf{Description}} \\
\midrule
\endhead
\midrule
\endfoot
\bottomrule
\endlastfoot
MaxDegree & 1x1 & double & - & Truncation degree, 1..coeff.nMax; default coeff.nMax \\
Rotation & 3x3\textbar{}\allowbreak{}3x3xK & double & - & Maps Earth-fixed vectors to the frame B is returned in (default []: B stays Earth-fixed) \\
\end{longtable}

\subsubsection*{OUTPUT ARGUMENTS}
\begin{longtable}{@{}>{\footnotesize\raggedright\arraybackslash\ttfamily}p{3.00cm}>{\footnotesize\raggedright\arraybackslash}p{1.50cm}>{\footnotesize\raggedright\arraybackslash}p{2.20cm}>{\footnotesize\raggedright\arraybackslash}p{1.50cm}>{\footnotesize\raggedright\arraybackslash}p{6.40cm}@{}}
\toprule
{\footnotesize\textbf{Name}} & {\footnotesize\textbf{Size}} & {\footnotesize\textbf{Type}} & {\footnotesize\textbf{Units}} & {\footnotesize\textbf{Description}} \\
\midrule
\endfirsthead
\toprule
{\footnotesize\textbf{Name}} & {\footnotesize\textbf{Size}} & {\footnotesize\textbf{Type}} & {\footnotesize\textbf{Units}} & {\footnotesize\textbf{Description}} \\
\midrule
\endhead
\midrule
\endfoot
\bottomrule
\endlastfoot
B & 3xK & double & T & Field vector, Earth-fixed or rotated by 'Rotation' \\
info & 1x1 & struct & - & .itrs\_T (3xK Earth-fixed), .northEastDown\_T (3xK, geocentric local frame: X north, Y east, Z down), .intensity\_T (1xK), .maxDegree, .year (coeff.year), .radius\_m (1xK), .provenance \\
\end{longtable}

\subsubsection*{ERRORS}
\begin{description}[leftmargin=1.2em,style=nextline,itemsep=0pt]
\item[\texttt{NonGravForces:\allowbreak{}magneticField:\allowbreak{}inputCoefficients}] coeff has no g/h at an epoch (pass 'Epoch' to io.igrfCoefficients)
\item[\texttt{NonGravForces:\allowbreak{}magneticField:\allowbreak{}inputShape}] rItrs not 3xK, or Rotation not 3x3 or 3x3xK
\item[\texttt{NonGravForces:\allowbreak{}magneticField:\allowbreak{}inputRange}] \textbar{}r\textbar{} = 0, or MaxDegree outside 1..nMax
\item[\texttt{NonGravForces:\allowbreak{}magneticField:\allowbreak{}numericNaN}] NaN/Inf in rItrs, the coefficients or Rotation
\end{description}

\subsection{\texttt{NonGravForces.env.radiationSources}}
\label{api:env.radiationSources}
\noindent Source list \{dir, P, band\} in the body frame: Sun (with shadow) plus Earth albedo and longwave tiles.

\begin{lstlisting}
[src, info] = NonGravForces.env.radiationSources(epoch, r, q, options)
[src, info] = NonGravForces.env.radiationSources(epoch, r, q, options, 'Ephemeris', eph)
[src, info] = NonGravForces.env.radiationSources(epoch, r, q, options, 'Ceres', cer, 'RotationItrsToGcrs', R)
\end{lstlisting}

\subsubsection*{INPUT ARGUMENTS}
\begin{longtable}{@{}>{\footnotesize\raggedright\arraybackslash\ttfamily}p{3.00cm}>{\footnotesize\raggedright\arraybackslash}p{1.50cm}>{\footnotesize\raggedright\arraybackslash}p{2.20cm}>{\footnotesize\raggedright\arraybackslash}p{1.50cm}>{\footnotesize\raggedright\arraybackslash}p{6.40cm}@{}}
\toprule
{\footnotesize\textbf{Name}} & {\footnotesize\textbf{Size}} & {\footnotesize\textbf{Type}} & {\footnotesize\textbf{Units}} & {\footnotesize\textbf{Description}} \\
\midrule
\endfirsthead
\toprule
{\footnotesize\textbf{Name}} & {\footnotesize\textbf{Size}} & {\footnotesize\textbf{Type}} & {\footnotesize\textbf{Units}} & {\footnotesize\textbf{Description}} \\
\midrule
\endhead
\midrule
\endfoot
\bottomrule
\endlastfoot
epoch & 1x1 & struct & - & referenceSystem epoch struct, K entries \\
r & 3xK & double & m & Satellite position, GCRS \\
q\_\allowbreak{}body2gcrs & 4xK & double & - & Unit quaternions BODY-\textgreater{}GCRS, scalar first \\
options & 1x1 & NonGravForces.\allowbreak{}Options & - & srp.*, erp.*, sun.source are read \\
\end{longtable}

\subsubsection*{NAME-VALUE ARGUMENTS}
\begin{longtable}{@{}>{\footnotesize\raggedright\arraybackslash\ttfamily}p{3.00cm}>{\footnotesize\raggedright\arraybackslash}p{1.50cm}>{\footnotesize\raggedright\arraybackslash}p{2.20cm}>{\footnotesize\raggedright\arraybackslash}p{1.50cm}>{\footnotesize\raggedright\arraybackslash}p{6.40cm}@{}}
\toprule
{\footnotesize\textbf{Name}} & {\footnotesize\textbf{Size}} & {\footnotesize\textbf{Type}} & {\footnotesize\textbf{Units}} & {\footnotesize\textbf{Description}} \\
\midrule
\endfirsthead
\toprule
{\footnotesize\textbf{Name}} & {\footnotesize\textbf{Size}} & {\footnotesize\textbf{Type}} & {\footnotesize\textbf{Units}} & {\footnotesize\textbf{Description}} \\
\midrule
\endhead
\midrule
\endfoot
\bottomrule
\endlastfoot
Ephemeris & 1x1 & struct & - & referenceSystem.readEphemeris result; required for sun.source "rs.de440". Default []. \\
Sun & 3xK & double & m & Caller-supplied Earth -\textgreater{} Sun vector (GCRS) that REPLACES the ephemeris (info.sunSource "caller"): for tests and external ephemerides. Default []. \\
Moon & 3xK & double & m & Caller-supplied Earth -\textgreater{} Moon vector, used with 'Sun' when srp.lunarEclipse is true. Default []. \\
EclipticNormal & 3xK & double & - & Caller-supplied unit ecliptic normal, needed with 'Sun' when srp.eclipse is solaarsCF (the ephemeris path measures it as r\_sun x v\_sun). Default []. \\
Ceres & 1x1 & struct & - & io.importCERES result; erp.backend ceres only. Default []. \\
Climatology & 1x1 & struct & - & .albedo and .longwave io.importGroopsGrid results; erp.backend ceresClimatology only. Default []. \\
RotationItrsToGcrs & 3x3xK & double & - & ITRS-\textgreater{}GCRS rotation; erp.backend ceres and ceresClimatology only. Default []. \\
\end{longtable}

\subsubsection*{OUTPUT ARGUMENTS}
\begin{longtable}{@{}>{\footnotesize\raggedright\arraybackslash\ttfamily}p{3.00cm}>{\footnotesize\raggedright\arraybackslash}p{1.50cm}>{\footnotesize\raggedright\arraybackslash}p{2.20cm}>{\footnotesize\raggedright\arraybackslash}p{1.50cm}>{\footnotesize\raggedright\arraybackslash}p{6.40cm}@{}}
\toprule
{\footnotesize\textbf{Name}} & {\footnotesize\textbf{Size}} & {\footnotesize\textbf{Type}} & {\footnotesize\textbf{Units}} & {\footnotesize\textbf{Description}} \\
\midrule
\endfirsthead
\toprule
{\footnotesize\textbf{Name}} & {\footnotesize\textbf{Size}} & {\footnotesize\textbf{Type}} & {\footnotesize\textbf{Units}} & {\footnotesize\textbf{Description}} \\
\midrule
\endhead
\midrule
\endfoot
\bottomrule
\endlastfoot
src & 1x1 & struct & - & .dir (3xMxK, body, unit, satellite -\textgreater{} source), .P (MxK, Pa), .band (1xM string "vis"\textbar{}"ir"), .nu (1xK Sun shadow factor), .M, .K, .frame "body", .meta (see info). \\
info & 1x1 & struct & - & .sunSource (Options value or "caller"), .sunDistance\_m (1xK), .sunUnit\_gcrs (3xK), .eclipseModel, .eclipse (eclipseFactor's info), .lunarEclipse, .erpBackend, .grid, .tiles (earthRadiance struct or []), .tsi\_W\_m2, .earthRadius\_m, .speedOfLight\_m\_s, .provenance. \\
\end{longtable}

\subsubsection*{ERRORS}
\begin{description}[leftmargin=1.2em,style=nextline,itemsep=0pt]
\item[\texttt{NonGravForces:\allowbreak{}radiationSources:\allowbreak{}inputShape}] epoch not an epoch struct, r not 3xK, q not 4xK, or K differs from epoch; a caller 'Sun' not a finite 3xK, or given without 'EclipticNormal' (solaarsCF) or without 'Moon' (srp.lunarEclipse)
\item[\texttt{NonGravForces:\allowbreak{}radiationSources:\allowbreak{}numericNaN}] NaN/Inf in r or q
\item[\texttt{NonGravForces:\allowbreak{}radiationSources:\allowbreak{}depKernel}] sun.source rs.de440 without 'Ephemeris'
\item[\texttt{NonGravForces:\allowbreak{}radiationSources:\allowbreak{}inputCeres}] 'Ceres', 'Climatology' or 'RotationItrsToGcrs' given with erp.backend none
\item[\texttt{NonGravForces:\allowbreak{}earthRadiance:\allowbreak{}inputCeres}] ceres inputs missing, or given to another backend
\item[\texttt{NonGravForces:\allowbreak{}earthRadiance:\allowbreak{}inputClimatology}] ceresClimatology inputs missing, or given to another backend
\item[\texttt{NonGravForces:\allowbreak{}eclipseFactor:\allowbreak{}*, NonGravForces:\allowbreak{}earthRadiance:\allowbreak{}*}] from the callees
\item[\texttt{referenceSystem:\allowbreak{}*}] propagated unwrapped
\end{description}

\section{L3 \texttt{geom} --- panel geometry, attitude and occlusion}

\subsection{\texttt{NonGravForces.geom.MacroModel}}
\label{api:geom.MacroModel}
\noindent Panel ("macro") geometry of a spacecraft, frozen and validated once.

\begin{lstlisting}
mm = NonGravForces.geom.MacroModel(panels, 'Name', "GRACE-C", 'Inertia', I)
s  = mm.silhouetteDefect()     % 3x1; zero for any closed body
\end{lstlisting}

\subsubsection*{METHODS}
\begin{lstlisting}
s = silhouetteDefect(obj)  sum_k A_k n_k, the area-weighted sum of the outward normals, 3x1
                           in the body frame. It is ZERO for ANY closed body, convex or not
                           (the divergence theorem), so it is the one-line check that a panel
                           set closes: a non-zero value means the areas do not enclose a
                           volume -- a missing face, a normal pointing inward, or printed
                           areas that do not close on each other (the GRACE/GRACE-FO tables:
                           0.14 m^2 along z). It is a DIAGNOSTIC and no force reads it;
                           geom.graceFO's and geom.grace's tests assert its exact value.
\end{lstlisting}
\begin{lstlisting}
PANELS is a struct array, one element per surface element, with fields
  name        string     label
  kind        string     "plate" (the only kind so far; "cylinder" and
                         "sphere" are reserved until a kernel consumes them)
  area_m2     1x1 double area (the authority for every force)
  normal      3x1 double outward unit normal, body frame
  centroid_m  3x1 double application point, body frame (torques,
                         self-shadowing); NaN if unknown
  absVis, specVis, diffVis   visible band, must sum to 1
  emisIR, specIR, diffIR     infrared band, must sum to 1
  group       string     "bus" | "array" | "boom" (informational)
(a per-panel accommodation coefficient `alpha` was required until 2026-09-15 and read
by nothing; it is no longer required and, when present, ignored: HANDOVER O-23 closed
by removal. The accommodation is the environment's: kernel.accommodation.)
and OPTIONALLY (added 2026-09-07, LP-4)
  vertices_m  3xV double corners of the plate polygon, body frame, wound
                         counter-clockwise about the outward normal, V >= 3;
                         used ONLY by geom.selfShadow (occlusion). Absent or
                         empty: no ray casting possible for this model.
  heatCapacity_J_m2K 1x1 double heat capacity per unit area of the wall
                         behind the surface; used ONLY by geom.thermalState
                         "transient". NaN if unknown.
  internalShare 1x1 double share of the internal dissipation radiated through
                         this plate; given on every plate or on none, >= 0,
                         summing to 1 (geom.thermalState 'InternalShare').
\end{lstlisting}
NAME-VALUE 'Mass\_kg' (added 2026-09-15, LP-C F10): the spacecraft mass a Model divides by when it is given no 'Mass' of its own; NaN (the default) means unknown, and a Model built on such a macro model refuses to guess.

Validation at construction is the ONLY validation: kernels trust the packed arrays. Energy closure \textbar{}a+s+d-1\textbar{} \textless{}= 1e-6 in both bands, unit normals to 1e-5 (the handbook prints sin/cos 50 deg to six digits, whose norm is 1 - 7.3e-7, measured): the packed N holds the normals RENORMALISED to unit length (decided 2026-09-07, HANDOVER O-2: the printed digits are a truncation of a unit vector, so one renormalisation here restores the authority's intent and no kernel renormalises), while panels(k).normal keeps the value as given; positive areas, finite centroids or all-NaN. Vertices, when given: planar to 1 mm about the normal, polygon area within 5 \% of area\_m2 (the polygons are derived from exterior dimensions, the areas are the printed authority), and wound counter-clockwise (a clockwise polygon is refused, never flipped), and a finite centroid\_m within 5 cm of the polygon's area centroid (R1).


\subsubsection*{PROPERTIES (read-only)}
\begin{lstlisting}
name        string
n           double     number of elements
panels      struct     as given, validated
A           1xN double areas
N           3xN double normals
C           3xN double centroids (NaN allowed)
optVis      3xN double [abs; spec; diff] visible
optIR       3xN double [emis; spec; diff] infrared
V           1xN cell   vertices (3xV each) or empty cells
hasVertices logical    every panel carries vertices
heatCapacity 1xN double J/(m^2 K), NaN where unknown
internalShare 1xN double internal dissipation shares, NaN where not given
mass_kg     1x1 double spacecraft mass, NaN if not given
inertia     3x3 double kg m^2 (NaN if not given)
centreOfMass    3x1 double m, body frame (origin = ACC proof mass for GRACE-FO)
proofMassOffset 3x1 double m, CG -> proof mass (CoM-coupling term, LP-7)
reference   string     source of the numbers
\end{lstlisting}

\subsubsection*{NAME-VALUE ARGUMENTS}
The properties above are read-only once built; these are how they are SET at construction. The names are PascalCase where the properties they fill are camelCase, and three of them were described nowhere in this help until 2026-09-18 -- the property was documented and the way to fill it was not, which is a gap a reader cannot close by guessing.

\begin{longtable}{@{}>{\footnotesize\raggedright\arraybackslash\ttfamily}p{3.00cm}>{\footnotesize\raggedright\arraybackslash}p{1.50cm}>{\footnotesize\raggedright\arraybackslash}p{2.20cm}>{\footnotesize\raggedright\arraybackslash}p{1.50cm}>{\footnotesize\raggedright\arraybackslash}p{6.40cm}@{}}
\toprule
{\footnotesize\textbf{Name}} & {\footnotesize\textbf{Size}} & {\footnotesize\textbf{Type}} & {\footnotesize\textbf{Units}} & {\footnotesize\textbf{Description}} \\
\midrule
\endfirsthead
\toprule
{\footnotesize\textbf{Name}} & {\footnotesize\textbf{Size}} & {\footnotesize\textbf{Type}} & {\footnotesize\textbf{Units}} & {\footnotesize\textbf{Description}} \\
\midrule
\endhead
\midrule
\endfoot
\bottomrule
\endlastfoot
Name & 1x1 & string & - & Model label, e.g. "GRACE-C"; default "" \\
Inertia & 3x3 & double & kg m\textasciicircum{}2 & Inertia tensor about the CoM, body frame; default NaN(3). Required by kernel.bodyTorques' gravity-gradient term. \\
CentreOfMass & 3x1 & double & m & CoM in the body frame, whose ORIGIN for GRACE-FO is the accelerometer proof mass; default zeros. Read by the torque reduction as the point moments are taken about. \\
ProofMassOffset & 3x1 & double & m & CG -\textgreater{} proof mass, body frame; default zeros. This is the `dr` of sensor.Electrostatic's CoM-coupling term, and it is NOT the negative of CentreOfMass in general: the two are independent statements about the same geometry, and a model that carries one need not carry the other. \\
Reference & 1x1 & string & - & Provenance of the numbers, e.g. "JPL D-56935 Tab.4"; default "". Carried into provenance, never computed with. \\
Mass\_\allowbreak{}kg & 1x1 & double & kg & Spacecraft mass; default NaN (a Model then needs 'Mass'). MAS1B is COLD-GAS TANK mass, not this: see io.massModel. \\
\end{longtable}

\subsubsection*{ERRORS}
\begin{description}[leftmargin=1.2em,style=nextline,itemsep=0pt]
\item[\texttt{NonGravForces:\allowbreak{}MacroModel:\allowbreak{}inputPanels}] missing field; an area not finite and positive; a heatCapacity\_J\_m2K \textless{}= 0; internalShare or Mass\_kg invalid
\item[\texttt{NonGravForces:\allowbreak{}MacroModel:\allowbreak{}inputEnergy}] a+s+d \textasciitilde{}= 1 in a band
\item[\texttt{NonGravForces:\allowbreak{}MacroModel:\allowbreak{}inputNormal}] non-unit normal
\item[\texttt{NonGravForces:\allowbreak{}MacroModel:\allowbreak{}inputVertices}] vertices not 3xV, not planar, area off by \textgreater{} 5 \%, wound clockwise, or centroid\_m \textgreater{} 5 cm from the polygon centroid
\end{description}

\subsection{\texttt{NonGravForces.geom.graceFO}}
\label{api:geom.graceFO}
\noindent GRACE-FO macro model from JPL D-56935 (Wen et al. 2019), Tables 4 and 5.

\begin{lstlisting}
mm = NonGravForces.geom.graceFO("C")
mm = NonGravForces.geom.graceFO("D", 'SolarArrayOperating', false)
\end{lstlisting}

\subsubsection*{INPUT ARGUMENTS}
\begin{longtable}{@{}>{\footnotesize\raggedright\arraybackslash\ttfamily}p{3.00cm}>{\footnotesize\raggedright\arraybackslash}p{1.50cm}>{\footnotesize\raggedright\arraybackslash}p{2.20cm}>{\footnotesize\raggedright\arraybackslash}p{1.50cm}>{\footnotesize\raggedright\arraybackslash}p{6.40cm}@{}}
\toprule
{\footnotesize\textbf{Name}} & {\footnotesize\textbf{Size}} & {\footnotesize\textbf{Type}} & {\footnotesize\textbf{Units}} & {\footnotesize\textbf{Description}} \\
\midrule
\endfirsthead
\toprule
{\footnotesize\textbf{Name}} & {\footnotesize\textbf{Size}} & {\footnotesize\textbf{Type}} & {\footnotesize\textbf{Units}} & {\footnotesize\textbf{Description}} \\
\midrule
\endhead
\midrule
\endfoot
\bottomrule
\endlastfoot
satellite & 1x1 & string & - & "C" (FM1, 601.214 kg) or "D" (FM2, 601.209 kg) \\
\end{longtable}

\subsubsection*{NAME-VALUE ARGUMENTS}
\begin{longtable}{@{}>{\footnotesize\raggedright\arraybackslash\ttfamily}p{3.00cm}>{\footnotesize\raggedright\arraybackslash}p{1.50cm}>{\footnotesize\raggedright\arraybackslash}p{2.20cm}>{\footnotesize\raggedright\arraybackslash}p{7.90cm}@{}}
\toprule
{\footnotesize\textbf{Name}} & {\footnotesize\textbf{Size}} & {\footnotesize\textbf{Type}} & {\footnotesize\textbf{Description}} \\
\midrule
\endfirsthead
\toprule
{\footnotesize\textbf{Name}} & {\footnotesize\textbf{Size}} & {\footnotesize\textbf{Type}} & {\footnotesize\textbf{Description}} \\
\midrule
\endhead
\midrule
\endfoot
\bottomrule
\endlastfoot
SolarArrayOperating & 1x1 & logical & true: absVis 0.65 / diffVis 0.30 (closes). false: absVis 0.72 with diffVis 0.23 so that the budget closes (the printed 0.30 would sum to 1.07 - a handbook defect). \\
\end{longtable}

\subsubsection*{OUTPUT ARGUMENTS}
\begin{longtable}{@{}>{\footnotesize\raggedright\arraybackslash\ttfamily}p{3.00cm}>{\footnotesize\raggedright\arraybackslash}p{1.50cm}>{\footnotesize\raggedright\arraybackslash}p{10.10cm}@{}}
\toprule
{\footnotesize\textbf{Name}} & {\footnotesize\textbf{Size}} & {\footnotesize\textbf{Description}} \\
\midrule
\endfirsthead
\toprule
{\footnotesize\textbf{Name}} & {\footnotesize\textbf{Size}} & {\footnotesize\textbf{Description}} \\
\midrule
\endhead
\midrule
\endfoot
\bottomrule
\endlastfoot
mm & 1x1 & NonGravForces.geom.MacroModel with .inertia (Table 4), .centreOfMass and .proofMassOffset (Table 4, CG and CG -\textgreater{} ACC), .V (plate polygons), .heatCapacity, .mass\_kg (the handbook mass of that satellite, NonGravForces.const.graceFO) \\
\end{longtable}

\subsubsection*{ERRORS}
\begin{description}[leftmargin=1.2em,style=nextline,itemsep=0pt]
\item[\texttt{MATLAB:\allowbreak{}validators:\allowbreak{}mustBeMember}] satellite is not "C" or "D" (arguments block)
\end{description}
plus the contract errors of geom.MacroModel, which validates the table.


\subsection{\texttt{NonGravForces.geom.grace}}
\label{api:geom.grace}
\noindent GRACE-A/B macro model from the GRACE Product Specification Document (Bettadpur 2012), Sec. I.5.3.

\begin{lstlisting}
mm = NonGravForces.geom.grace("A")
mm = NonGravForces.geom.grace("B", 'SolarArrayOperating', false)
\end{lstlisting}

\subsubsection*{INPUT ARGUMENTS}
\begin{longtable}{@{}>{\footnotesize\raggedright\arraybackslash\ttfamily}p{3.00cm}>{\footnotesize\raggedright\arraybackslash}p{1.50cm}>{\footnotesize\raggedright\arraybackslash}p{2.20cm}>{\footnotesize\raggedright\arraybackslash}p{1.50cm}>{\footnotesize\raggedright\arraybackslash}p{6.40cm}@{}}
\toprule
{\footnotesize\textbf{Name}} & {\footnotesize\textbf{Size}} & {\footnotesize\textbf{Type}} & {\footnotesize\textbf{Units}} & {\footnotesize\textbf{Description}} \\
\midrule
\endfirsthead
\toprule
{\footnotesize\textbf{Name}} & {\footnotesize\textbf{Size}} & {\footnotesize\textbf{Type}} & {\footnotesize\textbf{Units}} & {\footnotesize\textbf{Description}} \\
\midrule
\endhead
\midrule
\endfoot
\bottomrule
\endlastfoot
satellite & 1x1 & string & - & "A" (GRACE-1, FM-1) or "B" (GRACE-2, FM-2) \\
\end{longtable}

\subsubsection*{NAME-VALUE ARGUMENTS}
\begin{longtable}{@{}>{\footnotesize\raggedright\arraybackslash\ttfamily}p{3.00cm}>{\footnotesize\raggedright\arraybackslash}p{1.50cm}>{\footnotesize\raggedright\arraybackslash}p{2.20cm}>{\footnotesize\raggedright\arraybackslash}p{7.90cm}@{}}
\toprule
{\footnotesize\textbf{Name}} & {\footnotesize\textbf{Size}} & {\footnotesize\textbf{Type}} & {\footnotesize\textbf{Description}} \\
\midrule
\endfirsthead
\toprule
{\footnotesize\textbf{Name}} & {\footnotesize\textbf{Size}} & {\footnotesize\textbf{Type}} & {\footnotesize\textbf{Description}} \\
\midrule
\endhead
\midrule
\endfoot
\bottomrule
\endlastfoot
SolarArrayOperating & 1x1 & logical & true: absVis 0.65 / diffVis 0.30 (closes). false: absVis 0.72 with diffVis 0.23 so that the budget closes (the printed 0.30 would sum to 1.07 - a handbook defect). \\
\end{longtable}

\subsubsection*{OUTPUT ARGUMENTS}
\begin{longtable}{@{}>{\footnotesize\raggedright\arraybackslash\ttfamily}p{3.00cm}>{\footnotesize\raggedright\arraybackslash}p{1.50cm}>{\footnotesize\raggedright\arraybackslash}p{10.10cm}@{}}
\toprule
{\footnotesize\textbf{Name}} & {\footnotesize\textbf{Size}} & {\footnotesize\textbf{Description}} \\
\midrule
\endfirsthead
\toprule
{\footnotesize\textbf{Name}} & {\footnotesize\textbf{Size}} & {\footnotesize\textbf{Description}} \\
\midrule
\endhead
\midrule
\endfoot
\bottomrule
\endlastfoot
mm & 1x1 & NonGravForces.geom.MacroModel with .V (plate polygons), .heatCapacity, .mass\_kg (the launch mass, NonGravForces.const.grace), .inertia NaN(3) \\
\end{longtable}

\subsubsection*{ERRORS}
\begin{description}[leftmargin=1.2em,style=nextline,itemsep=0pt]
\item[\texttt{MATLAB:\allowbreak{}validators:\allowbreak{}mustBeMember}] satellite is not "A" or "B" (arguments block)
\end{description}
plus the contract errors of geom.MacroModel, which validates the table.


\subsection{\texttt{NonGravForces.geom.boxWing}}
\label{api:geom.boxWing}
\noindent Generic box-wing macro model: a rectangular bus of six plates and optional solar-array plates.

\begin{lstlisting}
mm = NonGravForces.geom.boxWing([Lx Ly Lz], optics)
mm = NonGravForces.geom.boxWing([Lx Ly Lz], optics, 'Array', arr, 'Name', "MySat", 'Mass_kg', 480)
\end{lstlisting}

\subsubsection*{INPUT ARGUMENTS}
\begin{longtable}{@{}>{\footnotesize\raggedright\arraybackslash\ttfamily}p{3.00cm}>{\footnotesize\raggedright\arraybackslash}p{1.50cm}>{\footnotesize\raggedright\arraybackslash}p{2.20cm}>{\footnotesize\raggedright\arraybackslash}p{1.50cm}>{\footnotesize\raggedright\arraybackslash}p{6.40cm}@{}}
\toprule
{\footnotesize\textbf{Name}} & {\footnotesize\textbf{Size}} & {\footnotesize\textbf{Type}} & {\footnotesize\textbf{Units}} & {\footnotesize\textbf{Description}} \\
\midrule
\endfirsthead
\toprule
{\footnotesize\textbf{Name}} & {\footnotesize\textbf{Size}} & {\footnotesize\textbf{Type}} & {\footnotesize\textbf{Units}} & {\footnotesize\textbf{Description}} \\
\midrule
\endhead
\midrule
\endfoot
\bottomrule
\endlastfoot
dims & 1x3 & double & m & [Lx Ly Lz] edge lengths of the bus, positive \\
optics & 1x1 & struct & - & .bus and (if 'Array' is given) .array, each with fields vis (1x3 [abs spec diff], sums to 1) and ir (1x3 [emis spec diff], sums to 1) \\
\end{longtable}

\subsubsection*{NAME-VALUE ARGUMENTS}
\begin{longtable}{@{}>{\footnotesize\raggedright\arraybackslash\ttfamily}p{3.00cm}>{\footnotesize\raggedright\arraybackslash}p{1.50cm}>{\footnotesize\raggedright\arraybackslash}p{2.20cm}>{\footnotesize\raggedright\arraybackslash}p{1.50cm}>{\footnotesize\raggedright\arraybackslash}p{6.40cm}@{}}
\toprule
{\footnotesize\textbf{Name}} & {\footnotesize\textbf{Size}} & {\footnotesize\textbf{Type}} & {\footnotesize\textbf{Units}} & {\footnotesize\textbf{Description}} \\
\midrule
\endfirsthead
\toprule
{\footnotesize\textbf{Name}} & {\footnotesize\textbf{Size}} & {\footnotesize\textbf{Type}} & {\footnotesize\textbf{Units}} & {\footnotesize\textbf{Description}} \\
\midrule
\endhead
\midrule
\endfoot
\bottomrule
\endlastfoot
Array & 1xP & struct & - & per array plate: .area\_m2, .normal (3x1 unit), .centroid\_m (3x1), optional .name; default none \\
Name & 1x1 & string & - & Model name (default "boxWing") \\
Mass\_\allowbreak{}kg & 1x1 & double & kg & Spacecraft mass; default NaN (unknown: a Model then needs 'Mass') \\
HeatCapacity\_\allowbreak{}J\_\allowbreak{}m2K & 1x2 & double & J/\allowbreak{}(m\textasciicircum{}2 K) & [bus array] heat capacity per area for thermalState transient; default [NaN NaN] \\
\end{longtable}

\subsubsection*{OUTPUT ARGUMENTS}
\begin{longtable}{@{}>{\footnotesize\raggedright\arraybackslash\ttfamily}p{3.00cm}>{\footnotesize\raggedright\arraybackslash}p{1.50cm}>{\footnotesize\raggedright\arraybackslash}p{2.20cm}>{\footnotesize\raggedright\arraybackslash}p{1.50cm}>{\footnotesize\raggedright\arraybackslash}p{6.40cm}@{}}
\toprule
{\footnotesize\textbf{Name}} & {\footnotesize\textbf{Size}} & {\footnotesize\textbf{Type}} & {\footnotesize\textbf{Units}} & {\footnotesize\textbf{Description}} \\
\midrule
\endfirsthead
\toprule
{\footnotesize\textbf{Name}} & {\footnotesize\textbf{Size}} & {\footnotesize\textbf{Type}} & {\footnotesize\textbf{Units}} & {\footnotesize\textbf{Description}} \\
\midrule
\endhead
\midrule
\endfoot
\bottomrule
\endlastfoot
mm & 1x1 & NonGravForces.\allowbreak{}geom.\allowbreak{}MacroModel &  & 6 + P plates, polygons on the bus plates \\
\end{longtable}

\subsubsection*{ERRORS}
\begin{description}[leftmargin=1.2em,style=nextline,itemsep=0pt]
\item[\texttt{NonGravForces:\allowbreak{}boxWing:\allowbreak{}inputDims}] dims not three positive finite lengths
\item[\texttt{NonGravForces:\allowbreak{}boxWing:\allowbreak{}inputOptics}] optics lacks .bus.vis/.bus.ir (or .array.vis/.array.ir when arrays are given), or a vis/ir triple does not have three elements
\item[\texttt{NonGravForces:\allowbreak{}boxWing:\allowbreak{}inputArray}] an array entry lacks area\_m2/normal/centroid\_m, its area is not a positive scalar, or its normal is not 3-element unit (1e-5)
\item[\texttt{NonGravForces:\allowbreak{}MacroModel:\allowbreak{}*}] from the constructor: inputEnergy when a vis/ir triple does not sum to 1, inputPanels for a non-positive HeatCapacity\_J\_m2K or Mass\_kg
\end{description}

\subsection{\texttt{NonGravForces.geom.selfShadow}}
\label{api:geom.selfShadow}
\noindent Unobstructed fraction of every plate toward body-frame directions: ray cast or cube-map table.

\begin{lstlisting}
[f, info] = NonGravForces.geom.selfShadow(macro, dir_body)                       % ray cast
[f, info] = NonGravForces.geom.selfShadow(macro, dir_body, 'Method', "lut", 'LutResolution', 16)
lut = NonGravForces.geom.selfShadow(macro, [], 'Method', "lut")                   % table only
f = NonGravForces.geom.selfShadow(macro, dir_body, 'Method', "lut", 'Lut', lut)   % reuse a table
\end{lstlisting}

\subsubsection*{INPUT ARGUMENTS}
\begin{longtable}{@{}>{\footnotesize\raggedright\arraybackslash\ttfamily}p{3.00cm}>{\footnotesize\raggedright\arraybackslash}p{1.50cm}>{\footnotesize\raggedright\arraybackslash}p{2.20cm}>{\footnotesize\raggedright\arraybackslash}p{1.50cm}>{\footnotesize\raggedright\arraybackslash}p{6.40cm}@{}}
\toprule
{\footnotesize\textbf{Name}} & {\footnotesize\textbf{Size}} & {\footnotesize\textbf{Type}} & {\footnotesize\textbf{Units}} & {\footnotesize\textbf{Description}} \\
\midrule
\endfirsthead
\toprule
{\footnotesize\textbf{Name}} & {\footnotesize\textbf{Size}} & {\footnotesize\textbf{Type}} & {\footnotesize\textbf{Units}} & {\footnotesize\textbf{Description}} \\
\midrule
\endhead
\midrule
\endfoot
\bottomrule
\endlastfoot
macro & 1x1 & NonGravForces.\allowbreak{}geom.\allowbreak{}MacroModel & - & With vertices on every plate (hasVertices) \\
dir\_\allowbreak{}body & 3xK & double & - & Unit directions from the body, body frame; [] with Method "lut" returns only the table \\
\end{longtable}

\subsubsection*{NAME-VALUE ARGUMENTS}
\begin{longtable}{@{}>{\footnotesize\raggedright\arraybackslash\ttfamily}p{3.00cm}>{\footnotesize\raggedright\arraybackslash}p{1.50cm}>{\footnotesize\raggedright\arraybackslash}p{2.20cm}>{\footnotesize\raggedright\arraybackslash}p{1.50cm}>{\footnotesize\raggedright\arraybackslash}p{6.40cm}@{}}
\toprule
{\footnotesize\textbf{Name}} & {\footnotesize\textbf{Size}} & {\footnotesize\textbf{Type}} & {\footnotesize\textbf{Units}} & {\footnotesize\textbf{Description}} \\
\midrule
\endfirsthead
\toprule
{\footnotesize\textbf{Name}} & {\footnotesize\textbf{Size}} & {\footnotesize\textbf{Type}} & {\footnotesize\textbf{Units}} & {\footnotesize\textbf{Description}} \\
\midrule
\endhead
\midrule
\endfoot
\bottomrule
\endlastfoot
Method & 1x1 & string & - & "raycast" (default) \textbar{} "lut" \\
Samples & 1x1 & double & - & grid points per plate edge for the ray cast (default 40) \\
LutResolution & 1x1 & double & - & cells per cube-face edge, integer in [1, 256]; default: the resolution of 'Lut' if given, else 16 (5.6 deg) \\
Lut & 1x1 & struct & - & a previously returned table for this macro model; its own sample count is kept and reported in info.samples (Samples is not compared with it) \\
Plate & 1x1 & double & - & evaluate this one plate only (index into the macro model): f is then 1xK; default [] (every plate) \\
\end{longtable}

\subsubsection*{OUTPUT ARGUMENTS}
\begin{longtable}{@{}>{\footnotesize\raggedright\arraybackslash\ttfamily}p{3.00cm}>{\footnotesize\raggedright\arraybackslash}p{1.50cm}>{\footnotesize\raggedright\arraybackslash}p{2.20cm}>{\footnotesize\raggedright\arraybackslash}p{1.50cm}>{\footnotesize\raggedright\arraybackslash}p{6.40cm}@{}}
\toprule
{\footnotesize\textbf{Name}} & {\footnotesize\textbf{Size}} & {\footnotesize\textbf{Type}} & {\footnotesize\textbf{Units}} & {\footnotesize\textbf{Description}} \\
\midrule
\endfirsthead
\toprule
{\footnotesize\textbf{Name}} & {\footnotesize\textbf{Size}} & {\footnotesize\textbf{Type}} & {\footnotesize\textbf{Units}} & {\footnotesize\textbf{Description}} \\
\midrule
\endhead
\midrule
\endfoot
\bottomrule
\endlastfoot
f & NxK & double & - & Unobstructed fraction per plate and direction, in [0, 1]; or, with dir\_body = [] and Method "lut", the table struct itself \\
info & 1x1 & struct & - & .method, .samples, .lut (struct .f N x 6(n+1)\textasciicircum{}2, .nodes 3 x 6(n+1)\textasciicircum{}2 node directions face by face, .resolution, .samples, .macroName, .geometry; only for "lut"), .provenance (built only when info is requested) \\
\end{longtable}

\subsubsection*{ERRORS}
\begin{description}[leftmargin=1.2em,style=nextline,itemsep=0pt]
\item[\texttt{NonGravForces:\allowbreak{}selfShadow:\allowbreak{}inputMacro}] not a MacroModel, or without vertices on every plate
\item[\texttt{NonGravForces:\allowbreak{}selfShadow:\allowbreak{}inputShape}] dir\_body not 3xK, not unit, or not finite; 'Plate' not a plate index
\item[\texttt{NonGravForces:\allowbreak{}selfShadow:\allowbreak{}inputResolution}] LutResolution not an integer in [1, 256]
\item[\texttt{NonGravForces:\allowbreak{}selfShadow:\allowbreak{}inputLut}] 'Lut' built for another geometry or resolution (or not a table)
\end{description}

\subsection{\texttt{NonGravForces.geom.flowShadow}}
\label{api:geom.flowShadow}
\noindent Self-shadowing of a free-molecular gas stream: flux-weighted occlusion ratios per plate, epoch and species.

\begin{lstlisting}
S = NonGravForces.geom.flowShadow(macro, vrel_body, T_K, 'Method', "lut", 'Lut', lut)
S = NonGravForces.geom.flowShadow(macro, vrel_body, T_K)                 % exact ray cast (slow)
S = NonGravForces.geom.flowShadow(macro, vrel_body, [], 'Lut', lut, 'Method', "lut")   % hyperthermal only
S = NonGravForces.geom.flowShadow(macro, vrel_body, T_K, 'Method', "table", 'Lut', lut)  % tabulated "lut"
tb = NonGravForces.geom.flowShadow(macro, [], [], 'Method', "table", 'Lut', lut)       % the table itself
[S, info] = NonGravForces.geom.flowShadow(___, 'Integrals', true)
\end{lstlisting}

\subsubsection*{INPUT ARGUMENTS}
\begin{longtable}{@{}>{\footnotesize\raggedright\arraybackslash\ttfamily}p{3.00cm}>{\footnotesize\raggedright\arraybackslash}p{1.50cm}>{\footnotesize\raggedright\arraybackslash}p{2.20cm}>{\footnotesize\raggedright\arraybackslash}p{1.50cm}>{\footnotesize\raggedright\arraybackslash}p{6.40cm}@{}}
\toprule
{\footnotesize\textbf{Name}} & {\footnotesize\textbf{Size}} & {\footnotesize\textbf{Type}} & {\footnotesize\textbf{Units}} & {\footnotesize\textbf{Description}} \\
\midrule
\endfirsthead
\toprule
{\footnotesize\textbf{Name}} & {\footnotesize\textbf{Size}} & {\footnotesize\textbf{Type}} & {\footnotesize\textbf{Units}} & {\footnotesize\textbf{Description}} \\
\midrule
\endhead
\midrule
\endfoot
\bottomrule
\endlastfoot
macro & 1x1 & NonGravForces.\allowbreak{}geom.\allowbreak{}MacroModel & - & With vertices on every plate \\
vrel\_\allowbreak{}body & 3xK & double & m/\allowbreak{}s & Satellite velocity relative to the gas, body frame, \textbar{}v\textbar{} \textgreater{} 0 \\
T\_\allowbreak{}K & 1xK & double & K & Gas kinetic temperature at the satellite, \textgreater{} 0; [] for the hyperthermal limit only (S.beam) \\
\end{longtable}

\subsubsection*{NAME-VALUE ARGUMENTS}
\begin{longtable}{@{}>{\footnotesize\raggedright\arraybackslash\ttfamily}p{3.00cm}>{\footnotesize\raggedright\arraybackslash}p{1.50cm}>{\footnotesize\raggedright\arraybackslash}p{2.20cm}>{\footnotesize\raggedright\arraybackslash}p{1.50cm}>{\footnotesize\raggedright\arraybackslash}p{6.40cm}@{}}
\toprule
{\footnotesize\textbf{Name}} & {\footnotesize\textbf{Size}} & {\footnotesize\textbf{Type}} & {\footnotesize\textbf{Units}} & {\footnotesize\textbf{Description}} \\
\midrule
\endfirsthead
\toprule
{\footnotesize\textbf{Name}} & {\footnotesize\textbf{Size}} & {\footnotesize\textbf{Type}} & {\footnotesize\textbf{Units}} & {\footnotesize\textbf{Description}} \\
\midrule
\endhead
\midrule
\endfoot
\bottomrule
\endlastfoot
Method & 1x1 & string & - & "raycast" (default, exact to the sampling) \textbar{} "lut" \textbar{} "table" (the "lut" ratios read from a table; see THE TABLE) \\
Samples & 1x1 & double & - & ray-cast samples per plate edge (default 40) \\
Lut & 1x1 & struct & - & a geom.selfShadow table for this macro model (required for "lut" and "table") \\
Quadrature & 1x2 & double & - & [nodes per chi segment, nodes per arc], default [12 16] \\
Integrals & 1x1 & logical & - & also return the unobstructed coefficient integrals in info (default false) \\
\end{longtable}

\subsubsection*{OUTPUT ARGUMENTS}
\begin{longtable}{@{}>{\footnotesize\raggedright\arraybackslash\ttfamily}p{3.00cm}>{\footnotesize\raggedright\arraybackslash}p{1.50cm}>{\footnotesize\raggedright\arraybackslash}p{2.20cm}>{\footnotesize\raggedright\arraybackslash}p{1.50cm}>{\footnotesize\raggedright\arraybackslash}p{6.40cm}@{}}
\toprule
{\footnotesize\textbf{Name}} & {\footnotesize\textbf{Size}} & {\footnotesize\textbf{Type}} & {\footnotesize\textbf{Units}} & {\footnotesize\textbf{Description}} \\
\midrule
\endfirsthead
\toprule
{\footnotesize\textbf{Name}} & {\footnotesize\textbf{Size}} & {\footnotesize\textbf{Type}} & {\footnotesize\textbf{Units}} & {\footnotesize\textbf{Description}} \\
\midrule
\endhead
\midrule
\endfoot
\bottomrule
\endlastfoot
S & 1x1 & struct & - & .beam NxK (occlusion along v\_rel), .number .normal .tangential NxKx8 (species in const.speciesName order; [] without T\_K), .speedRatio 8xK, .kind "flow", .tableFallback (1xK logical: the epochs Method "table" computed exactly because they lie outside the table; [] for the other methods). With vrel\_body = [] and T\_K = [] under Method "table": the table itself (.n .t0 .h .ns .s0 .hs .halfWidthTan .shaded .quadrature .plus .minus, rows plate x species x \{number normal tangential\}, only the plates .shaded that the Lut ever occludes) \\
info & 1x1 & struct & - & .method, .quadrature, .directionsPerEpoch, .integrals (struct CN CpInc CtauInc NxKx8, with 'Integrals'), .provenance; under Method "table" also .table (the grid without its values) and .fallbackEpochs (count of epochs computed exactly) \\
\end{longtable}

\subsubsection*{ERRORS}
\begin{description}[leftmargin=1.2em,style=nextline,itemsep=0pt]
\item[\texttt{NonGravForces:\allowbreak{}flowShadow:\allowbreak{}inputMacro}] not a MacroModel, or without vertices on every plate
\item[\texttt{NonGravForces:\allowbreak{}flowShadow:\allowbreak{}inputShape}] vrel\_body not 3xK or not finite, T\_K not 1xK or [], Quadrature not positive integers
\item[\texttt{NonGravForces:\allowbreak{}flowShadow:\allowbreak{}inputRange}] \textbar{}vrel\textbar{} = 0 or T\_K \textless{}= 0
\item[\texttt{NonGravForces:\allowbreak{}flowShadow:\allowbreak{}inputLut}] Method "lut" or "table" without a 'Lut'
\item[\texttt{NonGravForces:\allowbreak{}flowShadow:\allowbreak{}unsupportedIntegrals}] 'Integrals' under Method "table" (they are not tabulated)
\item[\texttt{NonGravForces:\allowbreak{}selfShadow:\allowbreak{}*}] propagated from the occlusion look-up
\end{description}

\subsection{\texttt{NonGravForces.geom.thermalState}}
\label{api:geom.thermalState}
\noindent Per-panel surface temperatures from the radiation source list: steady state or transient lumped wall.

\begin{lstlisting}
[T, info] = NonGravForces.geom.thermalState(macro, src)                          % steady state
[T, info] = NonGravForces.geom.thermalState(macro, src, 'Model', "transient", 'Epoch', ep, 'T0_K', T0)
[T, info] = NonGravForces.geom.thermalState(macro, src, 'InternalPower_W', 250, 'InternalShare', s)
\end{lstlisting}

\subsubsection*{INPUT ARGUMENTS}
\begin{longtable}{@{}>{\footnotesize\raggedright\arraybackslash\ttfamily}p{3.00cm}>{\footnotesize\raggedright\arraybackslash}p{1.50cm}>{\footnotesize\raggedright\arraybackslash}p{2.20cm}>{\footnotesize\raggedright\arraybackslash}p{1.50cm}>{\footnotesize\raggedright\arraybackslash}p{6.40cm}@{}}
\toprule
{\footnotesize\textbf{Name}} & {\footnotesize\textbf{Size}} & {\footnotesize\textbf{Type}} & {\footnotesize\textbf{Units}} & {\footnotesize\textbf{Description}} \\
\midrule
\endfirsthead
\toprule
{\footnotesize\textbf{Name}} & {\footnotesize\textbf{Size}} & {\footnotesize\textbf{Type}} & {\footnotesize\textbf{Units}} & {\footnotesize\textbf{Description}} \\
\midrule
\endhead
\midrule
\endfoot
\bottomrule
\endlastfoot
macro & 1x1 & NonGravForces.\allowbreak{}geom.\allowbreak{}MacroModel & - & Optics (optVis(1,:), optIR(1,:)), areas, heat capacities \\
src & 1x1 & struct & - & env.radiationSources output: .dir 3xMxK, .P MxK, .band 1xM \\
\end{longtable}

\subsubsection*{NAME-VALUE ARGUMENTS}
\begin{longtable}{@{}>{\footnotesize\raggedright\arraybackslash\ttfamily}p{3.00cm}>{\footnotesize\raggedright\arraybackslash}p{1.50cm}>{\footnotesize\raggedright\arraybackslash}p{2.20cm}>{\footnotesize\raggedright\arraybackslash}p{1.50cm}>{\footnotesize\raggedright\arraybackslash}p{6.40cm}@{}}
\toprule
{\footnotesize\textbf{Name}} & {\footnotesize\textbf{Size}} & {\footnotesize\textbf{Type}} & {\footnotesize\textbf{Units}} & {\footnotesize\textbf{Description}} \\
\midrule
\endfirsthead
\toprule
{\footnotesize\textbf{Name}} & {\footnotesize\textbf{Size}} & {\footnotesize\textbf{Type}} & {\footnotesize\textbf{Units}} & {\footnotesize\textbf{Description}} \\
\midrule
\endhead
\midrule
\endfoot
\bottomrule
\endlastfoot
Model & 1x1 & string & - & "steadyState" (default) \textbar{} "transient" \\
Exposed & NxKxM & double & - & Unobstructed fraction of every plate toward every source, from geom.selfShadow (default: none = 1); N x K when the list has one source \\
InternalPower\_\allowbreak{}W & 1x1 & double & W & Total internal dissipation radiated through the panels (default 0) \\
InternalShare & 1xN & double & - & Share of the internal power per panel (\textgreater{}= 0, sums to 1); REQUIRED when InternalPower\_W \textgreater{} 0 (rule R5), ignored otherwise; macro.internalShare is NOT read here \\
Epoch & 1x1 & struct & - & referenceSystem epoch struct, K entries (transient) \\
T0\_\allowbreak{}K & 1xN & double & K & Initial temperatures (transient) \\
\end{longtable}

\subsubsection*{OUTPUT ARGUMENTS}
\begin{longtable}{@{}>{\footnotesize\raggedright\arraybackslash\ttfamily}p{3.00cm}>{\footnotesize\raggedright\arraybackslash}p{1.50cm}>{\footnotesize\raggedright\arraybackslash}p{2.20cm}>{\footnotesize\raggedright\arraybackslash}p{1.50cm}>{\footnotesize\raggedright\arraybackslash}p{6.40cm}@{}}
\toprule
{\footnotesize\textbf{Name}} & {\footnotesize\textbf{Size}} & {\footnotesize\textbf{Type}} & {\footnotesize\textbf{Units}} & {\footnotesize\textbf{Description}} \\
\midrule
\endfirsthead
\toprule
{\footnotesize\textbf{Name}} & {\footnotesize\textbf{Size}} & {\footnotesize\textbf{Type}} & {\footnotesize\textbf{Units}} & {\footnotesize\textbf{Description}} \\
\midrule
\endhead
\midrule
\endfoot
\bottomrule
\endlastfoot
T & NxK & double & K & Surface temperature per panel and epoch \\
info & 1x1 & struct & - & .model, .absorbed\_W\_m2 (NxK q\_abs), .internal\_W\_m2 (NxK), .emitted\_W\_m2 (NxK eps sigma T\textasciicircum{}4), .heatCapacity\_J\_m2K (1xN), .substeps (1x(K-1), transient), .timeConstant\_s (NxK, transient, linearised C/(4 eps sigma T\textasciicircum{}3)), .provenance \\
\end{longtable}

\subsubsection*{ERRORS}
\begin{description}[leftmargin=1.2em,style=nextline,itemsep=0pt]
\item[\texttt{NonGravForces:\allowbreak{}thermalState:\allowbreak{}inputMacro}] not a MacroModel; transient with a NaN or non-positive heat capacity
\item[\texttt{NonGravForces:\allowbreak{}thermalState:\allowbreak{}inputSources}] src lacks dir/P/band or its sizes disagree
\item[\texttt{NonGravForces:\allowbreak{}thermalState:\allowbreak{}inputShape}] Exposed not NxKxM in [0,1]; transient with T0\_K not N finite values \textgreater{}= 0, or Epoch not an epoch struct of K strictly increasing entries
\item[\texttt{NonGravForces:\allowbreak{}thermalState:\allowbreak{}inputInternal}] InternalPower\_W \textgreater{} 0 without a share summing to 1
\item[\texttt{NonGravForces:\allowbreak{}thermalState:\allowbreak{}numericNaN}] NaN/Inf in the sources or the result
\end{description}

\subsection{\texttt{NonGravForces.geom.AttitudeProvider}}
\label{api:geom.AttitudeProvider}
\noindent BODY-\textgreater{}GCRS quaternions at query epochs: SCA1B nodes interpolated by referenceSystem, or a nominal attitude from the state.

\begin{lstlisting}
ap = NonGravForces.geom.AttitudeProvider(epochNodes, qNodes)                 % from io.readSCA1B
ap = NonGravForces.geom.AttitudeProvider(epochNodes, qNodes, 'Method', "squad")
q  = ap.at(epoch)                                                            % 4xK BODY->GCRS
q  = NonGravForces.geom.AttitudeProvider.nominal(r, v, xAlongVelocity)       % 4xK from the state
ap = NonGravForces.geom.AttitudeProvider.fromItsg(files)                     % TU Graz's attitude product
\end{lstlisting}

\subsubsection*{CONSTRUCTOR INPUTS}
\begin{longtable}{@{}>{\footnotesize\raggedright\arraybackslash\ttfamily}p{3.00cm}>{\footnotesize\raggedright\arraybackslash}p{1.50cm}>{\footnotesize\raggedright\arraybackslash}p{2.20cm}>{\footnotesize\raggedright\arraybackslash}p{1.50cm}>{\footnotesize\raggedright\arraybackslash}p{6.40cm}@{}}
\toprule
{\footnotesize\textbf{Name}} & {\footnotesize\textbf{Size}} & {\footnotesize\textbf{Type}} & {\footnotesize\textbf{Units}} & {\footnotesize\textbf{Description}} \\
\midrule
\endfirsthead
\toprule
{\footnotesize\textbf{Name}} & {\footnotesize\textbf{Size}} & {\footnotesize\textbf{Type}} & {\footnotesize\textbf{Units}} & {\footnotesize\textbf{Description}} \\
\midrule
\endhead
\midrule
\endfoot
\bottomrule
\endlastfoot
epochNodes & 1x1 & struct & - & referenceSystem epoch struct, M entries, strictly increasing \\
qNodes & 4xM & double & - & BODY-\textgreater{}GCRS quaternions at the nodes (io.readSCA1B output) \\
\end{longtable}

\subsubsection*{NAME-VALUE}
\begin{longtable}{@{}>{\footnotesize\raggedright\arraybackslash\ttfamily}p{3.00cm}>{\footnotesize\raggedright\arraybackslash}p{1.50cm}>{\footnotesize\raggedright\arraybackslash}p{2.20cm}>{\footnotesize\raggedright\arraybackslash}p{1.50cm}>{\footnotesize\raggedright\arraybackslash}p{6.40cm}@{}}
\toprule
{\footnotesize\textbf{Name}} & {\footnotesize\textbf{Size}} & {\footnotesize\textbf{Type}} & {\footnotesize\textbf{Units}} & {\footnotesize\textbf{Description}} \\
\midrule
\endfirsthead
\toprule
{\footnotesize\textbf{Name}} & {\footnotesize\textbf{Size}} & {\footnotesize\textbf{Type}} & {\footnotesize\textbf{Units}} & {\footnotesize\textbf{Description}} \\
\midrule
\endhead
\midrule
\endfoot
\bottomrule
\endlastfoot
Method & 1x1 & string & - & "slerp" (default) \textbar{} "squad" \\
\end{longtable}

\subsubsection*{PROPERTIES (read-only)}
\begin{lstlisting}
source       string   "nodes", or "itsg" from fromItsg
method       string
epochNodes   struct
qNodes       4xM double, as given (sign continuity is the interpolator's)
nodeCount    double
scale        string   time scale of the nodes; queries are converted to it
span         1x2 double decimal day numbers of the first and last node, in .scale
convention   struct   referenceSystem.attitudeConvention()
provenance   struct
\end{lstlisting}

\subsubsection*{METHODS}
\begin{lstlisting}
q = at(epoch)              4xK unit quaternions BODY->GCRS at the query epochs
[q, info] = at(epoch)      info from interpolateQuaternions plus .converted (logical)
q = nominal(r, v, xAlongVelocity)   static, 4xK from a 3xK state and a logical
ap = fromItsg(files, 'Method', m)   static, a provider over TU Graz's attitude files (one satellite, in
                                    any order: sorted, an epoch held twice with the same rotation kept
                                    once), conjugated once; source "itsg"
\end{lstlisting}

\subsubsection*{ERRORS}
\begin{description}[leftmargin=1.2em,style=nextline,itemsep=0pt]
\item[\texttt{NonGravForces:\allowbreak{}AttitudeProvider:\allowbreak{}inputNodes}] qNodes not 4xM, epochNodes not an epoch struct with M entries, M \textless{} 2
\item[\texttt{NonGravForces:\allowbreak{}AttitudeProvider:\allowbreak{}inputEpoch}] at(): not an epoch struct
\item[\texttt{NonGravForces:\allowbreak{}AttitudeProvider:\allowbreak{}inputShape}] nominal(): r or v not 3xK, or degenerate (r x v = 0)
\item[\texttt{NonGravForces:\allowbreak{}AttitudeProvider:\allowbreak{}dataItsg}] fromItsg(): a file is not a 4-column quaternion product, or two files hold one epoch with different attitudes
\item[\texttt{gracefoReading:\allowbreak{}*}] fromItsg(): propagated unwrapped
\item[\texttt{referenceSystem:\allowbreak{}interpolateQuaternions:\allowbreak{}*}] propagated unwrapped (queryOutOfRange in particular)
\item[\texttt{referenceSystem:\allowbreak{}*}] propagated unwrapped
\end{description}

\subsection{\texttt{NonGravForces.geom.checkAttitudeSense}}
\label{api:geom.checkAttitudeSense}
\noindent Physics test of a BODY-\textgreater{}GCRS quaternion series: body +Z must point to nadir, body +X along the flight line.

\begin{lstlisting}
[ok, report] = NonGravForces.geom.checkAttitudeSense(q, r, v)
[ok, report] = NonGravForces.geom.checkAttitudeSense(q, r, v, 'NadirTolerance_deg', 5)
\end{lstlisting}

\subsubsection*{INPUT ARGUMENTS}
\begin{longtable}{@{}>{\footnotesize\raggedright\arraybackslash\ttfamily}p{3.00cm}>{\footnotesize\raggedright\arraybackslash}p{1.50cm}>{\footnotesize\raggedright\arraybackslash}p{2.20cm}>{\footnotesize\raggedright\arraybackslash}p{1.50cm}>{\footnotesize\raggedright\arraybackslash}p{6.40cm}@{}}
\toprule
{\footnotesize\textbf{Name}} & {\footnotesize\textbf{Size}} & {\footnotesize\textbf{Type}} & {\footnotesize\textbf{Units}} & {\footnotesize\textbf{Description}} \\
\midrule
\endfirsthead
\toprule
{\footnotesize\textbf{Name}} & {\footnotesize\textbf{Size}} & {\footnotesize\textbf{Type}} & {\footnotesize\textbf{Units}} & {\footnotesize\textbf{Description}} \\
\midrule
\endhead
\midrule
\endfoot
\bottomrule
\endlastfoot
q\_\allowbreak{}body2gcrs & 4xK & double & - & Quaternions to test, scalar first \\
r & 3xK & double & m & GCRS position at the same epochs \\
v & 3xK & double & m/\allowbreak{}s & GCRS velocity (flight-line sense) \\
\end{longtable}

\subsubsection*{NAME-VALUE ARGUMENTS}
\begin{longtable}{@{}>{\footnotesize\raggedright\arraybackslash\ttfamily}p{3.00cm}>{\footnotesize\raggedright\arraybackslash}p{1.50cm}>{\footnotesize\raggedright\arraybackslash}p{2.20cm}>{\footnotesize\raggedright\arraybackslash}p{1.50cm}>{\footnotesize\raggedright\arraybackslash}p{6.40cm}@{}}
\toprule
{\footnotesize\textbf{Name}} & {\footnotesize\textbf{Size}} & {\footnotesize\textbf{Type}} & {\footnotesize\textbf{Units}} & {\footnotesize\textbf{Description}} \\
\midrule
\endfirsthead
\toprule
{\footnotesize\textbf{Name}} & {\footnotesize\textbf{Size}} & {\footnotesize\textbf{Type}} & {\footnotesize\textbf{Units}} & {\footnotesize\textbf{Description}} \\
\midrule
\endhead
\midrule
\endfoot
\bottomrule
\endlastfoot
NadirTolerance\_\allowbreak{}deg & 1x1 & double & deg & Largest acceptable nadir angle (default 5: GRACE-FO pointing offsets are \textasciitilde{}1 deg, attitude-control deadbands smaller) \\
\end{longtable}

\subsubsection*{OUTPUT ARGUMENTS}
\begin{longtable}{@{}>{\footnotesize\raggedright\arraybackslash\ttfamily}p{3.00cm}>{\footnotesize\raggedright\arraybackslash}p{1.50cm}>{\footnotesize\raggedright\arraybackslash}p{2.20cm}>{\footnotesize\raggedright\arraybackslash}p{1.50cm}>{\footnotesize\raggedright\arraybackslash}p{6.40cm}@{}}
\toprule
{\footnotesize\textbf{Name}} & {\footnotesize\textbf{Size}} & {\footnotesize\textbf{Type}} & {\footnotesize\textbf{Units}} & {\footnotesize\textbf{Description}} \\
\midrule
\endfirsthead
\toprule
{\footnotesize\textbf{Name}} & {\footnotesize\textbf{Size}} & {\footnotesize\textbf{Type}} & {\footnotesize\textbf{Units}} & {\footnotesize\textbf{Description}} \\
\midrule
\endhead
\midrule
\endfoot
\bottomrule
\endlastfoot
ok & 1x1 & logical & - & every epoch's nadir angle is within tolerance \\
report & 1x1 & struct & - & .nadirAngle\_deg (1xK, body +Z vs -r/\textbar{}r\textbar{}), .nadirAngleConjugate\_deg (1xK, same with q*), .alongTrackAngle\_deg (1xK, body +X vs v/\textbar{}v\textbar{}), .xSense ("alongVelocity" \textbar{} "againstVelocity" \textbar{} "mixed", from the along-track angles), .maxNadirAngle\_deg, .tolerance\_deg, .K, .provenance \\
\end{longtable}

\subsubsection*{ERRORS}
\begin{description}[leftmargin=1.2em,style=nextline,itemsep=0pt]
\item[\texttt{NonGravForces:\allowbreak{}checkAttitudeSense:\allowbreak{}inputShape}] q not 4xK, r or v not 3xK, or NaN/Inf
\item[\texttt{referenceSystem:\allowbreak{}*}] propagated unwrapped
\end{description}

\subsection{\texttt{NonGravForces.geom.attitudeQuality}}
\label{api:geom.attitudeQuality}
\noindent Attitude-quality diagnostic: star-camera rates against the IMU gyros, attitude against the nominal frame, 1 and 2 cpr.

\begin{lstlisting}
Q = NonGravForces.geom.attitudeQuality(epochSca, qSca, epochImu, angleImu, gyroId)
Q = NonGravForces.geom.attitudeQuality(..., 'Orbit', struct('r', r, 'v', v), 'AlongVelocity', false)
Q = NonGravForces.geom.attitudeQuality(..., 'Period_s', 5640)
\end{lstlisting}

\subsubsection*{INPUT ARGUMENTS}
\begin{longtable}{@{}>{\footnotesize\raggedright\arraybackslash\ttfamily}p{3.00cm}>{\footnotesize\raggedright\arraybackslash}p{1.50cm}>{\footnotesize\raggedright\arraybackslash}p{2.20cm}>{\footnotesize\raggedright\arraybackslash}p{1.50cm}>{\footnotesize\raggedright\arraybackslash}p{6.40cm}@{}}
\toprule
{\footnotesize\textbf{Name}} & {\footnotesize\textbf{Size}} & {\footnotesize\textbf{Type}} & {\footnotesize\textbf{Units}} & {\footnotesize\textbf{Description}} \\
\midrule
\endfirsthead
\toprule
{\footnotesize\textbf{Name}} & {\footnotesize\textbf{Size}} & {\footnotesize\textbf{Type}} & {\footnotesize\textbf{Units}} & {\footnotesize\textbf{Description}} \\
\midrule
\endhead
\midrule
\endfoot
\bottomrule
\endlastfoot
epochSca & 1x1 & struct & - & referenceSystem epoch struct, K entries (io.readSCA1B), strictly increasing \\
qSca & 4xK & double & - & BODY-\textgreater{}GCRS quaternions at epochSca \\
epochImu & 1x1 & struct & - & referenceSystem epoch struct, N entries (io.readIMU1B), any scale \\
angleImu & 1xN & double & rad & accumulated gyro angle of each record \\
gyroId & 1xN & double & - & gyro of each record \\
\end{longtable}

\subsubsection*{NAME-VALUE ARGUMENTS}
\begin{longtable}{@{}>{\footnotesize\raggedright\arraybackslash\ttfamily}p{3.00cm}>{\footnotesize\raggedright\arraybackslash}p{1.50cm}>{\footnotesize\raggedright\arraybackslash}p{2.20cm}>{\footnotesize\raggedright\arraybackslash}p{1.50cm}>{\footnotesize\raggedright\arraybackslash}p{6.40cm}@{}}
\toprule
{\footnotesize\textbf{Name}} & {\footnotesize\textbf{Size}} & {\footnotesize\textbf{Type}} & {\footnotesize\textbf{Units}} & {\footnotesize\textbf{Description}} \\
\midrule
\endfirsthead
\toprule
{\footnotesize\textbf{Name}} & {\footnotesize\textbf{Size}} & {\footnotesize\textbf{Type}} & {\footnotesize\textbf{Units}} & {\footnotesize\textbf{Description}} \\
\midrule
\endhead
\midrule
\endfoot
\bottomrule
\endlastfoot
Orbit & 1x1 & struct & m,m/\allowbreak{}s & .r .v, 3xK GCRS at epochSca; default [] (no nominal comparison) \\
AlongVelocity & 1x1 & logical & - & body +X along the velocity for the nominal frame; default false (GRACE-C on 2020-01-01; GRACE-D flies true) \\
Period\_\allowbreak{}s & 1x1 & double & s & orbital period for the cpr fits; default NaN (from Orbit, else none) \\
MaxGap\_\allowbreak{}s & 1x1 & double & s & longest SCA1B interval and gyro sample spacing used; default 2 \\
JumpThreshold\_\allowbreak{}rad & 1x1 & double & rad & a gyro step larger than this is a reset or wrap: the gyro series is split there and no interval spans it; default deg2rad(1) \\
\end{longtable}

\subsubsection*{OUTPUT ARGUMENTS}
\begin{longtable}{@{}>{\footnotesize\raggedright\arraybackslash\ttfamily}p{3.00cm}>{\footnotesize\raggedright\arraybackslash}p{1.50cm}>{\footnotesize\raggedright\arraybackslash}p{2.20cm}>{\footnotesize\raggedright\arraybackslash}p{1.50cm}>{\footnotesize\raggedright\arraybackslash}p{6.40cm}@{}}
\toprule
{\footnotesize\textbf{Name}} & {\footnotesize\textbf{Size}} & {\footnotesize\textbf{Type}} & {\footnotesize\textbf{Units}} & {\footnotesize\textbf{Description}} \\
\midrule
\endfirsthead
\toprule
{\footnotesize\textbf{Name}} & {\footnotesize\textbf{Size}} & {\footnotesize\textbf{Type}} & {\footnotesize\textbf{Units}} & {\footnotesize\textbf{Description}} \\
\midrule
\endhead
\midrule
\endfoot
\bottomrule
\endlastfoot
Q & 1x1 & struct & - & .t\_s (1xM mid-interval times, s from the first SCA1B epoch), .omega\_rad\_s (3xM SCA1B mean body rates), .gyros (1xG ids), and per gyro in the struct array .gyro(g): .id, .rate\_rad\_s (1xM, NaN where not covered), .axis (3x1 unit, SRF), .scale, .bias\_rad\_s (the minimum-norm fit; see .condition), .condition (2-norm condition number of the column-normalised design), .residual\_rad\_s (1xM), .rms\_rad\_s, .rmsRate\_rad\_s, .n, .cpr (1x2 amplitudes at 1 and 2 cpr, rad/s; NaN without a period); .period\_s; with Orbit: .nominal.t\_s (1xK), .nominal.angles\_rad (3xK roll pitch yaw), .nominal.mean\_rad, .nominal.rms\_rad, .nominal.cpr (3x2 rad) \\
\end{longtable}

\subsubsection*{ERRORS}
\begin{description}[leftmargin=1.2em,style=nextline,itemsep=0pt]
\item[\texttt{NonGravForces:\allowbreak{}attitudeQuality:\allowbreak{}inputShape}] sizes disagree, fewer than 2 SCA1B epochs, or an epoch is not a referenceSystem epoch struct
\item[\texttt{NonGravForces:\allowbreak{}attitudeQuality:\allowbreak{}inputOrbit}] Orbit is not .r/.v 3xK at the SCA1B epochs
\item[\texttt{NonGravForces:\allowbreak{}attitudeQuality:\allowbreak{}dataOverlap}] no SCA1B interval is covered by a gyro, or a gyro has fewer than 10 covered intervals (the fit needs them)
\item[\texttt{NonGravForces:\allowbreak{}AttitudeProvider:\allowbreak{}*, referenceSystem:\allowbreak{}*}] propagated unwrapped
\end{description}

\section{L4 \texttt{kernel} --- per-panel forces, their derivatives and their reductions}

\subsection{\texttt{NonGravForces.kernel.gsiCoefficients}}
\label{api:kernel.gsiCoefficients}
\noindent Normal and tangential free-molecular force coefficients of a flat plate: sentman, maxwell, constantCd.

\begin{lstlisting}
[Cp, Ctau] = NonGravForces.kernel.gsiCoefficients("sentman", s, theta, alpha, Tw_K, Tinf_K)
[Cp, Ctau] = NonGravForces.kernel.gsiCoefficients("maxwell", s, theta, alpha, Tw_K, Tinf_K, 'SpecularFraction', 0.1)
[Cp, Ctau] = NonGravForces.kernel.gsiCoefficients("constantCd", [], theta, [], [], [], 'Cd', 2.2)
[Cp, Ctau, info] = NonGravForces.kernel.gsiCoefficients(___)
\end{lstlisting}

\subsubsection*{INPUT ARGUMENTS}
\begin{longtable}{@{}>{\footnotesize\raggedright\arraybackslash\ttfamily}p{3.00cm}>{\footnotesize\raggedright\arraybackslash}p{1.50cm}>{\footnotesize\raggedright\arraybackslash}p{2.20cm}>{\footnotesize\raggedright\arraybackslash}p{1.50cm}>{\footnotesize\raggedright\arraybackslash}p{6.40cm}@{}}
\toprule
{\footnotesize\textbf{Name}} & {\footnotesize\textbf{Size}} & {\footnotesize\textbf{Type}} & {\footnotesize\textbf{Units}} & {\footnotesize\textbf{Description}} \\
\midrule
\endfirsthead
\toprule
{\footnotesize\textbf{Name}} & {\footnotesize\textbf{Size}} & {\footnotesize\textbf{Type}} & {\footnotesize\textbf{Units}} & {\footnotesize\textbf{Description}} \\
\midrule
\endhead
\midrule
\endfoot
\bottomrule
\endlastfoot
model & 1x1 & string & - & "sentman" \textbar{} "maxwell" \textbar{} "constantCd" \textbar{} "cll" \textbar{} "schamberg" \\
s & any & double & - & Speed ratio U / sqrt(2 k T / m) \textgreater{} 0 (ignored by constantCd) \\
theta & any & double & rad & Angle between the plate normal and the relative velocity v, in [0, pi] \\
alpha & any & double & - & Energy accommodation coefficient in [0, 1] \\
Tw\_\allowbreak{}K & any & double & K & Wall temperature \textgreater{} 0 \\
Tinf\_\allowbreak{}K & any & double & K & Gas kinetic temperature \textgreater{} 0 \\
\end{longtable}

\noindent All arrays must be compatible for implicit expansion (e.g. theta NxK, s 1xK).

\subsubsection*{NAME-VALUE ARGUMENTS}
\begin{longtable}{@{}>{\footnotesize\raggedright\arraybackslash\ttfamily}p{3.00cm}>{\footnotesize\raggedright\arraybackslash}p{1.50cm}>{\footnotesize\raggedright\arraybackslash}p{2.20cm}>{\footnotesize\raggedright\arraybackslash}p{1.50cm}>{\footnotesize\raggedright\arraybackslash}p{6.40cm}@{}}
\toprule
{\footnotesize\textbf{Name}} & {\footnotesize\textbf{Size}} & {\footnotesize\textbf{Type}} & {\footnotesize\textbf{Units}} & {\footnotesize\textbf{Description}} \\
\midrule
\endfirsthead
\toprule
{\footnotesize\textbf{Name}} & {\footnotesize\textbf{Size}} & {\footnotesize\textbf{Type}} & {\footnotesize\textbf{Units}} & {\footnotesize\textbf{Description}} \\
\midrule
\endhead
\midrule
\endfoot
\bottomrule
\endlastfoot
Cd & 1x1 & double & - & Drag coefficient for constantCd, \textgreater{} 0 (required by that model) \\
SpecularFraction & 1x1 & double & - & Specular fraction for maxwell, in [0, 1] (required by that model) \\
Derivatives & 1x1 & logical & - & Also return the analytic derivatives of the coefficients with respect to alpha and to the speed ratio s, in info (default false). They live here, beside the closed form they differentiate, so the two cannot drift apart (rule R1); kernel.partials assembles them. \\
Provenance & 1x1 & logical & - & Build info.provenance when info is requested (default true); a kernel that calls this per species passes false, because the record costs more than the coefficients (HANDOVER LP-C, finding F12). \\
\end{longtable}

\subsubsection*{OUTPUT ARGUMENTS}
\begin{longtable}{@{}>{\footnotesize\raggedright\arraybackslash\ttfamily}p{3.00cm}>{\footnotesize\raggedright\arraybackslash}p{1.50cm}>{\footnotesize\raggedright\arraybackslash}p{2.20cm}>{\footnotesize\raggedright\arraybackslash}p{1.50cm}>{\footnotesize\raggedright\arraybackslash}p{6.40cm}@{}}
\toprule
{\footnotesize\textbf{Name}} & {\footnotesize\textbf{Size}} & {\footnotesize\textbf{Type}} & {\footnotesize\textbf{Units}} & {\footnotesize\textbf{Description}} \\
\midrule
\endfirsthead
\toprule
{\footnotesize\textbf{Name}} & {\footnotesize\textbf{Size}} & {\footnotesize\textbf{Type}} & {\footnotesize\textbf{Units}} & {\footnotesize\textbf{Description}} \\
\midrule
\endhead
\midrule
\endfoot
\bottomrule
\endlastfoot
Cp & expanded & double & - & Normal coefficient (\textgreater{}= 0) \\
Ctau & expanded & double & - & Tangential coefficient (\textgreater{}= 0) \\
info & 1x1 & struct & - & .model, .CD, .CL (same size as Cp), .CpIncident and .CpReemitted (Cp = CpIncident + CpReemitted: the incident normal momentum, including a specular fraction, and the diffusely re-emitted one, which is proportional to the incident NUMBER flux; geom.flowShadow scales the two differently), .vRatio (re-emission speed ratio; empty for constantCd), .provenance. With 'Derivatives': .dCp\_dalpha, .dCtau\_dalpha (identically zero: alpha enters only the re-emitted normal momentum), .dCp\_ds, .dCtau\_ds, .dCpIncident\_ds, .dCpReemitted\_ds, all the size of Cp. \\
\end{longtable}

\subsubsection*{ERRORS}
\begin{description}[leftmargin=1.2em,style=nextline,itemsep=0pt]
\item[\texttt{NonGravForces:\allowbreak{}gsiCoefficients:\allowbreak{}unsupportedModel}] cll or schamberg (O-17)
\item[\texttt{NonGravForces:\allowbreak{}gsiCoefficients:\allowbreak{}inputShape}] arrays not compatible for expansion
\item[\texttt{NonGravForces:\allowbreak{}gsiCoefficients:\allowbreak{}inputRange}] s \textless{}= 0, theta outside [0, pi], alpha outside [0, 1], Tw or Tinf \textless{}= 0
\item[\texttt{NonGravForces:\allowbreak{}gsiCoefficients:\allowbreak{}inputCd}] constantCd without a positive Cd
\item[\texttt{NonGravForces:\allowbreak{}gsiCoefficients:\allowbreak{}inputSpecular}] maxwell without a SpecularFraction in [0, 1]
\item[\texttt{NonGravForces:\allowbreak{}gsiCoefficients:\allowbreak{}numericNaN}] NaN or Inf in an input
\end{description}

\subsection{\texttt{NonGravForces.kernel.accommodation}}
\label{api:kernel.accommodation}
\noindent Energy accommodation of the gas-surface interaction per epoch: fixed, either Pilinski fit, the Walker/Mehta adsorbate mixture, or SESAM.

\begin{lstlisting}
acc = NonGravForces.kernel.accommodation("fixed", species, T_K, 'Alpha', 0.85)
acc = NonGravForces.kernel.accommodation("pilinski2010", species, T_K)
acc = NonGravForces.kernel.accommodation("pilinski2011", species, T_K)
acc = NonGravForces.kernel.accommodation("walker2014", species, T_K, 'SurfaceMass_amu', 263.3223)
acc = NonGravForces.kernel.accommodation("sesam", species, T_K, 'RelativeSpeed_ms', 7600)
[acc, info] = NonGravForces.kernel.accommodation(___)
\end{lstlisting}

\subsubsection*{INPUT ARGUMENTS}
\begin{longtable}{@{}>{\footnotesize\raggedright\arraybackslash\ttfamily}p{3.00cm}>{\footnotesize\raggedright\arraybackslash}p{1.50cm}>{\footnotesize\raggedright\arraybackslash}p{2.20cm}>{\footnotesize\raggedright\arraybackslash}p{1.50cm}>{\footnotesize\raggedright\arraybackslash}p{6.40cm}@{}}
\toprule
{\footnotesize\textbf{Name}} & {\footnotesize\textbf{Size}} & {\footnotesize\textbf{Type}} & {\footnotesize\textbf{Units}} & {\footnotesize\textbf{Description}} \\
\midrule
\endfirsthead
\toprule
{\footnotesize\textbf{Name}} & {\footnotesize\textbf{Size}} & {\footnotesize\textbf{Type}} & {\footnotesize\textbf{Units}} & {\footnotesize\textbf{Description}} \\
\midrule
\endhead
\midrule
\endfoot
\bottomrule
\endlastfoot
model & 1x1 & string & - & "fixed" \textbar{} "pilinski2010" \textbar{} "pilinski2011" \textbar{} "walker2014" \textbar{} "sesam" \\
species & 8xK & double & 1/\allowbreak{}m\textasciicircum{}3 & Number densities in const.speciesName order (env.density) \\
T\_\allowbreak{}K & 1xK & double & K & Gas kinetic temperature at the satellite (density's T(2,:)) \\
\end{longtable}

\subsubsection*{NAME-VALUE ARGUMENTS}
\begin{longtable}{@{}>{\footnotesize\raggedright\arraybackslash\ttfamily}p{3.00cm}>{\footnotesize\raggedright\arraybackslash}p{1.50cm}>{\footnotesize\raggedright\arraybackslash}p{2.20cm}>{\footnotesize\raggedright\arraybackslash}p{1.50cm}>{\footnotesize\raggedright\arraybackslash}p{6.40cm}@{}}
\toprule
{\footnotesize\textbf{Name}} & {\footnotesize\textbf{Size}} & {\footnotesize\textbf{Type}} & {\footnotesize\textbf{Units}} & {\footnotesize\textbf{Description}} \\
\midrule
\endfirsthead
\toprule
{\footnotesize\textbf{Name}} & {\footnotesize\textbf{Size}} & {\footnotesize\textbf{Type}} & {\footnotesize\textbf{Units}} & {\footnotesize\textbf{Description}} \\
\midrule
\endhead
\midrule
\endfoot
\bottomrule
\endlastfoot
Alpha & 1x1\textbar{}\allowbreak{}1xK & double & - & Coefficient in [0, 1]; required by "fixed", refused by the others \\
SurfaceMass\_\allowbreak{}amu & 1x1 & double & amu & Surface lattice particle mass, \textgreater{} 0; required by "walker2014", refused by "sesam" (which fixes 65), ignored by "fixed" and "pilinski20xx" \\
RelativeSpeed\_\allowbreak{}ms & 1x1\textbar{}\allowbreak{}1xK & double & m/\allowbreak{}s & Satellite speed relative to the gas, \textgreater{} 0; required by "sesam", refused by the others \\
\end{longtable}

\subsubsection*{OUTPUT ARGUMENTS}
\begin{longtable}{@{}>{\footnotesize\raggedright\arraybackslash\ttfamily}p{3.00cm}>{\footnotesize\raggedright\arraybackslash}p{1.50cm}>{\footnotesize\raggedright\arraybackslash}p{2.20cm}>{\footnotesize\raggedright\arraybackslash}p{1.50cm}>{\footnotesize\raggedright\arraybackslash}p{6.40cm}@{}}
\toprule
{\footnotesize\textbf{Name}} & {\footnotesize\textbf{Size}} & {\footnotesize\textbf{Type}} & {\footnotesize\textbf{Units}} & {\footnotesize\textbf{Description}} \\
\midrule
\endfirsthead
\toprule
{\footnotesize\textbf{Name}} & {\footnotesize\textbf{Size}} & {\footnotesize\textbf{Type}} & {\footnotesize\textbf{Units}} & {\footnotesize\textbf{Description}} \\
\midrule
\endhead
\midrule
\endfoot
\bottomrule
\endlastfoot
acc & 1x1 & struct & - & .model; .alpha 1xK (fixed, pilinski20xx; for walker2014 and sesam the adsorbate state's, 1); .coverage 1xK (walker2014, sesam: theta; [] for a single state); .alphaSubstrate 1xK (walker2014, sesam; [] otherwise). The EFFECTIVE coefficient of a two-state surface is theta + (1 - theta) alphaSubstrate; the kernel mixes the two states' forces rather than their coefficients, so this unit does not form it. \\
info & 1x1 & struct & - & .K, .pressureO\_Pa (walker2014, sesam), .meanMass\_amu (walker2014, sesam), .insideFit (1xK logical, EVERY model: whether each epoch lies inside the range the model was fitted over -- alpha in 0.85-1.00 for pilinski20xx, P\_O \textgreater{}= 1e-5 Pa for walker2014, true everywhere for fixed, which extrapolates nothing, and for sesam, whose source states no domain this project can transcribe), .stickingCoefficient .langmuir\_invPa .speedRatioO (sesam), .provenance \\
\end{longtable}

\subsubsection*{WARNINGS}
\begin{description}[leftmargin=1.2em,style=nextline,itemsep=0pt]
\item[\texttt{NonGravForces:\allowbreak{}accommodation:\allowbreak{}outsideFit}] a computed model was evaluated outside the range it was fitted in. Raised ONCE PER CALL, naming how many of the K epochs and which range; silence it with warning('off', 'NonGravForces:accommodation:outsideFit') when the extrapolation is intended. Owner ruling of 2026-09-17 (HANDOVER O-42): warn rather than refuse, because extrapolating a fitted model knowingly is legitimate and refusing would forbid it. This was the first unit after dependencyReport allowed to warn (app.fitCalibration and app.Calibration have since joined it); a warning is not printing, since it carries an identifier and can be silenced. IT WILL FIRE ON GRACE-FO. At 520 km in solar minimum P\_O is \textasciitilde{}1.4e-8 Pa, three orders below walker2014's fitted domain, and the adsorption models return 0.11-0.17 where the fitted literature reports 0.82-0.93 (measured 2026-09-17). That is the extrapolation this warning exists to name, and it is why all three presets take a fixed alpha instead.
\end{description}

\subsubsection*{ERRORS}
\begin{description}[leftmargin=1.2em,style=nextline,itemsep=0pt]
\item[\texttt{NonGravForces:\allowbreak{}accommodation:\allowbreak{}inputShape}] species not 8xK or T\_K not 1xK with the same K
\item[\texttt{NonGravForces:\allowbreak{}accommodation:\allowbreak{}inputAlpha}] fixed without 'Alpha', Alpha outside [0, 1] or not 1x1/1xK, or Alpha given to another model
\item[\texttt{NonGravForces:\allowbreak{}accommodation:\allowbreak{}inputSurfaceMass}] walker2014 without a positive 'SurfaceMass\_amu', or a mass given to sesam
\item[\texttt{NonGravForces:\allowbreak{}accommodation:\allowbreak{}inputSpeed}] sesam without a positive 'RelativeSpeed\_ms', or a speed given to another model
\item[\texttt{NonGravForces:\allowbreak{}accommodation:\allowbreak{}numericNaN}] NaN/Inf in Alpha, species or T\_K (species and T\_K are screened for EVERY model, fixed included)
\item[\texttt{NonGravForces:\allowbreak{}accommodation:\allowbreak{}inputRange}] a negative density or a non-positive temperature (every model but fixed)
\end{description}

\subsection{\texttt{NonGravForces.kernel.dragPanelForce}}
\label{api:kernel.dragPanelForce}
\noindent Per-panel free-molecular gas force in the body frame, species by species.

\begin{lstlisting}
[F, info] = NonGravForces.kernel.dragPanelForce(macro, rho, species, T_K, vrel, 'Model', "sentman", 'Alpha', 0.85, 'Tw_K', 300)
[F, info] = NonGravForces.kernel.dragPanelForce(macro, rho, [], [], vrel, 'Model', "constantCd", 'Cd', 2.2)
[F, info] = NonGravForces.kernel.dragPanelForce(___, 'Shadow', geom.flowShadow(...))
[F, info] = NonGravForces.kernel.dragPanelForce(___, 'Alpha', kernel.accommodation("walker2014", ...))
\end{lstlisting}

\subsubsection*{INPUT ARGUMENTS}
\begin{longtable}{@{}>{\footnotesize\raggedright\arraybackslash\ttfamily}p{3.00cm}>{\footnotesize\raggedright\arraybackslash}p{1.50cm}>{\footnotesize\raggedright\arraybackslash}p{2.20cm}>{\footnotesize\raggedright\arraybackslash}p{1.50cm}>{\footnotesize\raggedright\arraybackslash}p{6.40cm}@{}}
\toprule
{\footnotesize\textbf{Name}} & {\footnotesize\textbf{Size}} & {\footnotesize\textbf{Type}} & {\footnotesize\textbf{Units}} & {\footnotesize\textbf{Description}} \\
\midrule
\endfirsthead
\toprule
{\footnotesize\textbf{Name}} & {\footnotesize\textbf{Size}} & {\footnotesize\textbf{Type}} & {\footnotesize\textbf{Units}} & {\footnotesize\textbf{Description}} \\
\midrule
\endhead
\midrule
\endfoot
\bottomrule
\endlastfoot
macro & 1x1 & MacroModel &  & Panel model (areas A, normals N in body frame) \\
rho & 1xK & double & kg/\allowbreak{}m\textasciicircum{}3 & Total mass density (env.density); used by constantCd and reported for the species-sum check \\
species & 8xK & double & 1/\allowbreak{}m\textasciicircum{}3 & Number densities in const.speciesName order (env.density); may be [] for constantCd \\
T\_\allowbreak{}K & 1xK & double & K & Gas kinetic temperature at the satellite (density's T(2,:)); may be [] for constantCd \\
vrel\_\allowbreak{}body & 3xK & double & m/\allowbreak{}s & Satellite velocity relative to the gas, body frame, \textbar{}v\textbar{} \textgreater{} 0 \\
\end{longtable}

\subsubsection*{NAME-VALUE ARGUMENTS}
\begin{longtable}{@{}>{\footnotesize\raggedright\arraybackslash\ttfamily}p{3.00cm}>{\footnotesize\raggedright\arraybackslash}p{1.50cm}>{\footnotesize\raggedright\arraybackslash}p{2.20cm}>{\footnotesize\raggedright\arraybackslash}p{1.50cm}>{\footnotesize\raggedright\arraybackslash}p{6.40cm}@{}}
\toprule
{\footnotesize\textbf{Name}} & {\footnotesize\textbf{Size}} & {\footnotesize\textbf{Type}} & {\footnotesize\textbf{Units}} & {\footnotesize\textbf{Description}} \\
\midrule
\endfirsthead
\toprule
{\footnotesize\textbf{Name}} & {\footnotesize\textbf{Size}} & {\footnotesize\textbf{Type}} & {\footnotesize\textbf{Units}} & {\footnotesize\textbf{Description}} \\
\midrule
\endhead
\midrule
\endfoot
\bottomrule
\endlastfoot
Model & 1x1 & string & - & "sentman" \textbar{} "maxwell" \textbar{} "constantCd" (required; no default) \\
Alpha & 1x1\textbar{}\allowbreak{}1xK & double\textbar{}\allowbreak{}struct & - & Accommodation coefficient, or a kernel.accommodation struct; required by sentman/maxwell \\
Tw\_\allowbreak{}K & 1x1\textbar{}\allowbreak{}1xK & double & K & Wall temperature, required by sentman/maxwell \\
Cd & 1x1 & double & - & Drag coefficient, required by constantCd \\
SpecularFraction & 1x1 & double & - & Required by maxwell \\
Shadow & NxK & double\textbar{}\allowbreak{}struct & - & Self-shadowing: a geom.flowShadow struct, or one fraction per plate and epoch; default none \\
\end{longtable}

\subsubsection*{OUTPUT ARGUMENTS}
\begin{longtable}{@{}>{\footnotesize\raggedright\arraybackslash\ttfamily}p{3.00cm}>{\footnotesize\raggedright\arraybackslash}p{1.50cm}>{\footnotesize\raggedright\arraybackslash}p{2.20cm}>{\footnotesize\raggedright\arraybackslash}p{1.50cm}>{\footnotesize\raggedright\arraybackslash}p{6.40cm}@{}}
\toprule
{\footnotesize\textbf{Name}} & {\footnotesize\textbf{Size}} & {\footnotesize\textbf{Type}} & {\footnotesize\textbf{Units}} & {\footnotesize\textbf{Description}} \\
\midrule
\endfirsthead
\toprule
{\footnotesize\textbf{Name}} & {\footnotesize\textbf{Size}} & {\footnotesize\textbf{Type}} & {\footnotesize\textbf{Units}} & {\footnotesize\textbf{Description}} \\
\midrule
\endhead
\midrule
\endfoot
\bottomrule
\endlastfoot
F & 3xNxK & double & N & Force on each plate, body frame \\
info & 1x1 & struct & - & .gamma (NxK), .theta\_rad (NxK), .speed\_m\_s (1xK), .speedRatio (8xK, NaN for constantCd), .Cp .Ctau (NxK, mass-weighted mixture coefficients per unit rho, shadowing included), .rhoSpecies\_kg\_m3 (1xK, sum of n\_j m\_j), .model, .shadow ("none" \textbar{} "fraction" \textbar{} "flow"), .provenance \\
\end{longtable}

\subsubsection*{ERRORS}
\begin{description}[leftmargin=1.2em,style=nextline,itemsep=0pt]
\item[\texttt{NonGravForces:\allowbreak{}dragPanelForce:\allowbreak{}inputMacro}] macro is not a NonGravForces.geom.MacroModel
\item[\texttt{NonGravForces:\allowbreak{}dragPanelForce:\allowbreak{}inputModel}] 'Model' missing or not in the set, or its required options absent
\item[\texttt{NonGravForces:\allowbreak{}dragPanelForce:\allowbreak{}inputShape}] rho, species, T\_K, vrel\_body, Alpha or Shadow with the wrong shape
\item[\texttt{NonGravForces:\allowbreak{}dragPanelForce:\allowbreak{}inputRange}] \textbar{}v\textbar{} = 0, rho \textless{} 0, T\_K \textless{}= 0, Tw \textless{}= 0, a fraction or a number/normal ratio outside [0, 1] (the tangential ratio is signed and need only be finite), or an accommodation struct's alpha/coverage/alphaSubstrate outside [0, 1]
\item[\texttt{NonGravForces:\allowbreak{}dragPanelForce:\allowbreak{}numericNaN}] NaN/Inf in an input
\item[\texttt{NonGravForces:\allowbreak{}gsiCoefficients:\allowbreak{}*}] propagated from the coefficient model (e.g. inputCd, inputSpecular)
\item[\texttt{NonGravForces:\allowbreak{}accommodation:\allowbreak{}*}] propagated when a numeric Alpha is validated
\end{description}

\subsection{\texttt{NonGravForces.kernel.radiationPanelForce}}
\label{api:kernel.radiationPanelForce}
\noindent Per-panel radiation force for a source list (Sun, Earth tiles) in the body frame.

\begin{lstlisting}
[F, info] = NonGravForces.kernel.radiationPanelForce(macro, src)
[F, info] = NonGravForces.kernel.radiationPanelForce(macro, src, 'Shadow', f)
\end{lstlisting}

\subsubsection*{INPUT ARGUMENTS}
\begin{longtable}{@{}>{\footnotesize\raggedright\arraybackslash\ttfamily}p{3.00cm}>{\footnotesize\raggedright\arraybackslash}p{1.50cm}>{\footnotesize\raggedright\arraybackslash}p{2.20cm}>{\footnotesize\raggedright\arraybackslash}p{1.50cm}>{\footnotesize\raggedright\arraybackslash}p{6.40cm}@{}}
\toprule
{\footnotesize\textbf{Name}} & {\footnotesize\textbf{Size}} & {\footnotesize\textbf{Type}} & {\footnotesize\textbf{Units}} & {\footnotesize\textbf{Description}} \\
\midrule
\endfirsthead
\toprule
{\footnotesize\textbf{Name}} & {\footnotesize\textbf{Size}} & {\footnotesize\textbf{Type}} & {\footnotesize\textbf{Units}} & {\footnotesize\textbf{Description}} \\
\midrule
\endhead
\midrule
\endfoot
\bottomrule
\endlastfoot
macro & 1x1 & MacroModel &  & Panel model (A, N, optVis [abs;spec;diff], optIR [emis;spec;diff]) \\
src & 1x1 & struct & - & Source list: .dir 3xMxK unit vectors satellite -\textgreater{} source (body frame), .P MxK pressure [Pa], .band 1xM string "vis"\textbar{}"ir" (env.radiationSources) \\
\end{longtable}

\subsubsection*{NAME-VALUE ARGUMENTS}
\begin{longtable}{@{}>{\footnotesize\raggedright\arraybackslash\ttfamily}p{3.00cm}>{\footnotesize\raggedright\arraybackslash}p{1.50cm}>{\footnotesize\raggedright\arraybackslash}p{2.20cm}>{\footnotesize\raggedright\arraybackslash}p{1.50cm}>{\footnotesize\raggedright\arraybackslash}p{6.40cm}@{}}
\toprule
{\footnotesize\textbf{Name}} & {\footnotesize\textbf{Size}} & {\footnotesize\textbf{Type}} & {\footnotesize\textbf{Units}} & {\footnotesize\textbf{Description}} \\
\midrule
\endfirsthead
\toprule
{\footnotesize\textbf{Name}} & {\footnotesize\textbf{Size}} & {\footnotesize\textbf{Type}} & {\footnotesize\textbf{Units}} & {\footnotesize\textbf{Description}} \\
\midrule
\endhead
\midrule
\endfoot
\bottomrule
\endlastfoot
Shadow & NxKxM & double & - & Self-shadow fraction per plate, epoch and source (geom.selfShadow for each source direction), default all 1 (no self-shadowing) \\
\end{longtable}

\subsubsection*{OUTPUT ARGUMENTS}
\begin{longtable}{@{}>{\footnotesize\raggedright\arraybackslash\ttfamily}p{3.00cm}>{\footnotesize\raggedright\arraybackslash}p{1.50cm}>{\footnotesize\raggedright\arraybackslash}p{2.20cm}>{\footnotesize\raggedright\arraybackslash}p{1.50cm}>{\footnotesize\raggedright\arraybackslash}p{6.40cm}@{}}
\toprule
{\footnotesize\textbf{Name}} & {\footnotesize\textbf{Size}} & {\footnotesize\textbf{Type}} & {\footnotesize\textbf{Units}} & {\footnotesize\textbf{Description}} \\
\midrule
\endfirsthead
\toprule
{\footnotesize\textbf{Name}} & {\footnotesize\textbf{Size}} & {\footnotesize\textbf{Type}} & {\footnotesize\textbf{Units}} & {\footnotesize\textbf{Description}} \\
\midrule
\endhead
\midrule
\endfoot
\bottomrule
\endlastfoot
F & 3xNxK & double & N & Force on each plate, body frame, summed over the sources \\
info & 1x1 & struct & - & .absorbed\_W\_m2 (NxK, a P c gamma f summed over sources, f the Shadow fraction: the flux thermalState needs), .litCount (NxK, sources with gamma \textgreater{} 0, counted before the Shadow), .M, .K, .bands, .provenance \\
\end{longtable}

\subsubsection*{ERRORS}
\begin{description}[leftmargin=1.2em,style=nextline,itemsep=0pt]
\item[\texttt{NonGravForces:\allowbreak{}radiationPanelForce:\allowbreak{}inputMacro}] macro is not a MacroModel
\item[\texttt{NonGravForces:\allowbreak{}radiationPanelForce:\allowbreak{}inputSources}] src lacks .dir/.P/.band, shapes disagree, band not vis\textbar{}ir, or a direction is not a unit vector (1e-9)
\item[\texttt{NonGravForces:\allowbreak{}radiationPanelForce:\allowbreak{}inputShape}] Shadow not NxKxM
\item[\texttt{NonGravForces:\allowbreak{}radiationPanelForce:\allowbreak{}inputRange}] P \textless{} 0 or Shadow outside [0, 1]
\item[\texttt{NonGravForces:\allowbreak{}radiationPanelForce:\allowbreak{}numericNaN}] NaN/Inf in src or Shadow
\end{description}

\subsection{\texttt{NonGravForces.kernel.thermalRecoil}}
\label{api:kernel.thermalRecoil}
\noindent Per-panel recoil force of Lambertian thermal emission in the body frame.

\begin{lstlisting}
[F, info] = NonGravForces.kernel.thermalRecoil(macro, T_K)
\end{lstlisting}

\subsubsection*{INPUT ARGUMENTS}
\begin{longtable}{@{}>{\footnotesize\raggedright\arraybackslash\ttfamily}p{3.00cm}>{\footnotesize\raggedright\arraybackslash}p{1.50cm}>{\footnotesize\raggedright\arraybackslash}p{2.20cm}>{\footnotesize\raggedright\arraybackslash}p{1.50cm}>{\footnotesize\raggedright\arraybackslash}p{6.40cm}@{}}
\toprule
{\footnotesize\textbf{Name}} & {\footnotesize\textbf{Size}} & {\footnotesize\textbf{Type}} & {\footnotesize\textbf{Units}} & {\footnotesize\textbf{Description}} \\
\midrule
\endfirsthead
\toprule
{\footnotesize\textbf{Name}} & {\footnotesize\textbf{Size}} & {\footnotesize\textbf{Type}} & {\footnotesize\textbf{Units}} & {\footnotesize\textbf{Description}} \\
\midrule
\endhead
\midrule
\endfoot
\bottomrule
\endlastfoot
macro & 1x1 & MacroModel &  & Panel model (A, N, optIR(1,:) = emissivity) \\
T\_\allowbreak{}K & NxK & double & K & Surface temperature per plate and epoch, \textgreater{}= 0 (an unlit steady-state plate without internal power sits at 0 K and emits nothing) \\
\end{longtable}

\subsubsection*{OUTPUT ARGUMENTS}
\begin{longtable}{@{}>{\footnotesize\raggedright\arraybackslash\ttfamily}p{3.00cm}>{\footnotesize\raggedright\arraybackslash}p{1.50cm}>{\footnotesize\raggedright\arraybackslash}p{2.20cm}>{\footnotesize\raggedright\arraybackslash}p{1.50cm}>{\footnotesize\raggedright\arraybackslash}p{6.40cm}@{}}
\toprule
{\footnotesize\textbf{Name}} & {\footnotesize\textbf{Size}} & {\footnotesize\textbf{Type}} & {\footnotesize\textbf{Units}} & {\footnotesize\textbf{Description}} \\
\midrule
\endfirsthead
\toprule
{\footnotesize\textbf{Name}} & {\footnotesize\textbf{Size}} & {\footnotesize\textbf{Type}} & {\footnotesize\textbf{Units}} & {\footnotesize\textbf{Description}} \\
\midrule
\endhead
\midrule
\endfoot
\bottomrule
\endlastfoot
F & 3xNxK & double & N & Recoil force on each plate, body frame \\
info & 1x1 & struct & - & .emitted\_W\_m2 (NxK, eps sigma T\textasciicircum{}4), .emittedPower\_W (1xK, total), .provenance \\
\end{longtable}

\subsubsection*{ERRORS}
\begin{description}[leftmargin=1.2em,style=nextline,itemsep=0pt]
\item[\texttt{NonGravForces:\allowbreak{}thermalRecoil:\allowbreak{}inputMacro}] macro is not a MacroModel
\item[\texttt{NonGravForces:\allowbreak{}thermalRecoil:\allowbreak{}inputShape}] T\_K is not NxK
\item[\texttt{NonGravForces:\allowbreak{}thermalRecoil:\allowbreak{}inputRange}] a negative temperature
\item[\texttt{NonGravForces:\allowbreak{}thermalRecoil:\allowbreak{}numericNaN}] NaN/Inf in T\_K
\end{description}

\subsection{\texttt{NonGravForces.kernel.assemble}}
\label{api:kernel.assemble}
\noindent Reduce per-panel forces to the total force and the torque about the centre of mass.

\begin{lstlisting}
[F, tau] = NonGravForces.kernel.assemble(Fk, macro)
[F, tau, info] = NonGravForces.kernel.assemble(Fk, macro, 'CentreOfMass', r_cm)
\end{lstlisting}

\subsubsection*{INPUT ARGUMENTS}
\begin{longtable}{@{}>{\footnotesize\raggedright\arraybackslash\ttfamily}p{3.00cm}>{\footnotesize\raggedright\arraybackslash}p{1.50cm}>{\footnotesize\raggedright\arraybackslash}p{2.20cm}>{\footnotesize\raggedright\arraybackslash}p{1.50cm}>{\footnotesize\raggedright\arraybackslash}p{6.40cm}@{}}
\toprule
{\footnotesize\textbf{Name}} & {\footnotesize\textbf{Size}} & {\footnotesize\textbf{Type}} & {\footnotesize\textbf{Units}} & {\footnotesize\textbf{Description}} \\
\midrule
\endfirsthead
\toprule
{\footnotesize\textbf{Name}} & {\footnotesize\textbf{Size}} & {\footnotesize\textbf{Type}} & {\footnotesize\textbf{Units}} & {\footnotesize\textbf{Description}} \\
\midrule
\endhead
\midrule
\endfoot
\bottomrule
\endlastfoot
Fk & 3xNxK & double & N & Per-panel forces, body frame (kernel.dragPanelForce, radiationPanelForce, thermalRecoil, or their sum) \\
macro & 1x1 & MacroModel &  & Panel model whose centroids C (3xN) are the lever arms \\
\end{longtable}

\subsubsection*{NAME-VALUE ARGUMENTS}
\begin{longtable}{@{}>{\footnotesize\raggedright\arraybackslash\ttfamily}p{3.00cm}>{\footnotesize\raggedright\arraybackslash}p{1.50cm}>{\footnotesize\raggedright\arraybackslash}p{2.20cm}>{\footnotesize\raggedright\arraybackslash}p{1.50cm}>{\footnotesize\raggedright\arraybackslash}p{6.40cm}@{}}
\toprule
{\footnotesize\textbf{Name}} & {\footnotesize\textbf{Size}} & {\footnotesize\textbf{Type}} & {\footnotesize\textbf{Units}} & {\footnotesize\textbf{Description}} \\
\midrule
\endfirsthead
\toprule
{\footnotesize\textbf{Name}} & {\footnotesize\textbf{Size}} & {\footnotesize\textbf{Type}} & {\footnotesize\textbf{Units}} & {\footnotesize\textbf{Description}} \\
\midrule
\endhead
\midrule
\endfoot
\bottomrule
\endlastfoot
CentreOfMass & 3x1 & double & m & Reference point of the torque, default macro.centreOfMass \\
\end{longtable}

\subsubsection*{OUTPUT ARGUMENTS}
\begin{longtable}{@{}>{\footnotesize\raggedright\arraybackslash\ttfamily}p{3.00cm}>{\footnotesize\raggedright\arraybackslash}p{1.50cm}>{\footnotesize\raggedright\arraybackslash}p{2.20cm}>{\footnotesize\raggedright\arraybackslash}p{1.50cm}>{\footnotesize\raggedright\arraybackslash}p{6.40cm}@{}}
\toprule
{\footnotesize\textbf{Name}} & {\footnotesize\textbf{Size}} & {\footnotesize\textbf{Type}} & {\footnotesize\textbf{Units}} & {\footnotesize\textbf{Description}} \\
\midrule
\endfirsthead
\toprule
{\footnotesize\textbf{Name}} & {\footnotesize\textbf{Size}} & {\footnotesize\textbf{Type}} & {\footnotesize\textbf{Units}} & {\footnotesize\textbf{Description}} \\
\midrule
\endhead
\midrule
\endfoot
\bottomrule
\endlastfoot
F & 3xK & double & N & Total force, body frame \\
tau & 3xK & double & N m & Torque about CentreOfMass, body frame \\
info & 1x1 & struct & - & .lever\_m (3xN, c\_k - r\_cm), .N, .K, .provenance \\
\end{longtable}

\subsubsection*{ERRORS}
\begin{description}[leftmargin=1.2em,style=nextline,itemsep=0pt]
\item[\texttt{NonGravForces:\allowbreak{}assemble:\allowbreak{}inputMacro}] macro is not a MacroModel
\item[\texttt{NonGravForces:\allowbreak{}assemble:\allowbreak{}inputShape}] Fk not 3xNxK with N = macro.n, or CentreOfMass not 3x1
\item[\texttt{NonGravForces:\allowbreak{}assemble:\allowbreak{}numericNaN}] NaN/Inf in Fk or CentreOfMass, or non-finite panel centroids
\end{description}

\subsection{\texttt{NonGravForces.kernel.bodyTorques}}
\label{api:kernel.bodyTorques}
\noindent Gravity-gradient and magnetic torques on the body, in the body frame; both opt-in.

\begin{lstlisting}
[tau, info] = NonGravForces.kernel.bodyTorques(macro, r_body, 'GravityGradient', true)
[tau, info] = NonGravForces.kernel.bodyTorques(macro, r_body, 'Magnetic', true, 'B_body', B, 'Dipole_Am2', m)
\end{lstlisting}

\subsubsection*{INPUT ARGUMENTS}
\begin{longtable}{@{}>{\footnotesize\raggedright\arraybackslash\ttfamily}p{3.00cm}>{\footnotesize\raggedright\arraybackslash}p{1.50cm}>{\footnotesize\raggedright\arraybackslash}p{2.20cm}>{\footnotesize\raggedright\arraybackslash}p{1.50cm}>{\footnotesize\raggedright\arraybackslash}p{6.40cm}@{}}
\toprule
{\footnotesize\textbf{Name}} & {\footnotesize\textbf{Size}} & {\footnotesize\textbf{Type}} & {\footnotesize\textbf{Units}} & {\footnotesize\textbf{Description}} \\
\midrule
\endfirsthead
\toprule
{\footnotesize\textbf{Name}} & {\footnotesize\textbf{Size}} & {\footnotesize\textbf{Type}} & {\footnotesize\textbf{Units}} & {\footnotesize\textbf{Description}} \\
\midrule
\endhead
\midrule
\endfoot
\bottomrule
\endlastfoot
macro & 1x1 & MacroModel &  & Panel model with a finite .inertia (needed by the gravity gradient) \\
r\_\allowbreak{}body & 3xK & double & m & Geocentric position rotated into the body frame (R' r\_gcrs) \\
\end{longtable}

\subsubsection*{NAME-VALUE ARGUMENTS}
\begin{longtable}{@{}>{\footnotesize\raggedright\arraybackslash\ttfamily}p{3.00cm}>{\footnotesize\raggedright\arraybackslash}p{1.50cm}>{\footnotesize\raggedright\arraybackslash}p{2.20cm}>{\footnotesize\raggedright\arraybackslash}p{1.50cm}>{\footnotesize\raggedright\arraybackslash}p{6.40cm}@{}}
\toprule
{\footnotesize\textbf{Name}} & {\footnotesize\textbf{Size}} & {\footnotesize\textbf{Type}} & {\footnotesize\textbf{Units}} & {\footnotesize\textbf{Description}} \\
\midrule
\endfirsthead
\toprule
{\footnotesize\textbf{Name}} & {\footnotesize\textbf{Size}} & {\footnotesize\textbf{Type}} & {\footnotesize\textbf{Units}} & {\footnotesize\textbf{Description}} \\
\midrule
\endhead
\midrule
\endfoot
\bottomrule
\endlastfoot
GravityGradient & 1x1 & logical & - & default false \\
Magnetic & 1x1 & logical & - & default false \\
B\_\allowbreak{}body & 3xK & double & T & Magnetic field in the body frame (IGRF via env.magneticField or MAG1B), required when Magnetic is true \\
Dipole\_\allowbreak{}Am2 & 3x1\textbar{}\allowbreak{}3xK & double & A m\textasciicircum{}2 & Residual plus commanded dipole in the body frame, required when Magnetic is true \\
GM & 1x1 & double & m\textasciicircum{}3/\allowbreak{}s\textasciicircum{}2 & default referenceSystem.iersConstants().gmEarth \\
\end{longtable}

\subsubsection*{OUTPUT ARGUMENTS}
\begin{longtable}{@{}>{\footnotesize\raggedright\arraybackslash\ttfamily}p{3.00cm}>{\footnotesize\raggedright\arraybackslash}p{1.50cm}>{\footnotesize\raggedright\arraybackslash}p{2.20cm}>{\footnotesize\raggedright\arraybackslash}p{1.50cm}>{\footnotesize\raggedright\arraybackslash}p{6.40cm}@{}}
\toprule
{\footnotesize\textbf{Name}} & {\footnotesize\textbf{Size}} & {\footnotesize\textbf{Type}} & {\footnotesize\textbf{Units}} & {\footnotesize\textbf{Description}} \\
\midrule
\endfirsthead
\toprule
{\footnotesize\textbf{Name}} & {\footnotesize\textbf{Size}} & {\footnotesize\textbf{Type}} & {\footnotesize\textbf{Units}} & {\footnotesize\textbf{Description}} \\
\midrule
\endhead
\midrule
\endfoot
\bottomrule
\endlastfoot
tau & 3xK & double & N m & Sum of the enabled torques, body frame \\
info & 1x1 & struct & - & .gravityGradient (3xK), .magnetic (3xK), .enabled (struct of logicals), .GM, .provenance \\
\end{longtable}

\subsubsection*{ERRORS}
\begin{description}[leftmargin=1.2em,style=nextline,itemsep=0pt]
\item[\texttt{NonGravForces:\allowbreak{}bodyTorques:\allowbreak{}inputMacro}] macro is not a MacroModel, or its inertia is not finite (gravity gradient requested)
\item[\texttt{NonGravForces:\allowbreak{}bodyTorques:\allowbreak{}inputShape}] r\_body not 3xK; B\_body not 3xK; Dipole\_Am2 not 3x1 or 3xK
\item[\texttt{NonGravForces:\allowbreak{}bodyTorques:\allowbreak{}inputMagnetic}] Magnetic requested without B\_body or Dipole\_Am2
\item[\texttt{NonGravForces:\allowbreak{}bodyTorques:\allowbreak{}inputRange}] \textbar{}r\textbar{} = 0 or GM \textless{}= 0
\item[\texttt{NonGravForces:\allowbreak{}bodyTorques:\allowbreak{}numericNaN}] NaN/Inf in an input
\end{description}

\subsection{\texttt{NonGravForces.kernel.partials}}
\label{api:kernel.partials}
\noindent Analytic derivatives of the per-panel kernel forces with respect to one scalar parameter.

\begin{lstlisting}
dF = NonGravForces.kernel.partials("drag", "alpha", macro, 'Rho', rho, 'Species', sp, ...
                                   'T_K', T, 'Vrel_body', v, 'Model', "sentman", 'Alpha', a, 'Tw_K', tw)
dF = NonGravForces.kernel.partials("drag", "speed", macro, ___)
dF = NonGravForces.kernel.partials("radiation", "pressureScale", macro, 'Sources', src)
dF = NonGravForces.kernel.partials("thermal", "temperature", macro, 'Temperature_K', T)
\end{lstlisting}

\subsubsection*{INPUT ARGUMENTS}
\begin{longtable}{@{}>{\footnotesize\raggedright\arraybackslash\ttfamily}p{3.00cm}>{\footnotesize\raggedright\arraybackslash}p{1.50cm}>{\footnotesize\raggedright\arraybackslash}p{2.20cm}>{\footnotesize\raggedright\arraybackslash}p{1.50cm}>{\footnotesize\raggedright\arraybackslash}p{6.40cm}@{}}
\toprule
{\footnotesize\textbf{Name}} & {\footnotesize\textbf{Size}} & {\footnotesize\textbf{Type}} & {\footnotesize\textbf{Units}} & {\footnotesize\textbf{Description}} \\
\midrule
\endfirsthead
\toprule
{\footnotesize\textbf{Name}} & {\footnotesize\textbf{Size}} & {\footnotesize\textbf{Type}} & {\footnotesize\textbf{Units}} & {\footnotesize\textbf{Description}} \\
\midrule
\endhead
\midrule
\endfoot
\bottomrule
\endlastfoot
what & 1x1 & string & - & "drag" \textbar{} "radiation" \textbar{} "thermal" \\
wrt & 1x1 & string & - & the parameter, from the list above \\
macro & 1x1 & MacroModel &  & the panel model the force was computed with \\
\end{longtable}

\subsubsection*{NAME-VALUE ARGUMENTS}
\begin{longtable}{@{}>{\footnotesize\raggedright\arraybackslash\ttfamily}p{3.00cm}>{\footnotesize\raggedright\arraybackslash}p{1.50cm}>{\footnotesize\raggedright\arraybackslash}p{2.20cm}>{\footnotesize\raggedright\arraybackslash}p{1.50cm}>{\footnotesize\raggedright\arraybackslash}p{6.40cm}@{}}
\toprule
{\footnotesize\textbf{Name}} & {\footnotesize\textbf{Size}} & {\footnotesize\textbf{Type}} & {\footnotesize\textbf{Units}} & {\footnotesize\textbf{Description}} \\
\midrule
\endfirsthead
\toprule
{\footnotesize\textbf{Name}} & {\footnotesize\textbf{Size}} & {\footnotesize\textbf{Type}} & {\footnotesize\textbf{Units}} & {\footnotesize\textbf{Description}} \\
\midrule
\endhead
\midrule
\endfoot
\bottomrule
\endlastfoot
Rho & 1xK & double & kg/\allowbreak{}m\textasciicircum{}3 & total mass density (drag) \\
Species & 8xK & double & 1/\allowbreak{}m\textasciicircum{}3 & number densities (drag; not for constantCd) \\
T\_\allowbreak{}K & 1xK & double & K & gas kinetic temperature (drag) \\
Vrel\_\allowbreak{}body & 3xK & double & m/\allowbreak{}s & relative velocity, body frame (drag) \\
Model & 1x1 & string & - & "sentman" \textbar{} "maxwell" \textbar{} "constantCd" (drag) \\
Alpha & 1x1\textbar{}\allowbreak{}1xK & double\textbar{}\allowbreak{}struct & - & accommodation coefficient or kernel.accommodation struct (drag, kinetic models) \\
Tw\_\allowbreak{}K & 1x1\textbar{}\allowbreak{}1xK & double & K & wall temperature (drag, kinetic models) \\
Cd & 1x1 & double & - & drag coefficient (drag, constantCd) \\
SpecularFraction & 1x1 & double & - & specular fraction (drag, maxwell) \\
Shadow & NxK & double\textbar{}\allowbreak{}struct & - & geom.flowShadow struct or one fraction per plate and epoch (drag), held fixed; default none \\
Sources & 1x1 & struct & - & source list (radiation), as env.radiationSources \\
RadShadow & NxKxM & double & - & self-shadow fractions (radiation), default all 1 \\
Temperature\_\allowbreak{}K & NxK & double & K & panel temperatures (thermal) \\
\end{longtable}

\subsubsection*{OUTPUT ARGUMENTS}
\begin{longtable}{@{}>{\footnotesize\raggedright\arraybackslash\ttfamily}p{3.00cm}>{\footnotesize\raggedright\arraybackslash}p{1.50cm}>{\footnotesize\raggedright\arraybackslash}p{2.20cm}>{\footnotesize\raggedright\arraybackslash}p{1.50cm}>{\footnotesize\raggedright\arraybackslash}p{6.40cm}@{}}
\toprule
{\footnotesize\textbf{Name}} & {\footnotesize\textbf{Size}} & {\footnotesize\textbf{Type}} & {\footnotesize\textbf{Units}} & {\footnotesize\textbf{Description}} \\
\midrule
\endfirsthead
\toprule
{\footnotesize\textbf{Name}} & {\footnotesize\textbf{Size}} & {\footnotesize\textbf{Type}} & {\footnotesize\textbf{Units}} & {\footnotesize\textbf{Description}} \\
\midrule
\endhead
\midrule
\endfoot
\bottomrule
\endlastfoot
dF & 3xNxK & double & N per unit of the parameter & derivative of the per-panel force \\
info & 1x1 & struct & - & .what, .wrt, .unit (the parameter's unit as a string), .K, .N, .provenance \\
\end{longtable}

\subsubsection*{ERRORS}
\begin{description}[leftmargin=1.2em,style=nextline,itemsep=0pt]
\item[\texttt{NonGravForces:\allowbreak{}partials:\allowbreak{}inputMacro}] macro is not a MacroModel
\item[\texttt{NonGravForces:\allowbreak{}partials:\allowbreak{}unsupportedWrt}] the parameter is not defined for this force (alpha of a two-state surface included)
\item[\texttt{NonGravForces:\allowbreak{}partials:\allowbreak{}inputMissing}] a name-value this combination needs was not given
\item[\texttt{NonGravForces:\allowbreak{}partials:\allowbreak{}inputShape}] an input has the wrong shape
\item[\texttt{NonGravForces:\allowbreak{}gsiCoefficients:\allowbreak{}*}] propagated from the coefficient model
\item[\texttt{NonGravForces:\allowbreak{}dragPanelForce:\allowbreak{}*, :\allowbreak{}radiationPanelForce:\allowbreak{}*, :\allowbreak{}thermalRecoil:\allowbreak{}*, :\allowbreak{}accommodation:\allowbreak{}*}] propagated: the force each derivative is formed from is evaluated (and its inputs validated) by that unit
\end{description}

\subsection{\texttt{NonGravForces.kernel.singleEpoch}}
\label{api:kernel.singleEpoch}
\noindent Compiled single-epoch kernel: gas + radiation + thermal recoil, forces and torques in one call.

\begin{lstlisting}
[F, tau] = NonGravForces.kernel.singleEpoch(macro, gas, src, Tpanel_K, 'Model', "sentman", 'Alpha', 0.85, 'Tw_K', 300)
[F, tau] = NonGravForces.kernel.singleEpoch(macro, gas, src, [], 'Model', "constantCd", 'Cd', 2.2)
[F, tau, info] = NonGravForces.kernel.singleEpoch(___, 'CdScale', 1.03, 'CrScale', 0.97)
\end{lstlisting}

\subsubsection*{INPUT ARGUMENTS}
\begin{longtable}{@{}>{\footnotesize\raggedright\arraybackslash\ttfamily}p{3.00cm}>{\footnotesize\raggedright\arraybackslash}p{1.50cm}>{\footnotesize\raggedright\arraybackslash}p{2.20cm}>{\footnotesize\raggedright\arraybackslash}p{1.50cm}>{\footnotesize\raggedright\arraybackslash}p{6.40cm}@{}}
\toprule
{\footnotesize\textbf{Name}} & {\footnotesize\textbf{Size}} & {\footnotesize\textbf{Type}} & {\footnotesize\textbf{Units}} & {\footnotesize\textbf{Description}} \\
\midrule
\endfirsthead
\toprule
{\footnotesize\textbf{Name}} & {\footnotesize\textbf{Size}} & {\footnotesize\textbf{Type}} & {\footnotesize\textbf{Units}} & {\footnotesize\textbf{Description}} \\
\midrule
\endhead
\midrule
\endfoot
\bottomrule
\endlastfoot
macro & 1x1 & MacroModel & - & Panel model: areas, unit normals, centroids, optics \\
gas & 1x1 & struct & - & .rho 1x1 kg/m\textasciicircum{}3 total mass density (used by constantCd) .species 8x1 1/m\textasciicircum{}3 number densities in const.speciesName order, or [] for constantCd .T\_K 1x1 K gas kinetic temperature (density's T(2)). READ ONLY by sentman and maxwell: constantCd neither reads nor validates it, so the NaN that env.density("exponential") returns by design is legal .vrel\_body 3x1 m/s satellite velocity relative to the gas, body frame, \textbar{}v\textbar{} \textgreater{} 0 .shadow Nx1 self-shadow fraction of each plate for the gas stream in [0,1], or [] for all 1 \\
src & 1x1 & struct & - & .dir 3xM unit vectors satellite -\textgreater{} source, body frame .P Mx1 Pa radiation pressure (the eclipse factor is already inside the Sun's) .band 1xM string "vis" \textbar{} "ir" .shadow NxM per plate and source in [0,1], or [] for all 1 M = 0 is legal and gives no radiation force \\
Tpanel\_\allowbreak{}K & Nx1 & double & K & Plate temperatures (geom.thermalState), \textgreater{}= 0, or [] for no thermal recoil (Options.thermal.model == "none") \\
\end{longtable}

\subsubsection*{NAME-VALUE ARGUMENTS}
\begin{longtable}{@{}>{\footnotesize\raggedright\arraybackslash\ttfamily}p{3.00cm}>{\footnotesize\raggedright\arraybackslash}p{1.50cm}>{\footnotesize\raggedright\arraybackslash}p{2.20cm}>{\footnotesize\raggedright\arraybackslash}p{1.50cm}>{\footnotesize\raggedright\arraybackslash}p{6.40cm}@{}}
\toprule
{\footnotesize\textbf{Name}} & {\footnotesize\textbf{Size}} & {\footnotesize\textbf{Type}} & {\footnotesize\textbf{Units}} & {\footnotesize\textbf{Description}} \\
\midrule
\endfirsthead
\toprule
{\footnotesize\textbf{Name}} & {\footnotesize\textbf{Size}} & {\footnotesize\textbf{Type}} & {\footnotesize\textbf{Units}} & {\footnotesize\textbf{Description}} \\
\midrule
\endhead
\midrule
\endfoot
\bottomrule
\endlastfoot
Model & 1x1 & string & - & "sentman" \textbar{} "maxwell" \textbar{} "constantCd" (required; no default) \\
Alpha & 1x1 & double & - & Accommodation coefficient in [0,1]; required by sentman/maxwell \\
Tw\_\allowbreak{}K & 1x1 & double & K & Wall temperature \textgreater{} 0; required by sentman/maxwell \\
Cd & 1x1 & double & - & Drag coefficient \textgreater{} 0; required by constantCd \\
SpecularFraction & 1x1 & double & - & In [0,1]; required by maxwell \\
CdScale & 1x1 & double & - & Factor on the gas force (Options.drag.cdScale), default 1 \\
CrScale & 1x1 & double & - & Factor on the radiation force (Options.srp.crScale), default 1 \\
CentreOfMass & 3x1 & double & m & Torque reference point, default macro.centreOfMass \\
\end{longtable}

\subsubsection*{OUTPUT ARGUMENTS}
\begin{longtable}{@{}>{\footnotesize\raggedright\arraybackslash\ttfamily}p{3.00cm}>{\footnotesize\raggedright\arraybackslash}p{1.50cm}>{\footnotesize\raggedright\arraybackslash}p{2.20cm}>{\footnotesize\raggedright\arraybackslash}p{1.50cm}>{\footnotesize\raggedright\arraybackslash}p{6.40cm}@{}}
\toprule
{\footnotesize\textbf{Name}} & {\footnotesize\textbf{Size}} & {\footnotesize\textbf{Type}} & {\footnotesize\textbf{Units}} & {\footnotesize\textbf{Description}} \\
\midrule
\endfirsthead
\toprule
{\footnotesize\textbf{Name}} & {\footnotesize\textbf{Size}} & {\footnotesize\textbf{Type}} & {\footnotesize\textbf{Units}} & {\footnotesize\textbf{Description}} \\
\midrule
\endhead
\midrule
\endfoot
\bottomrule
\endlastfoot
F & 3x3 & double & N & Total force, body frame; columns drag \textbar{} radiation \textbar{} thermal \\
tau & 3x3 & double & N m & Torque about CentreOfMass, same column order \\
info & 1x1 & struct & - & .Fk 3xNx3 per-panel forces, pages in the same order .gamma Nx1 n.v of the gas stream, .theta\_rad Nx1 .speedRatio 8x1 per species (NaN for constantCd) .Cp .Ctau Nx1 mixture coefficients per unit rho .absorbed\_W\_m2 Nx1 (before CrScale), .emitted\_W\_m2 Nx1 .model, .backend "mex", .kernelVersion, .provenance \\
\end{longtable}

\subsubsection*{ERRORS}
\begin{description}[leftmargin=1.2em,style=nextline,itemsep=0pt]
\item[\texttt{NonGravForces:\allowbreak{}singleEpoch:\allowbreak{}depMissing}] the MEX is not built on this machine
\item[\texttt{NonGravForces:\allowbreak{}singleEpoch:\allowbreak{}inputMacro}] macro is not a NonGravForces.geom.MacroModel
\item[\texttt{NonGravForces:\allowbreak{}singleEpoch:\allowbreak{}inputModel}] 'Model' missing, unknown, or missing its required options
\item[\texttt{NonGravForces:\allowbreak{}singleEpoch:\allowbreak{}inputShape}] a field of gas or src, Tpanel\_K or CentreOfMass has the wrong shape
\item[\texttt{NonGravForces:\allowbreak{}singleEpoch:\allowbreak{}inputSources}] src lacks dir/P/band, a band is not vis\textbar{}ir, or a direction is not a unit vector
\item[\texttt{NonGravForces:\allowbreak{}singleEpoch:\allowbreak{}inputRange}] \textbar{}v\textbar{} = 0, rho \textless{} 0, T \textless{}= 0 (sentman/maxwell), alpha outside [0,1], Tw\_K \textless{}= 0, P \textless{} 0, Tpanel\_K \textless{} 0, a shadow outside [0,1]
\item[\texttt{NonGravForces:\allowbreak{}singleEpoch:\allowbreak{}numericNaN}] NaN or Inf in an input the model reads, or a non-finite CentreOfMass or panel centroid (no torque can be formed)
\item[\texttt{NonGravForces:\allowbreak{}kernelGateway:\allowbreak{}*}] raised by the C gateway; reaching one means a contract above let something through
\end{description}

\section{L5 \texttt{Model} --- the single point of access}

\subsection{\texttt{NonGravForces.Model}}
\label{api:NonGravForces.Model}
\noindent Single point of access: non-gravitational forces, torques and partials.

\begin{lstlisting}
M = NonGravForces.Model(macro, options, 'Eop', eop, 'SpaceWeather', sw)
M = NonGravForces.Model(macro, options, 'Attitude', att, 'Mass', 601.2, 'Eop', eop, 'SpaceWeather', sw)
[a, info]        = M.accel(epoch, r_gcrs, v_gcrs, q_body2gcrs)   % 3xK m/s^2, body frame
[tau, info]      = M.torque(epoch, r, v, q)                      % 3xK N m, body frame
[a, tau, info]   = M.evaluate(epoch, r, v, q)                    % both from ONE environment pass
T                = M.series(epoch, r, v, q)                      % timetable over chunks
[J, names, info] = M.partials(epoch, r, v, q, 'WithRespectTo', ["r","v","cdScale"])
B                = M.budget(epoch, r, v, q)                      % per-effect table
\end{lstlisting}

\subsubsection*{METHOD NAME-VALUE ARGUMENTS}
Every method takes 'Frame', "body" (default) or "gcrs". These are the rest, and none of them appeared in this help until 2026-09-18: a caller could pass them, the code validated them, and the class that is this toolbox's single point of access described them nowhere.

\begin{longtable}{@{}>{\footnotesize\raggedright\arraybackslash\ttfamily}p{1.80cm}>{\footnotesize\raggedright\arraybackslash}p{3.00cm}>{\footnotesize\raggedright\arraybackslash}p{1.50cm}>{\footnotesize\raggedright\arraybackslash}p{2.20cm}>{\footnotesize\raggedright\arraybackslash}p{1.50cm}>{\footnotesize\raggedright\arraybackslash}p{4.60cm}@{}}
\toprule
{\footnotesize\textbf{Method}} & {\footnotesize\textbf{Name}} & {\footnotesize\textbf{Size}} & {\footnotesize\textbf{Type}} & {\footnotesize\textbf{Units}} & {\footnotesize\textbf{Description}} \\
\midrule
\endfirsthead
\toprule
{\footnotesize\textbf{Method}} & {\footnotesize\textbf{Name}} & {\footnotesize\textbf{Size}} & {\footnotesize\textbf{Type}} & {\footnotesize\textbf{Units}} & {\footnotesize\textbf{Description}} \\
\midrule
\endhead
\midrule
\endfoot
\bottomrule
\endlastfoot
torque, evaluate & B\_\allowbreak{}body & 3xK & double & T & Magnetic field at the satellite, BODY frame. Required by the magnetic torque and by nothing else; default [] and the term is omitted. Get it from env.magneticField, rotated to the body. \\
torque, evaluate & Dipole\_\allowbreak{}Am2 & 3x1 & double & A m\textasciicircum{}2 & Residual magnetic dipole of the spacecraft, body frame; default []. The torque is m x B, so BOTH of these are needed or neither is: one alone leaves the term out silently. torques.magnetic must also be true. \\
accel, evaluate & T0\_\allowbreak{}K & Nx1 & double & K & Initial panel temperatures for a TRANSIENT thermal model (the standard preset), one per panel; default [] and the steady state at the call's FIRST epoch is the start, so pass the last column of info.temperature\_K to carry the walls from one call into the next (series does this between its chunks). Ignored by thermal.model "none" and "steadyState". \\
series & Quantities & 1xQ & string & - & "accel" (default), "torque", or both. Selects which columns the timetable carries; a torque column costs the torque reduction, which accel alone does not pay for. \\
partials & Method & 1x1 & string & - & "" (default: the Options' partials.method), "analytic" or "finiteDiff". analytic uses kernel.partials' closed forms; finiteDiff is a central difference in the parameter, for cross-checking the analytic form. A complex step is NOT offered and is refused by name: MATLAB's erf, erfc and every norm in the chain discard the imaginary part (measured 2026-09-08). \\
partials & WithRespectTo & 1xP & string & - & Any of "r", "v" (default ["r" "v"]), "cdScale", "crScale", "alpha". alpha is refused unless drag.accommodation is "fixed" (every other model computes its own), and under finiteDiff when it sits on the edge of [0, 1]. \\
\end{longtable}
The Model owns the packed macro model, the attitude provider, the mass, the Earth-orientation, space-weather, ephemeris and CERES sources, the resolved Options and the self-shadow table. It computes forces and torques and nothing else: time, frames and attitude kinematics are referenceSystem, reading is gracefoReading, integration is gjToolbox (through app.DynamicsSource).

ONE KERNEL (principle P2). accel and torque are two reductions of the same per-panel force array F\_k: F = sum\_k F\_k and tau = sum\_k (c\_k - r\_cm) x F\_k, computed once per epoch. budget is the same array kept per effect instead of summed, and it adds back to accel bitwise. evaluate returns both reductions from one pass; series uses it when both are asked for.

\begin{lstlisting}
WHAT IT COMPOSES, in order, per epoch:
  1. the relative velocity            env.wind (co-rotation from referenceSystem)
  2. the geodetic height              referenceSystem.cartesianToGeodetic on the ITRS position
  3. density and composition          env.density with the space-weather drivers, times the observed
                                      daily scale when drag.densityScale is "swarm" (env.densityScale,
                                      observed), "driver" (env.densityDriverScale, predicted) or
                                      "combined" (env.densityScale with 'Fallback', the standard preset)
  4. the radiation source list        env.radiationSources (Sun, Earth tiles, eclipse)
  5. occlusion                        geom.selfShadow per radiation source; geom.flowShadow for the gas
                                      (Method "table" when geom.flowShadow is "table", the standard preset)
  6. panel temperatures               geom.thermalState with the per-source occlusion (steady state,
                                      or integrated over the call's epochs for "transient")
  7. accommodation                    kernel.accommodation (fixed, pilinski2010/2011, walker2014, sesam)
  8. per-panel forces                 kernel.dragPanelForce, radiationPanelForce, thermalRecoil
  9. reductions                       kernel.assemble, and kernel.bodyTorques for the torque
Every step is the unit that owns that fact; the Model adds no physics. Until
2026-09-15 steps 5 and 6 fed a radiation-beam fraction to the gas (deleting the
edge-on shear) and the Sun's fraction to every source's heating (deleting the
Earth heating of panels facing away from the Sun): HANDOVER LP-C, F1 and F2.
\end{lstlisting}
FRAME. Results are BODY frame by default (principle P4: the body frame is what an accelerometer reads and what the sensor layer maps from). Pass 'Frame', "gcrs" for the inertial frame, which is what an integrator wants.


\subsubsection*{CONSTRUCTOR INPUTS}
\begin{longtable}{@{}>{\footnotesize\raggedright\arraybackslash\ttfamily}p{3.00cm}>{\footnotesize\raggedright\arraybackslash}p{1.50cm}>{\footnotesize\raggedright\arraybackslash}p{2.20cm}>{\footnotesize\raggedright\arraybackslash}p{1.50cm}>{\footnotesize\raggedright\arraybackslash}p{6.40cm}@{}}
\toprule
{\footnotesize\textbf{Name}} & {\footnotesize\textbf{Size}} & {\footnotesize\textbf{Type}} & {\footnotesize\textbf{Units}} & {\footnotesize\textbf{Description}} \\
\midrule
\endfirsthead
\toprule
{\footnotesize\textbf{Name}} & {\footnotesize\textbf{Size}} & {\footnotesize\textbf{Type}} & {\footnotesize\textbf{Units}} & {\footnotesize\textbf{Description}} \\
\midrule
\endhead
\midrule
\endfoot
\bottomrule
\endlastfoot
macro & 1x1 & NonGravForces.\allowbreak{}geom.\allowbreak{}MacroModel & - & Packed panel geometry \\
options & 1x1 & NonGravForces.\allowbreak{}Options & - & Resolved options (default "standard", which evaluates only with 'DensityScale') \\
\end{longtable}

\subsubsection*{CONSTRUCTOR NAME-VALUE ARGUMENTS}
\begin{longtable}{@{}>{\footnotesize\raggedright\arraybackslash\ttfamily}p{3.00cm}>{\footnotesize\raggedright\arraybackslash}p{1.50cm}>{\footnotesize\raggedright\arraybackslash}p{2.20cm}>{\footnotesize\raggedright\arraybackslash}p{1.50cm}>{\footnotesize\raggedright\arraybackslash}p{6.40cm}@{}}
\toprule
{\footnotesize\textbf{Name}} & {\footnotesize\textbf{Size}} & {\footnotesize\textbf{Type}} & {\footnotesize\textbf{Units}} & {\footnotesize\textbf{Description}} \\
\midrule
\endfirsthead
\toprule
{\footnotesize\textbf{Name}} & {\footnotesize\textbf{Size}} & {\footnotesize\textbf{Type}} & {\footnotesize\textbf{Units}} & {\footnotesize\textbf{Description}} \\
\midrule
\endhead
\midrule
\endfoot
\bottomrule
\endlastfoot
Attitude & 1x1 & AttitudeProvider & - & BODY-\textgreater{}GCRS quaternions, used when a method is called without q. Default []. \\
Mass & 1x1 & double\textbar{}\allowbreak{}struct & kg & A constant mass, or an io.massModel result. Default: the macro model's mass\_kg (geom.graceFO carries the handbook value); a macro model without one is REFUSED, never given GRACE-C's (LP-C F10). \\
Eop & 1x1 & struct & - & referenceSystem.readEop result or a file path. REQUIRED to evaluate: the geodetic height needs the ITRS rotation and the co-rotation its rate. \\
SpaceWeather & 1x1 & struct & - & io.spaceWeather result. Required by msis00, msis21 and dtm2020, and by drag.wind "hwm14". \\
Ephemeris & 1x1 & struct & - & referenceSystem.readEphemeris result, required when sun.source is "rs.de440". \\
Exponential & 1x1 & struct & - & .rho0\_kg\_m3 .h0\_m .H\_m, required when drag.density is "exponential". \\
Ceres & 1x1 & struct & - & io.importCERES result, required when erp.backend is "ceres". \\
Climatology & 1x1 & struct & - & .albedo and .longwave io.importGroopsGrid results, required when erp.backend is "ceresClimatology". \\
DensityScale & 1x1 & struct & - & env.densityScale result (drag.densityScale "swarm"; with 'Fallback' "combined") or env.densityDriverScale result ("driver"), required by those (so by the standard preset) and refused otherwise; its kind must match the option and it must have been formed for this Model's drag.density. Density and every species density are multiplied by S.at(epoch, geodetic height), so the gas force scales exactly (it is linear in them). \\
\end{longtable}

\subsubsection*{OUTPUT ARGUMENTS (accel / torque / evaluate)}
\begin{longtable}{@{}>{\footnotesize\raggedright\arraybackslash\ttfamily}p{3.00cm}>{\footnotesize\raggedright\arraybackslash}p{1.50cm}>{\footnotesize\raggedright\arraybackslash}p{2.20cm}>{\footnotesize\raggedright\arraybackslash}p{7.90cm}@{}}
\toprule
{\footnotesize\textbf{Name}} & {\footnotesize\textbf{Size}} & {\footnotesize\textbf{Type}} & {\footnotesize\textbf{Description}} \\
\midrule
\endfirsthead
\toprule
{\footnotesize\textbf{Name}} & {\footnotesize\textbf{Size}} & {\footnotesize\textbf{Type}} & {\footnotesize\textbf{Description}} \\
\midrule
\endhead
\midrule
\endfoot
\bottomrule
\endlastfoot
a, tau & 3xK & double & acceleration [m/s\textasciicircum{}2] or torque [N m] \\
info & 1x1 & struct & .frame, .K, .drag .radiation .thermal (3xK each, the same frame; torques from torque), .bodyTorque (torque and evaluate), .torque (evaluate only: .drag .radiation .thermal torques), .rho (1xK), .nu (1xK shadow factor), .mass\_kg (1xK), .vrel\_body, .sources (number M of radiation sources), .temperature\_K (NxK), .flowShadow, .accommodation, .lla, .densityScale (1xK, ones with drag.densityScale "none"), .backends, .provenance \\
OUTPUT & (series & /\allowbreak{} budget /\allowbreak{} partials) &  \\
T & timetable &  & one row per epoch: ax ay az [m/s\textasciicircum{}2] and/or tx ty tz [N m] \\
B & table &  & rows drag, radiation, thermal, total; rms\_x..rms\_z, max\_x..max\_z, rms\_norm [m/s\textasciicircum{}2] \\
J & 3xPxK & double & d a / d parameter, one column per entry of names \\
names & 1xP & string & "rx" "ry" "rz" "vx" "vy" "vz" and the scalar parameters \\
info & 1x1 & struct & .method .frame .stateStep .analytic .K .provenance \\
\end{longtable}

\subsubsection*{ERRORS}
\begin{description}[leftmargin=1.2em,style=nextline,itemsep=0pt]
\item[\texttt{NonGravForces:\allowbreak{}Model:\allowbreak{}inputMacro}] macro is not a MacroModel
\item[\texttt{NonGravForces:\allowbreak{}Model:\allowbreak{}inputOptions}] options is not an Options object, or its choices contradict each other (alpha beside a non-fixed accommodation, maxwell without a specular fraction, walker2014 without a surface mass, internal power without shares, geom.flowShadow table beside selfShadow raycast, a 'DensityScale' with drag.densityScale "none", of the other kind, or formed for another drag.density)
\item[\texttt{NonGravForces:\allowbreak{}Model:\allowbreak{}inputMass}] no 'Mass' and no mass on the macro model, or 'Mass' not a positive mass / massModel
\item[\texttt{NonGravForces:\allowbreak{}Model:\allowbreak{}inputState}] epoch, r, v or q disagree in size, or are not finite, or the attitude provider returned the wrong number of quaternions
\item[\texttt{NonGravForces:\allowbreak{}Model:\allowbreak{}inputAttitude}] no q was given and the Model has no attitude provider
\item[\texttt{NonGravForces:\allowbreak{}Model:\allowbreak{}dependEop}] evaluated without 'Eop'
\item[\texttt{NonGravForces:\allowbreak{}Model:\allowbreak{}dependWeather}] evaluated without 'SpaceWeather' or 'Exponential' for the density backend, or without 'SpaceWeather' for drag.wind hwm14
\item[\texttt{NonGravForces:\allowbreak{}Model:\allowbreak{}dependDensityScale}] evaluated with drag.densityScale "swarm", "driver" or "combined" and no 'DensityScale'
\item[\texttt{NonGravForces:\allowbreak{}Model:\allowbreak{}dependCeres}] evaluated with erp.backend ceres and no 'Ceres', or ceresClimatology and no 'Climatology'
\item[\texttt{NonGravForces:\allowbreak{}Model:\allowbreak{}unsupported}] a method option this Model cannot honour (partials of alpha for a computed accommodation or, by finiteDiff, at the edge of [0, 1]; an unknown 'WithRespectTo' entry or 'Method'; transient thermal under series.useParallel; unknown 'Quantities')
\item[\texttt{NonGravForces:\allowbreak{}\textless{}unit\textgreater{}:\allowbreak{}*}] propagated unwrapped from the unit that refused
\end{description}

\section{L6 \texttt{sensor} --- instrument models: true dynamics mapped to what one instrument reads}

\subsection{\texttt{NonGravForces.sensor.Sensor}}
\label{api:sensor.Sensor}
\noindent Abstract measurement model: maps true non-gravitational dynamics to an observable.

\subsubsection*{ABSTRACT METHODS (every subclass implements)}
\begin{lstlisting}
y          = observe(obj, dyn)           dyn: struct .a .tau .omega .omegaDot
                                         .epoch; y: struct .value (MxK)
                                         .epoch .frame .mask (logical MxK)
[J, names] = partials(obj, dyn)          d(y)/d(theta), M*K x P, theta named
S          = noiseModel(obj, f_Hz)       one-sided ASD [unit/sqrt(Hz)] at f
n          = numParameters(obj)
\end{lstlisting}

\subsubsection*{PROPERTIES}
\begin{lstlisting}
name        string   instrument label
axes        string   e.g. ["x","y","z"] in the sensor frame
sampleRate  double   Hz
\end{lstlisting}

\subsubsection*{ERRORS}
Subclasses use NonGravForces.fail("\textless{}Subclass\textgreater{}", ...).


\subsection{\texttt{NonGravForces.sensor.Electrostatic}}
\label{api:sensor.Electrostatic}
\noindent Electrostatic accelerometer: bias, 3x3 scale, proof-mass offset coupling; L1B CRN filter by default.

\begin{lstlisting}
s = NonGravForces.sensor.Electrostatic(bias, S, dr)                       % CRN-filtered, 10 Hz input
s = NonGravForces.sensor.Electrostatic(bias, S, dr, "crn", 'OutputRate', 1) % as ACT1B: 1 Hz samples
s = NonGravForces.sensor.Electrostatic(bias, S, dr, [])                   % unfiltered, any rate
s = NonGravForces.sensor.Electrostatic(bias, S, dr, kernel, 'Rotation', R, 'SampleRate', fs)
y = s.observe(dyn);  [J, names] = s.partials(dyn);  asd = s.noiseModel(f)
\end{lstlisting}

\subsubsection*{CONSTRUCTOR ARGUMENTS}
\begin{longtable}{@{}>{\footnotesize\raggedright\arraybackslash\ttfamily}p{3.00cm}>{\footnotesize\raggedright\arraybackslash}p{1.50cm}>{\footnotesize\raggedright\arraybackslash}p{2.20cm}>{\footnotesize\raggedright\arraybackslash}p{1.50cm}>{\footnotesize\raggedright\arraybackslash}p{6.40cm}@{}}
\toprule
{\footnotesize\textbf{Name}} & {\footnotesize\textbf{Size}} & {\footnotesize\textbf{Type}} & {\footnotesize\textbf{Units}} & {\footnotesize\textbf{Description}} \\
\midrule
\endfirsthead
\toprule
{\footnotesize\textbf{Name}} & {\footnotesize\textbf{Size}} & {\footnotesize\textbf{Type}} & {\footnotesize\textbf{Units}} & {\footnotesize\textbf{Description}} \\
\midrule
\endhead
\midrule
\endfoot
\bottomrule
\endlastfoot
bias & 3x1 & double & m/\allowbreak{}s\textasciicircum{}2 & Additive bias in the accelerometer frame; default zeros \\
S & 3x3 & double & - & Scale matrix; a 3x1 is taken as its diagonal; default eye(3) \\
dr & 3x1 & double & m & Proof-mass offset from the CoM, body frame; default zeros \\
kernel & 1x1 & string & - & "crn" (default): timeSeriesPreparation.crnKernel("ACC") 1x1 struct - a timeSeriesPreparation.crnKernel result (e.g. other options); its .fs must equal 'SampleRate' 1xL double - any FIR, L odd, sum \textasciitilde{}= 0, applied at the caller's rate; [] = no filter \\
\end{longtable}

\subsubsection*{NAME-VALUE ARGUMENTS}
\begin{longtable}{@{}>{\footnotesize\raggedright\arraybackslash\ttfamily}p{3.00cm}>{\footnotesize\raggedright\arraybackslash}p{1.50cm}>{\footnotesize\raggedright\arraybackslash}p{2.20cm}>{\footnotesize\raggedright\arraybackslash}p{1.50cm}>{\footnotesize\raggedright\arraybackslash}p{6.40cm}@{}}
\toprule
{\footnotesize\textbf{Name}} & {\footnotesize\textbf{Size}} & {\footnotesize\textbf{Type}} & {\footnotesize\textbf{Units}} & {\footnotesize\textbf{Description}} \\
\midrule
\endfirsthead
\toprule
{\footnotesize\textbf{Name}} & {\footnotesize\textbf{Size}} & {\footnotesize\textbf{Type}} & {\footnotesize\textbf{Units}} & {\footnotesize\textbf{Description}} \\
\midrule
\endhead
\midrule
\endfoot
\bottomrule
\endlastfoot
Rotation & 3x3 & double & - & Body -\textgreater{} accelerometer frame; default eye(3) \\
SampleRate & 1x1 & double & Hz & Rate of dyn; must equal a CRN kernel's fs; default 10 \\
OutputRate & 1x1 & double & Hz & Rate of y: every (SampleRate/OutputRate)-th filtered sample from the first, an integer ratio; default 'SampleRate' (no sampling). 1 reproduces ACT1B's rate \\
NoiseAsd & 3x1 & double & m/\allowbreak{}s\textasciicircum{}2/\allowbreak{}rtHz & White floor per axis; default [1e-10; 1e-9; 1e-10], the GRACE-FO order with the cross-track axis ten times worse (Kornfeld et al. 2019) \\
NoiseCorner & 1x1 & double & Hz & Below it the ASD rises as 1/f; default 5e-3 \\
Name & 1x1 & string & - & Instrument label; default "electrostatic" \\
\end{longtable}

\subsubsection*{METHODS}
\begin{lstlisting}
y          = observe(obj, dyn)     dyn: .a 3xK body, .omega 3xK body, .omegaDot 3xK body,
                                   .epoch (1xK s, datetime, or a referenceSystem epoch with
                                   .day/.sec; REQUIRED with a CRN kernel, optional otherwise),
                                   optional .gravityGradient 3x3xK.
                                   y: .value 3xM (accelerometer frame, M = ceil(K/stride)),
                                   .mask 3xM logical, .frame, .epoch (sampled like .value),
                                   .info (.K, .filtered, .kernel "crn ACC" | "crn <preset>" |
                                   "user" | "none", .edgeSamples, .outputStride, .gradient
                                   "none" | "supplied", .provenance)
[J, names] = partials(obj, dyn)    d(y)/d(theta), (3M)x15, theta = [bias; S(:); dr].
                                   Row (k-1)*3+i is axis i of epoch k: axis fastest
asd        = noiseModel(obj, f_Hz) 3xN one-sided amplitude spectral density
n          = numParameters(obj)    15
S          = Electrostatic.skew(v) static: the 3x3 matrix with S w = v x w
\end{lstlisting}

\subsubsection*{ERRORS}
\begin{description}[leftmargin=1.2em,style=nextline,itemsep=0pt]
\item[\texttt{NonGravForces:\allowbreak{}Electrostatic:\allowbreak{}inputShape}] an argument or a dyn field has the wrong size, 'OutputRate' does not divide 'SampleRate' by a whole number, 'NoiseAsd' is negative, or noiseModel's f\_Hz is not positive and finite
\item[\texttt{NonGravForces:\allowbreak{}Electrostatic:\allowbreak{}inputKernel}] the kernel is empty-summed, of even length, not a vector, text other than "crn", a struct without .F/.fs, or a CRN kernel whose .fs differs from 'SampleRate'
\item[\texttt{NonGravForces:\allowbreak{}Electrostatic:\allowbreak{}depMissing}] "crn" asked for and timeSeriesPreparation is not on the path
\item[\texttt{NonGravForces:\allowbreak{}Electrostatic:\allowbreak{}inputDyn}] dyn lacks a, omega or omegaDot, or their epoch counts differ; with a CRN kernel: dyn.epoch missing, of the wrong length or type, or not on the kernel's grid
\item[\texttt{NonGravForces:\allowbreak{}Electrostatic:\allowbreak{}numericNaN}] an input contains NaN or Inf
\end{description}

\subsection{\texttt{NonGravForces.sensor.AtomInterferometer}}
\label{api:sensor.AtomInterferometer}
\noindent Cold-atom interferometer: phase = k\_eff a T\textasciicircum{}2, with dead time, rotation and the 2 pi ambiguity.

\begin{lstlisting}
s = NonGravForces.sensor.AtomInterferometer(kEff, T, cycle, axis)
s = NonGravForces.sensor.AtomInterferometer(kEff, T, cycle, axis, 'Rotation', true, 'SampleRate', 10)
y = s.observe(dyn);  [J, names] = s.partials(dyn);  asd = s.noiseModel(f)
\end{lstlisting}

\subsubsection*{CONSTRUCTOR ARGUMENTS}
\begin{longtable}{@{}>{\footnotesize\raggedright\arraybackslash\ttfamily}p{3.00cm}>{\footnotesize\raggedright\arraybackslash}p{1.50cm}>{\footnotesize\raggedright\arraybackslash}p{2.20cm}>{\footnotesize\raggedright\arraybackslash}p{1.50cm}>{\footnotesize\raggedright\arraybackslash}p{6.40cm}@{}}
\toprule
{\footnotesize\textbf{Name}} & {\footnotesize\textbf{Size}} & {\footnotesize\textbf{Type}} & {\footnotesize\textbf{Units}} & {\footnotesize\textbf{Description}} \\
\midrule
\endfirsthead
\toprule
{\footnotesize\textbf{Name}} & {\footnotesize\textbf{Size}} & {\footnotesize\textbf{Type}} & {\footnotesize\textbf{Units}} & {\footnotesize\textbf{Description}} \\
\midrule
\endhead
\midrule
\endfoot
\bottomrule
\endlastfoot
kEff & 1x1 & double & rad/\allowbreak{}m & Effective wave vector, \textgreater{} 0. Two-photon Raman on the rubidium D2 line is 4 pi / 780.24 nm = 1.610e7 rad/m \\
T & 1x1 & double & s & Interrogation time (half the pulse separation), \textgreater{} 0 \\
cycle & 1x1 & double & s & Cycle time, \textgreater{}= 2 T. The dead time is cycle - 2 T \\
axis & 3x1 & double & - & Raman beam direction in the body frame; normalised on entry \\
\end{longtable}

\subsubsection*{NAME-VALUE ARGUMENTS}
\begin{longtable}{@{}>{\footnotesize\raggedright\arraybackslash\ttfamily}p{3.00cm}>{\footnotesize\raggedright\arraybackslash}p{1.50cm}>{\footnotesize\raggedright\arraybackslash}p{2.20cm}>{\footnotesize\raggedright\arraybackslash}p{1.50cm}>{\footnotesize\raggedright\arraybackslash}p{6.40cm}@{}}
\toprule
{\footnotesize\textbf{Name}} & {\footnotesize\textbf{Size}} & {\footnotesize\textbf{Type}} & {\footnotesize\textbf{Units}} & {\footnotesize\textbf{Description}} \\
\midrule
\endfirsthead
\toprule
{\footnotesize\textbf{Name}} & {\footnotesize\textbf{Size}} & {\footnotesize\textbf{Type}} & {\footnotesize\textbf{Units}} & {\footnotesize\textbf{Description}} \\
\midrule
\endhead
\midrule
\endfoot
\bottomrule
\endlastfoot
Rotation & 1x1 & logical & - & Include the Coriolis phase; default false. When true, dyn must carry .omega and .atomVelocity \\
SampleRate & 1x1 & double & Hz & Sample rate of dyn, \textgreater{} 0; default 1. Sets the grid the triangular weighting and the cycle are counted on \\
PhaseBias & 1x1 & double & rad & Additive phase offset; default 0 \\
GainError & 1x1 & double & - & Relative error of k\_eff T\textasciicircum{}2; default 0 \\
PhaseAsd & 1x1 & double & rad/\allowbreak{}rtHz & White phase noise floor; default 1e-3 \\
Name & 1x1 & string & - & Instrument label; default "atomInterferometer" \\
\end{longtable}

\subsubsection*{METHODS}
\begin{lstlisting}
y          = observe(obj, dyn)     dyn: .a 3xK body [m/s^2]; when 'Rotation', .omega 3xK
                                   [rad/s] and .atomVelocity 3xK [m/s], the atom cloud
                                   relative to the instrument, body frame.
                                   dyn.velocity (retired 2026-09-15) is refused.
                                   y: .value 1xK full phase [rad], .mask 1xK,
                                   .frame "beam", .epoch (dyn.epoch if given, else []),
                                   .info (.wrapped .fringe .weighted .coriolis
                                   .cycleEpochs .deadTime_s .provenance)
[J, names] = partials(obj, dyn)    d(phase)/d([phaseBias; gainError]), Kx2
asd        = noiseModel(obj, f_Hz) 1xN phase ASD [rad/sqrt(Hz)], 'PhaseAsd' up to the
                                   Nyquist frequency of the cycle rate, 1/(2 cycle), and
                                   Inf above it (nothing is measured there); divide by
                                   k_eff T^2 for the equivalent acceleration ASD
n          = numParameters(obj)    2
\end{lstlisting}

\subsubsection*{ERRORS}
\begin{description}[leftmargin=1.2em,style=nextline,itemsep=0pt]
\item[\texttt{NonGravForces:\allowbreak{}AtomInterferometer:\allowbreak{}inputShape}] axis or dyn.a has the wrong size, or noiseModel's f\_Hz is not positive and finite
\item[\texttt{NonGravForces:\allowbreak{}AtomInterferometer:\allowbreak{}inputRange}] kEff \textless{}= 0, T \textless{}= 0, cycle \textless{} 2T, SampleRate \textless{}= 0, PhaseAsd \textless{} 0, or the axis has zero length
\item[\texttt{NonGravForces:\allowbreak{}AtomInterferometer:\allowbreak{}inputDyn}] dyn lacks a, or lacks omega/atomVelocity when 'Rotation' is on, or their epoch counts differ, or it carries the retired field velocity
\item[\texttt{NonGravForces:\allowbreak{}AtomInterferometer:\allowbreak{}numericNaN}] an input contains NaN or Inf
\end{description}

\subsection{\texttt{NonGravForces.sensor.DragFree}}
\label{api:sensor.DragFree}
\noindent Residual acceleration left by a drag-free control loop: sensitivity high-pass and actuator saturation.

\begin{lstlisting}
s = NonGravForces.sensor.DragFree(bandwidth_Hz, order)
s = NonGravForces.sensor.DragFree(1e-3, 2, 'Authority', 1e-6, 'SampleRate', 10)
y = s.observe(dyn);  [J, names] = s.partials(dyn);  asd = s.noiseModel(f)
\end{lstlisting}

\subsubsection*{CONSTRUCTOR ARGUMENTS}
\begin{longtable}{@{}>{\footnotesize\raggedright\arraybackslash\ttfamily}p{3.00cm}>{\footnotesize\raggedright\arraybackslash}p{1.50cm}>{\footnotesize\raggedright\arraybackslash}p{2.20cm}>{\footnotesize\raggedright\arraybackslash}p{1.50cm}>{\footnotesize\raggedright\arraybackslash}p{6.40cm}@{}}
\toprule
{\footnotesize\textbf{Name}} & {\footnotesize\textbf{Size}} & {\footnotesize\textbf{Type}} & {\footnotesize\textbf{Units}} & {\footnotesize\textbf{Description}} \\
\midrule
\endfirsthead
\toprule
{\footnotesize\textbf{Name}} & {\footnotesize\textbf{Size}} & {\footnotesize\textbf{Type}} & {\footnotesize\textbf{Units}} & {\footnotesize\textbf{Description}} \\
\midrule
\endhead
\midrule
\endfoot
\bottomrule
\endlastfoot
bandwidth\_\allowbreak{}Hz & 1x1 & double & Hz & Control bandwidth f\_c, \textgreater{} 0. Below it the loop rejects \\
order & 1x1 & double & - & Number of first-order sections, integer \textgreater{}= 1; default 1 \\
\end{longtable}

\subsubsection*{NAME-VALUE ARGUMENTS}
\begin{longtable}{@{}>{\footnotesize\raggedright\arraybackslash\ttfamily}p{3.00cm}>{\footnotesize\raggedright\arraybackslash}p{1.50cm}>{\footnotesize\raggedright\arraybackslash}p{2.20cm}>{\footnotesize\raggedright\arraybackslash}p{1.50cm}>{\footnotesize\raggedright\arraybackslash}p{6.40cm}@{}}
\toprule
{\footnotesize\textbf{Name}} & {\footnotesize\textbf{Size}} & {\footnotesize\textbf{Type}} & {\footnotesize\textbf{Units}} & {\footnotesize\textbf{Description}} \\
\midrule
\endfirsthead
\toprule
{\footnotesize\textbf{Name}} & {\footnotesize\textbf{Size}} & {\footnotesize\textbf{Type}} & {\footnotesize\textbf{Units}} & {\footnotesize\textbf{Description}} \\
\midrule
\endhead
\midrule
\endfoot
\bottomrule
\endlastfoot
Authority & 1x1 & double & m/\allowbreak{}s\textasciicircum{}2 & Maximum acceleration the actuators can command, \textgreater{} 0; default Inf (an ideal loop) \\
SampleRate & 1x1 & double & Hz & Sample rate of dyn, \textgreater{} 0; default 1. Must exceed 2 f\_c or the bandwidth is not on the grid \\
Bias & 3x1 & double & m/\allowbreak{}s\textasciicircum{}2 & Additive residual offset; default zeros \\
GainError & 1x1 & double & - & Relative error of the residual; default 0 \\
ThrusterAsd & 3x1 & double & m/\allowbreak{}s\textasciicircum{}2/\allowbreak{}rtHz & Actuation noise floor injected inside the bandwidth, \textgreater{}= 0; a scalar applies to all three axes; default 1e-11 per axis \\
Name & 1x1 & string & - & Label; default "dragFree" \\
\end{longtable}

\subsubsection*{METHODS}
\begin{lstlisting}
y          = observe(obj, dyn)     dyn: .a 3xK non-gravitational acceleration, body frame.
                                   y: .value 3xK residual, .mask 3xK, .frame "body",
                                   .epoch (dyn.epoch if given, else []), .info (.commanded
                                   .saturated .settleSamples .gainAtBandwidth .provenance)
[J, names] = partials(obj, dyn)    d(residual)/d([bias; gainError]), (3K)x4, axis fastest
asd        = noiseModel(obj, f_Hz) 3xN residual ASD: thruster noise shaped by the
                                   complementary sensitivity, i.e. injected where the loop
                                   is active and absent where it is not
n          = numParameters(obj)    4
m          = magnitudeResponse(obj, f_Hz)   |S(f)| of the cascade, evaluated on the unit
                                   circle from the difference-equation coefficients, so a
                                   caller (and a test) can ask what this loop does to a
                                   frequency without running a series through it
\end{lstlisting}

\subsubsection*{ERRORS}
\begin{description}[leftmargin=1.2em,style=nextline,itemsep=0pt]
\item[\texttt{NonGravForces:\allowbreak{}DragFree:\allowbreak{}inputShape}] Bias, ThrusterAsd or dyn.a has the wrong size, or noiseModel's f\_Hz is not positive and finite
\item[\texttt{NonGravForces:\allowbreak{}DragFree:\allowbreak{}inputRange}] bandwidth \textless{}= 0, order \textless{} 1 or not an integer, Authority \textless{}= 0, ThrusterAsd \textless{} 0, SampleRate \textless{}= 2 f\_c
\item[\texttt{NonGravForces:\allowbreak{}DragFree:\allowbreak{}inputDyn}] dyn lacks a
\item[\texttt{NonGravForces:\allowbreak{}DragFree:\allowbreak{}numericNaN}] an input contains NaN or Inf
\end{description}

\section{L8 \texttt{app} --- applications built from the layers below}

\subsection{\texttt{NonGravForces.app.calibrate}}
\label{api:app.calibrate}
\noindent Calibrate an accelerometer arc against a NonGravForces Model or a reference arc, through timeSeriesPreparation.run and lsqvce.fit.

\begin{lstlisting}
res = NonGravForces.app.calibrate(acc, M, 'Orbit', orb, 'Attitude', att)   % a priori from the Model
res = NonGravForces.app.calibrate(acc, refArc)                              % an external reference
res = NonGravForces.app.calibrate(acc, M, 'Orbit', orb, 'Fit', {'BiasBasis', "pspline", 'Scale', "full"})
res = NonGravForces.app.calibrate(accDays, M, 'Orbit', orbDays)            % N days, one scale (N-09)
\end{lstlisting}

\subsubsection*{INPUT ARGUMENTS}
\begin{longtable}{@{}>{\footnotesize\raggedright\arraybackslash\ttfamily}p{3.00cm}>{\footnotesize\raggedright\arraybackslash}p{1.50cm}>{\footnotesize\raggedright\arraybackslash}p{2.20cm}>{\footnotesize\raggedright\arraybackslash}p{1.50cm}>{\footnotesize\raggedright\arraybackslash}p{6.40cm}@{}}
\toprule
{\footnotesize\textbf{Name}} & {\footnotesize\textbf{Size}} & {\footnotesize\textbf{Type}} & {\footnotesize\textbf{Units}} & {\footnotesize\textbf{Description}} \\
\midrule
\endfirsthead
\toprule
{\footnotesize\textbf{Name}} & {\footnotesize\textbf{Size}} & {\footnotesize\textbf{Type}} & {\footnotesize\textbf{Units}} & {\footnotesize\textbf{Description}} \\
\midrule
\endhead
\midrule
\endfoot
\bottomrule
\endlastfoot
acc & 1xN & timeSeriesPreparation.\allowbreak{}Arc & m/\allowbreak{}s\textasciicircum{}2 & Observation, three axes in m/s2 (e.g. ACT1B, lin\_accl\_x/y/z); its product selects the preset \\
reference & 1x1 & NonGravForces.\allowbreak{}Model or Arc & m/\allowbreak{}s\textasciicircum{}2 & The a priori: a Model (needs 'Orbit'), or an Arc with acc's three columns on the raw grid \\
\end{longtable}

\subsubsection*{NAME-VALUE ARGUMENTS}
\begin{longtable}{@{}>{\footnotesize\raggedright\arraybackslash\ttfamily}p{3.00cm}>{\footnotesize\raggedright\arraybackslash}p{1.50cm}>{\footnotesize\raggedright\arraybackslash}p{2.20cm}>{\footnotesize\raggedright\arraybackslash}p{1.50cm}>{\footnotesize\raggedright\arraybackslash}p{6.40cm}@{}}
\toprule
{\footnotesize\textbf{Name}} & {\footnotesize\textbf{Size}} & {\footnotesize\textbf{Type}} & {\footnotesize\textbf{Units}} & {\footnotesize\textbf{Description}} \\
\midrule
\endfirsthead
\toprule
{\footnotesize\textbf{Name}} & {\footnotesize\textbf{Size}} & {\footnotesize\textbf{Type}} & {\footnotesize\textbf{Units}} & {\footnotesize\textbf{Description}} \\
\midrule
\endhead
\midrule
\endfoot
\bottomrule
\endlastfoot
Orbit & 1xN & struct & - & .epoch (referenceSystem epoch, GPS), .r and .v (3xK, GCRS, m and m/s), e.g. from io.readGNV1B; required with a Model. It should cover the arc plus half a CRN window; it is clipped to its overlap with that span and refused only without any overlap \\
Attitude & 1x1 & object & - & geom.AttitudeProvider, evaluated at the orbit epochs; [] (default) uses the Model's own provider \\
Config & 1x1 & object & - & timeSeriesPreparation.Config; default Config(acc) (of the first arc in a campaign); run.filterReference is set true either way \\
QualFlags & 1x1 & struct & - & Passed to run's mask stage; default none \\
Thrusters & 1x1 & struct & - & THR1B Arc, passed to run's mask stage; default none \\
Jumps & 1x1 & logical & - & Turn run's "jump" events into bias steps; default false (they are reported in diagnostics.jumps either way) \\
VifMax & 1x1 & double & - & Attribution's VIF gate; default 10 \\
Fit & 1xN & cell & - & Name-value pairs for app.fitCalibration (BiasBasis, Scale, ScaleBasis for a time-varying scale, Noise, NBoot, ...); default \{\} \\
\end{longtable}

\subsubsection*{OUTPUT ARGUMENTS}
\begin{longtable}{@{}>{\footnotesize\raggedright\arraybackslash\ttfamily}p{3.00cm}>{\footnotesize\raggedright\arraybackslash}p{1.50cm}>{\footnotesize\raggedright\arraybackslash}p{2.20cm}>{\footnotesize\raggedright\arraybackslash}p{1.50cm}>{\footnotesize\raggedright\arraybackslash}p{6.40cm}@{}}
\toprule
{\footnotesize\textbf{Name}} & {\footnotesize\textbf{Size}} & {\footnotesize\textbf{Type}} & {\footnotesize\textbf{Units}} & {\footnotesize\textbf{Description}} \\
\midrule
\endfirsthead
\toprule
{\footnotesize\textbf{Name}} & {\footnotesize\textbf{Size}} & {\footnotesize\textbf{Type}} & {\footnotesize\textbf{Units}} & {\footnotesize\textbf{Description}} \\
\midrule
\endhead
\midrule
\endfoot
\bottomrule
\endlastfoot
res & 1x1 & struct & - & timeSeriesPreparation.run's fields (arc, reference, diagnostics, correlation, flags, config, provenance), diagnostics.fit (fitCalibration's info), diagnostics.rawReference (the reference on the raw grid, as fitted) and diagnostics.jumps (1xJ int64, qc's jump epochs in GPS ns, fitted only with 'Jumps' true), then .calibration (app.Calibration) and .attribution (.effects 1xE string, .rms ExE3 a priori RMS of the filtered effects over the fitted rows, .share, .relativeCorrection Ex3, .sigma Ex3, .vif Ex3, .verdict 1x3 "resolved" \textbar{} "unresolved" \textbar{} "unavailable", .note) \\
\end{longtable}

\subsubsection*{ERRORS}
\begin{description}[leftmargin=1.2em,style=nextline,itemsep=0pt]
\item[\texttt{NonGravForces:\allowbreak{}calibrate:\allowbreak{}inputArc}] acc is not a three-axis Arc in m/s2
\item[\texttt{NonGravForces:\allowbreak{}calibrate:\allowbreak{}inputReference}] reference is neither a Model nor such an Arc, one arc with several references, or a campaign's reference array whose length is neither 1 nor N
\item[\texttt{NonGravForces:\allowbreak{}calibrate:\allowbreak{}inputOrbit}] a Model without 'Orbit', a malformed Orbit (fields, 3xK shapes, increasing epochs; a campaign's not 1 or 1xN), or an orbit that does not overlap the arc
\item[\texttt{NonGravForces:\allowbreak{}calibrate:\allowbreak{}inputConfig}] the Config names no CRN preset
\item[\texttt{NonGravForces:\allowbreak{}calibrate:\allowbreak{}dataEmpty}] a campaign in which run refused every day
\item[\texttt{NonGravForces:\allowbreak{}calibrate:\allowbreak{}unsupportedMultiArc}] 'Jumps' true in a campaign
\item[\texttt{NonGravForces:\allowbreak{}fitCalibration:\allowbreak{}*}] the estimate
\item[\texttt{NonGravForces:\allowbreak{}Model:\allowbreak{}*}] the Model evaluation
\item[\texttt{timeSeriesPreparation:\allowbreak{}*, lsqvce:\allowbreak{}*, referenceSystem:\allowbreak{}*}] propagated unwrapped
\end{description}

\subsection{\texttt{NonGravForces.app.DynamicsSource}}
\label{api:app.DynamicsSource}
\noindent The non-gravitational acceleration as gjToolbox's frozen integrator interface, with a native C entry.

\begin{lstlisting}
src = NonGravForces.app.DynamicsSource(model, 'Epoch', tai)
src = NonGravForces.app.DynamicsSource(model, 'Epoch', tai, 'Gravity', gravitySource)
src = NonGravForces.app.DynamicsSource(model, 'Epoch', tai, 'Parameters', ["cdScale" "crScale"])
src = NonGravForces.app.DynamicsSource(model, 'Epoch', tai, 'ThermalOrbit', orb)   % thermal.model transient
sol = gjToolbox.propagate(src, [0 5400], r0, v0, gjToolbox.Options('Epoch', tai, 'Step', 5, 'Binding', "native"))
\end{lstlisting}

\subsubsection*{CONSTRUCTOR ARGUMENTS}
\begin{longtable}{@{}>{\footnotesize\raggedright\arraybackslash\ttfamily}p{3.00cm}>{\footnotesize\raggedright\arraybackslash}p{1.50cm}>{\footnotesize\raggedright\arraybackslash}p{2.20cm}>{\footnotesize\raggedright\arraybackslash}p{1.50cm}>{\footnotesize\raggedright\arraybackslash}p{6.40cm}@{}}
\toprule
{\footnotesize\textbf{Name}} & {\footnotesize\textbf{Size}} & {\footnotesize\textbf{Type}} & {\footnotesize\textbf{Units}} & {\footnotesize\textbf{Description}} \\
\midrule
\endfirsthead
\toprule
{\footnotesize\textbf{Name}} & {\footnotesize\textbf{Size}} & {\footnotesize\textbf{Type}} & {\footnotesize\textbf{Units}} & {\footnotesize\textbf{Description}} \\
\midrule
\endhead
\midrule
\endfoot
\bottomrule
\endlastfoot
model & 1x1 & NonGravForces.\allowbreak{}Model & - & Configured and complete (attitude, EOP, space weather, CERES, mass) for the arc to be integrated \\
\end{longtable}

\subsubsection*{NAME-VALUE ARGUMENTS}
\begin{longtable}{@{}>{\footnotesize\raggedright\arraybackslash\ttfamily}p{3.00cm}>{\footnotesize\raggedright\arraybackslash}p{1.50cm}>{\footnotesize\raggedright\arraybackslash}p{2.20cm}>{\footnotesize\raggedright\arraybackslash}p{1.50cm}>{\footnotesize\raggedright\arraybackslash}p{6.40cm}@{}}
\toprule
{\footnotesize\textbf{Name}} & {\footnotesize\textbf{Size}} & {\footnotesize\textbf{Type}} & {\footnotesize\textbf{Units}} & {\footnotesize\textbf{Description}} \\
\midrule
\endfirsthead
\toprule
{\footnotesize\textbf{Name}} & {\footnotesize\textbf{Size}} & {\footnotesize\textbf{Type}} & {\footnotesize\textbf{Units}} & {\footnotesize\textbf{Description}} \\
\midrule
\endhead
\midrule
\endfoot
\bottomrule
\endlastfoot
Epoch & 1x1 & datetime & TAI & The run epoch, unzoned, as gjToolbox.Options carries it. Required (default NaT, refused); a context Epoch must equal it. \\
Attitude & 1x1 & string & - & "provider" (default, the Model's) or "nominal" \\
AlongVelocity & 1x1 & logical & - & Nominal attitude: +X along the velocity (default true) \\
Parameters & 1xP & string & - & Estimated parameters: any of "cdScale", "crScale"; default none \\
Native & 1x1 & logical & - & Publish the compiled entry (default true) \\
Gravity & 1x1 & object & - & Another gjToolbox.DynamicsSource added before this one; default [] \\
Margin\_\allowbreak{}s & 1x1 & double & s & Table margin beyond the stencil span (default 3600) \\
TableStep\_\allowbreak{}s & 1x1 & double & s & Table spacing (default 60) \\
ThermalOrbit & 1x1 & struct & - & thermal.model transient only, and required by it: the reference orbit (.epoch, a referenceSystem epoch struct with K increasing entries; .r and .v, 3xK in GCRS, m and m/s; e.g. the three outputs of io.readGNV1B packed into a struct) along which prepare() integrates the wall temperatures; it must cover the tables' span and should start well before it (the walls start from the steady state at its first epoch). Default struct() \\
\end{longtable}

\subsubsection*{PROPERTIES (beyond the interface's)}
\begin{lstlisting}
model, epoch, attitudeSource, alongVelocity, parameters, native, gravity, margin_s,
tableStep_s, thermalOrbit, tables (the prepared tables: .nT, .nb, .breakpoints, .handle,
.tPanel, .scale, ...), StateSize (3), ParameterRule, Vectorized (true)
\end{lstlisting}

\subsubsection*{METHODS}
\begin{lstlisting}
prepare(context)   Bind to one run: tables, the native entry, the breakpoints, the ParameterSpec
accel(t, R, V)     3xm accelerations in GCRS, from the entry (native) or from Model.accel
hasGradient()      true only when Parameters were declared
gradient(t, r, v)  Gr, Gv by central differences of this source; Gp analytic from Model.partials
release()          Free the native context (called under onCleanup by gjToolbox.Dynamics)
\end{lstlisting}

\subsubsection*{ERRORS}
\begin{description}[leftmargin=1.2em,style=nextline,itemsep=0pt]
\item[\texttt{NonGravForces:\allowbreak{}DynamicsSource:\allowbreak{}inputModel}] model is not a scalar NonGravForces.Model; its options cannot be served (raycast; a non-slerp attitude provider; transient thermal without 'ThermalOrbit' or with 'Native' false); the Model lacks what its options need (Eop, the attitude provider, SpaceWeather, Ceres, DensityScale, Exponential); 'Parameters' outside cdScale/crScale; 'Margin\_s' \textless{} 0 or 'TableStep\_s' \textless{}= 0; accel() given a non-scalar t or mismatched R, V; diagnostics() without Native
\item[\texttt{NonGravForces:\allowbreak{}DynamicsSource:\allowbreak{}unsupportedOption}] erp.backend ceresClimatology or drag.density dtm2020, which have no native evaluator; use Model.accel or series instead
\item[\texttt{NonGravForces:\allowbreak{}DynamicsSource:\allowbreak{}inputEpoch}] no 'Epoch' (or a zoned one), or a context Epoch that disagrees
\item[\texttt{NonGravForces:\allowbreak{}DynamicsSource:\allowbreak{}inputGravity}] 'Gravity' is not a gjToolbox.DynamicsSource, or (Native) publishes no NativeEntry
\item[\texttt{NonGravForces:\allowbreak{}DynamicsSource:\allowbreak{}depMissing}] the rhs MEX is not built (Native), or sun.source rs.de440 without the Model's 'Ephemeris'. Without gjToolbox the class does not load at all (MATLAB's own error)
\item[\texttt{NonGravForces:\allowbreak{}DynamicsSource:\allowbreak{}dataRange}] the attitude nodes or 'ThermalOrbit' do not cover the tables' span (span plus the margin); a space-weather file that does not cover it raises io.spaceWeather's own error
\item[\texttt{NonGravForces:\allowbreak{}rhsGateway:\allowbreak{}dataRange}] native evaluation outside its domain: t outside the tables, a UTC day without a density scale, a height too far outside Swarm's bracket, hwm14 below the ellipsoid
\item[\texttt{NonGravForces:\allowbreak{}rhsGateway:\allowbreak{}numericModel}] native evaluation failed otherwise (e.g. the gravity entry)
\item[\texttt{gjToolbox:\allowbreak{}dynamics:\allowbreak{}notPrepared}] accel() before prepare()
\end{description}

\subsection{\texttt{NonGravForces.app.forceBudget}}
\label{api:app.forceBudget}
\noindent Per-effect per-axis acceleration budget over an arc, in chunks, for one or several models.

\begin{lstlisting}
B = NonGravForces.app.forceBudget(model, epoch, r, v, q)
B = NonGravForces.app.forceBudget({fast, standard}, epoch, r, v, q, 'Labels', ["fast" "standard"])
[B, info] = NonGravForces.app.forceBudget(___, 'Chunk', 50, 'Frame', "gcrs")
\end{lstlisting}

\subsubsection*{INPUT ARGUMENTS}
\begin{longtable}{@{}>{\footnotesize\raggedright\arraybackslash\ttfamily}p{3.00cm}>{\footnotesize\raggedright\arraybackslash}p{1.50cm}>{\footnotesize\raggedright\arraybackslash}p{2.20cm}>{\footnotesize\raggedright\arraybackslash}p{1.50cm}>{\footnotesize\raggedright\arraybackslash}p{6.40cm}@{}}
\toprule
{\footnotesize\textbf{Name}} & {\footnotesize\textbf{Size}} & {\footnotesize\textbf{Type}} & {\footnotesize\textbf{Units}} & {\footnotesize\textbf{Description}} \\
\midrule
\endfirsthead
\toprule
{\footnotesize\textbf{Name}} & {\footnotesize\textbf{Size}} & {\footnotesize\textbf{Type}} & {\footnotesize\textbf{Units}} & {\footnotesize\textbf{Description}} \\
\midrule
\endhead
\midrule
\endfoot
\bottomrule
\endlastfoot
models & 1x1\textbar{}\allowbreak{}1xM & Model & - & One NonGravForces.Model, or a cell array of them. Every or cell model must accept the same epochs; they may differ in preset, options, macro model and mass \\
epoch & 1x1 & struct & - & referenceSystem epoch struct with K entries \\
r & 3xK & double & m & Position, GCRS \\
v & 3xK & double & m/\allowbreak{}s & Velocity, GCRS \\
q & 4xK & double & - & BODY-\textgreater{}GCRS quaternions, or [] to take them from each model's attitude provider \\
\end{longtable}

\subsubsection*{NAME-VALUE ARGUMENTS}
\begin{longtable}{@{}>{\footnotesize\raggedright\arraybackslash\ttfamily}p{3.00cm}>{\footnotesize\raggedright\arraybackslash}p{1.50cm}>{\footnotesize\raggedright\arraybackslash}p{2.20cm}>{\footnotesize\raggedright\arraybackslash}p{1.50cm}>{\footnotesize\raggedright\arraybackslash}p{6.40cm}@{}}
\toprule
{\footnotesize\textbf{Name}} & {\footnotesize\textbf{Size}} & {\footnotesize\textbf{Type}} & {\footnotesize\textbf{Units}} & {\footnotesize\textbf{Description}} \\
\midrule
\endfirsthead
\toprule
{\footnotesize\textbf{Name}} & {\footnotesize\textbf{Size}} & {\footnotesize\textbf{Type}} & {\footnotesize\textbf{Units}} & {\footnotesize\textbf{Description}} \\
\midrule
\endhead
\midrule
\endfoot
\bottomrule
\endlastfoot
Frame & 1x1 & string & - & "body" (default) or "gcrs" \\
Chunk & 1x1 & double & - & Epochs per evaluation, \textgreater{}= 1; default the first model's Options.series.chunk; rounded to an integer. For a The result does not depend on it beyond the rounding named under ACCURACY, transient walls included \\
Labels & 1xM & string & - & One label per model, default the model's preset name with an index appended when two share it \\
\end{longtable}

\subsubsection*{OUTPUT ARGUMENTS}
\begin{longtable}{@{}>{\footnotesize\raggedright\arraybackslash\ttfamily}p{3.00cm}>{\footnotesize\raggedright\arraybackslash}p{1.50cm}>{\footnotesize\raggedright\arraybackslash}p{2.20cm}>{\footnotesize\raggedright\arraybackslash}p{1.50cm}>{\footnotesize\raggedright\arraybackslash}p{6.40cm}@{}}
\toprule
{\footnotesize\textbf{Name}} & {\footnotesize\textbf{Size}} & {\footnotesize\textbf{Type}} & {\footnotesize\textbf{Units}} & {\footnotesize\textbf{Description}} \\
\midrule
\endfirsthead
\toprule
{\footnotesize\textbf{Name}} & {\footnotesize\textbf{Size}} & {\footnotesize\textbf{Type}} & {\footnotesize\textbf{Units}} & {\footnotesize\textbf{Description}} \\
\midrule
\endhead
\midrule
\endfoot
\bottomrule
\endlastfoot
B & (4M)x10 & table & - & Four rows per model - drag, radiation, thermal, total - with .model .effect .rms\_x .rms\_y .rms\_z .max\_x .max\_y .max\_z .rms\_norm .share (rms\_norm over the model's total rms\_norm; 1 on the total row, and 0 throughout for a null arc rather than 0/0). VariableUnits are m/s\textasciicircum{}2 on every numeric column but share \\
info & 1x1 & struct & - & .K, .chunks, .chunkSize, .frame, .labels, .provenance \\
\end{longtable}

\subsubsection*{ERRORS}
\begin{description}[leftmargin=1.2em,style=nextline,itemsep=0pt]
\item[\texttt{NonGravForces:\allowbreak{}forceBudget:\allowbreak{}inputModels}] models is not a Model or a cell array of Models
\item[\texttt{NonGravForces:\allowbreak{}forceBudget:\allowbreak{}inputShape}] r, v or q do not match the epoch count, or Labels is not 1xM
\item[\texttt{NonGravForces:\allowbreak{}forceBudget:\allowbreak{}inputRange}] Chunk \textless{} 1 or not finite
\item[\texttt{NonGravForces:\allowbreak{}Model:\allowbreak{}*}] propagated from the evaluation (missing Eop, attitude, weather)
\end{description}

\subsection{\texttt{NonGravForces.app.densityInversion}}
\label{api:app.densityInversion}
\noindent Thermosphere density from calibrated non-gravitational accelerations.

\begin{lstlisting}
T = NonGravForces.app.densityInversion(model, epoch, r, v, q, aMeasured)
[T, info] = NonGravForces.app.densityInversion(___, 'Frame', "body", 'Chunk', 200)
\end{lstlisting}

\subsubsection*{INPUT ARGUMENTS}
\begin{longtable}{@{}>{\footnotesize\raggedright\arraybackslash\ttfamily}p{3.00cm}>{\footnotesize\raggedright\arraybackslash}p{1.50cm}>{\footnotesize\raggedright\arraybackslash}p{2.20cm}>{\footnotesize\raggedright\arraybackslash}p{1.50cm}>{\footnotesize\raggedright\arraybackslash}p{6.40cm}@{}}
\toprule
{\footnotesize\textbf{Name}} & {\footnotesize\textbf{Size}} & {\footnotesize\textbf{Type}} & {\footnotesize\textbf{Units}} & {\footnotesize\textbf{Description}} \\
\midrule
\endfirsthead
\toprule
{\footnotesize\textbf{Name}} & {\footnotesize\textbf{Size}} & {\footnotesize\textbf{Type}} & {\footnotesize\textbf{Units}} & {\footnotesize\textbf{Description}} \\
\midrule
\endhead
\midrule
\endfoot
\bottomrule
\endlastfoot
model & 1x1 & Model & - & A configured NonGravForces.Model \\
epoch & 1x1 & struct & - & referenceSystem epoch struct with K entries \\
r & 3xK & double & m & Position, GCRS \\
v & 3xK & double & m/\allowbreak{}s & Velocity, GCRS \\
q & 4xK & double & - & BODY-\textgreater{}GCRS quaternions, or [] for the model's provider \\
aMeasured & 3xK & double & m/\allowbreak{}s\textasciicircum{}2 & Calibrated non-gravitational acceleration, in 'Frame' \\
\end{longtable}

\subsubsection*{NAME-VALUE ARGUMENTS}
\begin{longtable}{@{}>{\footnotesize\raggedright\arraybackslash\ttfamily}p{3.00cm}>{\footnotesize\raggedright\arraybackslash}p{1.50cm}>{\footnotesize\raggedright\arraybackslash}p{2.20cm}>{\footnotesize\raggedright\arraybackslash}p{1.50cm}>{\footnotesize\raggedright\arraybackslash}p{6.40cm}@{}}
\toprule
{\footnotesize\textbf{Name}} & {\footnotesize\textbf{Size}} & {\footnotesize\textbf{Type}} & {\footnotesize\textbf{Units}} & {\footnotesize\textbf{Description}} \\
\midrule
\endfirsthead
\toprule
{\footnotesize\textbf{Name}} & {\footnotesize\textbf{Size}} & {\footnotesize\textbf{Type}} & {\footnotesize\textbf{Units}} & {\footnotesize\textbf{Description}} \\
\midrule
\endhead
\midrule
\endfoot
\bottomrule
\endlastfoot
Frame & 1x1 & string & - & "body" (default) or "gcrs"; aMeasured is read in it \\
Chunk & 1x1 & double & - & Epochs per evaluation, \textgreater{}= 1; default the model's Options.series.chunk; rounded to an integer. Changes the result by rounding only, transient walls included (see ACCURACY) \\
MinDrag & 1x1 & double & m/\allowbreak{}s\textasciicircum{}2 & Below this modelled drag magnitude the epoch is marked unusable instead of dividing by it; default 1e-12 \\
\end{longtable}

\subsubsection*{OUTPUT ARGUMENTS}
\begin{longtable}{@{}>{\footnotesize\raggedright\arraybackslash\ttfamily}p{3.00cm}>{\footnotesize\raggedright\arraybackslash}p{1.50cm}>{\footnotesize\raggedright\arraybackslash}p{2.20cm}>{\footnotesize\raggedright\arraybackslash}p{1.50cm}>{\footnotesize\raggedright\arraybackslash}p{6.40cm}@{}}
\toprule
{\footnotesize\textbf{Name}} & {\footnotesize\textbf{Size}} & {\footnotesize\textbf{Type}} & {\footnotesize\textbf{Units}} & {\footnotesize\textbf{Description}} \\
\midrule
\endfirsthead
\toprule
{\footnotesize\textbf{Name}} & {\footnotesize\textbf{Size}} & {\footnotesize\textbf{Type}} & {\footnotesize\textbf{Units}} & {\footnotesize\textbf{Description}} \\
\midrule
\endhead
\midrule
\endfoot
\bottomrule
\endlastfoot
T & Kx7 & timetable & - & .rho\_kg\_m3 inverted density, .rho\_model\_kg\_m3 the backend's own, .scale their ratio, .dragProjected\_m\_s2 the measurement along u, .dragModel\_m\_s2 the model along u, .residualPerp\_m\_s2 the norm of what is left across u, .usable logical (false where the modelled drag is below MinDrag, where rho\_model is zero, or where the epoch was not finite). rho is NaN wherever usable is false \\
info & 1x1 & struct & - & .K .chunks .chunkSize .frame .usable (count) .minDrag .provenance \\
\end{longtable}

\subsubsection*{ERRORS}
\begin{description}[leftmargin=1.2em,style=nextline,itemsep=0pt]
\item[\texttt{NonGravForces:\allowbreak{}densityInversion:\allowbreak{}inputModel}] model is not a NonGravForces.Model
\item[\texttt{NonGravForces:\allowbreak{}densityInversion:\allowbreak{}inputShape}] r, v, q or aMeasured do not match the epoch count
\item[\texttt{NonGravForces:\allowbreak{}densityInversion:\allowbreak{}inputRange}] Chunk \textless{} 1, or MinDrag \textless{} 0, or not finite
\item[\texttt{NonGravForces:\allowbreak{}densityInversion:\allowbreak{}numericNaN}] aMeasured contains NaN or Inf
\item[\texttt{NonGravForces:\allowbreak{}Model:\allowbreak{}*}] propagated from the evaluation
\end{description}

\subsection{\texttt{NonGravForces.app.syntheticPair}}
\label{api:app.syntheticPair}
\noindent Truth generator: a 10 Hz model arc and its CRN-filtered 1 Hz observation, with every planted value.

\begin{lstlisting}
[obs, ref, truth] = NonGravForces.app.syntheticPair()
[obs, ref, truth] = NonGravForces.app.syntheticPair('Days', 0.25, 'Seed', 7, 'Scale', S, 'Quadratic', k2)
\end{lstlisting}

\subsubsection*{INPUT ARGUMENTS}
None.


\subsubsection*{NAME-VALUE ARGUMENTS}
\begin{longtable}{@{}>{\footnotesize\raggedright\arraybackslash\ttfamily}p{3.00cm}>{\footnotesize\raggedright\arraybackslash}p{1.50cm}>{\footnotesize\raggedright\arraybackslash}p{2.20cm}>{\footnotesize\raggedright\arraybackslash}p{1.50cm}>{\footnotesize\raggedright\arraybackslash}p{6.40cm}@{}}
\toprule
{\footnotesize\textbf{Name}} & {\footnotesize\textbf{Size}} & {\footnotesize\textbf{Type}} & {\footnotesize\textbf{Units}} & {\footnotesize\textbf{Description}} \\
\midrule
\endfirsthead
\toprule
{\footnotesize\textbf{Name}} & {\footnotesize\textbf{Size}} & {\footnotesize\textbf{Type}} & {\footnotesize\textbf{Units}} & {\footnotesize\textbf{Description}} \\
\midrule
\endhead
\midrule
\endfoot
\bottomrule
\endlastfoot
Days & 1x1 & double & day & Arc length, \textgreater{} 0; default 1 \\
Start & 1x1 & int64 & ns & First epoch, GPS ns past 2000-01-01 12:00:00; default 2020-01-01 00:00:18 GPS \\
Seed & 1x1 & double & - & Seed of every random draw; default 1 \\
Scale & 3x3 & double & - & True scale matrix S (a 3x1 is its diagonal); default diag([0.96 0.98 0.95]) \\
RegressorFcn & 1x1 & function\_\allowbreak{}handle & - & Measured series z(ts) (n x K, ts in s since Start) whose planted coupling raw += z * RegressorCoupling'; default none (N-11a) \\
RegressorCoupling & 3xn & double & m/\allowbreak{}s\textasciicircum{}2/\allowbreak{}unit & The true coefficients, one column per column of z; default zeros(3, 0) \\
ScaleTemperature & 3x3 & double & 1/\allowbreak{}unit & True temperature dependence of the scale (N-11b), driven by the FIRST column z1 of RegressorFcn: S(t) = Scale + ScaleTemperature * (z1(t) - Tref); default zeros \\
ScaleTemperatureReference & 1x1 & double & unit & Tref; default [] = the mean of z1 over the raw grid, which is the reference fitCalibration uses on one arc \\
ScaleDrift & 3x3 & double & 1/\allowbreak{}day & True linear drift of S, so S(t) = Scale + ScaleDrift * (t - Start) / 1 day; default zeros \\
Bias & 3x1 & double & m/\allowbreak{}s\textasciicircum{}2 & True bias at Start; default [2e-7; -1e-7; 5e-8] \\
BiasDrift & 3x1 & double & m/\allowbreak{}s\textasciicircum{}2/\allowbreak{}day & True linear bias drift; default [1e-9; 0; -5e-10] \\
Quadratic & 3x1 & double & s\textasciicircum{}2/\allowbreak{}m & True K2 per axis; default zeros \\
NoiseAsd & 3x1 & double & m/\allowbreak{}s\textasciicircum{}2/\allowbreak{}rtHz & White noise floor per axis; default [1e-10; 1e-9; 1e-10] \\
NoiseCorner & 1x1 & double & Hz & 1/f corner of the noise; default 5e-3 \\
Spikes & 1x1 & double & - & Number of planted 1 s spikes; default 3 \\
SpikeAmplitude & 1x1 & double & m/\allowbreak{}s\textasciicircum{}2 & Their amplitude; default 5e-8 \\
Gap\_\allowbreak{}s & 1x1 & double & s & Length of the one planted gap, 0 for none; default 30 \\
OrbitPeriod\_\allowbreak{}s & 1x1 & double & s & Orbital period; default 5640 (\textasciitilde{}490 km) \\
\end{longtable}

\subsubsection*{OUTPUT ARGUMENTS}
\begin{longtable}{@{}>{\footnotesize\raggedright\arraybackslash\ttfamily}p{3.00cm}>{\footnotesize\raggedright\arraybackslash}p{1.50cm}>{\footnotesize\raggedright\arraybackslash}p{2.20cm}>{\footnotesize\raggedright\arraybackslash}p{1.50cm}>{\footnotesize\raggedright\arraybackslash}p{6.40cm}@{}}
\toprule
{\footnotesize\textbf{Name}} & {\footnotesize\textbf{Size}} & {\footnotesize\textbf{Type}} & {\footnotesize\textbf{Units}} & {\footnotesize\textbf{Description}} \\
\midrule
\endfirsthead
\toprule
{\footnotesize\textbf{Name}} & {\footnotesize\textbf{Size}} & {\footnotesize\textbf{Type}} & {\footnotesize\textbf{Units}} & {\footnotesize\textbf{Description}} \\
\midrule
\endhead
\midrule
\endfoot
\bottomrule
\endlastfoot
obs & 1x1 & timeSeriesPreparation.\allowbreak{}Arc & m/\allowbreak{}s\textasciicircum{}2 & CRN-filtered observation at 1 Hz, product ACT1B, columns lin\_accl\_x/y/z; NaN (and the layer crnWindow) where the filter window met an edge or the gap \\
ref & 1x1 & timeSeriesPreparation.\allowbreak{}Arc & m/\allowbreak{}s\textasciicircum{}2 & The model at 10 Hz, same column names \\
truth & 1x1 & struct & - & .scale (at Start) .scaleDrift\_per\_day .regressors (Arc of the series on the 10 Hz grid, columns reg1..regK, or []) .regressorCoupling .scaleTemperature (3x3) .scaleTemperatureReference .bias0 .biasDrift\_per\_day .quadratic .bias (3xK1 at obs.t) .noise (.asd .corner .seeds) .spikes (table: t int64, axis, amplitude) .gap (int64 [first last] on the 10 Hz grid, or empty) .kernel (crnKernel struct) .orbitPeriod\_s .seed .provenance \\
\end{longtable}

\subsubsection*{ERRORS}
\begin{description}[leftmargin=1.2em,style=nextline,itemsep=0pt]
\item[\texttt{NonGravForces:\allowbreak{}syntheticPair:\allowbreak{}depMissing}] timeSeriesPreparation is not on the path
\item[\texttt{NonGravForces:\allowbreak{}syntheticPair:\allowbreak{}inputShape}] Scale, ScaleDrift, Bias, BiasDrift, Quadratic or NoiseAsd of the wrong size, or not finite; ScaleTemperature not a finite 3x3, or nonzero without RegressorFcn; RegressorFcn not returning one row per raw sample, or RegressorCoupling not 3x(its columns)
\item[\texttt{NonGravForces:\allowbreak{}syntheticPair:\allowbreak{}inputRange}] the arc is too short for the CRN window plus the spikes and the gap, or the spikes find no room clear of each other
\item[\texttt{timeSeriesPreparation:\allowbreak{}*}] from crnKernel, crnApply, synthTimeSeries, Arc
\end{description}

\subsection{\texttt{NonGravForces.app.Calibration}}
\label{api:app.Calibration}
\noindent An accelerometer calibration: scale, time-varying bias, uncertainty, gates and provenance, applicable to an arc.

\subsubsection*{CONSTRUCTOR ARGUMENTS}
\begin{longtable}{@{}>{\footnotesize\raggedright\arraybackslash\ttfamily}p{3.00cm}>{\footnotesize\raggedright\arraybackslash}p{1.50cm}>{\footnotesize\raggedright\arraybackslash}p{2.20cm}>{\footnotesize\raggedright\arraybackslash}p{1.50cm}>{\footnotesize\raggedright\arraybackslash}p{6.40cm}@{}}
\toprule
{\footnotesize\textbf{Name}} & {\footnotesize\textbf{Size}} & {\footnotesize\textbf{Type}} & {\footnotesize\textbf{Units}} & {\footnotesize\textbf{Description}} \\
\midrule
\endfirsthead
\toprule
{\footnotesize\textbf{Name}} & {\footnotesize\textbf{Size}} & {\footnotesize\textbf{Type}} & {\footnotesize\textbf{Units}} & {\footnotesize\textbf{Description}} \\
\midrule
\endhead
\midrule
\endfoot
\bottomrule
\endlastfoot
fit & 1x1 & struct & - & .scale 3x3; .bias struct (.basis, .t0 int64, .knots int64 1xQ, optional .jumps int64 1xJ, .coef 3xP with P including the J steps, optional .cov 1x3 cell of PxP); optional .covScale 9x9 (cov of scale(:)), .quadratic 3x1, .uncertainty struct, .gates struct, .convention "observation"\textbar{}"correction", .provenance struct; optional .scaleCurve struct (.basis, .t0, .knots, .tRef int64, .coef 9xQ, .cov 1x9 cell of QxQ, [] where not estimated) for a time-varying scale; optional .regressors struct (.names 1xK, .units 1xK, .mean 1xK, .coef 3xK, .cov 1x3 cell of KxK); optional .scaleTemperature struct (.name, .units, .reference, .coef 3x3, .cov 9x9) \\
\end{longtable}

\subsubsection*{NAME-VALUE ARGUMENTS}
None.


\subsubsection*{PROPERTIES (read-only)}
\begin{lstlisting}
scale        3x3  double  -      S; for a time-varying scale, S(scaleCurve.tRef)
scaleCurve   1x1  struct  -      .basis ("constant" when the scale is not a curve) .t0
                                 .knots .tRef .coef (9xQ) .cov (1x9 cell)
covScale     9x9  double  -      covariance of scale(:), NaN where not estimated
bias         1x1  struct  m/s^2  .basis .t0 .knots .jumps .coef (3xP) .cov (1x3 cell)
quadratic    3x1  double  s^2/m  K2 per axis; NaN = not estimated
regressors   1x1  struct  -      .names .units .mean (1xK) .coef (3xK, m/s^2 per unit) .cov
                                 (1x3 cell, KxK); K = 0 without measured regressors
scaleTemperature 1x1 struct -    .name ("" without) .units .reference T0 .coef (3x3, dS/dT, 1 per
                                 unit, 0 where not estimated) .cov (9x9, cov of coef(:), NaN
                                 where not estimated)
uncertainty  1x1  struct  -      as the estimator reports it (sigma, inflation, EIV bound, ...)
gates        1x1  struct  -      observability and preflight verdicts
convention   1x1  string  -      "observation" (y = b + S a) or "correction" (a = b + S y)
provenance   1x1  struct  -      NonGravForces.provenance of the fit, plus the estimator's
\end{lstlisting}

\subsubsection*{METHODS}
\begin{lstlisting}
b = biasAt(obj, t)          t Kx1 int64 GPS ns; b 3xK in m/s^2
S = scaleAt(obj, t)         3x3xK; the constant S repeated when not time-varying
s = scaleSigmaAt(obj, t)    3x3xK pointwise sigma of S(t); NaN where not estimated
tf = isTimeVarying(obj)     scaleCurve.basis is not "constant"
tf = isTemperatureDependent(obj)   scaleTemperature.name is not ""
S = scaleAtTemperature(obj, T)   3x3xK: scale + coef (T - reference), T 1xK in the series' unit
r = regressorTermAt(obj, t, R)   3xK: sum_k c_k (z_k(t) - m_k) from the measured series R
T = temperatureAt(obj, t, R)     1xK: the series the scale depends on, at int64 epochs t
Z = Calibration.interpolateSeries(R, col, t)   columns col of Arc R, linearly at int64 t (NaN
                            outside); the one rule of the estimator and of apply
d = drift(obj, alpha)       .statistic .dof .pValue (3x3, NaN where not tested)
                            .significant (3x3 logical) .alpha; alpha default 0.05
tf = isDriftSignificant(obj, alpha)   drift(obj, alpha).significant
out = apply(obj, arc, 'Regressors', R)   R needed when regressors.names is not empty;
out = apply(obj, arc)       arc: timeSeriesPreparation.Arc with 3 columns in m/s2; out: the
                            calibrated arc (derive: same epochs, provenance appended)
inv = invert(obj)           the other convention, first-order covariance
T = summaryTable(obj)       one row per scale entry and bias coefficient, plus one per
                            scale-curve coefficient when time-varying and one per estimated
                            temperature coefficient (dS_ij_dT)
save(obj, file)             versioned MAT file (schemaVersion "4"; "1" to "3" files load)
cal = Calibration.load(file)
B = Calibration.basis(bias, t)   KxP basis matrix at int64 epochs
\end{lstlisting}

\subsubsection*{OUTPUT ARGUMENTS}
\begin{longtable}{@{}>{\footnotesize\raggedright\arraybackslash\ttfamily}p{3.00cm}>{\footnotesize\raggedright\arraybackslash}p{1.50cm}>{\footnotesize\raggedright\arraybackslash}p{10.10cm}@{}}
\toprule
{\footnotesize\textbf{Name}} & {\footnotesize\textbf{Size}} & {\footnotesize\textbf{Description}} \\
\midrule
\endfirsthead
\toprule
{\footnotesize\textbf{Name}} & {\footnotesize\textbf{Size}} & {\footnotesize\textbf{Description}} \\
\midrule
\endhead
\midrule
\endfoot
\bottomrule
\endlastfoot
cal & 1x1 & NonGravForces.app.Calibration \\
\end{longtable}

\subsubsection*{ERRORS}
\begin{description}[leftmargin=1.2em,style=nextline,itemsep=0pt]
\item[\texttt{NonGravForces:\allowbreak{}Calibration:\allowbreak{}inputFit}] the fit struct lacks a field, or one has the wrong size
\item[\texttt{NonGravForces:\allowbreak{}Calibration:\allowbreak{}inputBasis}] unknown basis, knots or jumps not increasing, too few knots, or epochs that are not int64 GPS ns
\item[\texttt{NonGravForces:\allowbreak{}Calibration:\allowbreak{}inputArc}] apply: not a 3-column arc in m/s2
\item[\texttt{NonGravForces:\allowbreak{}Calibration:\allowbreak{}unsupportedQuadratic}] apply with a non-zero K2
\item[\texttt{NonGravForces:\allowbreak{}Calibration:\allowbreak{}dataSchema}] load: not a calibration file, or no schemaVersion
\item[\texttt{NonGravForces:\allowbreak{}Calibration:\allowbreak{}numericSingular}] invert/apply: S is singular to working precision (at some epoch, for a time-varying scale)
\item[\texttt{NonGravForces:\allowbreak{}Calibration:\allowbreak{}unsupportedTimeVarying}] invert of a time-varying scale
\item[\texttt{NonGravForces:\allowbreak{}Calibration:\allowbreak{}inputRegressors}] apply without the regressor series the calibration carries (or the temperature its scale depends on), or an Arc that lacks one of their columns
\item[\texttt{NonGravForces:\allowbreak{}Calibration:\allowbreak{}unsupportedTemperatureScale}] invert of a temperature-dependent scale, or one combined with a time-varying scale
\end{description}

\subsubsection*{WARNINGS}
\begin{description}[leftmargin=1.2em,style=nextline,itemsep=0pt]
\item[\texttt{NonGravForces:\allowbreak{}Calibration:\allowbreak{}newerSchema}] load: written by a newer schema; read field by field
\item[\texttt{NonGravForces:\allowbreak{}Calibration:\allowbreak{}formalDrift}] drift: the covariance is the formal white-noise one (NBoot 0, Noise "white"), which a CRN-filtered residual violates; the test is then far too liberal
\end{description}

\subsection{\texttt{NonGravForces.app.fitCalibration}}
\label{api:app.fitCalibration}
\noindent Estimate an accelerometer calibration from a filtered observation arc and a raw-rate model arc, through lsqvce.fit.

\begin{lstlisting}
cal = NonGravForces.app.fitCalibration(obs, ref)
[cal, info] = NonGravForces.app.fitCalibration(obs, ref, 'Scale', "full", 'BiasBasis', "pspline", 'NBoot', 0)
cal = NonGravForces.app.fitCalibration(obs, ref, 'ScaleBasis', "linear")   % S(t), N-03a
cal = NonGravForces.app.fitCalibration(obsDays, refDays)                  % 1xN: one scale, N biases (N-09)
\end{lstlisting}

\subsubsection*{INPUT ARGUMENTS}
\begin{longtable}{@{}>{\footnotesize\raggedright\arraybackslash\ttfamily}p{3.00cm}>{\footnotesize\raggedright\arraybackslash}p{1.50cm}>{\footnotesize\raggedright\arraybackslash}p{2.20cm}>{\footnotesize\raggedright\arraybackslash}p{1.50cm}>{\footnotesize\raggedright\arraybackslash}p{6.40cm}@{}}
\toprule
{\footnotesize\textbf{Name}} & {\footnotesize\textbf{Size}} & {\footnotesize\textbf{Type}} & {\footnotesize\textbf{Units}} & {\footnotesize\textbf{Description}} \\
\midrule
\endfirsthead
\toprule
{\footnotesize\textbf{Name}} & {\footnotesize\textbf{Size}} & {\footnotesize\textbf{Type}} & {\footnotesize\textbf{Units}} & {\footnotesize\textbf{Description}} \\
\midrule
\endhead
\midrule
\endfoot
\bottomrule
\endlastfoot
obs & 1xN & timeSeriesPreparation.\allowbreak{}Arc & m/\allowbreak{}s\textasciicircum{}2 & Observation, three axes, units m/s2, already through the L1B chain (e.g. ACT1B); its epochs lie on ref's grid \\
ref & 1xN & timeSeriesPreparation.\allowbreak{}Arc & m/\allowbreak{}s\textasciicircum{}2 & Model accelerations, three axes, units m/s2, on the kernel's raw grid around obs (10 Hz for "ACC"); with 'Kernel' [] on obs's epochs \\
\end{longtable}

\subsubsection*{NAME-VALUE ARGUMENTS}
\begin{longtable}{@{}>{\footnotesize\raggedright\arraybackslash\ttfamily}p{3.00cm}>{\footnotesize\raggedright\arraybackslash}p{1.50cm}>{\footnotesize\raggedright\arraybackslash}p{2.20cm}>{\footnotesize\raggedright\arraybackslash}p{1.50cm}>{\footnotesize\raggedright\arraybackslash}p{6.40cm}@{}}
\toprule
{\footnotesize\textbf{Name}} & {\footnotesize\textbf{Size}} & {\footnotesize\textbf{Type}} & {\footnotesize\textbf{Units}} & {\footnotesize\textbf{Description}} \\
\midrule
\endfirsthead
\toprule
{\footnotesize\textbf{Name}} & {\footnotesize\textbf{Size}} & {\footnotesize\textbf{Type}} & {\footnotesize\textbf{Units}} & {\footnotesize\textbf{Description}} \\
\midrule
\endhead
\midrule
\endfoot
\bottomrule
\endlastfoot
Kernel & 1x1 & string & - & "crn" (default): crnKernel("ACC"); a crnKernel struct (fields F, fs, Nf, Nh, Tf, preset); or [] for an unfiltered observation \\
BiasBasis & 1x1 & string & - & "constant" \textbar{} "linear" (default) \textbar{} "hat" \textbar{} "pspline" \\
BiasKnotSpacing\_\allowbreak{}s & 1x1 & double & s & Knot spacing of "hat" and "pspline"; default 10800. Below the kernel span Tf it warns \\
Scale & 3x3 & string & - & "diag" (default) \textbar{} "full" \textbar{} a 3x3 logical mask \\
ScaleBasis & 1x1 & string & - & "constant" (default) \textbar{} "linear" \textbar{} "hat" \textbar{} "pspline": the time dependence of every estimated entry \\
ScaleKnotSpacing\_\allowbreak{}s & 1x1 & double & s & Knot spacing of a "hat" or "pspline" scale; default 43200, at least the bias's when both are knotted \\
ScaleMode & 1x1 & string & - & With several arcs: "common" (default) or "perArc" \\
Regressors & 1xN & Arc & - & Measured series, one Arc of K columns (one per obs arc, or one covering all), interpolated linearly but never across a gap in a series (Calibration.interpolateSeries); default none \\
ScaleTemperature & 1x1 & string & - & Name of one 'Regressors' column the scale depends on (N-11b): S(T) = S + S1 (T - T0); default "" (none) \\
Span & 1x2 & int64 & ns & [first last] epoch the bias and scale bases must cover; default [] = the first and last usable observation epoch. Cross-validation passes the WHOLE arc's span, so that a B-spline is defined over a held-out block (N-03b) \\
FixCrossTrack & 1x1 & logical & - & Hold S\_yy at its a-priori value; default true \\
ScaleApriori & 3x3 & double & - & Values of the entries not estimated; default eye(3) \\
Quadratic & 1x1 & logical & - & Estimate K2 per axis; default false \\
Noise & 1x1 & string & - & "white" (default) \textbar{} "estimate" \textbar{} an lsqvce covariance spec struct \\
NBoot & 1x1 & double & - & Bootstrap replicates, 0 for none; default 200 \\
Seed & 1x1 & double & - & Bootstrap seed; default 0 \\
VifMax & 1x1 & double & - & VIF gate for a row with several model columns; default 10 \\
MinWeightFraction & 1x1 & double & - & Minimum-data gate of observability; default 0.5 \\
Jumps & 1xJ & int64 & ns & Epochs (GPS ns) of steps in the bias, one Heaviside column each, filtered like the basis; default none \\
Robust & 1x1 & string & - & lsqvce M-estimation: "none" (default) \textbar{} "huber" \textbar{} "igg3" \textbar{} "bisquare"; a single arc's bootstrap refits with it, the day bootstrap keeps its final weights \\
\end{longtable}

\subsubsection*{OUTPUT ARGUMENTS}
\begin{longtable}{@{}>{\footnotesize\raggedright\arraybackslash\ttfamily}p{3.00cm}>{\footnotesize\raggedright\arraybackslash}p{1.50cm}>{\footnotesize\raggedright\arraybackslash}p{2.20cm}>{\footnotesize\raggedright\arraybackslash}p{1.50cm}>{\footnotesize\raggedright\arraybackslash}p{6.40cm}@{}}
\toprule
{\footnotesize\textbf{Name}} & {\footnotesize\textbf{Size}} & {\footnotesize\textbf{Type}} & {\footnotesize\textbf{Units}} & {\footnotesize\textbf{Description}} \\
\midrule
\endfirsthead
\toprule
{\footnotesize\textbf{Name}} & {\footnotesize\textbf{Size}} & {\footnotesize\textbf{Type}} & {\footnotesize\textbf{Units}} & {\footnotesize\textbf{Description}} \\
\midrule
\endhead
\midrule
\endfoot
\bottomrule
\endlastfoot
cal & 1xN & NonGravForces.\allowbreak{}app.\allowbreak{}Calibration & - & One per arc. Convention "observation"; scaleCurve when 'ScaleBasis' is not "constant"; covScale and bias cov from the bootstrap when NBoot \textgreater{} 0, else formal; .uncertainty (.sigmaScaleFormal, .sigmaScaleBootstrap, .inflation 3x3, .residualRms 1x3 m/s\textasciicircum{}2, .eiv, .noise, .nBoot, .stopReason); .gates (.vif, .cond, .maxAbsCorr, .weightFraction, .vifMax, .verdict) \\
info & 1x1 & struct & - & .residual (Arc at obs epochs, NaN where not fitted) .names (1x3 cell of column names) .rows (Nx3 logical) .Qxx (1x3 cell, the unit-weight cofactor of each axis, in the unscaled parameters) .results (1x3 cell of lsqvce.Results) .observability (1x3 cell) .prediction (Arc at obs epochs: the fitted observation, fixed scale entries included, at EVERY row whose design is finite, fitted or not; NaN elsewhere). With several arcs in common mode: .mode "common", .residual (1xN cell of Arc), .rows (1xN cell), .names, .Qxx and .results per axis; in perArc mode: .mode, .arcs (1xN cell of single-arc infos) \\
\end{longtable}

\subsubsection*{ERRORS}
\begin{description}[leftmargin=1.2em,style=nextline,itemsep=0pt]
\item[\texttt{NonGravForces:\allowbreak{}fitCalibration:\allowbreak{}depMissing}] lsqvce or timeSeriesPreparation is not on the path
\item[\texttt{NonGravForces:\allowbreak{}fitCalibration:\allowbreak{}inputArc}] obs or ref not a three-axis Arc in m/s2; one obs arc with several refs; several obs arcs with neither one ref nor one per arc, or not in time order, or overlapping
\item[\texttt{NonGravForces:\allowbreak{}fitCalibration:\allowbreak{}inputScale}] a mask that is neither "diag", "full" nor 3x3, or a non-finite 'ScaleApriori'
\item[\texttt{NonGravForces:\allowbreak{}fitCalibration:\allowbreak{}inputGrid}] 'Kernel' [] and an obs epoch ref does not carry
\item[\texttt{NonGravForces:\allowbreak{}fitCalibration:\allowbreak{}inputKernel}] a kernel that is neither "crn", a struct nor []
\item[\texttt{NonGravForces:\allowbreak{}fitCalibration:\allowbreak{}inputJumps}] 'Jumps' not int64
\item[\texttt{NonGravForces:\allowbreak{}fitCalibration:\allowbreak{}inputScaleBasis}] a knotted scale basis finer than the knotted bias basis
\item[\texttt{NonGravForces:\allowbreak{}fitCalibration:\allowbreak{}inputSpan}] 'Span' not an increasing int64 pair
\item[\texttt{NonGravForces:\allowbreak{}fitCalibration:\allowbreak{}inputRegressors}] 'Regressors' not an Arc (or one per arc), or a series with no finite value or no variation on the arc
\item[\texttt{NonGravForces:\allowbreak{}fitCalibration:\allowbreak{}unsupportedMultiArc}] common mode with a knotted bias, 'Jumps' or a time-varying scale
\item[\texttt{NonGravForces:\allowbreak{}fitCalibration:\allowbreak{}inputScaleTemperature}] 'ScaleTemperature' names no column of 'Regressors'
\item[\texttt{NonGravForces:\allowbreak{}fitCalibration:\allowbreak{}unsupportedScaleTemperature}] 'ScaleTemperature' beside a time-varying scale
\item[\texttt{NonGravForces:\allowbreak{}fitCalibration:\allowbreak{}numericCollinear}] the VIF gate refused a row with several model columns
\item[\texttt{NonGravForces:\allowbreak{}Calibration:\allowbreak{}inputBasis}] knots the basis refuses
\item[\texttt{timeSeriesPreparation:\allowbreak{}*}] crnApply, observability, bootstrap
\item[\texttt{lsqvce:\allowbreak{}*}] the solve
\end{description}

\subsubsection*{WARNINGS}
\begin{description}[leftmargin=1.2em,style=nextline,itemsep=0pt]
\item[\texttt{NonGravForces:\allowbreak{}fitCalibration:\allowbreak{}knotsBelowFilterSpan}] knot spacing below the kernel span Tf: the filtered basis columns are nearly dependent
\item[\texttt{NonGravForces:\allowbreak{}fitCalibration:\allowbreak{}fewArcs}] a bootstrap over fewer than 10 days
\item[\texttt{NonGravForces:\allowbreak{}fitCalibration:\allowbreak{}jumpOutsideFit}] a jump whose step is constant over the fitted rows; it is dropped from the bias
\end{description}

\subsection{\texttt{NonGravForces.app.compareCalibrations}}
\label{api:app.compareCalibrations}
\noindent Rank candidate calibration models by blocked cross-validation, per axis.

\begin{lstlisting}
T = NonGravForces.app.compareCalibrations(obs, ref, specs)
[T, info] = NonGravForces.app.compareCalibrations(obs, ref, specs, 'Names', names, 'Folds', 5, 'Fit', {'Scale', "full"})
\end{lstlisting}

\subsubsection*{INPUT ARGUMENTS}
\begin{longtable}{@{}>{\footnotesize\raggedright\arraybackslash\ttfamily}p{3.00cm}>{\footnotesize\raggedright\arraybackslash}p{1.50cm}>{\footnotesize\raggedright\arraybackslash}p{2.20cm}>{\footnotesize\raggedright\arraybackslash}p{1.50cm}>{\footnotesize\raggedright\arraybackslash}p{6.40cm}@{}}
\toprule
{\footnotesize\textbf{Name}} & {\footnotesize\textbf{Size}} & {\footnotesize\textbf{Type}} & {\footnotesize\textbf{Units}} & {\footnotesize\textbf{Description}} \\
\midrule
\endfirsthead
\toprule
{\footnotesize\textbf{Name}} & {\footnotesize\textbf{Size}} & {\footnotesize\textbf{Type}} & {\footnotesize\textbf{Units}} & {\footnotesize\textbf{Description}} \\
\midrule
\endhead
\midrule
\endfoot
\bottomrule
\endlastfoot
obs & 1x1 & timeSeriesPreparation.\allowbreak{}Arc & m/\allowbreak{}s\textasciicircum{}2 & Observation, as for app.fitCalibration \\
ref & 1x1 & timeSeriesPreparation.\allowbreak{}Arc & m/\allowbreak{}s\textasciicircum{}2 & Model on the kernel's raw grid, as for app.fitCalibration \\
specs & 1xC & cell & - & Candidate models: each a cell of app.fitCalibration name-value pairs (\{\} is the default model); C \textgreater{}= 1 \\
\end{longtable}

\subsubsection*{NAME-VALUE ARGUMENTS}
\begin{longtable}{@{}>{\footnotesize\raggedright\arraybackslash\ttfamily}p{3.00cm}>{\footnotesize\raggedright\arraybackslash}p{1.50cm}>{\footnotesize\raggedright\arraybackslash}p{2.20cm}>{\footnotesize\raggedright\arraybackslash}p{1.50cm}>{\footnotesize\raggedright\arraybackslash}p{6.40cm}@{}}
\toprule
{\footnotesize\textbf{Name}} & {\footnotesize\textbf{Size}} & {\footnotesize\textbf{Type}} & {\footnotesize\textbf{Units}} & {\footnotesize\textbf{Description}} \\
\midrule
\endfirsthead
\toprule
{\footnotesize\textbf{Name}} & {\footnotesize\textbf{Size}} & {\footnotesize\textbf{Type}} & {\footnotesize\textbf{Units}} & {\footnotesize\textbf{Description}} \\
\midrule
\endhead
\midrule
\endfoot
\bottomrule
\endlastfoot
Names & 1xC & string & - & A name per candidate; default "spec1" ... "specC" \\
Fit & 1xN & cell & - & Name-value pairs common to every candidate (a spec's own pair wins); 'Kernel' belongs here; default \{\} \\
Folds & 1x1 & double & - & Contiguous blocks, integer \textgreater{}= 2; default 5 \\
Gap & 1x1 & double & rows & Buffer on each side of a test block; default [] = the larger of crossval's automatic block length and the kernel span in rows \\
\end{longtable}

\subsubsection*{OUTPUT ARGUMENTS}
\begin{longtable}{@{}>{\footnotesize\raggedright\arraybackslash\ttfamily}p{3.00cm}>{\footnotesize\raggedright\arraybackslash}p{1.50cm}>{\footnotesize\raggedright\arraybackslash}p{2.20cm}>{\footnotesize\raggedright\arraybackslash}p{1.50cm}>{\footnotesize\raggedright\arraybackslash}p{6.40cm}@{}}
\toprule
{\footnotesize\textbf{Name}} & {\footnotesize\textbf{Size}} & {\footnotesize\textbf{Type}} & {\footnotesize\textbf{Units}} & {\footnotesize\textbf{Description}} \\
\midrule
\endfirsthead
\toprule
{\footnotesize\textbf{Name}} & {\footnotesize\textbf{Size}} & {\footnotesize\textbf{Type}} & {\footnotesize\textbf{Units}} & {\footnotesize\textbf{Description}} \\
\midrule
\endhead
\midrule
\endfoot
\bottomrule
\endlastfoot
T & Cx8 & table & - & One row per candidate: name (string), cvMse (1x3, m\textasciicircum{}2/s\textasciicircum{}4, mean over the folds of the test-row mean squared prediction error), cvMseSe (1x3, its standard error over the folds), rank (1x3, 1 = lowest cvMse on that axis), nParams (1x3), nRows (1x3, rows of the full-arc fit), aic and bic (1x3, from the full-arc fit; IC\_unreliable). T.Properties.UserData: .best (1x3 string, the rank-1 name per axis) .folds .gap .note \\
info & 1x1 & struct & - & .foldLoss (1xC cell of Folds x 3) .testRows (Folds x 2) .gap .calibrations (1xC cell of the full-arc app.Calibration) \\
\end{longtable}

\subsubsection*{ERRORS}
\begin{description}[leftmargin=1.2em,style=nextline,itemsep=0pt]
\item[\texttt{NonGravForces:\allowbreak{}compareCalibrations:\allowbreak{}depMissing}] timeSeriesPreparation is not on the path
\item[\texttt{NonGravForces:\allowbreak{}compareCalibrations:\allowbreak{}inputFolds}] a candidate's knot spacing does not exceed the held-out window, so some fold cannot fit it
\item[\texttt{NonGravForces:\allowbreak{}compareCalibrations:\allowbreak{}inputSpecs}] specs is not a non-empty cell of name-value cells, a spec sets 'Kernel', 'NBoot' or 'Span', or 'Names' does not have one name per spec
\item[\texttt{NonGravForces:\allowbreak{}fitCalibration:\allowbreak{}*}] a candidate the estimator refuses
\item[\texttt{timeSeriesPreparation:\allowbreak{}crossval:\allowbreak{}*}] folds that leave no training or test row
\item[\texttt{timeSeriesPreparation:\allowbreak{}blockLength:\allowbreak{}tooShort}] the automatic gap on too short an arc
\end{description}

\subsection{\texttt{NonGravForces.app.perRevolution}}
\label{api:app.perRevolution}
\noindent Accelerometer bias and diagonal scale fitted per orbital revolution, gated per window.

\begin{lstlisting}
T = NonGravForces.app.perRevolution(obs, ref)
[T, info] = NonGravForces.app.perRevolution(obs, ref, 'Period_s', 5640, 'NBoot', 50, 'Effects', effects)
\end{lstlisting}

\subsubsection*{INPUT ARGUMENTS}
\begin{longtable}{@{}>{\footnotesize\raggedright\arraybackslash\ttfamily}p{3.00cm}>{\footnotesize\raggedright\arraybackslash}p{1.50cm}>{\footnotesize\raggedright\arraybackslash}p{2.20cm}>{\footnotesize\raggedright\arraybackslash}p{1.50cm}>{\footnotesize\raggedright\arraybackslash}p{6.40cm}@{}}
\toprule
{\footnotesize\textbf{Name}} & {\footnotesize\textbf{Size}} & {\footnotesize\textbf{Type}} & {\footnotesize\textbf{Units}} & {\footnotesize\textbf{Description}} \\
\midrule
\endfirsthead
\toprule
{\footnotesize\textbf{Name}} & {\footnotesize\textbf{Size}} & {\footnotesize\textbf{Type}} & {\footnotesize\textbf{Units}} & {\footnotesize\textbf{Description}} \\
\midrule
\endhead
\midrule
\endfoot
\bottomrule
\endlastfoot
obs & 1x1 & timeSeriesPreparation.\allowbreak{}Arc & m/\allowbreak{}s\textasciicircum{}2 & Observation, as for app.fitCalibration \\
ref & 1x1 & timeSeriesPreparation.\allowbreak{}Arc & m/\allowbreak{}s\textasciicircum{}2 & Model on the kernel's raw grid, as for app.fitCalibration \\
\end{longtable}

\subsubsection*{NAME-VALUE ARGUMENTS}
\begin{longtable}{@{}>{\footnotesize\raggedright\arraybackslash\ttfamily}p{3.00cm}>{\footnotesize\raggedright\arraybackslash}p{1.50cm}>{\footnotesize\raggedright\arraybackslash}p{2.20cm}>{\footnotesize\raggedright\arraybackslash}p{1.50cm}>{\footnotesize\raggedright\arraybackslash}p{6.40cm}@{}}
\toprule
{\footnotesize\textbf{Name}} & {\footnotesize\textbf{Size}} & {\footnotesize\textbf{Type}} & {\footnotesize\textbf{Units}} & {\footnotesize\textbf{Description}} \\
\midrule
\endfirsthead
\toprule
{\footnotesize\textbf{Name}} & {\footnotesize\textbf{Size}} & {\footnotesize\textbf{Type}} & {\footnotesize\textbf{Units}} & {\footnotesize\textbf{Description}} \\
\midrule
\endhead
\midrule
\endfoot
\bottomrule
\endlastfoot
Period\_\allowbreak{}s & 1x1 & double & s & Window length, the orbital period; default 5640 (\textasciitilde{}490 km) \\
MinWindowFraction & 1x1 & double & - & Usable share of a window's rows below which it is gated; default 0.8 \\
NBoot & 1x1 & double & - & Bootstrap replicates per window, 0 for formal sigmas; default 50 \\
Seed & 1x1 & double & - & Bootstrap seed; window w uses Seed + w; default 0 \\
Fit & 1xN & cell & - & Further app.fitCalibration name-value pairs (Kernel, FixCrossTrack, Noise, Robust, ...); 'BiasBasis', 'ScaleBasis', 'Span', 'NBoot' and 'Seed' are this unit's \\
Effects & 1xE & struct & m/\allowbreak{}s\textasciicircum{}2 & .name (string) .arc (timeSeriesPreparation.Arc, 3 axes, any epochs); default none \\
\end{longtable}

\subsubsection*{OUTPUT ARGUMENTS}
\begin{longtable}{@{}>{\footnotesize\raggedright\arraybackslash\ttfamily}p{3.00cm}>{\footnotesize\raggedright\arraybackslash}p{1.50cm}>{\footnotesize\raggedright\arraybackslash}p{2.20cm}>{\footnotesize\raggedright\arraybackslash}p{1.50cm}>{\footnotesize\raggedright\arraybackslash}p{6.40cm}@{}}
\toprule
{\footnotesize\textbf{Name}} & {\footnotesize\textbf{Size}} & {\footnotesize\textbf{Type}} & {\footnotesize\textbf{Units}} & {\footnotesize\textbf{Description}} \\
\midrule
\endfirsthead
\toprule
{\footnotesize\textbf{Name}} & {\footnotesize\textbf{Size}} & {\footnotesize\textbf{Type}} & {\footnotesize\textbf{Units}} & {\footnotesize\textbf{Description}} \\
\midrule
\endhead
\midrule
\endfoot
\bottomrule
\endlastfoot
T & RxN & table & - & One row per window: window, t0, t1, tMid (int64 GPS ns), usable (share of the window's rows), scale, sigmaScale (Rx3, NaN where held or gated), bias, sigmaBias (Rx3, m/s\textasciicircum{}2), gate (Rx3 string: "pass", "warn" -- fitted with a VIF above VifMax on one model column -- or the reason it is NaN), and amp\_\textless{}name\textgreater{} (Rx3, m/s\textasciicircum{}2) per effect. UserData: .period\_s .nBoot .sigma ("bootstrap" or "formal") \\
info & 1x1 & struct & - & .calibrations (1xR cell of app.Calibration, [] where gated) .messages (Rx1 string, the gate's message where gated) \\
\end{longtable}

\subsubsection*{ERRORS}
\begin{description}[leftmargin=1.2em,style=nextline,itemsep=0pt]
\item[\texttt{NonGravForces:\allowbreak{}perRevolution:\allowbreak{}inputArc}] obs is not a three-column timeSeriesPreparation.Arc
\item[\texttt{NonGravForces:\allowbreak{}perRevolution:\allowbreak{}inputFit}] 'Fit' sets a pair this unit owns, or is not name-value
\item[\texttt{NonGravForces:\allowbreak{}perRevolution:\allowbreak{}inputEffects}] an effect without .name and a 3-axis .arc
\item[\texttt{NonGravForces:\allowbreak{}perRevolution:\allowbreak{}dataEmpty}] no usable observation epoch
\item[\texttt{NonGravForces:\allowbreak{}fitCalibration:\allowbreak{}*}] an estimator error that is not a gate
\item[\texttt{timeSeriesPreparation:\allowbreak{}*, lsqvce:\allowbreak{}*}] as app.fitCalibration
\end{description}

\subsection{\texttt{NonGravForces.app.attributeRevolutions}}
\label{api:app.attributeRevolutions}
\noindent Attribute per-revolution scale variations to drivers: density, geometry or sensor.

\begin{lstlisting}
A = NonGravForces.app.attributeRevolutions(T, 'Orbit', orb, 'SpaceWeather', sw)
[A, info] = NonGravForces.app.attributeRevolutions(T, 'Drivers', D, 'Families', struct('myDriver', "geometry"))
\end{lstlisting}

\subsubsection*{INPUT ARGUMENTS}
\begin{longtable}{@{}>{\footnotesize\raggedright\arraybackslash\ttfamily}p{3.00cm}>{\footnotesize\raggedright\arraybackslash}p{1.50cm}>{\footnotesize\raggedright\arraybackslash}p{2.20cm}>{\footnotesize\raggedright\arraybackslash}p{1.50cm}>{\footnotesize\raggedright\arraybackslash}p{6.40cm}@{}}
\toprule
{\footnotesize\textbf{Name}} & {\footnotesize\textbf{Size}} & {\footnotesize\textbf{Type}} & {\footnotesize\textbf{Units}} & {\footnotesize\textbf{Description}} \\
\midrule
\endfirsthead
\toprule
{\footnotesize\textbf{Name}} & {\footnotesize\textbf{Size}} & {\footnotesize\textbf{Type}} & {\footnotesize\textbf{Units}} & {\footnotesize\textbf{Description}} \\
\midrule
\endhead
\midrule
\endfoot
\bottomrule
\endlastfoot
T & RxN & table & - & app.perRevolution's table (t0, t1, tMid, scale, optional amp\_*) \\
\end{longtable}

\subsubsection*{NAME-VALUE ARGUMENTS}
\begin{longtable}{@{}>{\footnotesize\raggedright\arraybackslash\ttfamily}p{3.00cm}>{\footnotesize\raggedright\arraybackslash}p{1.50cm}>{\footnotesize\raggedright\arraybackslash}p{2.20cm}>{\footnotesize\raggedright\arraybackslash}p{1.50cm}>{\footnotesize\raggedright\arraybackslash}p{6.40cm}@{}}
\toprule
{\footnotesize\textbf{Name}} & {\footnotesize\textbf{Size}} & {\footnotesize\textbf{Type}} & {\footnotesize\textbf{Units}} & {\footnotesize\textbf{Description}} \\
\midrule
\endfirsthead
\toprule
{\footnotesize\textbf{Name}} & {\footnotesize\textbf{Size}} & {\footnotesize\textbf{Type}} & {\footnotesize\textbf{Units}} & {\footnotesize\textbf{Description}} \\
\midrule
\endhead
\midrule
\endfoot
\bottomrule
\endlastfoot
Orbit & 1x1 & struct & m & .epoch .r .v (GCRS, as io.readGNV1B gives them); adds betaPrime and eclipseFraction; default none \\
SpaceWeather & 1x1 & NonGravForces.\allowbreak{}io.\allowbreak{}spaceWeather & - & Adds f107, f107a, ap; default none \\
Drivers & RxN & table & - & Further drivers, one numeric column each; default none \\
Families & 1x1 & struct & - & Field = driver name, value = family; overrides and extends the defaults above \\
Significance & 1x1 & double & - & \textbar{}coef/sigma\textbar{} above which a driver counts; default 3 \\
MaxVif & 1x1 & double & - & regressDrivers' VIF gate; default 10 \\
NBoot & 1x1 & double & - & regressDrivers' replicates; default 1000 \\
Seed & 1x1 & double & - & regressDrivers' seed; default 0 \\
BlockLength & 1x1 & double & - & regressDrivers' block length in revolutions; default [] = automatic (needs \textgreater{}= 20 rows) \\
\end{longtable}

\subsubsection*{OUTPUT ARGUMENTS}
\begin{longtable}{@{}>{\footnotesize\raggedright\arraybackslash\ttfamily}p{3.00cm}>{\footnotesize\raggedright\arraybackslash}p{1.50cm}>{\footnotesize\raggedright\arraybackslash}p{2.20cm}>{\footnotesize\raggedright\arraybackslash}p{1.50cm}>{\footnotesize\raggedright\arraybackslash}p{6.40cm}@{}}
\toprule
{\footnotesize\textbf{Name}} & {\footnotesize\textbf{Size}} & {\footnotesize\textbf{Type}} & {\footnotesize\textbf{Units}} & {\footnotesize\textbf{Description}} \\
\midrule
\endfirsthead
\toprule
{\footnotesize\textbf{Name}} & {\footnotesize\textbf{Size}} & {\footnotesize\textbf{Type}} & {\footnotesize\textbf{Units}} & {\footnotesize\textbf{Description}} \\
\midrule
\endhead
\midrule
\endfoot
\bottomrule
\endlastfoot
A & Jx6 & table & - & One row per estimated per-revolution scale: parameter ("S\_xx", ...), verdict, reason, significant (string, the significant drivers), r2, nRows \\
info & 1x1 & struct & - & .drivers (R-row table of every driver offered) .families (struct) .dropped (string) .regression (1xJ cell of regressDrivers results, [] where it refused) \\
\end{longtable}

\subsubsection*{ERRORS}
\begin{description}[leftmargin=1.2em,style=nextline,itemsep=0pt]
\item[\texttt{NonGravForces:\allowbreak{}attributeRevolutions:\allowbreak{}inputTable}] T lacks t0, t1, tMid or scale
\item[\texttt{NonGravForces:\allowbreak{}attributeRevolutions:\allowbreak{}inputOrbit}] 'Orbit' without .epoch, .r, .v of matching size
\item[\texttt{NonGravForces:\allowbreak{}attributeRevolutions:\allowbreak{}inputDrivers}] 'Drivers' without one numeric row per revolution
\item[\texttt{NonGravForces:\allowbreak{}attributeRevolutions:\allowbreak{}dataNoDriver}] no driver varies over the usable revolutions
\item[\texttt{timeSeriesPreparation:\allowbreak{}regressDrivers:\allowbreak{}tooFewRows}] fewer revolutions than drivers + 2
\item[\texttt{timeSeriesPreparation:\allowbreak{}blockLength:\allowbreak{}tooShort}] fewer than 20 revolutions and no 'BlockLength'
\item[\texttt{referenceSystem:\allowbreak{}*, NonGravForces:\allowbreak{}spaceWeather:\allowbreak{}*}] from the geometry and the index lookup
\end{description}

\subsection{\texttt{NonGravForces.app.plotCalibration}}
\label{api:app.plotCalibration}
\noindent Figure of an accelerometer calibration: scale and bias with their bands, residual ACF.

\begin{lstlisting}
f = NonGravForces.app.plotCalibration(cal, t)
[f, files] = NonGravForces.app.plotCalibration(cal, t, 'Residual', info.residual, 'Visible', "off", 'ExportTo', folder)
\end{lstlisting}

\subsubsection*{INPUT ARGUMENTS}
\begin{longtable}{@{}>{\footnotesize\raggedright\arraybackslash\ttfamily}p{3.00cm}>{\footnotesize\raggedright\arraybackslash}p{1.50cm}>{\footnotesize\raggedright\arraybackslash}p{2.20cm}>{\footnotesize\raggedright\arraybackslash}p{1.50cm}>{\footnotesize\raggedright\arraybackslash}p{6.40cm}@{}}
\toprule
{\footnotesize\textbf{Name}} & {\footnotesize\textbf{Size}} & {\footnotesize\textbf{Type}} & {\footnotesize\textbf{Units}} & {\footnotesize\textbf{Description}} \\
\midrule
\endfirsthead
\toprule
{\footnotesize\textbf{Name}} & {\footnotesize\textbf{Size}} & {\footnotesize\textbf{Type}} & {\footnotesize\textbf{Units}} & {\footnotesize\textbf{Description}} \\
\midrule
\endhead
\midrule
\endfoot
\bottomrule
\endlastfoot
cal & 1x1 & NonGravForces.\allowbreak{}app.\allowbreak{}Calibration & - & The calibration to draw \\
t & Kx1 & int64 & ns & Epochs to draw the curves at, GPS ns \\
\end{longtable}

\subsubsection*{NAME-VALUE ARGUMENTS}
\begin{longtable}{@{}>{\footnotesize\raggedright\arraybackslash\ttfamily}p{3.00cm}>{\footnotesize\raggedright\arraybackslash}p{1.50cm}>{\footnotesize\raggedright\arraybackslash}p{2.20cm}>{\footnotesize\raggedright\arraybackslash}p{1.50cm}>{\footnotesize\raggedright\arraybackslash}p{6.40cm}@{}}
\toprule
{\footnotesize\textbf{Name}} & {\footnotesize\textbf{Size}} & {\footnotesize\textbf{Type}} & {\footnotesize\textbf{Units}} & {\footnotesize\textbf{Description}} \\
\midrule
\endfirsthead
\toprule
{\footnotesize\textbf{Name}} & {\footnotesize\textbf{Size}} & {\footnotesize\textbf{Type}} & {\footnotesize\textbf{Units}} & {\footnotesize\textbf{Description}} \\
\midrule
\endhead
\midrule
\endfoot
\bottomrule
\endlastfoot
Residual & 1x1 & timeSeriesPreparation.\allowbreak{}Arc & m/\allowbreak{}s\textasciicircum{}2 & app.fitCalibration's info.residual; default none \\
Order & 1x1 & double & - & AR order of the whitening; default 5 \\
MaxLag & 1x1 & double & - & ACF lags drawn; default 300 (samples) \\
Visible & 1x1 & string & - & "on" (default) or "off" \\
ExportTo & 1x1 & string & - & Folder: writes calibration.png at 200 dpi; default "" \\
\end{longtable}

\subsubsection*{OUTPUT ARGUMENTS}
\begin{longtable}{@{}>{\footnotesize\raggedright\arraybackslash\ttfamily}p{3.00cm}>{\footnotesize\raggedright\arraybackslash}p{1.50cm}>{\footnotesize\raggedright\arraybackslash}p{2.20cm}>{\footnotesize\raggedright\arraybackslash}p{1.50cm}>{\footnotesize\raggedright\arraybackslash}p{6.40cm}@{}}
\toprule
{\footnotesize\textbf{Name}} & {\footnotesize\textbf{Size}} & {\footnotesize\textbf{Type}} & {\footnotesize\textbf{Units}} & {\footnotesize\textbf{Description}} \\
\midrule
\endfirsthead
\toprule
{\footnotesize\textbf{Name}} & {\footnotesize\textbf{Size}} & {\footnotesize\textbf{Type}} & {\footnotesize\textbf{Units}} & {\footnotesize\textbf{Description}} \\
\midrule
\endhead
\midrule
\endfoot
\bottomrule
\endlastfoot
f & 1x1 & figure & - & The figure \\
files & 1xF & string & - & Files written: calibration.png in 'ExportTo', or a 1x0 string without it \\
\end{longtable}

\subsubsection*{ERRORS}
\begin{description}[leftmargin=1.2em,style=nextline,itemsep=0pt]
\item[\texttt{NonGravForces:\allowbreak{}plotCalibration:\allowbreak{}inputCalibration}] cal is not an app.Calibration
\item[\texttt{NonGravForces:\allowbreak{}plotCalibration:\allowbreak{}inputEpochs}] t is not a non-empty int64 vector
\item[\texttt{NonGravForces:\allowbreak{}plotCalibration:\allowbreak{}inputResidual}] 'Residual' is not a 3-column Arc
\item[\texttt{NonGravForces:\allowbreak{}plotCalibration:\allowbreak{}inputExport}] 'ExportTo' is not an existing folder
\item[\texttt{timeSeriesPreparation:\allowbreak{}*}] from tsAcf and tsFitAR (Yule-Walker)
\end{description}

\subsection{\texttt{NonGravForces.app.plotAttribution}}
\label{api:app.plotAttribution}
\noindent Figure of the per-revolution attribution: scales over time, scales against drivers, verdict.

\begin{lstlisting}
f = NonGravForces.app.plotAttribution(T, A, info)
[f, files] = NonGravForces.app.plotAttribution(T, A, info, 'MaxDrivers', 3, 'Visible', "off", 'ExportTo', folder)
\end{lstlisting}

\subsubsection*{INPUT ARGUMENTS}
\begin{longtable}{@{}>{\footnotesize\raggedright\arraybackslash\ttfamily}p{3.00cm}>{\footnotesize\raggedright\arraybackslash}p{1.50cm}>{\footnotesize\raggedright\arraybackslash}p{2.20cm}>{\footnotesize\raggedright\arraybackslash}p{1.50cm}>{\footnotesize\raggedright\arraybackslash}p{6.40cm}@{}}
\toprule
{\footnotesize\textbf{Name}} & {\footnotesize\textbf{Size}} & {\footnotesize\textbf{Type}} & {\footnotesize\textbf{Units}} & {\footnotesize\textbf{Description}} \\
\midrule
\endfirsthead
\toprule
{\footnotesize\textbf{Name}} & {\footnotesize\textbf{Size}} & {\footnotesize\textbf{Type}} & {\footnotesize\textbf{Units}} & {\footnotesize\textbf{Description}} \\
\midrule
\endhead
\midrule
\endfoot
\bottomrule
\endlastfoot
T & RxN & table & - & app.perRevolution's table (tMid, scale, optional sigmaScale) \\
A & Jx6 & table & - & app.attributeRevolutions' verdicts \\
info & 1x1 & struct & - & app.attributeRevolutions' info (.drivers, .regression) \\
\end{longtable}

\subsubsection*{NAME-VALUE ARGUMENTS}
\begin{longtable}{@{}>{\footnotesize\raggedright\arraybackslash\ttfamily}p{3.00cm}>{\footnotesize\raggedright\arraybackslash}p{1.50cm}>{\footnotesize\raggedright\arraybackslash}p{2.20cm}>{\footnotesize\raggedright\arraybackslash}p{1.50cm}>{\footnotesize\raggedright\arraybackslash}p{6.40cm}@{}}
\toprule
{\footnotesize\textbf{Name}} & {\footnotesize\textbf{Size}} & {\footnotesize\textbf{Type}} & {\footnotesize\textbf{Units}} & {\footnotesize\textbf{Description}} \\
\midrule
\endfirsthead
\toprule
{\footnotesize\textbf{Name}} & {\footnotesize\textbf{Size}} & {\footnotesize\textbf{Type}} & {\footnotesize\textbf{Units}} & {\footnotesize\textbf{Description}} \\
\midrule
\endhead
\midrule
\endfoot
\bottomrule
\endlastfoot
MaxDrivers & 1x1 & double & - & Driver panels per parameter, the strongest first; default 3 \\
Visible & 1x1 & string & - & "on" (default) or "off" \\
ExportTo & 1x1 & string & - & Folder: writes attribution.png at 200 dpi; default "" \\
\end{longtable}

\subsubsection*{OUTPUT ARGUMENTS}
\begin{longtable}{@{}>{\footnotesize\raggedright\arraybackslash\ttfamily}p{3.00cm}>{\footnotesize\raggedright\arraybackslash}p{1.50cm}>{\footnotesize\raggedright\arraybackslash}p{2.20cm}>{\footnotesize\raggedright\arraybackslash}p{1.50cm}>{\footnotesize\raggedright\arraybackslash}p{6.40cm}@{}}
\toprule
{\footnotesize\textbf{Name}} & {\footnotesize\textbf{Size}} & {\footnotesize\textbf{Type}} & {\footnotesize\textbf{Units}} & {\footnotesize\textbf{Description}} \\
\midrule
\endfirsthead
\toprule
{\footnotesize\textbf{Name}} & {\footnotesize\textbf{Size}} & {\footnotesize\textbf{Type}} & {\footnotesize\textbf{Units}} & {\footnotesize\textbf{Description}} \\
\midrule
\endhead
\midrule
\endfoot
\bottomrule
\endlastfoot
f & 1x1 & figure & - & The figure \\
files & 1xF & string & - & Files written: attribution.png in 'ExportTo', or a 1x0 string without it \\
\end{longtable}

\subsubsection*{ERRORS}
\begin{description}[leftmargin=1.2em,style=nextline,itemsep=0pt]
\item[\texttt{NonGravForces:\allowbreak{}plotAttribution:\allowbreak{}inputTable}] T lacks tMid or scale, or A lacks parameter or verdict
\item[\texttt{NonGravForces:\allowbreak{}plotAttribution:\allowbreak{}inputInfo}] info lacks drivers or regression, or does not match A
\item[\texttt{NonGravForces:\allowbreak{}plotAttribution:\allowbreak{}inputExport}] 'ExportTo' is not an existing folder
\end{description}


%=======================================================================================
\chapter{References}
% SHARED SECTION, \input by both NonGravForces_Theory_and_Implementation.tex and
% NonGravForces_UserGuide.tex. It carries no preamble and opens no document, so that the
% theory exists ONCE and both documents render the same words (R1).
\small
\begin{itemize}[nosep,leftmargin=*]
\item Bandikova T., McCullough C., Kruizinga G.L., Save H., Christophe B. (2019). GRACE accelerometer data transplant. Adv.\ Space Res.\ 64:623--644.
\item Behzadpour S., Mayer-Gürr T., Krauss S. (2021). GRACE Follow-On accelerometer data recovery. JGR Solid Earth 126:e2020JB021297.
\item Bettadpur S. (2012). GRACE product specification document, GRACE 327-720 (CSR-GR-03-02).
\item Bhattarai S. et al.\ (2019). Anisotropic Earth radiation pressure models. J.\ Geodesy.
\item Bowman B.R. et al.\ (2008). JB2008. AIAA 2008-6438.
\item Bruinsma S., Boniface C. (2021). DTM-2020. J.\ Space Weather Space Clim.\ 11:47.
\item Doornbos E. (2011). Thermospheric density and wind determination from satellite dynamics. PhD, TU Delft; Springer 2012.
\item Drob D.P. et al.\ (2015). HWM14. Earth and Space Science 2:301--319.
\item Emmert J.T. et al.\ (2021). NRLMSIS 2.0: a whole-atmosphere empirical model of temperature and neutral species densities. Earth and Space Science 8:e2020EA001321. \emph{The toolbox ships NRLMSIS 2.1, whose source is the same distribution; this is the model paper.}
\item Flury J., Bettadpur S., Tapley B.D. (2008). Precise accelerometry onboard GRACE. Adv.\ Space Res.\ 42:1414--1423.
\item Guo X., Zhang R. (2024). Performance of NRLMSIS 2.1 using GRACE-A and Swarm-C. Adv.\ Space Res.\ 74:2475.
\item Hu X. et al.\ (2022). Estimation of the centre of mass of GRACE-type satellites. Remote Sensing 14:4030.
\item IAGA (2024). IGRF-14. Zenodo 10.5281/zenodo.14012303.
\item Killick R., Fearnhead P., Eckley I.A. (2012). Optimal detection of changepoints with a linear computational cost. JASA 107:1590--1598.
\item Klinger B., Mayer-Gürr T. (2016). Accelerometer calibration within GRACE gravity field recovery. Adv.\ Space Res.\ 58:1597--1609.
\item Knocke P.C., Ries J.C., Tapley B.D. (1988). Earth radiation pressure effects on satellites. AIAA 88-4292.
\item Kornfeld R.P. et al.\ (2019). GRACE-FO: the extension of the legacy. J.\ Spacecraft Rockets.
\item Lévèque T. et al.\ (2021, 2022). Quantum accelerometers for space gravimetry; CARIOQA definition; rotation effects. J.\ Geodesy 95:15; arXiv:2211.01215; arXiv:2212.06686.
\item March G., Doornbos E.N., Visser P.N.A.M. (2019). High-fidelity geometry models for improving the consistency of CHAMP, GRACE, GOCE and Swarm thermospheric density data sets. Adv.\ Space Res.\ 63:213--238. DOI 10.1016/j.asr.2018.07.009. \emph{Geometry, not gas--surface interaction.}
\item March G., Visser P.N.A.M., Doornbos E.N., Siemes C. (2019). CHAMP and GOCE thermospheric wind characterization with improved gas-surface interactions modelling. Adv.\ Space Res.\ 64:1225--1242. DOI 10.1016/j.asr.2019.06.023.
\item March G., van den IJssel J., Siemes C., Visser P.N.A.M., Doornbos E.N., Pilinski M. (2021). Gas-surface interactions modelling influence on satellite aerodynamics and thermosphere mass density. J.\ Space Weather Space Clim.\ 11:54. DOI 10.1051/swsc/2021035. \emph{The fitted $\alpha_E=0.85$ every preset uses.}
\item McCullough C. et al.\ (2019). Calibrated GRACE-FO accelerometer data products (ACT). JPL D-103863.
\item Montenbruck O., Gill E. (2000). Satellite Orbits. Springer.
\item Moe K., Moe M.M. (2005). Gas--surface interactions and satellite drag coefficients. Planet.\ Space Sci.\ 53:793.
\item Peterseim N. et al.\ (2012). Magnetic torquer induced disturbing signals within GRACE accelerometer data. Adv.\ Space Res.
\item Picone J.M. et al.\ (2002). NRLMSISE-00. JGR Space Physics 107:1468.
\item Pilinski M.D., Argrow B.M., Palo S.E. (2010). Semi-empirical model for satellite energy-accommodation coefficients. J.\ Spacecraft Rockets 47(6):951--956. DOI 10.2514/1.49330.
\item Pilinski M.D., Argrow B.M., Palo S.E., Bowman B.R. (2013). Semi-empirical satellite accommodation model for spherical and randomly tumbling objects. J.\ Spacecraft Rockets 50(3):556--571. DOI 10.2514/1.A32348. \emph{SESAM; the source of \texttt{drag.accommodation = "sesam"}.}
\item Robertson R.V. (2015). Highly physical solar radiation pressure modelling with atmospheric effects (SOLAARS). PhD, Virginia Tech; Zhang et al.\ (2019).
\item Sentman L.H. (1961). Free molecule flow theory and its application to the determination of aerodynamic forces. Lockheed LMSC-448514.
\item Siemes C. et al.\ (2023). New thermosphere neutral mass density and crosswind datasets from CHAMP, GRACE and GRACE-FO. J.\ Space Weather Space Clim.\ 13:16.
\item van Helleputte T. (2009). GPS-based orbit determination and accelerometer calibration for LEO satellites. PhD, TU Delft.
\item Vielberg K., Kusche J. (2020). Extended forward and inverse modeling of radiation pressure accelerations for LEO satellites. J.\ Geodesy 94:43. DOI 10.1007/s00190-020-01368-6.
\item Wen H.Y. et al.\ (2019). GRACE-FO Level-1 data product user handbook. JPL D-56935.
\item Wöske F. (2019/2021). PhD thesis, University of Bremen (ZARM); Wöske F., Kato T., Rievers B., List M. (2019). GRACE accelerometer calibration by high-precision non-gravitational force modelling. Adv.\ Space Res.\ 63:1318--1335.
\item Walker A., Mehta P., Koller J. (2014). Drag coefficient model using the Cercignani--Lampis--Lord gas-surface interaction model. J.\ Spacecraft Rockets 51(5):1544--1563. DOI 10.2514/1.A32677. \emph{The adsorbate coverage law of \texttt{drag.accommodation = "walker2014"}; its closed forms are whole-body and normal-incidence only, which is why \texttt{drag.gsi = "cll"} is refused.}
\item Paul S.N., Sheridan P.L., Licata R.J., Mehta P.M. (2022). Comparison of local and global gas-surface interaction models. arXiv:2210.08364v2. \emph{Restates the Walker coverage law and the Mehta substrate coefficient in the form \texttt{kernel.accommodation} implements.}
\item H\l{}adczuk N.A., van den IJssel J., Kodikara T., Siemes C., Visser P. (2024). GRACE-FO radiation pressure modelling for accurate density and crosswind retrieval. Adv.\ Space Res.\ 73:2355--2373. \emph{Table 3 grades the thermal and geometry choices; the source of the 4.22 and 3.04 nm/s$^2$ figures quoted in the limitations.}
\item Zahzam N. et al.\ (2022). Hybrid electrostatic--atomic accelerometer for future space gravity missions. Remote Sensing 14:3273.
\end{itemize}


\end{document}

